\documentclass[11pt]{article}

    \usepackage[breakable]{tcolorbox}
    \usepackage{parskip} % Stop auto-indenting (to mimic markdown behaviour)
    

    % Basic figure setup, for now with no caption control since it's done
    % automatically by Pandoc (which extracts ![](path) syntax from Markdown).
    \usepackage{graphicx}
    % Maintain compatibility with old templates. Remove in nbconvert 6.0
    \let\Oldincludegraphics\includegraphics
    % Ensure that by default, figures have no caption (until we provide a
    % proper Figure object with a Caption API and a way to capture that
    % in the conversion process - todo).
    \usepackage{caption}
    \DeclareCaptionFormat{nocaption}{}
    \captionsetup{format=nocaption,aboveskip=0pt,belowskip=0pt}

    \usepackage{float}
    \floatplacement{figure}{H} % forces figures to be placed at the correct location
    \usepackage{xcolor} % Allow colors to be defined
    \usepackage{enumerate} % Needed for markdown enumerations to work
    \usepackage{geometry} % Used to adjust the document margins
    \usepackage{amsmath} % Equations
    \usepackage{amssymb} % Equations
    \usepackage{textcomp} % defines textquotesingle
    % Hack from http://tex.stackexchange.com/a/47451/13684:
    \AtBeginDocument{%
        \def\PYZsq{\textquotesingle}% Upright quotes in Pygmentized code
    }
    \usepackage{upquote} % Upright quotes for verbatim code
    \usepackage{eurosym} % defines \euro

    \usepackage{iftex}
    \ifPDFTeX
        \usepackage[T1]{fontenc}
        \IfFileExists{alphabeta.sty}{
              \usepackage{alphabeta}
          }{
              \usepackage[mathletters]{ucs}
              \usepackage[utf8x]{inputenc}
          }
    \else
        \usepackage{fontspec}
        \usepackage{unicode-math}
    \fi

    \usepackage{fancyvrb} % verbatim replacement that allows latex
    \usepackage{grffile} % extends the file name processing of package graphics
                         % to support a larger range
    \makeatletter % fix for old versions of grffile with XeLaTeX
    \@ifpackagelater{grffile}{2019/11/01}
    {
      % Do nothing on new versions
    }
    {
      \def\Gread@@xetex#1{%
        \IfFileExists{"\Gin@base".bb}%
        {\Gread@eps{\Gin@base.bb}}%
        {\Gread@@xetex@aux#1}%
      }
    }
    \makeatother
    \usepackage[Export]{adjustbox} % Used to constrain images to a maximum size
    \adjustboxset{max size={0.9\linewidth}{0.9\paperheight}}

    % The hyperref package gives us a pdf with properly built
    % internal navigation ('pdf bookmarks' for the table of contents,
    % internal cross-reference links, web links for URLs, etc.)
    \usepackage{hyperref}
    % The default LaTeX title has an obnoxious amount of whitespace. By default,
    % titling removes some of it. It also provides customization options.
    \usepackage{titling}
    \usepackage{longtable} % longtable support required by pandoc >1.10
    \usepackage{booktabs}  % table support for pandoc > 1.12.2
    \usepackage{array}     % table support for pandoc >= 2.11.3
    \usepackage{calc}      % table minipage width calculation for pandoc >= 2.11.1
    \usepackage[inline]{enumitem} % IRkernel/repr support (it uses the enumerate* environment)
    \usepackage[normalem]{ulem} % ulem is needed to support strikethroughs (\sout)
                                % normalem makes italics be italics, not underlines
    \usepackage{soul}      % strikethrough (\st) support for pandoc >= 3.0.0
    \usepackage{mathrsfs}
    

    
    % Colors for the hyperref package
    \definecolor{urlcolor}{rgb}{0,.145,.698}
    \definecolor{linkcolor}{rgb}{.71,0.21,0.01}
    \definecolor{citecolor}{rgb}{.12,.54,.11}

    % ANSI colors
    \definecolor{ansi-black}{HTML}{3E424D}
    \definecolor{ansi-black-intense}{HTML}{282C36}
    \definecolor{ansi-red}{HTML}{E75C58}
    \definecolor{ansi-red-intense}{HTML}{B22B31}
    \definecolor{ansi-green}{HTML}{00A250}
    \definecolor{ansi-green-intense}{HTML}{007427}
    \definecolor{ansi-yellow}{HTML}{DDB62B}
    \definecolor{ansi-yellow-intense}{HTML}{B27D12}
    \definecolor{ansi-blue}{HTML}{208FFB}
    \definecolor{ansi-blue-intense}{HTML}{0065CA}
    \definecolor{ansi-magenta}{HTML}{D160C4}
    \definecolor{ansi-magenta-intense}{HTML}{A03196}
    \definecolor{ansi-cyan}{HTML}{60C6C8}
    \definecolor{ansi-cyan-intense}{HTML}{258F8F}
    \definecolor{ansi-white}{HTML}{C5C1B4}
    \definecolor{ansi-white-intense}{HTML}{A1A6B2}
    \definecolor{ansi-default-inverse-fg}{HTML}{FFFFFF}
    \definecolor{ansi-default-inverse-bg}{HTML}{000000}

    % common color for the border for error outputs.
    \definecolor{outerrorbackground}{HTML}{FFDFDF}

    % commands and environments needed by pandoc snippets
    % extracted from the output of `pandoc -s`
    \providecommand{\tightlist}{%
      \setlength{\itemsep}{0pt}\setlength{\parskip}{0pt}}
    \DefineVerbatimEnvironment{Highlighting}{Verbatim}{commandchars=\\\{\}}
    % Add ',fontsize=\small' for more characters per line
    \newenvironment{Shaded}{}{}
    \newcommand{\KeywordTok}[1]{\textcolor[rgb]{0.00,0.44,0.13}{\textbf{{#1}}}}
    \newcommand{\DataTypeTok}[1]{\textcolor[rgb]{0.56,0.13,0.00}{{#1}}}
    \newcommand{\DecValTok}[1]{\textcolor[rgb]{0.25,0.63,0.44}{{#1}}}
    \newcommand{\BaseNTok}[1]{\textcolor[rgb]{0.25,0.63,0.44}{{#1}}}
    \newcommand{\FloatTok}[1]{\textcolor[rgb]{0.25,0.63,0.44}{{#1}}}
    \newcommand{\CharTok}[1]{\textcolor[rgb]{0.25,0.44,0.63}{{#1}}}
    \newcommand{\StringTok}[1]{\textcolor[rgb]{0.25,0.44,0.63}{{#1}}}
    \newcommand{\CommentTok}[1]{\textcolor[rgb]{0.38,0.63,0.69}{\textit{{#1}}}}
    \newcommand{\OtherTok}[1]{\textcolor[rgb]{0.00,0.44,0.13}{{#1}}}
    \newcommand{\AlertTok}[1]{\textcolor[rgb]{1.00,0.00,0.00}{\textbf{{#1}}}}
    \newcommand{\FunctionTok}[1]{\textcolor[rgb]{0.02,0.16,0.49}{{#1}}}
    \newcommand{\RegionMarkerTok}[1]{{#1}}
    \newcommand{\ErrorTok}[1]{\textcolor[rgb]{1.00,0.00,0.00}{\textbf{{#1}}}}
    \newcommand{\NormalTok}[1]{{#1}}

    % Additional commands for more recent versions of Pandoc
    \newcommand{\ConstantTok}[1]{\textcolor[rgb]{0.53,0.00,0.00}{{#1}}}
    \newcommand{\SpecialCharTok}[1]{\textcolor[rgb]{0.25,0.44,0.63}{{#1}}}
    \newcommand{\VerbatimStringTok}[1]{\textcolor[rgb]{0.25,0.44,0.63}{{#1}}}
    \newcommand{\SpecialStringTok}[1]{\textcolor[rgb]{0.73,0.40,0.53}{{#1}}}
    \newcommand{\ImportTok}[1]{{#1}}
    \newcommand{\DocumentationTok}[1]{\textcolor[rgb]{0.73,0.13,0.13}{\textit{{#1}}}}
    \newcommand{\AnnotationTok}[1]{\textcolor[rgb]{0.38,0.63,0.69}{\textbf{\textit{{#1}}}}}
    \newcommand{\CommentVarTok}[1]{\textcolor[rgb]{0.38,0.63,0.69}{\textbf{\textit{{#1}}}}}
    \newcommand{\VariableTok}[1]{\textcolor[rgb]{0.10,0.09,0.49}{{#1}}}
    \newcommand{\ControlFlowTok}[1]{\textcolor[rgb]{0.00,0.44,0.13}{\textbf{{#1}}}}
    \newcommand{\OperatorTok}[1]{\textcolor[rgb]{0.40,0.40,0.40}{{#1}}}
    \newcommand{\BuiltInTok}[1]{{#1}}
    \newcommand{\ExtensionTok}[1]{{#1}}
    \newcommand{\PreprocessorTok}[1]{\textcolor[rgb]{0.74,0.48,0.00}{{#1}}}
    \newcommand{\AttributeTok}[1]{\textcolor[rgb]{0.49,0.56,0.16}{{#1}}}
    \newcommand{\InformationTok}[1]{\textcolor[rgb]{0.38,0.63,0.69}{\textbf{\textit{{#1}}}}}
    \newcommand{\WarningTok}[1]{\textcolor[rgb]{0.38,0.63,0.69}{\textbf{\textit{{#1}}}}}


    % Define a nice break command that doesn't care if a line doesn't already
    % exist.
    \def\br{\hspace*{\fill} \\* }
    % Math Jax compatibility definitions
    \def\gt{>}
    \def\lt{<}
    \let\Oldtex\TeX
    \let\Oldlatex\LaTeX
    \renewcommand{\TeX}{\textrm{\Oldtex}}
    \renewcommand{\LaTeX}{\textrm{\Oldlatex}}
    % Document parameters
    % Document title
    \title{ModeloPrediccionAgropecuaria}
    
    
    
    
    
    
    
% Pygments definitions
\makeatletter
\def\PY@reset{\let\PY@it=\relax \let\PY@bf=\relax%
    \let\PY@ul=\relax \let\PY@tc=\relax%
    \let\PY@bc=\relax \let\PY@ff=\relax}
\def\PY@tok#1{\csname PY@tok@#1\endcsname}
\def\PY@toks#1+{\ifx\relax#1\empty\else%
    \PY@tok{#1}\expandafter\PY@toks\fi}
\def\PY@do#1{\PY@bc{\PY@tc{\PY@ul{%
    \PY@it{\PY@bf{\PY@ff{#1}}}}}}}
\def\PY#1#2{\PY@reset\PY@toks#1+\relax+\PY@do{#2}}

\@namedef{PY@tok@w}{\def\PY@tc##1{\textcolor[rgb]{0.73,0.73,0.73}{##1}}}
\@namedef{PY@tok@c}{\let\PY@it=\textit\def\PY@tc##1{\textcolor[rgb]{0.24,0.48,0.48}{##1}}}
\@namedef{PY@tok@cp}{\def\PY@tc##1{\textcolor[rgb]{0.61,0.40,0.00}{##1}}}
\@namedef{PY@tok@k}{\let\PY@bf=\textbf\def\PY@tc##1{\textcolor[rgb]{0.00,0.50,0.00}{##1}}}
\@namedef{PY@tok@kp}{\def\PY@tc##1{\textcolor[rgb]{0.00,0.50,0.00}{##1}}}
\@namedef{PY@tok@kt}{\def\PY@tc##1{\textcolor[rgb]{0.69,0.00,0.25}{##1}}}
\@namedef{PY@tok@o}{\def\PY@tc##1{\textcolor[rgb]{0.40,0.40,0.40}{##1}}}
\@namedef{PY@tok@ow}{\let\PY@bf=\textbf\def\PY@tc##1{\textcolor[rgb]{0.67,0.13,1.00}{##1}}}
\@namedef{PY@tok@nb}{\def\PY@tc##1{\textcolor[rgb]{0.00,0.50,0.00}{##1}}}
\@namedef{PY@tok@nf}{\def\PY@tc##1{\textcolor[rgb]{0.00,0.00,1.00}{##1}}}
\@namedef{PY@tok@nc}{\let\PY@bf=\textbf\def\PY@tc##1{\textcolor[rgb]{0.00,0.00,1.00}{##1}}}
\@namedef{PY@tok@nn}{\let\PY@bf=\textbf\def\PY@tc##1{\textcolor[rgb]{0.00,0.00,1.00}{##1}}}
\@namedef{PY@tok@ne}{\let\PY@bf=\textbf\def\PY@tc##1{\textcolor[rgb]{0.80,0.25,0.22}{##1}}}
\@namedef{PY@tok@nv}{\def\PY@tc##1{\textcolor[rgb]{0.10,0.09,0.49}{##1}}}
\@namedef{PY@tok@no}{\def\PY@tc##1{\textcolor[rgb]{0.53,0.00,0.00}{##1}}}
\@namedef{PY@tok@nl}{\def\PY@tc##1{\textcolor[rgb]{0.46,0.46,0.00}{##1}}}
\@namedef{PY@tok@ni}{\let\PY@bf=\textbf\def\PY@tc##1{\textcolor[rgb]{0.44,0.44,0.44}{##1}}}
\@namedef{PY@tok@na}{\def\PY@tc##1{\textcolor[rgb]{0.41,0.47,0.13}{##1}}}
\@namedef{PY@tok@nt}{\let\PY@bf=\textbf\def\PY@tc##1{\textcolor[rgb]{0.00,0.50,0.00}{##1}}}
\@namedef{PY@tok@nd}{\def\PY@tc##1{\textcolor[rgb]{0.67,0.13,1.00}{##1}}}
\@namedef{PY@tok@s}{\def\PY@tc##1{\textcolor[rgb]{0.73,0.13,0.13}{##1}}}
\@namedef{PY@tok@sd}{\let\PY@it=\textit\def\PY@tc##1{\textcolor[rgb]{0.73,0.13,0.13}{##1}}}
\@namedef{PY@tok@si}{\let\PY@bf=\textbf\def\PY@tc##1{\textcolor[rgb]{0.64,0.35,0.47}{##1}}}
\@namedef{PY@tok@se}{\let\PY@bf=\textbf\def\PY@tc##1{\textcolor[rgb]{0.67,0.36,0.12}{##1}}}
\@namedef{PY@tok@sr}{\def\PY@tc##1{\textcolor[rgb]{0.64,0.35,0.47}{##1}}}
\@namedef{PY@tok@ss}{\def\PY@tc##1{\textcolor[rgb]{0.10,0.09,0.49}{##1}}}
\@namedef{PY@tok@sx}{\def\PY@tc##1{\textcolor[rgb]{0.00,0.50,0.00}{##1}}}
\@namedef{PY@tok@m}{\def\PY@tc##1{\textcolor[rgb]{0.40,0.40,0.40}{##1}}}
\@namedef{PY@tok@gh}{\let\PY@bf=\textbf\def\PY@tc##1{\textcolor[rgb]{0.00,0.00,0.50}{##1}}}
\@namedef{PY@tok@gu}{\let\PY@bf=\textbf\def\PY@tc##1{\textcolor[rgb]{0.50,0.00,0.50}{##1}}}
\@namedef{PY@tok@gd}{\def\PY@tc##1{\textcolor[rgb]{0.63,0.00,0.00}{##1}}}
\@namedef{PY@tok@gi}{\def\PY@tc##1{\textcolor[rgb]{0.00,0.52,0.00}{##1}}}
\@namedef{PY@tok@gr}{\def\PY@tc##1{\textcolor[rgb]{0.89,0.00,0.00}{##1}}}
\@namedef{PY@tok@ge}{\let\PY@it=\textit}
\@namedef{PY@tok@gs}{\let\PY@bf=\textbf}
\@namedef{PY@tok@ges}{\let\PY@bf=\textbf\let\PY@it=\textit}
\@namedef{PY@tok@gp}{\let\PY@bf=\textbf\def\PY@tc##1{\textcolor[rgb]{0.00,0.00,0.50}{##1}}}
\@namedef{PY@tok@go}{\def\PY@tc##1{\textcolor[rgb]{0.44,0.44,0.44}{##1}}}
\@namedef{PY@tok@gt}{\def\PY@tc##1{\textcolor[rgb]{0.00,0.27,0.87}{##1}}}
\@namedef{PY@tok@err}{\def\PY@bc##1{{\setlength{\fboxsep}{\string -\fboxrule}\fcolorbox[rgb]{1.00,0.00,0.00}{1,1,1}{\strut ##1}}}}
\@namedef{PY@tok@kc}{\let\PY@bf=\textbf\def\PY@tc##1{\textcolor[rgb]{0.00,0.50,0.00}{##1}}}
\@namedef{PY@tok@kd}{\let\PY@bf=\textbf\def\PY@tc##1{\textcolor[rgb]{0.00,0.50,0.00}{##1}}}
\@namedef{PY@tok@kn}{\let\PY@bf=\textbf\def\PY@tc##1{\textcolor[rgb]{0.00,0.50,0.00}{##1}}}
\@namedef{PY@tok@kr}{\let\PY@bf=\textbf\def\PY@tc##1{\textcolor[rgb]{0.00,0.50,0.00}{##1}}}
\@namedef{PY@tok@bp}{\def\PY@tc##1{\textcolor[rgb]{0.00,0.50,0.00}{##1}}}
\@namedef{PY@tok@fm}{\def\PY@tc##1{\textcolor[rgb]{0.00,0.00,1.00}{##1}}}
\@namedef{PY@tok@vc}{\def\PY@tc##1{\textcolor[rgb]{0.10,0.09,0.49}{##1}}}
\@namedef{PY@tok@vg}{\def\PY@tc##1{\textcolor[rgb]{0.10,0.09,0.49}{##1}}}
\@namedef{PY@tok@vi}{\def\PY@tc##1{\textcolor[rgb]{0.10,0.09,0.49}{##1}}}
\@namedef{PY@tok@vm}{\def\PY@tc##1{\textcolor[rgb]{0.10,0.09,0.49}{##1}}}
\@namedef{PY@tok@sa}{\def\PY@tc##1{\textcolor[rgb]{0.73,0.13,0.13}{##1}}}
\@namedef{PY@tok@sb}{\def\PY@tc##1{\textcolor[rgb]{0.73,0.13,0.13}{##1}}}
\@namedef{PY@tok@sc}{\def\PY@tc##1{\textcolor[rgb]{0.73,0.13,0.13}{##1}}}
\@namedef{PY@tok@dl}{\def\PY@tc##1{\textcolor[rgb]{0.73,0.13,0.13}{##1}}}
\@namedef{PY@tok@s2}{\def\PY@tc##1{\textcolor[rgb]{0.73,0.13,0.13}{##1}}}
\@namedef{PY@tok@sh}{\def\PY@tc##1{\textcolor[rgb]{0.73,0.13,0.13}{##1}}}
\@namedef{PY@tok@s1}{\def\PY@tc##1{\textcolor[rgb]{0.73,0.13,0.13}{##1}}}
\@namedef{PY@tok@mb}{\def\PY@tc##1{\textcolor[rgb]{0.40,0.40,0.40}{##1}}}
\@namedef{PY@tok@mf}{\def\PY@tc##1{\textcolor[rgb]{0.40,0.40,0.40}{##1}}}
\@namedef{PY@tok@mh}{\def\PY@tc##1{\textcolor[rgb]{0.40,0.40,0.40}{##1}}}
\@namedef{PY@tok@mi}{\def\PY@tc##1{\textcolor[rgb]{0.40,0.40,0.40}{##1}}}
\@namedef{PY@tok@il}{\def\PY@tc##1{\textcolor[rgb]{0.40,0.40,0.40}{##1}}}
\@namedef{PY@tok@mo}{\def\PY@tc##1{\textcolor[rgb]{0.40,0.40,0.40}{##1}}}
\@namedef{PY@tok@ch}{\let\PY@it=\textit\def\PY@tc##1{\textcolor[rgb]{0.24,0.48,0.48}{##1}}}
\@namedef{PY@tok@cm}{\let\PY@it=\textit\def\PY@tc##1{\textcolor[rgb]{0.24,0.48,0.48}{##1}}}
\@namedef{PY@tok@cpf}{\let\PY@it=\textit\def\PY@tc##1{\textcolor[rgb]{0.24,0.48,0.48}{##1}}}
\@namedef{PY@tok@c1}{\let\PY@it=\textit\def\PY@tc##1{\textcolor[rgb]{0.24,0.48,0.48}{##1}}}
\@namedef{PY@tok@cs}{\let\PY@it=\textit\def\PY@tc##1{\textcolor[rgb]{0.24,0.48,0.48}{##1}}}

\def\PYZbs{\char`\\}
\def\PYZus{\char`\_}
\def\PYZob{\char`\{}
\def\PYZcb{\char`\}}
\def\PYZca{\char`\^}
\def\PYZam{\char`\&}
\def\PYZlt{\char`\<}
\def\PYZgt{\char`\>}
\def\PYZsh{\char`\#}
\def\PYZpc{\char`\%}
\def\PYZdl{\char`\$}
\def\PYZhy{\char`\-}
\def\PYZsq{\char`\'}
\def\PYZdq{\char`\"}
\def\PYZti{\char`\~}
% for compatibility with earlier versions
\def\PYZat{@}
\def\PYZlb{[}
\def\PYZrb{]}
\makeatother


    % For linebreaks inside Verbatim environment from package fancyvrb.
    \makeatletter
        \newbox\Wrappedcontinuationbox
        \newbox\Wrappedvisiblespacebox
        \newcommand*\Wrappedvisiblespace {\textcolor{red}{\textvisiblespace}}
        \newcommand*\Wrappedcontinuationsymbol {\textcolor{red}{\llap{\tiny$\m@th\hookrightarrow$}}}
        \newcommand*\Wrappedcontinuationindent {3ex }
        \newcommand*\Wrappedafterbreak {\kern\Wrappedcontinuationindent\copy\Wrappedcontinuationbox}
        % Take advantage of the already applied Pygments mark-up to insert
        % potential linebreaks for TeX processing.
        %        {, <, #, %, $, ' and ": go to next line.
        %        _, }, ^, &, >, - and ~: stay at end of broken line.
        % Use of \textquotesingle for straight quote.
        \newcommand*\Wrappedbreaksatspecials {%
            \def\PYGZus{\discretionary{\char`\_}{\Wrappedafterbreak}{\char`\_}}%
            \def\PYGZob{\discretionary{}{\Wrappedafterbreak\char`\{}{\char`\{}}%
            \def\PYGZcb{\discretionary{\char`\}}{\Wrappedafterbreak}{\char`\}}}%
            \def\PYGZca{\discretionary{\char`\^}{\Wrappedafterbreak}{\char`\^}}%
            \def\PYGZam{\discretionary{\char`\&}{\Wrappedafterbreak}{\char`\&}}%
            \def\PYGZlt{\discretionary{}{\Wrappedafterbreak\char`\<}{\char`\<}}%
            \def\PYGZgt{\discretionary{\char`\>}{\Wrappedafterbreak}{\char`\>}}%
            \def\PYGZsh{\discretionary{}{\Wrappedafterbreak\char`\#}{\char`\#}}%
            \def\PYGZpc{\discretionary{}{\Wrappedafterbreak\char`\%}{\char`\%}}%
            \def\PYGZdl{\discretionary{}{\Wrappedafterbreak\char`\$}{\char`\$}}%
            \def\PYGZhy{\discretionary{\char`\-}{\Wrappedafterbreak}{\char`\-}}%
            \def\PYGZsq{\discretionary{}{\Wrappedafterbreak\textquotesingle}{\textquotesingle}}%
            \def\PYGZdq{\discretionary{}{\Wrappedafterbreak\char`\"}{\char`\"}}%
            \def\PYGZti{\discretionary{\char`\~}{\Wrappedafterbreak}{\char`\~}}%
        }
        % Some characters . , ; ? ! / are not pygmentized.
        % This macro makes them "active" and they will insert potential linebreaks
        \newcommand*\Wrappedbreaksatpunct {%
            \lccode`\~`\.\lowercase{\def~}{\discretionary{\hbox{\char`\.}}{\Wrappedafterbreak}{\hbox{\char`\.}}}%
            \lccode`\~`\,\lowercase{\def~}{\discretionary{\hbox{\char`\,}}{\Wrappedafterbreak}{\hbox{\char`\,}}}%
            \lccode`\~`\;\lowercase{\def~}{\discretionary{\hbox{\char`\;}}{\Wrappedafterbreak}{\hbox{\char`\;}}}%
            \lccode`\~`\:\lowercase{\def~}{\discretionary{\hbox{\char`\:}}{\Wrappedafterbreak}{\hbox{\char`\:}}}%
            \lccode`\~`\?\lowercase{\def~}{\discretionary{\hbox{\char`\?}}{\Wrappedafterbreak}{\hbox{\char`\?}}}%
            \lccode`\~`\!\lowercase{\def~}{\discretionary{\hbox{\char`\!}}{\Wrappedafterbreak}{\hbox{\char`\!}}}%
            \lccode`\~`\/\lowercase{\def~}{\discretionary{\hbox{\char`\/}}{\Wrappedafterbreak}{\hbox{\char`\/}}}%
            \catcode`\.\active
            \catcode`\,\active
            \catcode`\;\active
            \catcode`\:\active
            \catcode`\?\active
            \catcode`\!\active
            \catcode`\/\active
            \lccode`\~`\~
        }
    \makeatother

    \let\OriginalVerbatim=\Verbatim
    \makeatletter
    \renewcommand{\Verbatim}[1][1]{%
        %\parskip\z@skip
        \sbox\Wrappedcontinuationbox {\Wrappedcontinuationsymbol}%
        \sbox\Wrappedvisiblespacebox {\FV@SetupFont\Wrappedvisiblespace}%
        \def\FancyVerbFormatLine ##1{\hsize\linewidth
            \vtop{\raggedright\hyphenpenalty\z@\exhyphenpenalty\z@
                \doublehyphendemerits\z@\finalhyphendemerits\z@
                \strut ##1\strut}%
        }%
        % If the linebreak is at a space, the latter will be displayed as visible
        % space at end of first line, and a continuation symbol starts next line.
        % Stretch/shrink are however usually zero for typewriter font.
        \def\FV@Space {%
            \nobreak\hskip\z@ plus\fontdimen3\font minus\fontdimen4\font
            \discretionary{\copy\Wrappedvisiblespacebox}{\Wrappedafterbreak}
            {\kern\fontdimen2\font}%
        }%

        % Allow breaks at special characters using \PYG... macros.
        \Wrappedbreaksatspecials
        % Breaks at punctuation characters . , ; ? ! and / need catcode=\active
        \OriginalVerbatim[#1,codes*=\Wrappedbreaksatpunct]%
    }
    \makeatother

    % Exact colors from NB
    \definecolor{incolor}{HTML}{303F9F}
    \definecolor{outcolor}{HTML}{D84315}
    \definecolor{cellborder}{HTML}{CFCFCF}
    \definecolor{cellbackground}{HTML}{F7F7F7}

    % prompt
    \makeatletter
    \newcommand{\boxspacing}{\kern\kvtcb@left@rule\kern\kvtcb@boxsep}
    \makeatother
    \newcommand{\prompt}[4]{
        {\ttfamily\llap{{\color{#2}[#3]:\hspace{3pt}#4}}\vspace{-\baselineskip}}
    }
    

    
    % Prevent overflowing lines due to hard-to-break entities
    \sloppy
    % Setup hyperref package
    \hypersetup{
      breaklinks=true,  % so long urls are correctly broken across lines
      colorlinks=true,
      urlcolor=urlcolor,
      linkcolor=linkcolor,
      citecolor=citecolor,
      }
    % Slightly bigger margins than the latex defaults
    
    \geometry{verbose,tmargin=1in,bmargin=1in,lmargin=1in,rmargin=1in}
    
    

\begin{document}
    
    \maketitle
    
    

    
    \hypertarget{mercado-justo}{%
\section{Mercado Justo:}\label{mercado-justo}}

\hypertarget{anuxe1lisis-y-predicciuxf3n-del-precio-nacional-del-mauxedz-blanco-en-muxe9xico-mediante-series-de-tiempo-y-variables-exuxf3genas}{%
\subsection{Análisis y Predicción del Precio Nacional del Maíz Blanco en
México mediante Series de Tiempo y Variables
Exógenas}\label{anuxe1lisis-y-predicciuxf3n-del-precio-nacional-del-mauxedz-blanco-en-muxe9xico-mediante-series-de-tiempo-y-variables-exuxf3genas}}

\hypertarget{licenciatura-en-ciencia-de-datos-para-negocios}{%
\subsubsection{Licenciatura en Ciencia de Datos para
Negocios}\label{licenciatura-en-ciencia-de-datos-para-negocios}}

Universidad Nacional Rosario Castellanos

\hypertarget{problema-prototuxedpico-5o-semestre}{%
\subsubsection{Problema Prototípico 5o
semestre}\label{problema-prototuxedpico-5o-semestre}}

Mercado Justo: Análisis y predicción de precios de productos
agropecuarios

\hypertarget{integrantes}{%
\subsubsection{Integrantes}\label{integrantes}}

\begin{itemize}
\tightlist
\item
  Marco Ramírez Arizpe
\item
  Eder Martínez Martínez
\end{itemize}

\hypertarget{objetivo-general}{%
\subsubsection{Objetivo General}\label{objetivo-general}}

Desarrollar un modelo predictivo basado en series de tiempo para estimar
el comportamiento futuro del precio nacional del maíz blanco en México,
incorporando variables económicas, climáticas y productivas que permitan
apoyar estrategias de comercialización más justas para pequeños
productores.

    \hypertarget{planteamiento-del-problema}{%
\section{1. Planteamiento del
Problema}\label{planteamiento-del-problema}}

El maíz blanco constituye uno de los principales cultivos estratégicos
de México debido a su relevancia económica, social y alimentaria.

Sin embargo, los pequeños productores enfrentan una alta volatilidad en
los precios de comercialización debido a factores como:

\begin{itemize}
\tightlist
\item
  Variaciones en la producción nacional.
\item
  Cambios climáticos.
\item
  Sequías.
\item
  Inflación.
\item
  Tipo de cambio.
\item
  Costos de transporte.
\item
  Precio internacional del maíz.
\item
  Costos energéticos.
\end{itemize}

Esta incertidumbre dificulta la toma de decisiones relacionadas con:

\begin{itemize}
\tightlist
\item
  Venta inmediata.
\item
  Almacenamiento.
\item
  Coberturas de precio.
\item
  Comercialización regional.
\end{itemize}

Por ello se propone construir un modelo predictivo basado en series de
tiempo que permita anticipar el comportamiento futuro del precio
nacional del maíz blanco.

    \hypertarget{objetivos}{%
\section{2. Objetivos}\label{objetivos}}

\hypertarget{objetivo-general}{%
\subsection{Objetivo General}\label{objetivo-general}}

Construir un modelo SARIMAX para pronosticar el precio nacional del maíz
blanco en México utilizando variables económicas, climáticas y
productivas.

\hypertarget{objetivos-especuxedficos}{%
\subsection{Objetivos Específicos}\label{objetivos-especuxedficos}}

\begin{itemize}
\tightlist
\item
  Integrar una base de datos histórica nacional.
\item
  Analizar tendencias y estacionalidad.
\item
  Identificar relaciones entre variables exógenas.
\item
  Construir modelos ARIMA, SARIMA y SARIMAX.
\item
  Comparar métricas de desempeño.
\item
  Generar recomendaciones para productores.
\item
  Analizar el impacto potencial sobre el precio de la tortilla.
\end{itemize}

    \hypertarget{metodologuxeda-crisp-dm}{%
\section{3. Metodología CRISP-DM}\label{metodologuxeda-crisp-dm}}

Para el desarrollo de este proyecto se emplea la metodología
\textbf{CRISP-DM (Cross Industry Standard Process for Data Mining)}, uno
de los marcos de trabajo más utilizados en proyectos de minería de datos
y ciencia de datos.

Esta metodología proporciona una estructura sistemática para transformar
datos en conocimiento útil mediante un proceso iterativo compuesto por
seis fases:

\hypertarget{comprensiuxf3n-del-negocio}{%
\subsection{1. Comprensión del
negocio}\label{comprensiuxf3n-del-negocio}}

Se identifica el problema de investigación: la volatilidad del precio
nacional del maíz blanco y su impacto sobre la toma de decisiones de
comercialización para pequeños productores agrícolas.

\hypertarget{comprensiuxf3n-de-los-datos}{%
\subsection{2. Comprensión de los
datos}\label{comprensiuxf3n-de-los-datos}}

Se recopilan y analizan fuentes históricas relacionadas con el mercado
del maíz, variables económicas, climáticas, productivas y logísticas que
pueden influir en la formación de precios.

\hypertarget{preparaciuxf3n-de-los-datos}{%
\subsection{3. Preparación de los
datos}\label{preparaciuxf3n-de-los-datos}}

Se integran múltiples fuentes de información, se corrigen
inconsistencias, se homologan frecuencias temporales y se construye una
base de datos consolidada para análisis.

\hypertarget{modelado}{%
\subsection{4. Modelado}\label{modelado}}

Se desarrollan modelos de series de tiempo ARIMA, SARIMA y SARIMAX con
el fin de identificar la alternativa con mejor capacidad predictiva.

\hypertarget{evaluaciuxf3n}{%
\subsection{5. Evaluación}\label{evaluaciuxf3n}}

Los modelos se comparan mediante métricas de error como MAE, RMSE y
MAPE, verificando su capacidad para representar adecuadamente el
comportamiento histórico de la serie.

\hypertarget{despliegue-y-generaciuxf3n-de-valor}{%
\subsection{6. Despliegue y generación de
valor}\label{despliegue-y-generaciuxf3n-de-valor}}

Los resultados se traducen en recomendaciones de comercialización para
productores agrícolas y se analiza la posible transmisión de precios
hacia productos derivados como la tortilla.

La siguiente figura resume el flujo metodológico seguido durante el
proyecto.

    \hypertarget{descripciuxf3n-general-del-dataset}{%
\section{4. Descripción General del
Dataset}\label{descripciuxf3n-general-del-dataset}}

La base de datos utilizada en este proyecto integra información
histórica proveniente de organismos gubernamentales, instituciones
financieras, organismos reguladores y fuentes especializadas
relacionadas con el mercado agroalimentario mexicano.

El objetivo de esta integración es construir una serie temporal
multivariable capaz de representar los principales factores que influyen
sobre el precio nacional del maíz blanco en México.

La información fue sometida a un proceso de integración, depuración y
transformación para homologar todas las variables a frecuencia mensual y
facilitar su análisis conjunto.

La base final incorpora variables pertenecientes a cuatro dimensiones
principales:

\hypertarget{variables-econuxf3micas}{%
\subsubsection{Variables económicas}\label{variables-econuxf3micas}}

\begin{itemize}
\tightlist
\item
  Tipo de cambio.
\item
  Inflación.
\item
  Precio internacional del maíz.
\end{itemize}

\hypertarget{variables-productivas}{%
\subsubsection{Variables productivas}\label{variables-productivas}}

\begin{itemize}
\tightlist
\item
  Producción nacional.
\item
  Demanda nacional.
\end{itemize}

\hypertarget{variables-climuxe1ticas}{%
\subsubsection{Variables climáticas}\label{variables-climuxe1ticas}}

\begin{itemize}
\tightlist
\item
  Temperatura.
\item
  Precipitación.
\item
  Severidad de sequía.
\end{itemize}

\hypertarget{variables-loguxedsticas}{%
\subsubsection{Variables logísticas}\label{variables-loguxedsticas}}

\begin{itemize}
\tightlist
\item
  Precio del diésel.
\item
  Costos de transporte.
\end{itemize}

La variable objetivo principal corresponde al precio nacional del maíz
blanco al mayoreo, mientras que las demás variables se consideran
candidatas a variables exógenas dentro de modelos de series de tiempo.

Adicionalmente, como extensión del análisis, se explorará la posible
transmisión de precios desde el mercado primario del maíz hacia el
precio de la tortilla como producto de consumo final.

    \hypertarget{diccionario-de-variables}{%
\section{5. Diccionario de Variables}\label{diccionario-de-variables}}

La siguiente tabla describe las variables incluidas en la base de datos
consolidada utilizada para el análisis y modelado del precio nacional
del maíz blanco. Todas las variables fueron homologadas a frecuencia
mensual y, cuando fue necesario, transformadas a unidades comparables
para facilitar su integración dentro del modelo SARIMAX.

\begin{longtable}[]{@{}
  >{\raggedright\arraybackslash}p{(\columnwidth - 4\tabcolsep) * \real{0.3333}}
  >{\raggedright\arraybackslash}p{(\columnwidth - 4\tabcolsep) * \real{0.3636}}
  >{\raggedright\arraybackslash}p{(\columnwidth - 4\tabcolsep) * \real{0.3030}}@{}}
\toprule\noalign{}
\begin{minipage}[b]{\linewidth}\raggedright
Variable
\end{minipage} & \begin{minipage}[b]{\linewidth}\raggedright
Descripción
\end{minipage} & \begin{minipage}[b]{\linewidth}\raggedright
Unidad
\end{minipage} \\
\midrule\noalign{}
\endhead
\bottomrule\noalign{}
\endlastfoot
Ano & Año calendario de observación & Año \\
Mes & Mes calendario de observación & Mes \\
Temperatura & Temperatura media mensual nacional & °C \\
Lluvia\_mm & Precipitación acumulada mensual nacional & mm \\
Precio\_diesel & Precio promedio mensual nacional del diésel automotriz
& MXN/L \\
Indice\_severidad & Índice nacional ponderado de severidad de sequía &
Índice de severidad ponderado \\
Precio\_maiz & Precio nacional mensual del maíz blanco al mayoreo &
MXN/Kg \\
INPC\_prom & Índice Nacional de Precios al Consumidor promedio mensual &
Índice \\
Produccion & Producción nacional mensual de maíz blanco & Kg \\
Precio\_prom & Precio internacional equivalente del maíz (CBOT ajustado)
& MXN/Kg \\
Tipo\_cambio & Tipo de cambio FIX promedio mensual & MXN/USD \\
Demanda\_m & Demanda nacional estimada de maíz blanco & Kg \\
Costo\_transp & Costo logístico nacional estimado por kilogramo
transportado & MXN/Kg \\
\end{longtable}

\hypertarget{variable-objetivo}{%
\subsection{Variable objetivo}\label{variable-objetivo}}

La variable objetivo principal del proyecto es:

\textbf{Precio\_maiz}

correspondiente al precio nacional mensual del maíz blanco al mayoreo en
México.

\hypertarget{variables-exuxf3genas-candidatas}{%
\subsection{Variables exógenas
candidatas}\label{variables-exuxf3genas-candidatas}}

Las variables restantes serán evaluadas como posibles variables
explicativas dentro de modelos SARIMAX, considerando su relevancia
económica, productiva, climática y logística sobre la formación del
precio del maíz.ndice \textbar{}

    \hypertarget{fuentes-de-informaciuxf3n-y-trazabilidad-de-datos}{%
\section{6. Fuentes de Información y Trazabilidad de
Datos}\label{fuentes-de-informaciuxf3n-y-trazabilidad-de-datos}}

La información utilizada en este proyecto proviene exclusivamente de
fuentes oficiales nacionales e internacionales. Cada variable fue
documentada para garantizar la trazabilidad, reproducibilidad y
transparencia del proceso analítico.

\begin{longtable}[]{@{}
  >{\raggedright\arraybackslash}p{(\columnwidth - 4\tabcolsep) * \real{0.3143}}
  >{\raggedright\arraybackslash}p{(\columnwidth - 4\tabcolsep) * \real{0.3429}}
  >{\raggedright\arraybackslash}p{(\columnwidth - 4\tabcolsep) * \real{0.3429}}@{}}
\toprule\noalign{}
\begin{minipage}[b]{\linewidth}\raggedright
Variable
\end{minipage} & \begin{minipage}[b]{\linewidth}\raggedright
Institución
\end{minipage} & \begin{minipage}[b]{\linewidth}\raggedright
Fuente principal
\end{minipage} \\
\midrule\noalign{}
\endhead
\bottomrule\noalign{}
\endlastfoot
Precio\_maiz & Secretaría de Economía (Sistema Nacional de Información e
Integración de Mercados - SNIIM) & Precios de maíz blanco al mayoreo \\
INPC\_prom & Instituto Nacional de Estadística y Geografía (INEGI) &
Índice Nacional de Precios al Consumidor \\
Tipo\_cambio & Banco de México (Banxico) & Tipo de cambio FIX \\
Precio\_prom & Chicago Mercantile Exchange Group (CME Group) y Seguridad
Alimentaria Mexicana (SEGALMEX) & Precio internacional equivalente del
maíz \\
Produccion & Servicio de Información Agroalimentaria y Pesquera (SIAP) &
Producción nacional de maíz blanco \\
Demanda\_m & Fideicomisos Instituidos en Relación con la Agricultura
(FIRA) y Centro de Estudios de las Finanzas Públicas (CEFP) & Consumo
nacional estimado \\
Costo\_transp & Agencia Reguladora del Transporte Ferroviario (ARTF) y
Dirección General de Autotransporte Federal (DGAF) & Costos logísticos
nacionales \\
Temperatura & Comisión Nacional del Agua (CONAGUA) mediante el Servicio
Meteorológico Nacional (SMN) & Temperaturas medias nacionales \\
Lluvia\_mm & Comisión Nacional del Agua (CONAGUA) mediante el Servicio
Meteorológico Nacional (SMN) & Precipitación acumulada nacional \\
Indice\_severidad & Comisión Nacional del Agua (CONAGUA) mediante el
Monitor de Sequía de México (MSM) & Índice de severidad de sequía \\
Precio\_diesel & Comisión Reguladora de Energía (CRE) y Procuraduría
Federal del Consumidor (PROFECO) & Precio nacional del diésel \\
\end{longtable}

Todas las fuentes fueron transformadas a frecuencia mensual y
homologadas a unidades comparables para su integración dentro de una
única base analítica destinada al modelado de series de tiempo.

    \hypertarget{justificaciuxf3n-de-variables-exuxf3genas}{%
\section{7. Justificación de Variables
Exógenas}\label{justificaciuxf3n-de-variables-exuxf3genas}}

El precio del maíz blanco no depende únicamente de su comportamiento
histórico. Existen factores económicos, productivos, climáticos y
logísticos que influyen directa o indirectamente en la formación de
precios.

Por esta razón se considera apropiado utilizar modelos SARIMAX, los
cuales permiten incorporar variables externas (variables exógenas) para
explicar una parte de la variabilidad observada en la serie objetivo.

Las variables seleccionadas fueron agrupadas en cuatro dimensiones
principales.

\hypertarget{variables-econuxf3micas}{%
\subsection{Variables económicas}\label{variables-econuxf3micas}}

\hypertarget{tipo-de-cambio-tipo_cambio}{%
\subsubsection{Tipo de cambio
(Tipo\_cambio)}\label{tipo-de-cambio-tipo_cambio}}

México participa activamente en mercados internacionales de granos. Las
fluctuaciones del tipo de cambio afectan los costos de importación,
exportación y la competitividad del maíz nacional frente a los mercados
internacionales.

\hypertarget{inflaciuxf3n-inpc_prom}{%
\subsubsection{Inflación (INPC\_prom)}\label{inflaciuxf3n-inpc_prom}}

La inflación refleja cambios generalizados en los precios de bienes y
servicios, incluyendo insumos agrícolas, transporte y costos operativos
asociados a la producción y comercialización del maíz.

\hypertarget{precio-internacional-del-mauxedz-precio_prom}{%
\subsubsection{Precio internacional del maíz
(Precio\_prom)}\label{precio-internacional-del-mauxedz-precio_prom}}

El precio internacional del maíz actúa como referencia para los mercados
nacionales debido a la integración comercial existente entre México y
otros países productores.

\hypertarget{variables-productivas}{%
\subsection{Variables productivas}\label{variables-productivas}}

\hypertarget{producciuxf3n-nacional-produccion}{%
\subsubsection{Producción nacional
(Produccion)}\label{producciuxf3n-nacional-produccion}}

La producción representa la oferta disponible en el mercado nacional.
Incrementos importantes en la oferta pueden generar presiones a la baja
sobre los precios.

\hypertarget{demanda-nacional-demanda_m}{%
\subsubsection{Demanda nacional
(Demanda\_m)}\label{demanda-nacional-demanda_m}}

La demanda representa el consumo potencial del grano para uso
alimentario, pecuario e industrial. Cambios en la demanda pueden generar
presiones alcistas sobre el precio.

\hypertarget{variables-climuxe1ticas}{%
\subsection{Variables climáticas}\label{variables-climuxe1ticas}}

\hypertarget{temperatura-temperatura}{%
\subsubsection{Temperatura
(Temperatura)}\label{temperatura-temperatura}}

Las condiciones térmicas afectan el desarrollo fisiológico de los
cultivos y pueden influir sobre los rendimientos agrícolas.

\hypertarget{precipitaciuxf3n-lluvia_mm}{%
\subsubsection{Precipitación
(Lluvia\_mm)}\label{precipitaciuxf3n-lluvia_mm}}

La disponibilidad de agua es uno de los factores más importantes para la
producción agrícola, especialmente en regiones dependientes de temporal.

\hypertarget{severidad-de-sequuxeda-indice_severidad}{%
\subsubsection{Severidad de sequía
(Indice\_severidad)}\label{severidad-de-sequuxeda-indice_severidad}}

Los episodios de sequía reducen la disponibilidad de agua y pueden
afectar significativamente la producción nacional de maíz.

\hypertarget{variables-loguxedsticas}{%
\subsection{Variables logísticas}\label{variables-loguxedsticas}}

\hypertarget{precio-del-diuxe9sel-precio_diesel}{%
\subsubsection{Precio del diésel
(Precio\_diesel)}\label{precio-del-diuxe9sel-precio_diesel}}

El diésel constituye uno de los principales costos asociados al
transporte de insumos y productos agrícolas.

\hypertarget{costos-de-transporte-costo_transp}{%
\subsubsection{Costos de transporte
(Costo\_transp)}\label{costos-de-transporte-costo_transp}}

Los costos logísticos impactan directamente el precio final observado en
los mercados mayoristas.

En conjunto, estas variables permiten representar distintos mecanismos
económicos y productivos que podrían influir en la evolución del precio
nacional del maíz blanco.

    \hypertarget{arquitectura-de-la-soluciuxf3n}{%
\section{8. Arquitectura de la
Solución}\label{arquitectura-de-la-soluciuxf3n}}

El presente proyecto integra múltiples fuentes de información
económicas, productivas, climáticas y logísticas con el objetivo de
modelar y pronosticar el comportamiento del precio nacional del maíz
blanco en México.

Para ello se diseñó una arquitectura analítica basada en PySpark que
permite almacenar, transformar, integrar y analizar grandes volúmenes de
información provenientes de distintas instituciones oficiales.

La solución sigue una estructura secuencial que parte de la recopilación
de datos históricos, continúa con procesos de limpieza e integración y
culmina con la construcción de modelos predictivos de series de tiempo.

\hypertarget{flujo-general-del-proyecto}{%
\subsection{Flujo general del
proyecto}\label{flujo-general-del-proyecto}}

Fuentes Oficiales ↓ Extracción de Datos (ETL) ↓ Homologación de
Variables ↓ Integración del Dataset Maestro ↓ Análisis Exploratorio
(EDA) ↓ Análisis de Correlaciones ↓ Pruebas de Estacionariedad ↓ Modelos
ARIMA ↓ Modelos SARIMA ↓ Modelos SARIMAX ↓ Evaluación y Validación ↓
Pronósticos ↓ Recomendaciones para Productores ↓ Análisis Complementario
sobre Precio de la Tortilla

\hypertarget{justificaciuxf3n-tecnoluxf3gica}{%
\subsection{Justificación
tecnológica}\label{justificaciuxf3n-tecnoluxf3gica}}

El proyecto fue desarrollado utilizando PySpark debido a su capacidad
para procesar grandes volúmenes de información de manera distribuida y
eficiente.

Aunque el conjunto de datos final utilizado para el modelado es
relativamente pequeño, el proceso de construcción implicó la integración
de múltiples fuentes heterogéneas con diferentes estructuras,
frecuencias y niveles de agregación.

La utilización de PySpark permite reproducir un flujo de trabajo
escalable y alineado con prácticas modernas de ingeniería de datos y
minería de datos.

\hypertarget{arquitectura-analuxedtica}{%
\subsection{Arquitectura analítica}\label{arquitectura-analuxedtica}}

La arquitectura implementada se compone de cuatro capas principales:

\hypertarget{capa-de-adquisiciuxf3n}{%
\subsubsection{Capa de adquisición}\label{capa-de-adquisiciuxf3n}}

Obtención de información histórica proveniente de organismos oficiales
nacionales e internacionales.

\hypertarget{capa-de-preparaciuxf3n}{%
\subsubsection{Capa de preparación}\label{capa-de-preparaciuxf3n}}

Procesos ETL para limpieza, transformación, homologación temporal e
integración de variables.

\hypertarget{capa-analuxedtica}{%
\subsubsection{Capa analítica}\label{capa-analuxedtica}}

Aplicación de técnicas de análisis exploratorio, estadística descriptiva
y modelado de series de tiempo.

\hypertarget{capa-de-soporte-a-la-decisiuxf3n}{%
\subsubsection{Capa de soporte a la
decisión}\label{capa-de-soporte-a-la-decisiuxf3n}}

Generación de pronósticos y recomendaciones orientadas a apoyar
estrategias de comercialización más justas para pequeños productores
agrícolas.

    \hypertarget{configuraciuxf3n-del-entorno-pyspark}{%
\section{9. Configuración del Entorno
PySpark}\label{configuraciuxf3n-del-entorno-pyspark}}

En esta sección se inicializa el entorno de trabajo, se importan las
librerías necesarias y se establece la configuración base para el
procesamiento de datos mediante PySpark.

    \begin{tcolorbox}[breakable, size=fbox, boxrule=1pt, pad at break*=1mm,colback=cellbackground, colframe=cellborder]
\prompt{In}{incolor}{1}{\boxspacing}
\begin{Verbatim}[commandchars=\\\{\}]
\PY{k+kn}{import} \PY{n+nn}{warnings}
\PY{n}{warnings}\PY{o}{.}\PY{n}{filterwarnings}\PY{p}{(}\PY{l+s+s2}{\PYZdq{}}\PY{l+s+s2}{ignore}\PY{l+s+s2}{\PYZdq{}}\PY{p}{)}

\PY{c+c1}{\PYZsh{} PySpark}
\PY{k+kn}{from} \PY{n+nn}{pyspark}\PY{n+nn}{.}\PY{n+nn}{sql} \PY{k+kn}{import} \PY{n}{SparkSession}
\PY{k+kn}{from} \PY{n+nn}{pyspark}\PY{n+nn}{.}\PY{n+nn}{sql}\PY{n+nn}{.}\PY{n+nn}{functions} \PY{k+kn}{import} \PY{o}{*}

\PY{c+c1}{\PYZsh{} Manipulación de datos}
\PY{k+kn}{import} \PY{n+nn}{pandas} \PY{k}{as} \PY{n+nn}{pd}
\PY{k+kn}{import} \PY{n+nn}{numpy} \PY{k}{as} \PY{n+nn}{np}

\PY{c+c1}{\PYZsh{} Visualización}
\PY{k+kn}{import} \PY{n+nn}{matplotlib}\PY{n+nn}{.}\PY{n+nn}{pyplot} \PY{k}{as} \PY{n+nn}{plt}
\PY{k+kn}{import} \PY{n+nn}{seaborn} \PY{k}{as} \PY{n+nn}{sns}

\PY{c+c1}{\PYZsh{} Estadística}
\PY{k+kn}{from} \PY{n+nn}{scipy} \PY{k+kn}{import} \PY{n}{stats}

\PY{c+c1}{\PYZsh{} Series de tiempo}
\PY{k+kn}{from} \PY{n+nn}{statsmodels}\PY{n+nn}{.}\PY{n+nn}{tsa}\PY{n+nn}{.}\PY{n+nn}{stattools} \PY{k+kn}{import} \PY{n}{adfuller}
\PY{k+kn}{from} \PY{n+nn}{statsmodels}\PY{n+nn}{.}\PY{n+nn}{tsa}\PY{n+nn}{.}\PY{n+nn}{seasonal} \PY{k+kn}{import} \PY{n}{seasonal\PYZus{}decompose}
\end{Verbatim}
\end{tcolorbox}

    \begin{tcolorbox}[breakable, size=fbox, boxrule=1pt, pad at break*=1mm,colback=cellbackground, colframe=cellborder]
\prompt{In}{incolor}{2}{\boxspacing}
\begin{Verbatim}[commandchars=\\\{\}]
\PY{k+kn}{from} \PY{n+nn}{pyspark}\PY{n+nn}{.}\PY{n+nn}{sql} \PY{k+kn}{import} \PY{n}{SparkSession}
\PY{k+kn}{from} \PY{n+nn}{pyspark}\PY{n+nn}{.}\PY{n+nn}{sql}\PY{n+nn}{.}\PY{n+nn}{functions} \PY{k+kn}{import} \PY{n}{col}\PY{p}{,} \PY{n}{to\PYZus{}date}\PY{p}{,} \PY{n}{year}\PY{p}{,} \PY{n}{month}\PY{p}{,} \PY{n}{isnan}\PY{p}{,} \PY{n}{count}\PY{p}{,} \PY{n}{when}
\PY{k+kn}{from} \PY{n+nn}{pyspark}\PY{n+nn}{.}\PY{n+nn}{sql}\PY{n+nn}{.}\PY{n+nn}{types} \PY{k+kn}{import} \PY{n}{DoubleType}\PY{p}{,} \PY{n}{IntegerType}\PY{p}{,} \PY{n}{StringType}\PY{p}{,} \PY{n}{DateType}

\PY{n}{spark} \PY{o}{=} \PY{p}{(}
    \PY{n}{SparkSession}\PY{o}{.}\PY{n}{builder}
    \PY{o}{.}\PY{n}{appName}\PY{p}{(}\PY{l+s+s2}{\PYZdq{}}\PY{l+s+s2}{Modelo\PYZus{}Agroprecuario\PYZus{}Prediccion\PYZus{}Precio}\PY{l+s+s2}{\PYZdq{}}\PY{p}{)}
    \PY{o}{.}\PY{n}{config}\PY{p}{(}\PY{l+s+s2}{\PYZdq{}}\PY{l+s+s2}{spark.sql.shuffle.partitions}\PY{l+s+s2}{\PYZdq{}}\PY{p}{,} \PY{l+s+s2}{\PYZdq{}}\PY{l+s+s2}{8}\PY{l+s+s2}{\PYZdq{}}\PY{p}{)}
    \PY{o}{.}\PY{n}{getOrCreate}\PY{p}{(}\PY{p}{)}
\PY{p}{)}

\PY{n}{spark}
\end{Verbatim}
\end{tcolorbox}

            \begin{tcolorbox}[breakable, size=fbox, boxrule=.5pt, pad at break*=1mm, opacityfill=0]
\prompt{Out}{outcolor}{2}{\boxspacing}
\begin{Verbatim}[commandchars=\\\{\}]
<pyspark.sql.session.SparkSession at 0x76c7837af690>
\end{Verbatim}
\end{tcolorbox}
        
    \hypertarget{carga-del-dataset-maestro}{%
\section{10. Carga del Dataset
Maestro}\label{carga-del-dataset-maestro}}

C configurado el entorno de trabajo, se procede a cargar la base de
datipal para los análisis exploratorios y los modeposterioresente.

El archivo contiene observaciones mensuales correspondientes al periodo
2000--2026 y fue construido a partir de la integración de múltiples
fuentes oficiales previamente documentadas.

    \begin{tcolorbox}[breakable, size=fbox, boxrule=1pt, pad at break*=1mm,colback=cellbackground, colframe=cellborder]
\prompt{In}{incolor}{3}{\boxspacing}
\begin{Verbatim}[commandchars=\\\{\}]
\PY{n}{ruta\PYZus{}dataset} \PY{o}{=} \PY{l+s+s2}{\PYZdq{}}\PY{l+s+s2}{/home/jovyan/Precio maiz/Data/Raw/panorama\PYZus{}maizblanco\PYZus{}2000\PYZus{}2026.csv}\PY{l+s+s2}{\PYZdq{}}

\PY{n}{df} \PY{o}{=} \PY{p}{(}
    \PY{n}{spark}\PY{o}{.}\PY{n}{read}
    \PY{o}{.}\PY{n}{option}\PY{p}{(}\PY{l+s+s2}{\PYZdq{}}\PY{l+s+s2}{header}\PY{l+s+s2}{\PYZdq{}}\PY{p}{,} \PY{l+s+s2}{\PYZdq{}}\PY{l+s+s2}{true}\PY{l+s+s2}{\PYZdq{}}\PY{p}{)}
    \PY{o}{.}\PY{n}{option}\PY{p}{(}\PY{l+s+s2}{\PYZdq{}}\PY{l+s+s2}{inferSchema}\PY{l+s+s2}{\PYZdq{}}\PY{p}{,} \PY{l+s+s2}{\PYZdq{}}\PY{l+s+s2}{true}\PY{l+s+s2}{\PYZdq{}}\PY{p}{)}
    \PY{o}{.}\PY{n}{csv}\PY{p}{(}\PY{n}{ruta\PYZus{}dataset}\PY{p}{)}
\PY{p}{)}
\end{Verbatim}
\end{tcolorbox}

    \hypertarget{inspecciuxf3n-inicial-del-dataset}{%
\section{11. Inspección Inicial del
Dataset}\label{inspecciuxf3n-inicial-del-dataset}}

Antes de iniciar el análisis exploratorio es necesario verificar la
estructura general de la base de datos.

En esta etapa se revisan:

\begin{itemize}
\tightlist
\item
  Dimensiones del conjunto de datos.
\item
  Tipos de datos.
\item
  Primeras observaciones.
\item
  Distribución temporal.
\item
  Calidad general de la información.
\end{itemize}

Lo anterior nos identificar posibles problemas que puedan afectar los
análisis posteriores.

    \hypertarget{dimensiones}{%
\subsection{11.1 Dimensiones}\label{dimensiones}}

    \begin{tcolorbox}[breakable, size=fbox, boxrule=1pt, pad at break*=1mm,colback=cellbackground, colframe=cellborder]
\prompt{In}{incolor}{4}{\boxspacing}
\begin{Verbatim}[commandchars=\\\{\}]
\PY{n}{filas} \PY{o}{=} \PY{n}{df}\PY{o}{.}\PY{n}{count}\PY{p}{(}\PY{p}{)}
\PY{n}{columnas} \PY{o}{=} \PY{n+nb}{len}\PY{p}{(}\PY{n}{df}\PY{o}{.}\PY{n}{columns}\PY{p}{)}

\PY{n+nb}{print}\PY{p}{(}\PY{l+s+sa}{f}\PY{l+s+s2}{\PYZdq{}}\PY{l+s+s2}{Número de registros: }\PY{l+s+si}{\PYZob{}}\PY{n}{filas}\PY{l+s+si}{\PYZcb{}}\PY{l+s+s2}{\PYZdq{}}\PY{p}{)}
\PY{n+nb}{print}\PY{p}{(}\PY{l+s+sa}{f}\PY{l+s+s2}{\PYZdq{}}\PY{l+s+s2}{Número de variables: }\PY{l+s+si}{\PYZob{}}\PY{n}{columnas}\PY{l+s+si}{\PYZcb{}}\PY{l+s+s2}{\PYZdq{}}\PY{p}{)}
\end{Verbatim}
\end{tcolorbox}

    \begin{Verbatim}[commandchars=\\\{\}]
Número de registros: 317
Número de variables: 13
    \end{Verbatim}

    \hypertarget{esquema-de-variables}{%
\subsection{11.2 Esquema de variables}\label{esquema-de-variables}}

    \begin{tcolorbox}[breakable, size=fbox, boxrule=1pt, pad at break*=1mm,colback=cellbackground, colframe=cellborder]
\prompt{In}{incolor}{5}{\boxspacing}
\begin{Verbatim}[commandchars=\\\{\}]
\PY{n}{df}\PY{o}{.}\PY{n}{printSchema}\PY{p}{(}\PY{p}{)}
\end{Verbatim}
\end{tcolorbox}

    \begin{Verbatim}[commandchars=\\\{\}]
root
 |-- Ano: integer (nullable = true)
 |-- Mes: integer (nullable = true)
 |-- Temperatura\_celsius\_prom\_nacional: double (nullable = true)
 |-- Lluvia\_mm\_prom\_nacional: double (nullable = true)
 |-- Precio\_diesel\_automotriz\_mxplitro\_prom\_nacional: double (nullable = true)
 |-- Indice\_severidad\_ponderada\_sequia\_prom\_nacional: double (nullable = true)
 |-- Precio\_maizblanco\_mxpkilo\_prom\_nacional: double (nullable = true)
 |-- INPC\_prom\_nacional: double (nullable = true)
 |-- Produccion\_maizblanco\_kg\_prom\_nacional: long (nullable = true)
 |-- Precio\_prom\_CBOT\_mxpkilo\_maiz\_internacional: double (nullable = true)
 |-- Tipo\_cambio\_mxpdolar\_prom\_nacional: double (nullable = true)
 |-- Demanda\_maizblanco\_kg\_prom\_nacional: long (nullable = true)
 |-- Costo\_transporte\_maiz\_mxpkilo\_prom\_nac: double (nullable = true)

    \end{Verbatim}

    \hypertarget{primeras-observaciones}{%
\subsection{11.3 Primeras observaciones}\label{primeras-observaciones}}

    \begin{tcolorbox}[breakable, size=fbox, boxrule=1pt, pad at break*=1mm,colback=cellbackground, colframe=cellborder]
\prompt{In}{incolor}{6}{\boxspacing}
\begin{Verbatim}[commandchars=\\\{\}]
\PY{n}{df}\PY{o}{.}\PY{n}{show}\PY{p}{(}\PY{l+m+mi}{10}\PY{p}{,} \PY{n}{truncate}\PY{o}{=}\PY{k+kc}{False}\PY{p}{)}
\end{Verbatim}
\end{tcolorbox}

    \begin{Verbatim}[commandchars=\\\{\}]
+----+---+---------------------------------+-----------------------+------------
-----------------------------------+--------------------------------------------
---+---------------------------------------+------------------+-----------------
---------------------+-------------------------------------------+--------------
--------------------+-----------------------------------+-----------------------
---------------+
|Ano |Mes|Temperatura\_celsius\_prom\_nacional|Lluvia\_mm\_prom\_nacional|Precio\_diese
l\_automotriz\_mxplitro\_prom\_nacional|Indice\_severidad\_ponderada\_sequia\_prom\_nacio
nal|Precio\_maizblanco\_mxpkilo\_prom\_nacional|INPC\_prom\_nacional|Produccion\_maizbl
anco\_kg\_prom\_nacional|Precio\_prom\_CBOT\_mxpkilo\_maiz\_internacional|Tipo\_cambio\_mx
pdolar\_prom\_nacional|Demanda\_maizblanco\_kg\_prom\_nacional|Costo\_transporte\_maiz\_m
xpkilo\_prom\_nac|
+----+---+---------------------------------+-----------------------+------------
-----------------------------------+--------------------------------------------
---+---------------------------------------+------------------+-----------------
---------------------+-------------------------------------------+--------------
--------------------+-----------------------------------+-----------------------
---------------+
|2000|1  |16.6                             |20.21                  |5.12
|43.12                                          |2.05
|43.201            |2850400000                            |0.92
|9.48                              |1625000000                         |0.22
|
|2000|2  |18.3                             |15.02                  |5.15
|44.28                                          |2.06
|43.587            |850200000                             |0.91
|9.44                              |1627000000                         |0.22
|
|2000|3  |21.8                             |18.95                  |5.19
|46.24                                          |2.09
|43.831            |410100000                             |0.89
|9.31                              |1630000000                         |0.23
|
|2000|4  |24.4                             |27.95                  |5.22
|56.12                                          |2.12
|44.081            |380600000                             |0.92
|9.38                              |1628000000                         |0.23
|
|2000|5  |26.7                             |56.54                  |5.26
|53.83                                          |2.14
|44.244            |1520100000                            |1.0
|9.52                              |1632000000                         |0.23
|
|2000|6  |26.4                             |176.81                 |5.29
|55.59                                          |2.11
|44.401            |3890300000                            |0.99
|9.81                              |1635000000                         |0.23
|
|2000|7  |25.0                             |264.17                 |5.33
|44.92                                          |2.15
|44.577            |2100500000                            |0.84
|9.47                              |1633000000                         |0.24
|
|2000|8  |25.0                             |238.68                 |5.36
|48.71                                          |2.13
|44.789            |620400000                             |0.74
|9.29                              |1636000000                         |0.24
|
|2000|9  |24.6                             |231.39                 |5.4
|37.54                                          |2.12
|45.122            |510300000                             |0.78
|9.36                              |1640000000                         |0.24
|
|2000|10 |22.7                             |122.11                 |5.43
|31.79                                          |2.08
|45.432            |980200000                             |0.84
|9.51                              |1638000000                         |0.24
|
+----+---+---------------------------------+-----------------------+------------
-----------------------------------+--------------------------------------------
---+---------------------------------------+------------------+-----------------
---------------------+-------------------------------------------+--------------
--------------------+-----------------------------------+-----------------------
---------------+
only showing top 10 rows

    \end{Verbatim}

    \hypertarget{estaduxedsticos-descriptivos-iniciales}{%
\subsection{11.4 Estadísticos descriptivos
iniciales}\label{estaduxedsticos-descriptivos-iniciales}}

    \begin{tcolorbox}[breakable, size=fbox, boxrule=1pt, pad at break*=1mm,colback=cellbackground, colframe=cellborder]
\prompt{In}{incolor}{7}{\boxspacing}
\begin{Verbatim}[commandchars=\\\{\}]
\PY{n}{df}\PY{o}{.}\PY{n}{describe}\PY{p}{(}\PY{p}{)}\PY{o}{.}\PY{n}{show}\PY{p}{(}\PY{p}{)}
\end{Verbatim}
\end{tcolorbox}

    \begin{Verbatim}[commandchars=\\\{\}]
+-------+------------------+------------------+---------------------------------
+-----------------------+-----------------------------------------------+-------
----------------------------------------+---------------------------------------
+------------------+--------------------------------------+---------------------
----------------------+----------------------------------+----------------------
-------------+--------------------------------------+
|summary|               Ano|               Mes|Temperatura\_celsius\_prom\_nacional
|Lluvia\_mm\_prom\_nacional|Precio\_diesel\_automotriz\_mxplitro\_prom\_nacional|Indice\_
severidad\_ponderada\_sequia\_prom\_nacional|Precio\_maizblanco\_mxpkilo\_prom\_nacional
|INPC\_prom\_nacional|Produccion\_maizblanco\_kg\_prom\_nacional|Precio\_prom\_CBOT\_mxpk
ilo\_maiz\_internacional|Tipo\_cambio\_mxpdolar\_prom\_nacional|Demanda\_maizblanco\_kg\_
prom\_nacional|Costo\_transporte\_maiz\_mxpkilo\_prom\_nac|
+-------+------------------+------------------+---------------------------------
+-----------------------+-----------------------------------------------+-------
----------------------------------------+---------------------------------------
+------------------+--------------------------------------+---------------------
----------------------+----------------------------------+----------------------
-------------+--------------------------------------+
|  count|               317|               317|
317|                    317|                                            317|
317|                                    317|               317|
317|                                        317|
317|                                317|                                   317|
|   mean|2012.7129337539432| 6.444794952681388|
22.56309148264984|      98.13753943217661|
13.167760252365934|                               88.7492113564669|
4.625993690851736| 84.87167823343849|                  2.0251514195583596E9|
2.6361829652996835|                14.754984227129338|
1.9758548895899053E9|                    0.7877917981072556|
| stddev|  7.64043146348037|3.4624038220302773|
3.480596381654182|      88.73611031051769|
7.400222351641308|                              78.50573016058736|
2.3722664561976554| 29.18961021554382|                  1.4433367533537838E9|
1.2935686598747538|                 4.052737051735518|
2.0396896733529457E8|                   0.46939785781345217|
|    min|              2000|                 1|
16.2|                  12.78|                                            4.7|
28.71|                                   2.05|            43.201|
380600000|                                       0.74|
9.04|                         1625000000|                                  0.22|
|    max|              2026|                12|
27.4|                  264.8|                                          28.36|
377.76|                                  10.65|           145.831|
4990200000|                                        6.7|
24.23|                         2324000000|
1.87|
+-------+------------------+------------------+---------------------------------
+-----------------------+-----------------------------------------------+-------
----------------------------------------+---------------------------------------
+------------------+--------------------------------------+---------------------
----------------------+----------------------------------+----------------------
-------------+--------------------------------------+

    \end{Verbatim}

    \hypertarget{control-de-calidad-de-los-datos}{%
\section{12. Control de Calidad de los
Datos}\label{control-de-calidad-de-los-datos}}

La calidad de los datos constituye un requisito fundamental para el
desarrollo de modelos predictivos confiables.

En esta etapa se verifica la existencia de valores faltantes, registros
duplicados y posibles inconsistencias estructurales dentro de la base de
datos.

    \hypertarget{valores-faltantes}{%
\subsection{12.1 Valores faltantes}\label{valores-faltantes}}

    \begin{tcolorbox}[breakable, size=fbox, boxrule=1pt, pad at break*=1mm,colback=cellbackground, colframe=cellborder]
\prompt{In}{incolor}{8}{\boxspacing}
\begin{Verbatim}[commandchars=\\\{\}]
\PY{k+kn}{from} \PY{n+nn}{pyspark}\PY{n+nn}{.}\PY{n+nn}{sql}\PY{n+nn}{.}\PY{n+nn}{functions} \PY{k+kn}{import} \PY{n}{col}\PY{p}{,} \PY{n}{count}\PY{p}{,} \PY{n}{when}

\PY{n}{df}\PY{o}{.}\PY{n}{select}\PY{p}{(}\PY{p}{[}
    \PY{n}{count}\PY{p}{(}
        \PY{n}{when}\PY{p}{(}\PY{n}{col}\PY{p}{(}\PY{n}{c}\PY{p}{)}\PY{o}{.}\PY{n}{isNull}\PY{p}{(}\PY{p}{)}\PY{p}{,} \PY{n}{c}\PY{p}{)}
    \PY{p}{)}\PY{o}{.}\PY{n}{alias}\PY{p}{(}\PY{n}{c}\PY{p}{)}
    \PY{k}{for} \PY{n}{c} \PY{o+ow}{in} \PY{n}{df}\PY{o}{.}\PY{n}{columns}
\PY{p}{]}\PY{p}{)}\PY{o}{.}\PY{n}{show}\PY{p}{(}\PY{n}{vertical}\PY{o}{=}\PY{k+kc}{True}\PY{p}{)}
\end{Verbatim}
\end{tcolorbox}

    \begin{Verbatim}[commandchars=\\\{\}]
-RECORD 0----------------------------------------------
 Ano                                             | 0
 Mes                                             | 0
 Temperatura\_celsius\_prom\_nacional               | 0
 Lluvia\_mm\_prom\_nacional                         | 0
 Precio\_diesel\_automotriz\_mxplitro\_prom\_nacional | 0
 Indice\_severidad\_ponderada\_sequia\_prom\_nacional | 0
 Precio\_maizblanco\_mxpkilo\_prom\_nacional         | 0
 INPC\_prom\_nacional                              | 0
 Produccion\_maizblanco\_kg\_prom\_nacional          | 0
 Precio\_prom\_CBOT\_mxpkilo\_maiz\_internacional     | 0
 Tipo\_cambio\_mxpdolar\_prom\_nacional              | 0
 Demanda\_maizblanco\_kg\_prom\_nacional             | 0
 Costo\_transporte\_maiz\_mxpkilo\_prom\_nac          | 0

    \end{Verbatim}

    \hypertarget{valores-duplicados}{%
\subsection{12.2 Valores duplicados}\label{valores-duplicados}}

    \begin{tcolorbox}[breakable, size=fbox, boxrule=1pt, pad at break*=1mm,colback=cellbackground, colframe=cellborder]
\prompt{In}{incolor}{9}{\boxspacing}
\begin{Verbatim}[commandchars=\\\{\}]
\PY{n}{duplicados} \PY{o}{=} \PY{n}{df}\PY{o}{.}\PY{n}{count}\PY{p}{(}\PY{p}{)} \PY{o}{\PYZhy{}} \PY{n}{df}\PY{o}{.}\PY{n}{dropDuplicates}\PY{p}{(}\PY{p}{)}\PY{o}{.}\PY{n}{count}\PY{p}{(}\PY{p}{)}

\PY{n+nb}{print}\PY{p}{(}\PY{l+s+sa}{f}\PY{l+s+s2}{\PYZdq{}}\PY{l+s+s2}{Registros duplicados: }\PY{l+s+si}{\PYZob{}}\PY{n}{duplicados}\PY{l+s+si}{\PYZcb{}}\PY{l+s+s2}{\PYZdq{}}\PY{p}{)}
\end{Verbatim}
\end{tcolorbox}

    \begin{Verbatim}[commandchars=\\\{\}]
Registros duplicados: 0
    \end{Verbatim}

    \hypertarget{hallazgos-del-control-de-calidad}{%
\subsubsection{Hallazgos del control de
calidad}\label{hallazgos-del-control-de-calidad}}

La revisión inicial del conjunto de datos permitió identificar los
siguientes resultados:

\begin{itemize}
\tightlist
\item
  La base de datos contiene 317 observaciones mensuales.
\item
  Se dispone de 13 variables relacionadas con factores económicos,
  productivos, climáticos y logísticos.
\item
  No se identificaron valores faltantes en ninguna de las variables
  analizadas.
\item
  No se encontraron registros duplicados.
\item
  Los tipos de datos fueron inferidos correctamente por PySpark.
\item
  La cobertura temporal abarca el periodo comprendido entre enero de
  2000 y mayo de 2026.
\end{itemize}

Los resultados anteriores indican que la base de datos presenta una
calidad adecuada para continuar con las etapas de análisis exploratorio
y modelado.

    \hypertarget{construcciuxf3n-de-la-variable-temporal}{%
\section{13. Construcción de la Variable
Temporal}\label{construcciuxf3n-de-la-variable-temporal}}

Las variables de año y mes se integran en una única variable temporal
denominada \textbf{fecha}, la cual permitirá ordenar cronológicamente
las observaciones y facilitar la aplicación de técnicas de análisis de
series de tiempo.

La construcción de esta variable es un paso fundamental para la
posterior aplicación de modelos ARIMA, SARIMA y SARIMAX.

    \begin{tcolorbox}[breakable, size=fbox, boxrule=1pt, pad at break*=1mm,colback=cellbackground, colframe=cellborder]
\prompt{In}{incolor}{10}{\boxspacing}
\begin{Verbatim}[commandchars=\\\{\}]
\PY{k+kn}{from} \PY{n+nn}{pyspark}\PY{n+nn}{.}\PY{n+nn}{sql}\PY{n+nn}{.}\PY{n+nn}{functions} \PY{k+kn}{import} \PY{n}{concat\PYZus{}ws}\PY{p}{,} \PY{n}{lpad}\PY{p}{,} \PY{n}{to\PYZus{}date}\PY{p}{,} \PY{n}{col}\PY{p}{,} \PY{n}{lit}

\PY{n}{df} \PY{o}{=} \PY{n}{df}\PY{o}{.}\PY{n}{withColumn}\PY{p}{(}
    \PY{l+s+s2}{\PYZdq{}}\PY{l+s+s2}{fecha}\PY{l+s+s2}{\PYZdq{}}\PY{p}{,}
    \PY{n}{to\PYZus{}date}\PY{p}{(}
        \PY{n}{concat\PYZus{}ws}\PY{p}{(}
            \PY{l+s+s2}{\PYZdq{}}\PY{l+s+s2}{\PYZhy{}}\PY{l+s+s2}{\PYZdq{}}\PY{p}{,}
            \PY{n}{col}\PY{p}{(}\PY{l+s+s2}{\PYZdq{}}\PY{l+s+s2}{Ano}\PY{l+s+s2}{\PYZdq{}}\PY{p}{)}\PY{p}{,}
            \PY{n}{lpad}\PY{p}{(}\PY{n}{col}\PY{p}{(}\PY{l+s+s2}{\PYZdq{}}\PY{l+s+s2}{Mes}\PY{l+s+s2}{\PYZdq{}}\PY{p}{)}\PY{o}{.}\PY{n}{cast}\PY{p}{(}\PY{l+s+s2}{\PYZdq{}}\PY{l+s+s2}{string}\PY{l+s+s2}{\PYZdq{}}\PY{p}{)}\PY{p}{,} \PY{l+m+mi}{2}\PY{p}{,} \PY{l+s+s2}{\PYZdq{}}\PY{l+s+s2}{0}\PY{l+s+s2}{\PYZdq{}}\PY{p}{)}\PY{p}{,}
            \PY{n}{lit}\PY{p}{(}\PY{l+s+s2}{\PYZdq{}}\PY{l+s+s2}{01}\PY{l+s+s2}{\PYZdq{}}\PY{p}{)}
        \PY{p}{)}
    \PY{p}{)}
\PY{p}{)}

\PY{n}{df}\PY{o}{.}\PY{n}{select}\PY{p}{(}\PY{l+s+s2}{\PYZdq{}}\PY{l+s+s2}{fecha}\PY{l+s+s2}{\PYZdq{}}\PY{p}{)}\PY{o}{.}\PY{n}{show}\PY{p}{(}\PY{l+m+mi}{10}\PY{p}{,} \PY{k+kc}{False}\PY{p}{)}
\end{Verbatim}
\end{tcolorbox}

    \begin{Verbatim}[commandchars=\\\{\}]
+----------+
|fecha     |
+----------+
|2000-01-01|
|2000-02-01|
|2000-03-01|
|2000-04-01|
|2000-05-01|
|2000-06-01|
|2000-07-01|
|2000-08-01|
|2000-09-01|
|2000-10-01|
+----------+
only showing top 10 rows

    \end{Verbatim}

    \begin{tcolorbox}[breakable, size=fbox, boxrule=1pt, pad at break*=1mm,colback=cellbackground, colframe=cellborder]
\prompt{In}{incolor}{11}{\boxspacing}
\begin{Verbatim}[commandchars=\\\{\}]
\PY{n}{df} \PY{o}{=} \PY{n}{df}\PY{o}{.}\PY{n}{orderBy}\PY{p}{(}\PY{l+s+s2}{\PYZdq{}}\PY{l+s+s2}{fecha}\PY{l+s+s2}{\PYZdq{}}\PY{p}{)}
\end{Verbatim}
\end{tcolorbox}

    \begin{tcolorbox}[breakable, size=fbox, boxrule=1pt, pad at break*=1mm,colback=cellbackground, colframe=cellborder]
\prompt{In}{incolor}{12}{\boxspacing}
\begin{Verbatim}[commandchars=\\\{\}]
\PY{n}{df}\PY{o}{.}\PY{n}{select}\PY{p}{(}
    \PY{l+s+s2}{\PYZdq{}}\PY{l+s+s2}{fecha}\PY{l+s+s2}{\PYZdq{}}\PY{p}{,}
    \PY{l+s+s2}{\PYZdq{}}\PY{l+s+s2}{Precio\PYZus{}maizblanco\PYZus{}mxpkilo\PYZus{}prom\PYZus{}nacional}\PY{l+s+s2}{\PYZdq{}}
\PY{p}{)}\PY{o}{.}\PY{n}{show}\PY{p}{(}\PY{l+m+mi}{10}\PY{p}{,} \PY{k+kc}{False}\PY{p}{)}
\end{Verbatim}
\end{tcolorbox}

    \begin{Verbatim}[commandchars=\\\{\}]
+----------+---------------------------------------+
|fecha     |Precio\_maizblanco\_mxpkilo\_prom\_nacional|
+----------+---------------------------------------+
|2000-01-01|2.05                                   |
|2000-02-01|2.06                                   |
|2000-03-01|2.09                                   |
|2000-04-01|2.12                                   |
|2000-05-01|2.14                                   |
|2000-06-01|2.11                                   |
|2000-07-01|2.15                                   |
|2000-08-01|2.13                                   |
|2000-09-01|2.12                                   |
|2000-10-01|2.08                                   |
+----------+---------------------------------------+
only showing top 10 rows

    \end{Verbatim}

    \hypertarget{anuxe1lisis-exploratorio-de-la-serie-objetivo}{%
\section{14. Análisis Exploratorio de la Serie
Objetivo}\label{anuxe1lisis-exploratorio-de-la-serie-objetivo}}

La variable objetivo del proyecto corresponde al precio nacional mensual
del maíz blanco al mayoreo.

Antes de construir modelos predictivos resulta necesario comprender el
comportamiento histórico de la serie, identificar tendencias de largo
plazo, posibles cambios estructurales y periodos de crecimiento o
contracción.

El análisis exploratorio permitirá establecer una primera aproximación
al comportamiento temporal del mercado nacional del maíz.

    \begin{tcolorbox}[breakable, size=fbox, boxrule=1pt, pad at break*=1mm,colback=cellbackground, colframe=cellborder]
\prompt{In}{incolor}{15}{\boxspacing}
\begin{Verbatim}[commandchars=\\\{\}]
\PY{n}{pdf} \PY{o}{=} \PY{p}{(}
    \PY{n}{df}
    \PY{o}{.}\PY{n}{select}\PY{p}{(}
        \PY{l+s+s2}{\PYZdq{}}\PY{l+s+s2}{fecha}\PY{l+s+s2}{\PYZdq{}}\PY{p}{,}
        \PY{l+s+s2}{\PYZdq{}}\PY{l+s+s2}{Precio\PYZus{}maizblanco\PYZus{}mxpkilo\PYZus{}prom\PYZus{}nacional}\PY{l+s+s2}{\PYZdq{}}
    \PY{p}{)}
    \PY{o}{.}\PY{n}{toPandas}\PY{p}{(}\PY{p}{)}
\PY{p}{)}

\PY{n}{pdf}\PY{o}{.}\PY{n}{head}\PY{p}{(}\PY{p}{)}
\end{Verbatim}
\end{tcolorbox}

            \begin{tcolorbox}[breakable, size=fbox, boxrule=.5pt, pad at break*=1mm, opacityfill=0]
\prompt{Out}{outcolor}{15}{\boxspacing}
\begin{Verbatim}[commandchars=\\\{\}]
        fecha  Precio\_maizblanco\_mxpkilo\_prom\_nacional
0  2000-01-01                                     2.05
1  2000-02-01                                     2.06
2  2000-03-01                                     2.09
3  2000-04-01                                     2.12
4  2000-05-01                                     2.14
\end{Verbatim}
\end{tcolorbox}
        
    \hypertarget{evoluciuxf3n-histuxf3rica-del-precio-del-mauxedz}{%
\subsection{14.2 Evolución histórica del precio del
maíz}\label{evoluciuxf3n-histuxf3rica-del-precio-del-mauxedz}}

    \begin{tcolorbox}[breakable, size=fbox, boxrule=1pt, pad at break*=1mm,colback=cellbackground, colframe=cellborder]
\prompt{In}{incolor}{16}{\boxspacing}
\begin{Verbatim}[commandchars=\\\{\}]
\PY{k+kn}{import} \PY{n+nn}{matplotlib}\PY{n+nn}{.}\PY{n+nn}{pyplot} \PY{k}{as} \PY{n+nn}{plt}

\PY{n}{plt}\PY{o}{.}\PY{n}{figure}\PY{p}{(}\PY{n}{figsize}\PY{o}{=}\PY{p}{(}\PY{l+m+mi}{14}\PY{p}{,}\PY{l+m+mi}{6}\PY{p}{)}\PY{p}{)}

\PY{n}{plt}\PY{o}{.}\PY{n}{plot}\PY{p}{(}
    \PY{n}{pdf}\PY{p}{[}\PY{l+s+s2}{\PYZdq{}}\PY{l+s+s2}{fecha}\PY{l+s+s2}{\PYZdq{}}\PY{p}{]}\PY{p}{,}
    \PY{n}{pdf}\PY{p}{[}\PY{l+s+s2}{\PYZdq{}}\PY{l+s+s2}{Precio\PYZus{}maizblanco\PYZus{}mxpkilo\PYZus{}prom\PYZus{}nacional}\PY{l+s+s2}{\PYZdq{}}\PY{p}{]}
\PY{p}{)}

\PY{n}{plt}\PY{o}{.}\PY{n}{title}\PY{p}{(}\PY{l+s+s2}{\PYZdq{}}\PY{l+s+s2}{Precio Nacional del Maíz Blanco (2000\PYZhy{}2026)}\PY{l+s+s2}{\PYZdq{}}\PY{p}{)}
\PY{n}{plt}\PY{o}{.}\PY{n}{xlabel}\PY{p}{(}\PY{l+s+s2}{\PYZdq{}}\PY{l+s+s2}{Fecha}\PY{l+s+s2}{\PYZdq{}}\PY{p}{)}
\PY{n}{plt}\PY{o}{.}\PY{n}{ylabel}\PY{p}{(}\PY{l+s+s2}{\PYZdq{}}\PY{l+s+s2}{MXN por Kg}\PY{l+s+s2}{\PYZdq{}}\PY{p}{)}

\PY{n}{plt}\PY{o}{.}\PY{n}{grid}\PY{p}{(}\PY{k+kc}{True}\PY{p}{)}

\PY{n}{plt}\PY{o}{.}\PY{n}{show}\PY{p}{(}\PY{p}{)}
\end{Verbatim}
\end{tcolorbox}

    \begin{center}
    \adjustimage{max size={0.9\linewidth}{0.9\paperheight}}{output_36_0.png}
    \end{center}
    { \hspace*{\fill} \\}
    
    \hypertarget{crecimiento-acumulado}{%
\subsection{14.3 Crecimiento acumulado}\label{crecimiento-acumulado}}

    \begin{tcolorbox}[breakable, size=fbox, boxrule=1pt, pad at break*=1mm,colback=cellbackground, colframe=cellborder]
\prompt{In}{incolor}{17}{\boxspacing}
\begin{Verbatim}[commandchars=\\\{\}]
\PY{n}{precio\PYZus{}inicial} \PY{o}{=} \PY{n}{pdf}\PY{p}{[}\PY{l+s+s2}{\PYZdq{}}\PY{l+s+s2}{Precio\PYZus{}maizblanco\PYZus{}mxpkilo\PYZus{}prom\PYZus{}nacional}\PY{l+s+s2}{\PYZdq{}}\PY{p}{]}\PY{o}{.}\PY{n}{iloc}\PY{p}{[}\PY{l+m+mi}{0}\PY{p}{]}
\PY{n}{precio\PYZus{}final} \PY{o}{=} \PY{n}{pdf}\PY{p}{[}\PY{l+s+s2}{\PYZdq{}}\PY{l+s+s2}{Precio\PYZus{}maizblanco\PYZus{}mxpkilo\PYZus{}prom\PYZus{}nacional}\PY{l+s+s2}{\PYZdq{}}\PY{p}{]}\PY{o}{.}\PY{n}{iloc}\PY{p}{[}\PY{o}{\PYZhy{}}\PY{l+m+mi}{1}\PY{p}{]}

\PY{n}{incremento} \PY{o}{=} \PY{p}{(}
    \PY{p}{(}\PY{n}{precio\PYZus{}final} \PY{o}{\PYZhy{}} \PY{n}{precio\PYZus{}inicial}\PY{p}{)}
    \PY{o}{/} \PY{n}{precio\PYZus{}inicial}
\PY{p}{)} \PY{o}{*} \PY{l+m+mi}{100}

\PY{n+nb}{print}\PY{p}{(}\PY{l+s+sa}{f}\PY{l+s+s2}{\PYZdq{}}\PY{l+s+s2}{Precio inicial: }\PY{l+s+si}{\PYZob{}}\PY{n}{precio\PYZus{}inicial}\PY{l+s+si}{:}\PY{l+s+s2}{.2f}\PY{l+s+si}{\PYZcb{}}\PY{l+s+s2}{\PYZdq{}}\PY{p}{)}
\PY{n+nb}{print}\PY{p}{(}\PY{l+s+sa}{f}\PY{l+s+s2}{\PYZdq{}}\PY{l+s+s2}{Precio final: }\PY{l+s+si}{\PYZob{}}\PY{n}{precio\PYZus{}final}\PY{l+s+si}{:}\PY{l+s+s2}{.2f}\PY{l+s+si}{\PYZcb{}}\PY{l+s+s2}{\PYZdq{}}\PY{p}{)}
\PY{n+nb}{print}\PY{p}{(}\PY{l+s+sa}{f}\PY{l+s+s2}{\PYZdq{}}\PY{l+s+s2}{Incremento acumulado: }\PY{l+s+si}{\PYZob{}}\PY{n}{incremento}\PY{l+s+si}{:}\PY{l+s+s2}{.2f}\PY{l+s+si}{\PYZcb{}}\PY{l+s+s2}{\PYZpc{}}\PY{l+s+s2}{\PYZdq{}}\PY{p}{)}
\end{Verbatim}
\end{tcolorbox}

    \begin{Verbatim}[commandchars=\\\{\}]
Precio inicial: 2.05
Precio final: 9.82
Incremento acumulado: 379.02\%
    \end{Verbatim}

    \hypertarget{distribuciuxf3n-del-precio-del-mauxedz-blanco}{%
\subsection{14.4 Distribución del precio del maíz
blanco}\label{distribuciuxf3n-del-precio-del-mauxedz-blanco}}

    \begin{tcolorbox}[breakable, size=fbox, boxrule=1pt, pad at break*=1mm,colback=cellbackground, colframe=cellborder]
\prompt{In}{incolor}{18}{\boxspacing}
\begin{Verbatim}[commandchars=\\\{\}]
\PY{n}{plt}\PY{o}{.}\PY{n}{figure}\PY{p}{(}\PY{n}{figsize}\PY{o}{=}\PY{p}{(}\PY{l+m+mi}{8}\PY{p}{,}\PY{l+m+mi}{5}\PY{p}{)}\PY{p}{)}

\PY{n}{plt}\PY{o}{.}\PY{n}{hist}\PY{p}{(}
    \PY{n}{pdf}\PY{p}{[}\PY{l+s+s2}{\PYZdq{}}\PY{l+s+s2}{Precio\PYZus{}maizblanco\PYZus{}mxpkilo\PYZus{}prom\PYZus{}nacional}\PY{l+s+s2}{\PYZdq{}}\PY{p}{]}\PY{p}{,}
    \PY{n}{bins}\PY{o}{=}\PY{l+m+mi}{20}
\PY{p}{)}

\PY{n}{plt}\PY{o}{.}\PY{n}{title}\PY{p}{(}\PY{l+s+s2}{\PYZdq{}}\PY{l+s+s2}{Distribución del Precio Nacional del Maíz}\PY{l+s+s2}{\PYZdq{}}\PY{p}{)}
\PY{n}{plt}\PY{o}{.}\PY{n}{xlabel}\PY{p}{(}\PY{l+s+s2}{\PYZdq{}}\PY{l+s+s2}{MXN/Kg}\PY{l+s+s2}{\PYZdq{}}\PY{p}{)}
\PY{n}{plt}\PY{o}{.}\PY{n}{ylabel}\PY{p}{(}\PY{l+s+s2}{\PYZdq{}}\PY{l+s+s2}{Frecuencia}\PY{l+s+s2}{\PYZdq{}}\PY{p}{)}

\PY{n}{plt}\PY{o}{.}\PY{n}{show}\PY{p}{(}\PY{p}{)}
\end{Verbatim}
\end{tcolorbox}

    \begin{center}
    \adjustimage{max size={0.9\linewidth}{0.9\paperheight}}{output_40_0.png}
    \end{center}
    { \hspace*{\fill} \\}
    
    \hypertarget{estaduxedsticos-descriptivos-de-la-variable-objetivo}{%
\subsection{14.5 Estadísticos descriptivos de la variable
objetivo}\label{estaduxedsticos-descriptivos-de-la-variable-objetivo}}

    \begin{tcolorbox}[breakable, size=fbox, boxrule=1pt, pad at break*=1mm,colback=cellbackground, colframe=cellborder]
\prompt{In}{incolor}{19}{\boxspacing}
\begin{Verbatim}[commandchars=\\\{\}]
\PY{n}{pdf}\PY{p}{[}\PY{l+s+s2}{\PYZdq{}}\PY{l+s+s2}{Precio\PYZus{}maizblanco\PYZus{}mxpkilo\PYZus{}prom\PYZus{}nacional}\PY{l+s+s2}{\PYZdq{}}\PY{p}{]}\PY{o}{.}\PY{n}{describe}\PY{p}{(}\PY{p}{)}
\end{Verbatim}
\end{tcolorbox}

            \begin{tcolorbox}[breakable, size=fbox, boxrule=.5pt, pad at break*=1mm, opacityfill=0]
\prompt{Out}{outcolor}{19}{\boxspacing}
\begin{Verbatim}[commandchars=\\\{\}]
count    317.000000
mean       4.625994
std        2.372266
min        2.050000
25\%        2.640000
50\%        4.160000
75\%        5.010000
max       10.650000
Name: Precio\_maizblanco\_mxpkilo\_prom\_nacional, dtype: float64
\end{Verbatim}
\end{tcolorbox}
        
    \hypertarget{interpretaciuxf3n-de-la-evoluciuxf3n-histuxf3rica-del-precio}{%
\subsubsection{Interpretación de la evolución histórica del
precio}\label{interpretaciuxf3n-de-la-evoluciuxf3n-histuxf3rica-del-precio}}

La evolución histórica del precio nacional del maíz blanco muestra una
tendencia general creciente a lo largo del periodo 2000--2026.

Pueden identificarse cuatro etapas principales:

\hypertarget{estabilidad-relativa}{%
\paragraph{2000--2006: Estabilidad
relativa}\label{estabilidad-relativa}}

Durante los primeros años del periodo analizado, el precio permaneció
relativamente estable alrededor de los 2.0 a 2.5 MXN por kilogramo, con
fluctuaciones moderadas y baja volatilidad.

\hypertarget{primera-etapa-de-crecimiento}{%
\paragraph{2007--2012: Primera etapa de
crecimiento}\label{primera-etapa-de-crecimiento}}

A partir de 2007 se observa un incremento importante en los precios,
coincidiendo con el periodo de volatilidad internacional de materias
primas agrícolas. y la llamada ``Crisis alimentaria mundial''. Durante
esta etapa el precio supera los 4 MXN por kilogramo.

\hypertarget{periodo-de-estabilizaciuxf3n}{%
\paragraph{2013--2020: Periodo de
estabilización}\label{periodo-de-estabilizaciuxf3n}}

Después del crecimiento inicial, el mercado entra en una fase
relativamente estable, con oscilaciones moderadas alrededor de los 4 a 5
MXN por kilogramo.

\hypertarget{choque-estructural-reciente}{%
\paragraph{2021--2024: Choque estructural
reciente}\label{choque-estructural-reciente}}

Entre 2021 y 2024 se presenta el mayor incremento de toda la serie
histórica, alcanzando niveles superiores a los 10 MXN por kilogramo.

Este comportamiento coincide temporalmente con factores como:

\begin{itemize}
\tightlist
\item
  Incrementos en los costos energéticos.
\item
  Disrupciones logísticas globales.
\item
  Presiones inflacionarias internacionales.
\item
  Incrementos en precios internacionales de granos.
\item
  Eventos climáticos adversos.
\end{itemize}

\hypertarget{nueva-etapa-de-equilibrio}{%
\paragraph{2024--2026: Nueva etapa de
equilibrio}\label{nueva-etapa-de-equilibrio}}

Después del máximo histórico observado en 2023--2024, el precio presenta
una moderada corrección y posteriormente se estabiliza alrededor de los
9 a 10 MXN por kilogramo.

En términos generales, la serie presenta evidencia visual de:

\begin{itemize}
\tightlist
\item
  Tendencia creciente.
\item
  Cambios estructurales.
\item
  Posible estacionalidad.
\item
  Incremento de volatilidad en los años recientes.
\end{itemize}

\hypertarget{sobre-la-distribuciuxf3n-del-precio}{%
\subsubsection{Sobre la distribución del
precio}\label{sobre-la-distribuciuxf3n-del-precio}}

La distribución observada no presenta una forma aproximadamente norm de
distribución gaussianaal.

Se identifican múltiples concentraciones de frecuencia asociadas a
distintos niveles históricos de precio, lo que sugiere la existencia de
cambios estructurales en el mercado a lo largo del periodo de estudio.

Este comportamiento es consistente con la presencia de tendencias de
largo plazo y posibles cambios de régimen económico, por lo que refuerza
la necesidad de utilizar modelos de series de tiempo capaces de capturar
dinámicas temporales complejas.

    \hypertarget{evoluciuxf3n-de-las-variables-exuxf3genas}{%
\section{Evolución de las variables
exógenas}\label{evoluciuxf3n-de-las-variables-exuxf3genas}}

    \begin{tcolorbox}[breakable, size=fbox, boxrule=1pt, pad at break*=1mm,colback=cellbackground, colframe=cellborder]
\prompt{In}{incolor}{24}{\boxspacing}
\begin{Verbatim}[commandchars=\\\{\}]
\PY{n}{pdf\PYZus{}completo} \PY{o}{=} \PY{p}{(}
    \PY{n}{df}
    \PY{o}{.}\PY{n}{orderBy}\PY{p}{(}\PY{l+s+s2}{\PYZdq{}}\PY{l+s+s2}{fecha}\PY{l+s+s2}{\PYZdq{}}\PY{p}{)}
    \PY{o}{.}\PY{n}{toPandas}\PY{p}{(}\PY{p}{)}
\PY{p}{)}

\PY{n}{pdf\PYZus{}completo}\PY{o}{.}\PY{n}{head}\PY{p}{(}\PY{p}{)}
\end{Verbatim}
\end{tcolorbox}

            \begin{tcolorbox}[breakable, size=fbox, boxrule=.5pt, pad at break*=1mm, opacityfill=0]
\prompt{Out}{outcolor}{24}{\boxspacing}
\begin{Verbatim}[commandchars=\\\{\}]
    Ano  Mes  Temperatura\_celsius\_prom\_nacional  Lluvia\_mm\_prom\_nacional  \textbackslash{}
0  2000    1                               16.6                    20.21
1  2000    2                               18.3                    15.02
2  2000    3                               21.8                    18.95
3  2000    4                               24.4                    27.95
4  2000    5                               26.7                    56.54

   Precio\_diesel\_automotriz\_mxplitro\_prom\_nacional  \textbackslash{}
0                                             5.12
1                                             5.15
2                                             5.19
3                                             5.22
4                                             5.26

   Indice\_severidad\_ponderada\_sequia\_prom\_nacional  \textbackslash{}
0                                            43.12
1                                            44.28
2                                            46.24
3                                            56.12
4                                            53.83

   Precio\_maizblanco\_mxpkilo\_prom\_nacional  INPC\_prom\_nacional  \textbackslash{}
0                                     2.05              43.201
1                                     2.06              43.587
2                                     2.09              43.831
3                                     2.12              44.081
4                                     2.14              44.244

   Produccion\_maizblanco\_kg\_prom\_nacional  \textbackslash{}
0                              2850400000
1                               850200000
2                               410100000
3                               380600000
4                              1520100000

   Precio\_prom\_CBOT\_mxpkilo\_maiz\_internacional  \textbackslash{}
0                                         0.92
1                                         0.91
2                                         0.89
3                                         0.92
4                                         1.00

   Tipo\_cambio\_mxpdolar\_prom\_nacional  Demanda\_maizblanco\_kg\_prom\_nacional  \textbackslash{}
0                                9.48                           1625000000
1                                9.44                           1627000000
2                                9.31                           1630000000
3                                9.38                           1628000000
4                                9.52                           1632000000

   Costo\_transporte\_maiz\_mxpkilo\_prom\_nac       fecha
0                                    0.22  2000-01-01
1                                    0.22  2000-02-01
2                                    0.23  2000-03-01
3                                    0.23  2000-04-01
4                                    0.23  2000-05-01
\end{Verbatim}
\end{tcolorbox}
        
    \begin{tcolorbox}[breakable, size=fbox, boxrule=1pt, pad at break*=1mm,colback=cellbackground, colframe=cellborder]
\prompt{In}{incolor}{26}{\boxspacing}
\begin{Verbatim}[commandchars=\\\{\}]
\PY{n}{variables} \PY{o}{=} \PY{p}{[}
    \PY{l+s+s2}{\PYZdq{}}\PY{l+s+s2}{Precio\PYZus{}maizblanco\PYZus{}mxpkilo\PYZus{}prom\PYZus{}nacional}\PY{l+s+s2}{\PYZdq{}}\PY{p}{,}
    \PY{l+s+s2}{\PYZdq{}}\PY{l+s+s2}{INPC\PYZus{}prom\PYZus{}nacional}\PY{l+s+s2}{\PYZdq{}}\PY{p}{,}
    \PY{l+s+s2}{\PYZdq{}}\PY{l+s+s2}{Tipo\PYZus{}cambio\PYZus{}mxpdolar\PYZus{}prom\PYZus{}nacional}\PY{l+s+s2}{\PYZdq{}}\PY{p}{,}
    \PY{l+s+s2}{\PYZdq{}}\PY{l+s+s2}{Precio\PYZus{}prom\PYZus{}CBOT\PYZus{}mxpkilo\PYZus{}maiz\PYZus{}internacional}\PY{l+s+s2}{\PYZdq{}}\PY{p}{,}
    \PY{l+s+s2}{\PYZdq{}}\PY{l+s+s2}{Precio\PYZus{}diesel\PYZus{}automotriz\PYZus{}mxplitro\PYZus{}prom\PYZus{}nacional}\PY{l+s+s2}{\PYZdq{}}\PY{p}{,}
    \PY{l+s+s2}{\PYZdq{}}\PY{l+s+s2}{Produccion\PYZus{}maizblanco\PYZus{}kg\PYZus{}prom\PYZus{}nacional}\PY{l+s+s2}{\PYZdq{}}\PY{p}{,}
    \PY{l+s+s2}{\PYZdq{}}\PY{l+s+s2}{Demanda\PYZus{}maizblanco\PYZus{}kg\PYZus{}prom\PYZus{}nacional}\PY{l+s+s2}{\PYZdq{}}\PY{p}{,}
    \PY{l+s+s2}{\PYZdq{}}\PY{l+s+s2}{Temperatura\PYZus{}celsius\PYZus{}prom\PYZus{}nacional}\PY{l+s+s2}{\PYZdq{}}\PY{p}{,}
    \PY{l+s+s2}{\PYZdq{}}\PY{l+s+s2}{Lluvia\PYZus{}mm\PYZus{}prom\PYZus{}nacional}\PY{l+s+s2}{\PYZdq{}}\PY{p}{,}
    \PY{l+s+s2}{\PYZdq{}}\PY{l+s+s2}{Indice\PYZus{}severidad\PYZus{}ponderada\PYZus{}sequia\PYZus{}prom\PYZus{}nacional}\PY{l+s+s2}{\PYZdq{}}\PY{p}{,}
    \PY{l+s+s2}{\PYZdq{}}\PY{l+s+s2}{Costo\PYZus{}transporte\PYZus{}maiz\PYZus{}mxpkilo\PYZus{}prom\PYZus{}nac}\PY{l+s+s2}{\PYZdq{}}
\PY{p}{]}

\PY{k}{for} \PY{n}{variable} \PY{o+ow}{in} \PY{n}{variables}\PY{p}{:}

    \PY{n}{plt}\PY{o}{.}\PY{n}{figure}\PY{p}{(}\PY{n}{figsize}\PY{o}{=}\PY{p}{(}\PY{l+m+mi}{12}\PY{p}{,}\PY{l+m+mi}{4}\PY{p}{)}\PY{p}{)}

    \PY{n}{plt}\PY{o}{.}\PY{n}{plot}\PY{p}{(}
        \PY{n}{pdf\PYZus{}completo}\PY{p}{[}\PY{l+s+s2}{\PYZdq{}}\PY{l+s+s2}{fecha}\PY{l+s+s2}{\PYZdq{}}\PY{p}{]}\PY{p}{,}
        \PY{n}{pdf\PYZus{}completo}\PY{p}{[}\PY{n}{variable}\PY{p}{]}
    \PY{p}{)}

    \PY{n}{plt}\PY{o}{.}\PY{n}{title}\PY{p}{(}\PY{n}{variable}\PY{p}{)}

    \PY{n}{plt}\PY{o}{.}\PY{n}{grid}\PY{p}{(}\PY{k+kc}{True}\PY{p}{)}

    \PY{n}{plt}\PY{o}{.}\PY{n}{show}\PY{p}{(}\PY{p}{)}
\end{Verbatim}
\end{tcolorbox}

    \begin{center}
    \adjustimage{max size={0.9\linewidth}{0.9\paperheight}}{output_46_0.png}
    \end{center}
    { \hspace*{\fill} \\}
    
    \begin{center}
    \adjustimage{max size={0.9\linewidth}{0.9\paperheight}}{output_46_1.png}
    \end{center}
    { \hspace*{\fill} \\}
    
    \begin{center}
    \adjustimage{max size={0.9\linewidth}{0.9\paperheight}}{output_46_2.png}
    \end{center}
    { \hspace*{\fill} \\}
    
    \begin{center}
    \adjustimage{max size={0.9\linewidth}{0.9\paperheight}}{output_46_3.png}
    \end{center}
    { \hspace*{\fill} \\}
    
    \begin{center}
    \adjustimage{max size={0.9\linewidth}{0.9\paperheight}}{output_46_4.png}
    \end{center}
    { \hspace*{\fill} \\}
    
    \begin{center}
    \adjustimage{max size={0.9\linewidth}{0.9\paperheight}}{output_46_5.png}
    \end{center}
    { \hspace*{\fill} \\}
    
    \begin{center}
    \adjustimage{max size={0.9\linewidth}{0.9\paperheight}}{output_46_6.png}
    \end{center}
    { \hspace*{\fill} \\}
    
    \begin{center}
    \adjustimage{max size={0.9\linewidth}{0.9\paperheight}}{output_46_7.png}
    \end{center}
    { \hspace*{\fill} \\}
    
    \begin{center}
    \adjustimage{max size={0.9\linewidth}{0.9\paperheight}}{output_46_8.png}
    \end{center}
    { \hspace*{\fill} \\}
    
    \begin{center}
    \adjustimage{max size={0.9\linewidth}{0.9\paperheight}}{output_46_9.png}
    \end{center}
    { \hspace*{\fill} \\}
    
    \begin{center}
    \adjustimage{max size={0.9\linewidth}{0.9\paperheight}}{output_46_10.png}
    \end{center}
    { \hspace*{\fill} \\}
    
    \hypertarget{sobre-la-evoluciuxf3n-de-las-variables-exuxf3genas}{%
\subsection{Sobre la evolución de las variables
exógenas}\label{sobre-la-evoluciuxf3n-de-las-variables-exuxf3genas}}

El análisis visual de las variables exógenas permite identificar
diversos patrones temporales que podrían influir en el comportamiento
del precio nacional del maíz blanco.

\hypertarget{variables-econuxf3micas}{%
\subsubsection{Variables económicas}\label{variables-econuxf3micas}}

\hypertarget{inpc}{%
\paragraph{INPC}\label{inpc}}

El Índice Nacional de Precios al Consumidor presenta una tendencia
creciente sostenida durante todo el periodo de estudio, reflejando la
inflación acumulada de la economía mexicana.

\hypertarget{tipo-de-cambio}{%
\paragraph{Tipo de cambio}\label{tipo-de-cambio}}

El tipo de cambio muestra episodios de volatilidad importantes,
particularmente durante las crisis financieras internacionales y los
periodos de incertidumbre económica observados entre 2015 y 2024.

\hypertarget{precio-internacional-del-mauxedz-cbot}{%
\paragraph{Precio internacional del maíz
(CBOT)}\label{precio-internacional-del-mauxedz-cbot}}

El precio internacional del maíz exhibe una dinámica similar a la
observada en el precio nacional, con incrementos relevantes durante los
periodos 2007-2013 y 2021-2023, sugiriendo una posible relación entre
ambos mercados.

\hypertarget{variables-loguxedsticas}{%
\subsubsection{Variables logísticas}\label{variables-loguxedsticas}}

\hypertarget{precio-del-diuxe9sel}{%
\paragraph{Precio del diésel}\label{precio-del-diuxe9sel}}

El precio del diésel presenta una tendencia creciente de largo plazo,
especialmente después de 2017, lo que podría influir sobre los costos de
transporte y comercialización.

\hypertarget{costos-de-transporte}{%
\paragraph{Costos de transporte}\label{costos-de-transporte}}

Los costos logísticos muestran una tendencia ascendente consistente
durante todo el periodo analizado.

\hypertarget{variables-productivas}{%
\subsubsection{Variables productivas}\label{variables-productivas}}

\hypertarget{producciuxf3n-nacional}{%
\paragraph{Producción nacional}\label{producciuxf3n-nacional}}

La producción nacional presenta una fuerte estacionalidad anual,
asociada a los ciclos agrícolas de siembra y cosecha.

\hypertarget{demanda-nacional}{%
\paragraph{Demanda nacional}\label{demanda-nacional}}

La demanda estimada muestra una tendencia creciente relativamente
estable, sin evidencia visual de estacionalidad significativa.

\hypertarget{variables-climuxe1ticas}{%
\subsubsection{Variables climáticas}\label{variables-climuxe1ticas}}

\hypertarget{temperatura}{%
\paragraph{Temperatura}\label{temperatura}}

La temperatura media nacional presenta una marcada estacionalidad anual,
consistente con los ciclos climáticos observados en México.

\hypertarget{precipitaciuxf3n}{%
\paragraph{Precipitación}\label{precipitaciuxf3n}}

La precipitación acumulada presenta patrones estacionales muy definidos,
con incrementos recurrentes durante la temporada de lluvias.

\hypertarget{uxedndice-de-severidad-de-sequuxeda}{%
\paragraph{Índice de severidad de
sequía}\label{uxedndice-de-severidad-de-sequuxeda}}

La severidad de sequía presenta episodios extremos concentrados en
determinados periodos, particularmente entre 2011-2013 y 2023-2025, lo
que sugiere posibles impactos sobre la producción agrícola y los precios
del maíz.

En conjunto, las variables analizadas muestran evidencia visual de
tendencias, estacionalidad y eventos extraordinarios, características
que justifican su evaluación dentro de modelos SARIMAX.

    \hypertarget{anuxe1lisis-de-correlaciones}{%
\section{16. Análisis de
Correlaciones}\label{anuxe1lisis-de-correlaciones}}

Una vez analizado el comportamiento individual de las variables, se
procede a evaluar las relaciones lineales existentes entre ellas.

La correlación permite identificar variables que presentan movimientos
similares o inversos respecto al precio nacional del maíz blanco.

Es importante señalar que la correlación no implica causalidad; sin
embargo, constituye una herramienta útil para identificar posibles
variables candidatas a incorporarse dentro del modelo SARIMAX.

    \hypertarget{matriz-de-correlaciones}{%
\subsection{16.1 Matriz de
correlaciones}\label{matriz-de-correlaciones}}

    \begin{tcolorbox}[breakable, size=fbox, boxrule=1pt, pad at break*=1mm,colback=cellbackground, colframe=cellborder]
\prompt{In}{incolor}{28}{\boxspacing}
\begin{Verbatim}[commandchars=\\\{\}]
\PY{n}{corr} \PY{o}{=} \PY{n}{pdf\PYZus{}completo}\PY{o}{.}\PY{n}{drop}\PY{p}{(}\PY{n}{columns}\PY{o}{=}\PY{p}{[}\PY{l+s+s2}{\PYZdq{}}\PY{l+s+s2}{fecha}\PY{l+s+s2}{\PYZdq{}}\PY{p}{]}\PY{p}{)}\PY{o}{.}\PY{n}{corr}\PY{p}{(}\PY{p}{)}

\PY{n}{corr}\PY{o}{.}\PY{n}{round}\PY{p}{(}\PY{l+m+mi}{3}\PY{p}{)}
\end{Verbatim}
\end{tcolorbox}

            \begin{tcolorbox}[breakable, size=fbox, boxrule=.5pt, pad at break*=1mm, opacityfill=0]
\prompt{Out}{outcolor}{28}{\boxspacing}
\begin{Verbatim}[commandchars=\\\{\}]
                                                   Ano    Mes  \textbackslash{}
Ano                                              1.000 -0.028
Mes                                             -0.028  1.000
Temperatura\_celsius\_prom\_nacional                0.063  0.101
Lluvia\_mm\_prom\_nacional                         -0.016  0.357
Precio\_diesel\_automotriz\_mxplitro\_prom\_nacional  0.968 -0.004
Indice\_severidad\_ponderada\_sequia\_prom\_nacional  0.343 -0.128
Precio\_maizblanco\_mxpkilo\_prom\_nacional          0.901 -0.020
INPC\_prom\_nacional                               0.981  0.001
Produccion\_maizblanco\_kg\_prom\_nacional           0.084  0.257
Precio\_prom\_CBOT\_mxpkilo\_maiz\_internacional      0.824 -0.029
Tipo\_cambio\_mxpdolar\_prom\_nacional               0.916  0.013
Demanda\_maizblanco\_kg\_prom\_nacional              0.999  0.006
Costo\_transporte\_maiz\_mxpkilo\_prom\_nac           0.965 -0.002

Temperatura\_celsius\_prom\_nacional  \textbackslash{}
Ano
0.063
Mes
0.101
Temperatura\_celsius\_prom\_nacional
1.000
Lluvia\_mm\_prom\_nacional
0.640
Precio\_diesel\_automotriz\_mxplitro\_prom\_nacional
0.068
Indice\_severidad\_ponderada\_sequia\_prom\_nacional
0.141
Precio\_maizblanco\_mxpkilo\_prom\_nacional
0.118
INPC\_prom\_nacional
0.063
Produccion\_maizblanco\_kg\_prom\_nacional
-0.292
Precio\_prom\_CBOT\_mxpkilo\_maiz\_internacional
0.050
Tipo\_cambio\_mxpdolar\_prom\_nacional
0.046
Demanda\_maizblanco\_kg\_prom\_nacional
0.067
Costo\_transporte\_maiz\_mxpkilo\_prom\_nac
0.067

                                                 Lluvia\_mm\_prom\_nacional  \textbackslash{}
Ano                                                               -0.016
Mes                                                                0.357
Temperatura\_celsius\_prom\_nacional                                  0.640
Lluvia\_mm\_prom\_nacional                                            1.000
Precio\_diesel\_automotriz\_mxplitro\_prom\_nacional                   -0.009
Indice\_severidad\_ponderada\_sequia\_prom\_nacional                   -0.029
Precio\_maizblanco\_mxpkilo\_prom\_nacional                            0.029
INPC\_prom\_nacional                                                -0.012
Produccion\_maizblanco\_kg\_prom\_nacional                            -0.078
Precio\_prom\_CBOT\_mxpkilo\_maiz\_internacional                       -0.026
Tipo\_cambio\_mxpdolar\_prom\_nacional                                -0.012
Demanda\_maizblanco\_kg\_prom\_nacional                               -0.004
Costo\_transporte\_maiz\_mxpkilo\_prom\_nac                            -0.011

Precio\_diesel\_automotriz\_mxplitro\_prom\_nacional  \textbackslash{}
Ano
0.968
Mes
-0.004
Temperatura\_celsius\_prom\_nacional
0.068
Lluvia\_mm\_prom\_nacional
-0.009
Precio\_diesel\_automotriz\_mxplitro\_prom\_nacional
1.000
Indice\_severidad\_ponderada\_sequia\_prom\_nacional
0.340
Precio\_maizblanco\_mxpkilo\_prom\_nacional
0.917
INPC\_prom\_nacional
0.985
Produccion\_maizblanco\_kg\_prom\_nacional
0.077
Precio\_prom\_CBOT\_mxpkilo\_maiz\_internacional
0.770
Tipo\_cambio\_mxpdolar\_prom\_nacional
0.909
Demanda\_maizblanco\_kg\_prom\_nacional
0.965
Costo\_transporte\_maiz\_mxpkilo\_prom\_nac
0.988

Indice\_severidad\_ponderada\_sequia\_prom\_nacional  \textbackslash{}
Ano
0.343
Mes
-0.128
Temperatura\_celsius\_prom\_nacional
0.141
Lluvia\_mm\_prom\_nacional
-0.029
Precio\_diesel\_automotriz\_mxplitro\_prom\_nacional
0.340
Indice\_severidad\_ponderada\_sequia\_prom\_nacional
1.000
Precio\_maizblanco\_mxpkilo\_prom\_nacional
0.524
INPC\_prom\_nacional
0.398
Produccion\_maizblanco\_kg\_prom\_nacional
0.025
Precio\_prom\_CBOT\_mxpkilo\_maiz\_internacional
0.302
Tipo\_cambio\_mxpdolar\_prom\_nacional
0.120
Demanda\_maizblanco\_kg\_prom\_nacional
0.337
Costo\_transporte\_maiz\_mxpkilo\_prom\_nac
0.396

Precio\_maizblanco\_mxpkilo\_prom\_nacional  \textbackslash{}
Ano
0.901
Mes
-0.020
Temperatura\_celsius\_prom\_nacional
0.118
Lluvia\_mm\_prom\_nacional
0.029
Precio\_diesel\_automotriz\_mxplitro\_prom\_nacional
0.917
Indice\_severidad\_ponderada\_sequia\_prom\_nacional
0.524
Precio\_maizblanco\_mxpkilo\_prom\_nacional
1.000
INPC\_prom\_nacional
0.955
Produccion\_maizblanco\_kg\_prom\_nacional
0.028
Precio\_prom\_CBOT\_mxpkilo\_maiz\_internacional
0.761
Tipo\_cambio\_mxpdolar\_prom\_nacional
0.740
Demanda\_maizblanco\_kg\_prom\_nacional
0.896
Costo\_transporte\_maiz\_mxpkilo\_prom\_nac
0.956

                                                 INPC\_prom\_nacional  \textbackslash{}
Ano                                                           0.981
Mes                                                           0.001
Temperatura\_celsius\_prom\_nacional                             0.063
Lluvia\_mm\_prom\_nacional                                      -0.012
Precio\_diesel\_automotriz\_mxplitro\_prom\_nacional               0.985
Indice\_severidad\_ponderada\_sequia\_prom\_nacional               0.398
Precio\_maizblanco\_mxpkilo\_prom\_nacional                       0.955
INPC\_prom\_nacional                                            1.000
Produccion\_maizblanco\_kg\_prom\_nacional                        0.079
Precio\_prom\_CBOT\_mxpkilo\_maiz\_internacional                   0.788
Tipo\_cambio\_mxpdolar\_prom\_nacional                            0.882
Demanda\_maizblanco\_kg\_prom\_nacional                           0.979
Costo\_transporte\_maiz\_mxpkilo\_prom\_nac                        0.996

Produccion\_maizblanco\_kg\_prom\_nacional  \textbackslash{}
Ano
0.084
Mes
0.257
Temperatura\_celsius\_prom\_nacional
-0.292
Lluvia\_mm\_prom\_nacional
-0.078
Precio\_diesel\_automotriz\_mxplitro\_prom\_nacional
0.077
Indice\_severidad\_ponderada\_sequia\_prom\_nacional
0.025
Precio\_maizblanco\_mxpkilo\_prom\_nacional
0.028
INPC\_prom\_nacional
0.079
Produccion\_maizblanco\_kg\_prom\_nacional
1.000
Precio\_prom\_CBOT\_mxpkilo\_maiz\_internacional
0.103
Tipo\_cambio\_mxpdolar\_prom\_nacional
0.107
Demanda\_maizblanco\_kg\_prom\_nacional
0.095
Costo\_transporte\_maiz\_mxpkilo\_prom\_nac
0.072

Precio\_prom\_CBOT\_mxpkilo\_maiz\_internacional  \textbackslash{}
Ano
0.824
Mes
-0.029
Temperatura\_celsius\_prom\_nacional
0.050
Lluvia\_mm\_prom\_nacional
-0.026
Precio\_diesel\_automotriz\_mxplitro\_prom\_nacional
0.770
Indice\_severidad\_ponderada\_sequia\_prom\_nacional
0.302
Precio\_maizblanco\_mxpkilo\_prom\_nacional
0.761
INPC\_prom\_nacional
0.788
Produccion\_maizblanco\_kg\_prom\_nacional
0.103
Precio\_prom\_CBOT\_mxpkilo\_maiz\_internacional
1.000
Tipo\_cambio\_mxpdolar\_prom\_nacional
0.752
Demanda\_maizblanco\_kg\_prom\_nacional
0.826
Costo\_transporte\_maiz\_mxpkilo\_prom\_nac
0.754

Tipo\_cambio\_mxpdolar\_prom\_nacional  \textbackslash{}
Ano
0.916
Mes
0.013
Temperatura\_celsius\_prom\_nacional
0.046
Lluvia\_mm\_prom\_nacional
-0.012
Precio\_diesel\_automotriz\_mxplitro\_prom\_nacional
0.909
Indice\_severidad\_ponderada\_sequia\_prom\_nacional
0.120
Precio\_maizblanco\_mxpkilo\_prom\_nacional
0.740
INPC\_prom\_nacional
0.882
Produccion\_maizblanco\_kg\_prom\_nacional
0.107
Precio\_prom\_CBOT\_mxpkilo\_maiz\_internacional
0.752
Tipo\_cambio\_mxpdolar\_prom\_nacional
1.000
Demanda\_maizblanco\_kg\_prom\_nacional
0.915
Costo\_transporte\_maiz\_mxpkilo\_prom\_nac
0.873

Demanda\_maizblanco\_kg\_prom\_nacional  \textbackslash{}
Ano
0.999
Mes
0.006
Temperatura\_celsius\_prom\_nacional
0.067
Lluvia\_mm\_prom\_nacional
-0.004
Precio\_diesel\_automotriz\_mxplitro\_prom\_nacional
0.965
Indice\_severidad\_ponderada\_sequia\_prom\_nacional
0.337
Precio\_maizblanco\_mxpkilo\_prom\_nacional
0.896
INPC\_prom\_nacional
0.979
Produccion\_maizblanco\_kg\_prom\_nacional
0.095
Precio\_prom\_CBOT\_mxpkilo\_maiz\_internacional
0.826
Tipo\_cambio\_mxpdolar\_prom\_nacional
0.915
Demanda\_maizblanco\_kg\_prom\_nacional
1.000
Costo\_transporte\_maiz\_mxpkilo\_prom\_nac
0.961

Costo\_transporte\_maiz\_mxpkilo\_prom\_nac
Ano
0.965
Mes
-0.002
Temperatura\_celsius\_prom\_nacional
0.067
Lluvia\_mm\_prom\_nacional
-0.011
Precio\_diesel\_automotriz\_mxplitro\_prom\_nacional
0.988
Indice\_severidad\_ponderada\_sequia\_prom\_nacional
0.396
Precio\_maizblanco\_mxpkilo\_prom\_nacional
0.956
INPC\_prom\_nacional
0.996
Produccion\_maizblanco\_kg\_prom\_nacional
0.072
Precio\_prom\_CBOT\_mxpkilo\_maiz\_internacional
0.754
Tipo\_cambio\_mxpdolar\_prom\_nacional
0.873
Demanda\_maizblanco\_kg\_prom\_nacional
0.961
Costo\_transporte\_maiz\_mxpkilo\_prom\_nac
1.000
\end{Verbatim}
\end{tcolorbox}
        
    \hypertarget{mapa-de-calor-de-correlaciones-de-variables}{%
\subsection{16.2 Mapa de calor de correlaciones de
variables}\label{mapa-de-calor-de-correlaciones-de-variables}}

    \begin{tcolorbox}[breakable, size=fbox, boxrule=1pt, pad at break*=1mm,colback=cellbackground, colframe=cellborder]
\prompt{In}{incolor}{29}{\boxspacing}
\begin{Verbatim}[commandchars=\\\{\}]
\PY{k+kn}{import} \PY{n+nn}{matplotlib}\PY{n+nn}{.}\PY{n+nn}{pyplot} \PY{k}{as} \PY{n+nn}{plt}

\PY{n}{plt}\PY{o}{.}\PY{n}{figure}\PY{p}{(}\PY{n}{figsize}\PY{o}{=}\PY{p}{(}\PY{l+m+mi}{12}\PY{p}{,}\PY{l+m+mi}{10}\PY{p}{)}\PY{p}{)}

\PY{n}{plt}\PY{o}{.}\PY{n}{imshow}\PY{p}{(}
    \PY{n}{corr}\PY{p}{,}
    \PY{n}{aspect}\PY{o}{=}\PY{l+s+s2}{\PYZdq{}}\PY{l+s+s2}{auto}\PY{l+s+s2}{\PYZdq{}}
\PY{p}{)}

\PY{n}{plt}\PY{o}{.}\PY{n}{colorbar}\PY{p}{(}\PY{n}{label}\PY{o}{=}\PY{l+s+s2}{\PYZdq{}}\PY{l+s+s2}{Correlación}\PY{l+s+s2}{\PYZdq{}}\PY{p}{)}

\PY{n}{plt}\PY{o}{.}\PY{n}{xticks}\PY{p}{(}
    \PY{n+nb}{range}\PY{p}{(}\PY{n+nb}{len}\PY{p}{(}\PY{n}{corr}\PY{o}{.}\PY{n}{columns}\PY{p}{)}\PY{p}{)}\PY{p}{,}
    \PY{n}{corr}\PY{o}{.}\PY{n}{columns}\PY{p}{,}
    \PY{n}{rotation}\PY{o}{=}\PY{l+m+mi}{90}
\PY{p}{)}

\PY{n}{plt}\PY{o}{.}\PY{n}{yticks}\PY{p}{(}
    \PY{n+nb}{range}\PY{p}{(}\PY{n+nb}{len}\PY{p}{(}\PY{n}{corr}\PY{o}{.}\PY{n}{columns}\PY{p}{)}\PY{p}{)}\PY{p}{,}
    \PY{n}{corr}\PY{o}{.}\PY{n}{columns}
\PY{p}{)}

\PY{n}{plt}\PY{o}{.}\PY{n}{title}\PY{p}{(}\PY{l+s+s2}{\PYZdq{}}\PY{l+s+s2}{Matriz de Correlaciones}\PY{l+s+s2}{\PYZdq{}}\PY{p}{)}

\PY{n}{plt}\PY{o}{.}\PY{n}{tight\PYZus{}layout}\PY{p}{(}\PY{p}{)}

\PY{n}{plt}\PY{o}{.}\PY{n}{show}\PY{p}{(}\PY{p}{)}
\end{Verbatim}
\end{tcolorbox}

    \begin{center}
    \adjustimage{max size={0.9\linewidth}{0.9\paperheight}}{output_52_0.png}
    \end{center}
    { \hspace*{\fill} \\}
    
    \hypertarget{correlaciuxf3n-con-la-variable-objetivo}{%
\subsection{16.3 Correlación con la variable
objetivo}\label{correlaciuxf3n-con-la-variable-objetivo}}

    \begin{tcolorbox}[breakable, size=fbox, boxrule=1pt, pad at break*=1mm,colback=cellbackground, colframe=cellborder]
\prompt{In}{incolor}{30}{\boxspacing}
\begin{Verbatim}[commandchars=\\\{\}]
\PY{n}{corr\PYZus{}maiz} \PY{o}{=} \PY{p}{(}
    \PY{n}{corr}\PY{p}{[}\PY{l+s+s2}{\PYZdq{}}\PY{l+s+s2}{Precio\PYZus{}maizblanco\PYZus{}mxpkilo\PYZus{}prom\PYZus{}nacional}\PY{l+s+s2}{\PYZdq{}}\PY{p}{]}
    \PY{o}{.}\PY{n}{sort\PYZus{}values}\PY{p}{(}\PY{n}{ascending}\PY{o}{=}\PY{k+kc}{False}\PY{p}{)}
\PY{p}{)}

\PY{n}{corr\PYZus{}maiz}
\end{Verbatim}
\end{tcolorbox}

            \begin{tcolorbox}[breakable, size=fbox, boxrule=.5pt, pad at break*=1mm, opacityfill=0]
\prompt{Out}{outcolor}{30}{\boxspacing}
\begin{Verbatim}[commandchars=\\\{\}]
Precio\_maizblanco\_mxpkilo\_prom\_nacional            1.000000
Costo\_transporte\_maiz\_mxpkilo\_prom\_nac             0.955602
INPC\_prom\_nacional                                 0.955140
Precio\_diesel\_automotriz\_mxplitro\_prom\_nacional    0.917139
Ano                                                0.900642
Demanda\_maizblanco\_kg\_prom\_nacional                0.896478
Precio\_prom\_CBOT\_mxpkilo\_maiz\_internacional        0.761170
Tipo\_cambio\_mxpdolar\_prom\_nacional                 0.739921
Indice\_severidad\_ponderada\_sequia\_prom\_nacional    0.524208
Temperatura\_celsius\_prom\_nacional                  0.118466
Lluvia\_mm\_prom\_nacional                            0.028926
Produccion\_maizblanco\_kg\_prom\_nacional             0.028324
Mes                                               -0.019659
Name: Precio\_maizblanco\_mxpkilo\_prom\_nacional, dtype: float64
\end{Verbatim}
\end{tcolorbox}
        
    \hypertarget{gruxe1fico-de-importancia-de-variables}{%
\subsection{Gráfico de importancia de
variables}\label{gruxe1fico-de-importancia-de-variables}}

    \begin{tcolorbox}[breakable, size=fbox, boxrule=1pt, pad at break*=1mm,colback=cellbackground, colframe=cellborder]
\prompt{In}{incolor}{31}{\boxspacing}
\begin{Verbatim}[commandchars=\\\{\}]
\PY{n}{corr\PYZus{}maiz}\PY{o}{.}\PY{n}{drop}\PY{p}{(}
    \PY{l+s+s2}{\PYZdq{}}\PY{l+s+s2}{Precio\PYZus{}maizblanco\PYZus{}mxpkilo\PYZus{}prom\PYZus{}nacional}\PY{l+s+s2}{\PYZdq{}}
\PY{p}{)}\PY{o}{.}\PY{n}{plot}\PY{p}{(}
    \PY{n}{kind}\PY{o}{=}\PY{l+s+s2}{\PYZdq{}}\PY{l+s+s2}{barh}\PY{l+s+s2}{\PYZdq{}}\PY{p}{,}
    \PY{n}{figsize}\PY{o}{=}\PY{p}{(}\PY{l+m+mi}{10}\PY{p}{,}\PY{l+m+mi}{6}\PY{p}{)}
\PY{p}{)}

\PY{n}{plt}\PY{o}{.}\PY{n}{title}\PY{p}{(}
    \PY{l+s+s2}{\PYZdq{}}\PY{l+s+s2}{Correlación con el Precio Nacional del Maíz}\PY{l+s+s2}{\PYZdq{}}
\PY{p}{)}

\PY{n}{plt}\PY{o}{.}\PY{n}{xlabel}\PY{p}{(}\PY{l+s+s2}{\PYZdq{}}\PY{l+s+s2}{Coeficiente de correlación}\PY{l+s+s2}{\PYZdq{}}\PY{p}{)}

\PY{n}{plt}\PY{o}{.}\PY{n}{grid}\PY{p}{(}\PY{k+kc}{True}\PY{p}{)}

\PY{n}{plt}\PY{o}{.}\PY{n}{show}\PY{p}{(}\PY{p}{)}
\end{Verbatim}
\end{tcolorbox}

    \begin{center}
    \adjustimage{max size={0.9\linewidth}{0.9\paperheight}}{output_56_0.png}
    \end{center}
    { \hspace*{\fill} \\}
    
    \hypertarget{interpretaciuxf3n-de-correlaciones}{%
\subsection{16.5 Interpretación de
correlaciones}\label{interpretaciuxf3n-de-correlaciones}}

La matriz de correlaciones muestra asociaciones positivas elevadas entre
el precio nacional del maíz y diversas variables económicas y
logísticas.

Las mayores correlaciones observadas corresponden a:

\begin{itemize}
\tightlist
\item
  Costos de transporte (0.956)
\item
  INPC (0.955)
\item
  Precio del diésel (0.917)
\item
  Demanda nacional (0.896)
\item
  Precio internacional del maíz (0.761)
\item
  Tipo de cambio (0.740)
\end{itemize}

Sin embargo, es importante señalar que muchas de estas variables
presentan tendencias crecientes de largo plazo, por lo que parte de las
correlaciones observadas podría estar asociada al paso del tiempo y no
necesariamente a una relación causal directa.

Por esta razón, la selección final de variables para el modelo SARIMAX
no se realizará únicamente con base en correlaciones simples, sino
considerando además criterios económicos, climáticos y estadísticos.

    \hypertarget{selecciuxf3n-preliminar-de-variables-exuxf3genas}{%
\section{17. Selección Preliminar de Variables
Exógenas}\label{selecciuxf3n-preliminar-de-variables-exuxf3genas}}

La matriz de correlaciones permitió identificar asociaciones lineales
entre las variables disponibles y el precio nacional del maíz blanco.

Sin embargo, las correlaciones simples no son suficientes para
seleccionar variables dentro de un modelo SARIMAX, ya que algunas
variables pueden contener información redundante.

Para evaluar este fenómeno se realizará un análisis de multicolinealidad
mediante el Factor de Inflación de la Varianza (VIF).

El objetivo es identificar variables que aporten información única al
modelo y reducir la presencia de relaciones lineales excesivas entre
predictores.

    \hypertarget{evaluaciuxf3n-de-multicolinealidad-mediante-vif}{%
\subsection{17.1 Evaluación de Multicolinealidad mediante
VIF}\label{evaluaciuxf3n-de-multicolinealidad-mediante-vif}}

Aunque la matriz de correlaciones permitió identificar asociaciones
lineales entre las variables explicativas y el precio nacional del maíz,
las correlaciones simples no permiten determinar si varias variables
están aportando información redundante al modelo.

Por esta razón se calculó el Factor de Inflación de la Varianza
(Variance Inflation Factor, VIF), una medida utilizada para cuantificar
el grado de multicolinealidad entre variables explicativas.

Matemáticamente, el VIF se define como:

\[
VIF_i = \frac{1}{1 - R_i^2}
\]

donde \((R_i^2)\) representa el coeficiente de determinación obtenido al
explicar la variable \((i)\) mediante las demás variables explicativas.

Valores elevados de VIF indican que una variable puede ser explicada en
gran medida por otras variables del conjunto, generando redundancia
informativa y potencial inestabilidad en los coeficientes del modelo.

\hypertarget{criterios-de-interpretaciuxf3n}{%
\subsubsection{Criterios de
interpretación}\label{criterios-de-interpretaciuxf3n}}

\begin{longtable}[]{@{}ll@{}}
\toprule\noalign{}
VIF & Interpretación \\
\midrule\noalign{}
\endhead
\bottomrule\noalign{}
\endlastfoot
1 & Sin multicolinealidad \\
1 -- 5 & Aceptable \\
5 -- 10 & Moderada \\
\textgreater{} 10 & Problemática \\
\end{longtable}

    \begin{tcolorbox}[breakable, size=fbox, boxrule=1pt, pad at break*=1mm,colback=cellbackground, colframe=cellborder]
\prompt{In}{incolor}{33}{\boxspacing}
\begin{Verbatim}[commandchars=\\\{\}]
\PY{n}{variables\PYZus{}vif} \PY{o}{=} \PY{p}{[}
    \PY{l+s+s2}{\PYZdq{}}\PY{l+s+s2}{Temperatura\PYZus{}celsius\PYZus{}prom\PYZus{}nacional}\PY{l+s+s2}{\PYZdq{}}\PY{p}{,}
    \PY{l+s+s2}{\PYZdq{}}\PY{l+s+s2}{Lluvia\PYZus{}mm\PYZus{}prom\PYZus{}nacional}\PY{l+s+s2}{\PYZdq{}}\PY{p}{,}
    \PY{l+s+s2}{\PYZdq{}}\PY{l+s+s2}{Precio\PYZus{}diesel\PYZus{}automotriz\PYZus{}mxplitro\PYZus{}prom\PYZus{}nacional}\PY{l+s+s2}{\PYZdq{}}\PY{p}{,}
    \PY{l+s+s2}{\PYZdq{}}\PY{l+s+s2}{Indice\PYZus{}severidad\PYZus{}ponderada\PYZus{}sequia\PYZus{}prom\PYZus{}nacional}\PY{l+s+s2}{\PYZdq{}}\PY{p}{,}
    \PY{l+s+s2}{\PYZdq{}}\PY{l+s+s2}{INPC\PYZus{}prom\PYZus{}nacional}\PY{l+s+s2}{\PYZdq{}}\PY{p}{,}
    \PY{l+s+s2}{\PYZdq{}}\PY{l+s+s2}{Produccion\PYZus{}maizblanco\PYZus{}kg\PYZus{}prom\PYZus{}nacional}\PY{l+s+s2}{\PYZdq{}}\PY{p}{,}
    \PY{l+s+s2}{\PYZdq{}}\PY{l+s+s2}{Precio\PYZus{}prom\PYZus{}CBOT\PYZus{}mxpkilo\PYZus{}maiz\PYZus{}internacional}\PY{l+s+s2}{\PYZdq{}}\PY{p}{,}
    \PY{l+s+s2}{\PYZdq{}}\PY{l+s+s2}{Tipo\PYZus{}cambio\PYZus{}mxpdolar\PYZus{}prom\PYZus{}nacional}\PY{l+s+s2}{\PYZdq{}}\PY{p}{,}
    \PY{l+s+s2}{\PYZdq{}}\PY{l+s+s2}{Demanda\PYZus{}maizblanco\PYZus{}kg\PYZus{}prom\PYZus{}nacional}\PY{l+s+s2}{\PYZdq{}}\PY{p}{,}
    \PY{l+s+s2}{\PYZdq{}}\PY{l+s+s2}{Costo\PYZus{}transporte\PYZus{}maiz\PYZus{}mxpkilo\PYZus{}prom\PYZus{}nac}\PY{l+s+s2}{\PYZdq{}}
\PY{p}{]}

\PY{n}{X} \PY{o}{=} \PY{n}{pdf\PYZus{}completo}\PY{p}{[}\PY{n}{variables\PYZus{}vif}\PY{p}{]}

\PY{n}{X}\PY{o}{.}\PY{n}{head}\PY{p}{(}\PY{p}{)}
\end{Verbatim}
\end{tcolorbox}

            \begin{tcolorbox}[breakable, size=fbox, boxrule=.5pt, pad at break*=1mm, opacityfill=0]
\prompt{Out}{outcolor}{33}{\boxspacing}
\begin{Verbatim}[commandchars=\\\{\}]
   Temperatura\_celsius\_prom\_nacional  Lluvia\_mm\_prom\_nacional  \textbackslash{}
0                               16.6                    20.21
1                               18.3                    15.02
2                               21.8                    18.95
3                               24.4                    27.95
4                               26.7                    56.54

   Precio\_diesel\_automotriz\_mxplitro\_prom\_nacional  \textbackslash{}
0                                             5.12
1                                             5.15
2                                             5.19
3                                             5.22
4                                             5.26

   Indice\_severidad\_ponderada\_sequia\_prom\_nacional  INPC\_prom\_nacional  \textbackslash{}
0                                            43.12              43.201
1                                            44.28              43.587
2                                            46.24              43.831
3                                            56.12              44.081
4                                            53.83              44.244

   Produccion\_maizblanco\_kg\_prom\_nacional  \textbackslash{}
0                              2850400000
1                               850200000
2                               410100000
3                               380600000
4                              1520100000

   Precio\_prom\_CBOT\_mxpkilo\_maiz\_internacional  \textbackslash{}
0                                         0.92
1                                         0.91
2                                         0.89
3                                         0.92
4                                         1.00

   Tipo\_cambio\_mxpdolar\_prom\_nacional  Demanda\_maizblanco\_kg\_prom\_nacional  \textbackslash{}
0                                9.48                           1625000000
1                                9.44                           1627000000
2                                9.31                           1630000000
3                                9.38                           1628000000
4                                9.52                           1632000000

   Costo\_transporte\_maiz\_mxpkilo\_prom\_nac
0                                    0.22
1                                    0.22
2                                    0.23
3                                    0.23
4                                    0.23
\end{Verbatim}
\end{tcolorbox}
        
    \hypertarget{cuxe1lculo-de-vif}{%
\subsection{17.2 Cálculo de VIF}\label{cuxe1lculo-de-vif}}

    \begin{tcolorbox}[breakable, size=fbox, boxrule=1pt, pad at break*=1mm,colback=cellbackground, colframe=cellborder]
\prompt{In}{incolor}{35}{\boxspacing}
\begin{Verbatim}[commandchars=\\\{\}]
\PY{k+kn}{from} \PY{n+nn}{statsmodels}\PY{n+nn}{.}\PY{n+nn}{stats}\PY{n+nn}{.}\PY{n+nn}{outliers\PYZus{}influence} \PY{k+kn}{import} \PY{n}{variance\PYZus{}inflation\PYZus{}factor}
\PY{k+kn}{import} \PY{n+nn}{pandas} \PY{k}{as} \PY{n+nn}{pd}

\PY{n}{vif} \PY{o}{=} \PY{n}{pd}\PY{o}{.}\PY{n}{DataFrame}\PY{p}{(}\PY{p}{)}

\PY{n}{vif}\PY{p}{[}\PY{l+s+s2}{\PYZdq{}}\PY{l+s+s2}{Variable}\PY{l+s+s2}{\PYZdq{}}\PY{p}{]} \PY{o}{=} \PY{n}{X}\PY{o}{.}\PY{n}{columns}

\PY{n}{vif}\PY{p}{[}\PY{l+s+s2}{\PYZdq{}}\PY{l+s+s2}{VIF}\PY{l+s+s2}{\PYZdq{}}\PY{p}{]} \PY{o}{=} \PY{p}{[}
    \PY{n}{variance\PYZus{}inflation\PYZus{}factor}\PY{p}{(}
        \PY{n}{X}\PY{o}{.}\PY{n}{values}\PY{p}{,}
        \PY{n}{i}
    \PY{p}{)}
    \PY{k}{for} \PY{n}{i} \PY{o+ow}{in} \PY{n+nb}{range}\PY{p}{(}\PY{n}{X}\PY{o}{.}\PY{n}{shape}\PY{p}{[}\PY{l+m+mi}{1}\PY{p}{]}\PY{p}{)}
\PY{p}{]}

\PY{n}{vif}\PY{o}{.}\PY{n}{sort\PYZus{}values}\PY{p}{(}
    \PY{n}{by}\PY{o}{=}\PY{l+s+s2}{\PYZdq{}}\PY{l+s+s2}{VIF}\PY{l+s+s2}{\PYZdq{}}\PY{p}{,}
    \PY{n}{ascending}\PY{o}{=}\PY{k+kc}{False}
\PY{p}{)}
\end{Verbatim}
\end{tcolorbox}

            \begin{tcolorbox}[breakable, size=fbox, boxrule=.5pt, pad at break*=1mm, opacityfill=0]
\prompt{Out}{outcolor}{35}{\boxspacing}
\begin{Verbatim}[commandchars=\\\{\}]
                                          Variable          VIF
4                               INPC\_prom\_nacional  5166.890214
9           Costo\_transporte\_maiz\_mxpkilo\_prom\_nac  1897.810881
8              Demanda\_maizblanco\_kg\_prom\_nacional  1474.861106
2  Precio\_diesel\_automotriz\_mxplitro\_prom\_nacional   308.751802
7               Tipo\_cambio\_mxpdolar\_prom\_nacional   138.959007
0                Temperatura\_celsius\_prom\_nacional    88.643332
6      Precio\_prom\_CBOT\_mxpkilo\_maiz\_internacional    21.312846
3  Indice\_severidad\_ponderada\_sequia\_prom\_nacional     4.137364
1                          Lluvia\_mm\_prom\_nacional     3.986003
5           Produccion\_maizblanco\_kg\_prom\_nacional     3.457557
\end{Verbatim}
\end{tcolorbox}
        
    \hypertarget{interpretaciuxf3n}{%
\subsubsection{Interpretación}\label{interpretaciuxf3n}}

Los resultados sugieren que varias variables económicas presentan
información altamente redundante.

Particularmente, el INPC, los costos de transporte, la demanda nacional
y el precio del diésel exhiben niveles de multicolinealidad
incompatibles con una especificación parsimoniosa del modelo.

Este comportamiento era esperable debido a que dichas variables
presentan tendencias crecientes de largo plazo y se encuentran
fuertemente asociadas entre sí.

Por el contrario, las variables relacionadas con la producción agrícola
y los factores climáticos presentan niveles de multicolinealidad
considerablemente menores, indicando que aportan información más
independiente al modelo.

\hypertarget{decisiuxf3n}{%
\subsubsection{Decisión}\label{decisiuxf3n}}

Con el objetivo de reducir la redundancia informativa y mejorar la
estabilidad estadística del modelo SARIMAX, se realizará una segunda
iteración del análisis VIF eliminando inicialmente las variables:

\begin{itemize}
\tightlist
\item
  INPC\_prom\_nacional
\item
  Demanda\_maizblanco\_kg\_prom\_nacional
\end{itemize}

Estas variables fueron seleccionadas para su eliminación preliminar
debido a que presentan los mayores valores de VIF observados y, además,
muestran una fuerte asociación con la tendencia temporal de largo plazo.

Posteriormente se recalculará el VIF para evaluar si persisten problemas
de multicolinealidad entre las variables restantes.

    \hypertarget{segunda-iteraciuxf3n-del-anuxe1lisis-vif}{%
\subsection{17.3 Segunda iteración del análisis
VIF}\label{segunda-iteraciuxf3n-del-anuxe1lisis-vif}}

Tras identificar niveles extremos de multicolinealidad en el conjunto
inicial de variables exógenas, se procedió a eliminar preliminarmente
las variables:

\begin{itemize}
\tightlist
\item
  INPC\_prom\_nacional
\item
  Demanda\_maizblanco\_kg\_prom\_nacional
\end{itemize}

La eliminación se fundamenta en sus elevados valores de VIF y en la
fuerte redundancia observada con otras variables económicas y
logísticas.

Posteriormente se recalculó el VIF utilizando únicamente las variables
restantes, con el objetivo de evaluar si la multicolinealidad disminuye
y obtener una especificación más parsimoniosa para el modelo SARIMAX.con
otras variables económicas y logísticas.

Posteriormente se recalculó el VIF utilizando únicamente las variables
restantes, con el objetivo de evaluar si la multicolinealidad disminuye
y obtener una especificación más parsimoniosa para el modelo SARIMAX.

    \begin{tcolorbox}[breakable, size=fbox, boxrule=1pt, pad at break*=1mm,colback=cellbackground, colframe=cellborder]
\prompt{In}{incolor}{37}{\boxspacing}
\begin{Verbatim}[commandchars=\\\{\}]
\PY{n}{variables\PYZus{}vif\PYZus{}2} \PY{o}{=} \PY{p}{[}
    \PY{l+s+s2}{\PYZdq{}}\PY{l+s+s2}{Temperatura\PYZus{}celsius\PYZus{}prom\PYZus{}nacional}\PY{l+s+s2}{\PYZdq{}}\PY{p}{,}
    \PY{l+s+s2}{\PYZdq{}}\PY{l+s+s2}{Lluvia\PYZus{}mm\PYZus{}prom\PYZus{}nacional}\PY{l+s+s2}{\PYZdq{}}\PY{p}{,}
    \PY{l+s+s2}{\PYZdq{}}\PY{l+s+s2}{Precio\PYZus{}diesel\PYZus{}automotriz\PYZus{}mxplitro\PYZus{}prom\PYZus{}nacional}\PY{l+s+s2}{\PYZdq{}}\PY{p}{,}
    \PY{l+s+s2}{\PYZdq{}}\PY{l+s+s2}{Indice\PYZus{}severidad\PYZus{}ponderada\PYZus{}sequia\PYZus{}prom\PYZus{}nacional}\PY{l+s+s2}{\PYZdq{}}\PY{p}{,}
    \PY{l+s+s2}{\PYZdq{}}\PY{l+s+s2}{Produccion\PYZus{}maizblanco\PYZus{}kg\PYZus{}prom\PYZus{}nacional}\PY{l+s+s2}{\PYZdq{}}\PY{p}{,}
    \PY{l+s+s2}{\PYZdq{}}\PY{l+s+s2}{Precio\PYZus{}prom\PYZus{}CBOT\PYZus{}mxpkilo\PYZus{}maiz\PYZus{}internacional}\PY{l+s+s2}{\PYZdq{}}\PY{p}{,}
    \PY{l+s+s2}{\PYZdq{}}\PY{l+s+s2}{Tipo\PYZus{}cambio\PYZus{}mxpdolar\PYZus{}prom\PYZus{}nacional}\PY{l+s+s2}{\PYZdq{}}\PY{p}{,}
    \PY{l+s+s2}{\PYZdq{}}\PY{l+s+s2}{Costo\PYZus{}transporte\PYZus{}maiz\PYZus{}mxpkilo\PYZus{}prom\PYZus{}nac}\PY{l+s+s2}{\PYZdq{}}
\PY{p}{]}

\PY{n}{X2} \PY{o}{=} \PY{n}{pdf\PYZus{}completo}\PY{p}{[}\PY{n}{variables\PYZus{}vif\PYZus{}2}\PY{p}{]}
\end{Verbatim}
\end{tcolorbox}

    \begin{tcolorbox}[breakable, size=fbox, boxrule=1pt, pad at break*=1mm,colback=cellbackground, colframe=cellborder]
\prompt{In}{incolor}{39}{\boxspacing}
\begin{Verbatim}[commandchars=\\\{\}]
\PY{k+kn}{from} \PY{n+nn}{statsmodels}\PY{n+nn}{.}\PY{n+nn}{stats}\PY{n+nn}{.}\PY{n+nn}{outliers\PYZus{}influence} \PY{k+kn}{import} \PY{n}{variance\PYZus{}inflation\PYZus{}factor}
\PY{k+kn}{import} \PY{n+nn}{pandas} \PY{k}{as} \PY{n+nn}{pd}

\PY{n}{vif2} \PY{o}{=} \PY{n}{pd}\PY{o}{.}\PY{n}{DataFrame}\PY{p}{(}\PY{p}{)}

\PY{n}{vif2}\PY{p}{[}\PY{l+s+s2}{\PYZdq{}}\PY{l+s+s2}{Variable}\PY{l+s+s2}{\PYZdq{}}\PY{p}{]} \PY{o}{=} \PY{n}{X2}\PY{o}{.}\PY{n}{columns}

\PY{n}{vif2}\PY{p}{[}\PY{l+s+s2}{\PYZdq{}}\PY{l+s+s2}{VIF}\PY{l+s+s2}{\PYZdq{}}\PY{p}{]} \PY{o}{=} \PY{p}{[}
    \PY{n}{variance\PYZus{}inflation\PYZus{}factor}\PY{p}{(}
        \PY{n}{X2}\PY{o}{.}\PY{n}{values}\PY{p}{,}
        \PY{n}{i}
    \PY{p}{)}
    \PY{k}{for} \PY{n}{i} \PY{o+ow}{in} \PY{n+nb}{range}\PY{p}{(}\PY{n}{X2}\PY{o}{.}\PY{n}{shape}\PY{p}{[}\PY{l+m+mi}{1}\PY{p}{]}\PY{p}{)}
\PY{p}{]}

\PY{n}{vif2} \PY{o}{=} \PY{p}{(}
    \PY{n}{vif2}
    \PY{o}{.}\PY{n}{sort\PYZus{}values}\PY{p}{(}
        \PY{n}{by}\PY{o}{=}\PY{l+s+s2}{\PYZdq{}}\PY{l+s+s2}{VIF}\PY{l+s+s2}{\PYZdq{}}\PY{p}{,}
        \PY{n}{ascending}\PY{o}{=}\PY{k+kc}{False}
    \PY{p}{)}
\PY{p}{)}

\PY{n}{vif2}
\end{Verbatim}
\end{tcolorbox}

            \begin{tcolorbox}[breakable, size=fbox, boxrule=.5pt, pad at break*=1mm, opacityfill=0]
\prompt{Out}{outcolor}{39}{\boxspacing}
\begin{Verbatim}[commandchars=\\\{\}]
                                          Variable         VIF
2  Precio\_diesel\_automotriz\_mxplitro\_prom\_nacional  263.827525
7           Costo\_transporte\_maiz\_mxpkilo\_prom\_nac  199.027719
6               Tipo\_cambio\_mxpdolar\_prom\_nacional   91.933225
0                Temperatura\_celsius\_prom\_nacional   35.104434
5      Precio\_prom\_CBOT\_mxpkilo\_maiz\_internacional   13.690862
3  Indice\_severidad\_ponderada\_sequia\_prom\_nacional    3.936483
1                          Lluvia\_mm\_prom\_nacional    3.320416
4           Produccion\_maizblanco\_kg\_prom\_nacional    3.051112
\end{Verbatim}
\end{tcolorbox}
        
    \hypertarget{resultados-de-la-segunda-iteraciuxf3n}{%
\subsubsection{Resultados de la segunda
iteración}\label{resultados-de-la-segunda-iteraciuxf3n}}

La eliminación del INPC y la demanda nacional redujo de manera
importante los niveles generales de multicolinealidad observados en el
conjunto de variables exógenas.

No obstante, persisten niveles elevados de VIF en las variables
relacionadas con costos logísticos y energéticos, particularmente en:

\begin{itemize}
\tightlist
\item
  Precio del diésel.
\item
  Costos de transporte.
\item
  Tipo de cambio.
\end{itemize}

Estos resultados sugieren la existencia de un bloque económico altamente
correlacionado que captura fenómenos similares asociados a costos de
operación y comercialización.

Dado que el precio del diésel y los costos de transporte representan
mecanismos económicos estrechamente relacionados, se evaluarán
especificaciones alternativas del modelo conservando únicamente una de
estas variables en iteraciones posteriores.

    \hypertarget{tercera-iteraciuxf3n-conservando-costo-de-transporte}{%
\subsection{17.4 Tercera iteración: (conservando costo de
transporte)}\label{tercera-iteraciuxf3n-conservando-costo-de-transporte}}

En esta iteración se conserva la variable de costos de transporte y se
elimina el precio del diésel.

La decisión se fundamenta en que los costos de transporte representan
una medida más cercana al costo efectivo de movilización del maíz y
podrían capturar indirectamente el efecto de los combustibles.

El objetivo es evaluar si la eliminación del precio del diésel reduce
significativamente los niveles de multicolinealidad observados.

    \begin{tcolorbox}[breakable, size=fbox, boxrule=1pt, pad at break*=1mm,colback=cellbackground, colframe=cellborder]
\prompt{In}{incolor}{40}{\boxspacing}
\begin{Verbatim}[commandchars=\\\{\}]
\PY{n}{variables\PYZus{}vif\PYZus{}3A} \PY{o}{=} \PY{p}{[}
    \PY{l+s+s2}{\PYZdq{}}\PY{l+s+s2}{Temperatura\PYZus{}celsius\PYZus{}prom\PYZus{}nacional}\PY{l+s+s2}{\PYZdq{}}\PY{p}{,}
    \PY{l+s+s2}{\PYZdq{}}\PY{l+s+s2}{Lluvia\PYZus{}mm\PYZus{}prom\PYZus{}nacional}\PY{l+s+s2}{\PYZdq{}}\PY{p}{,}
    \PY{l+s+s2}{\PYZdq{}}\PY{l+s+s2}{Indice\PYZus{}severidad\PYZus{}ponderada\PYZus{}sequia\PYZus{}prom\PYZus{}nacional}\PY{l+s+s2}{\PYZdq{}}\PY{p}{,}
    \PY{l+s+s2}{\PYZdq{}}\PY{l+s+s2}{Produccion\PYZus{}maizblanco\PYZus{}kg\PYZus{}prom\PYZus{}nacional}\PY{l+s+s2}{\PYZdq{}}\PY{p}{,}
    \PY{l+s+s2}{\PYZdq{}}\PY{l+s+s2}{Precio\PYZus{}prom\PYZus{}CBOT\PYZus{}mxpkilo\PYZus{}maiz\PYZus{}internacional}\PY{l+s+s2}{\PYZdq{}}\PY{p}{,}
    \PY{l+s+s2}{\PYZdq{}}\PY{l+s+s2}{Tipo\PYZus{}cambio\PYZus{}mxpdolar\PYZus{}prom\PYZus{}nacional}\PY{l+s+s2}{\PYZdq{}}\PY{p}{,}
    \PY{l+s+s2}{\PYZdq{}}\PY{l+s+s2}{Costo\PYZus{}transporte\PYZus{}maiz\PYZus{}mxpkilo\PYZus{}prom\PYZus{}nac}\PY{l+s+s2}{\PYZdq{}}
\PY{p}{]}

\PY{n}{X3A} \PY{o}{=} \PY{n}{pdf\PYZus{}completo}\PY{p}{[}\PY{n}{variables\PYZus{}vif\PYZus{}3A}\PY{p}{]}
\end{Verbatim}
\end{tcolorbox}

    \begin{tcolorbox}[breakable, size=fbox, boxrule=1pt, pad at break*=1mm,colback=cellbackground, colframe=cellborder]
\prompt{In}{incolor}{41}{\boxspacing}
\begin{Verbatim}[commandchars=\\\{\}]
\PY{n}{vif3A} \PY{o}{=} \PY{n}{pd}\PY{o}{.}\PY{n}{DataFrame}\PY{p}{(}\PY{p}{)}

\PY{n}{vif3A}\PY{p}{[}\PY{l+s+s2}{\PYZdq{}}\PY{l+s+s2}{Variable}\PY{l+s+s2}{\PYZdq{}}\PY{p}{]} \PY{o}{=} \PY{n}{X3A}\PY{o}{.}\PY{n}{columns}

\PY{n}{vif3A}\PY{p}{[}\PY{l+s+s2}{\PYZdq{}}\PY{l+s+s2}{VIF}\PY{l+s+s2}{\PYZdq{}}\PY{p}{]} \PY{o}{=} \PY{p}{[}
    \PY{n}{variance\PYZus{}inflation\PYZus{}factor}\PY{p}{(}
        \PY{n}{X3A}\PY{o}{.}\PY{n}{values}\PY{p}{,}
        \PY{n}{i}
    \PY{p}{)}
    \PY{k}{for} \PY{n}{i} \PY{o+ow}{in} \PY{n+nb}{range}\PY{p}{(}\PY{n}{X3A}\PY{o}{.}\PY{n}{shape}\PY{p}{[}\PY{l+m+mi}{1}\PY{p}{]}\PY{p}{)}
\PY{p}{]}

\PY{n}{vif3A} \PY{o}{=} \PY{n}{vif3A}\PY{o}{.}\PY{n}{sort\PYZus{}values}\PY{p}{(}
    \PY{n}{by}\PY{o}{=}\PY{l+s+s2}{\PYZdq{}}\PY{l+s+s2}{VIF}\PY{l+s+s2}{\PYZdq{}}\PY{p}{,}
    \PY{n}{ascending}\PY{o}{=}\PY{k+kc}{False}
\PY{p}{)}

\PY{n}{vif3A}
\end{Verbatim}
\end{tcolorbox}

            \begin{tcolorbox}[breakable, size=fbox, boxrule=.5pt, pad at break*=1mm, opacityfill=0]
\prompt{Out}{outcolor}{41}{\boxspacing}
\begin{Verbatim}[commandchars=\\\{\}]
                                          Variable        VIF
5               Tipo\_cambio\_mxpdolar\_prom\_nacional  74.035622
0                Temperatura\_celsius\_prom\_nacional  32.902811
6           Costo\_transporte\_maiz\_mxpkilo\_prom\_nac  22.812526
4      Precio\_prom\_CBOT\_mxpkilo\_maiz\_internacional  13.530468
2  Indice\_severidad\_ponderada\_sequia\_prom\_nacional   3.875038
1                          Lluvia\_mm\_prom\_nacional   3.287472
3           Produccion\_maizblanco\_kg\_prom\_nacional   3.021775
\end{Verbatim}
\end{tcolorbox}
        
    \hypertarget{tercera-iteraciuxf3n-conservando-diuxe9sel}{%
\subsection{17.5 Tercera iteración: (Conservando
Diésel)}\label{tercera-iteraciuxf3n-conservando-diuxe9sel}}

En esta iteración se conserva el precio del diésel y se elimina la
variable de costos de transporte.

La lógica detrás de esta decisión es evaluar si una variable primaria
observada directamente puede representar adecuadamente el efecto de los
costos energéticos sin introducir niveles excesivos de
multicolinealidad.

Los resultados serán comparados con los obtenidos en el Modelo A.

    \begin{tcolorbox}[breakable, size=fbox, boxrule=1pt, pad at break*=1mm,colback=cellbackground, colframe=cellborder]
\prompt{In}{incolor}{42}{\boxspacing}
\begin{Verbatim}[commandchars=\\\{\}]
\PY{n}{variables\PYZus{}vif\PYZus{}3B} \PY{o}{=} \PY{p}{[}
    \PY{l+s+s2}{\PYZdq{}}\PY{l+s+s2}{Temperatura\PYZus{}celsius\PYZus{}prom\PYZus{}nacional}\PY{l+s+s2}{\PYZdq{}}\PY{p}{,}
    \PY{l+s+s2}{\PYZdq{}}\PY{l+s+s2}{Lluvia\PYZus{}mm\PYZus{}prom\PYZus{}nacional}\PY{l+s+s2}{\PYZdq{}}\PY{p}{,}
    \PY{l+s+s2}{\PYZdq{}}\PY{l+s+s2}{Indice\PYZus{}severidad\PYZus{}ponderada\PYZus{}sequia\PYZus{}prom\PYZus{}nacional}\PY{l+s+s2}{\PYZdq{}}\PY{p}{,}
    \PY{l+s+s2}{\PYZdq{}}\PY{l+s+s2}{Produccion\PYZus{}maizblanco\PYZus{}kg\PYZus{}prom\PYZus{}nacional}\PY{l+s+s2}{\PYZdq{}}\PY{p}{,}
    \PY{l+s+s2}{\PYZdq{}}\PY{l+s+s2}{Precio\PYZus{}prom\PYZus{}CBOT\PYZus{}mxpkilo\PYZus{}maiz\PYZus{}internacional}\PY{l+s+s2}{\PYZdq{}}\PY{p}{,}
    \PY{l+s+s2}{\PYZdq{}}\PY{l+s+s2}{Tipo\PYZus{}cambio\PYZus{}mxpdolar\PYZus{}prom\PYZus{}nacional}\PY{l+s+s2}{\PYZdq{}}\PY{p}{,}
    \PY{l+s+s2}{\PYZdq{}}\PY{l+s+s2}{Precio\PYZus{}diesel\PYZus{}automotriz\PYZus{}mxplitro\PYZus{}prom\PYZus{}nacional}\PY{l+s+s2}{\PYZdq{}}
\PY{p}{]}

\PY{n}{X3B} \PY{o}{=} \PY{n}{pdf\PYZus{}completo}\PY{p}{[}\PY{n}{variables\PYZus{}vif\PYZus{}3B}\PY{p}{]}
\end{Verbatim}
\end{tcolorbox}

    \begin{tcolorbox}[breakable, size=fbox, boxrule=1pt, pad at break*=1mm,colback=cellbackground, colframe=cellborder]
\prompt{In}{incolor}{43}{\boxspacing}
\begin{Verbatim}[commandchars=\\\{\}]
\PY{n}{vif3B} \PY{o}{=} \PY{n}{pd}\PY{o}{.}\PY{n}{DataFrame}\PY{p}{(}\PY{p}{)}

\PY{n}{vif3B}\PY{p}{[}\PY{l+s+s2}{\PYZdq{}}\PY{l+s+s2}{Variable}\PY{l+s+s2}{\PYZdq{}}\PY{p}{]} \PY{o}{=} \PY{n}{X3B}\PY{o}{.}\PY{n}{columns}

\PY{n}{vif3B}\PY{p}{[}\PY{l+s+s2}{\PYZdq{}}\PY{l+s+s2}{VIF}\PY{l+s+s2}{\PYZdq{}}\PY{p}{]} \PY{o}{=} \PY{p}{[}
    \PY{n}{variance\PYZus{}inflation\PYZus{}factor}\PY{p}{(}
        \PY{n}{X3B}\PY{o}{.}\PY{n}{values}\PY{p}{,}
        \PY{n}{i}
    \PY{p}{)}
    \PY{k}{for} \PY{n}{i} \PY{o+ow}{in} \PY{n+nb}{range}\PY{p}{(}\PY{n}{X3B}\PY{o}{.}\PY{n}{shape}\PY{p}{[}\PY{l+m+mi}{1}\PY{p}{]}\PY{p}{)}
\PY{p}{]}

\PY{n}{vif3B} \PY{o}{=} \PY{n}{vif3B}\PY{o}{.}\PY{n}{sort\PYZus{}values}\PY{p}{(}
    \PY{n}{by}\PY{o}{=}\PY{l+s+s2}{\PYZdq{}}\PY{l+s+s2}{VIF}\PY{l+s+s2}{\PYZdq{}}\PY{p}{,}
    \PY{n}{ascending}\PY{o}{=}\PY{k+kc}{False}
\PY{p}{)}

\PY{n}{vif3B}
\end{Verbatim}
\end{tcolorbox}

            \begin{tcolorbox}[breakable, size=fbox, boxrule=.5pt, pad at break*=1mm, opacityfill=0]
\prompt{Out}{outcolor}{43}{\boxspacing}
\begin{Verbatim}[commandchars=\\\{\}]
                                          Variable        VIF
5               Tipo\_cambio\_mxpdolar\_prom\_nacional  88.831344
0                Temperatura\_celsius\_prom\_nacional  35.079632
6  Precio\_diesel\_automotriz\_mxplitro\_prom\_nacional  30.239869
4      Precio\_prom\_CBOT\_mxpkilo\_maiz\_internacional  13.670499
2  Indice\_severidad\_ponderada\_sequia\_prom\_nacional   3.575313
1                          Lluvia\_mm\_prom\_nacional   3.320321
3           Produccion\_maizblanco\_kg\_prom\_nacional   3.050887
\end{Verbatim}
\end{tcolorbox}
        
    \hypertarget{cuarta-iteraciuxf3n-eliminando-tipo-de-cambio}{%
\subsection{17.6 Cuarta iteración: (Eliminando tipo de
cambio)}\label{cuarta-iteraciuxf3n-eliminando-tipo-de-cambio}}

Vamos a evaluar la hipótesis:

¿El Tipo de Cambio está inflando los VIF de CBOT, Transporte y
Temperatura?

    \begin{tcolorbox}[breakable, size=fbox, boxrule=1pt, pad at break*=1mm,colback=cellbackground, colframe=cellborder]
\prompt{In}{incolor}{44}{\boxspacing}
\begin{Verbatim}[commandchars=\\\{\}]
\PY{n}{variables\PYZus{}vif\PYZus{}4} \PY{o}{=} \PY{p}{[}
    \PY{l+s+s2}{\PYZdq{}}\PY{l+s+s2}{Temperatura\PYZus{}celsius\PYZus{}prom\PYZus{}nacional}\PY{l+s+s2}{\PYZdq{}}\PY{p}{,}
    \PY{l+s+s2}{\PYZdq{}}\PY{l+s+s2}{Lluvia\PYZus{}mm\PYZus{}prom\PYZus{}nacional}\PY{l+s+s2}{\PYZdq{}}\PY{p}{,}
    \PY{l+s+s2}{\PYZdq{}}\PY{l+s+s2}{Indice\PYZus{}severidad\PYZus{}ponderada\PYZus{}sequia\PYZus{}prom\PYZus{}nacional}\PY{l+s+s2}{\PYZdq{}}\PY{p}{,}
    \PY{l+s+s2}{\PYZdq{}}\PY{l+s+s2}{Produccion\PYZus{}maizblanco\PYZus{}kg\PYZus{}prom\PYZus{}nacional}\PY{l+s+s2}{\PYZdq{}}\PY{p}{,}
    \PY{l+s+s2}{\PYZdq{}}\PY{l+s+s2}{Precio\PYZus{}prom\PYZus{}CBOT\PYZus{}mxpkilo\PYZus{}maiz\PYZus{}internacional}\PY{l+s+s2}{\PYZdq{}}\PY{p}{,}
    \PY{l+s+s2}{\PYZdq{}}\PY{l+s+s2}{Costo\PYZus{}transporte\PYZus{}maiz\PYZus{}mxpkilo\PYZus{}prom\PYZus{}nac}\PY{l+s+s2}{\PYZdq{}}
\PY{p}{]}

\PY{n}{X4} \PY{o}{=} \PY{n}{pdf\PYZus{}completo}\PY{p}{[}\PY{n}{variables\PYZus{}vif\PYZus{}4}\PY{p}{]}
\end{Verbatim}
\end{tcolorbox}

    \begin{tcolorbox}[breakable, size=fbox, boxrule=1pt, pad at break*=1mm,colback=cellbackground, colframe=cellborder]
\prompt{In}{incolor}{45}{\boxspacing}
\begin{Verbatim}[commandchars=\\\{\}]
\PY{n}{vif4} \PY{o}{=} \PY{n}{pd}\PY{o}{.}\PY{n}{DataFrame}\PY{p}{(}\PY{p}{)}

\PY{n}{vif4}\PY{p}{[}\PY{l+s+s2}{\PYZdq{}}\PY{l+s+s2}{Variable}\PY{l+s+s2}{\PYZdq{}}\PY{p}{]} \PY{o}{=} \PY{n}{X4}\PY{o}{.}\PY{n}{columns}

\PY{n}{vif4}\PY{p}{[}\PY{l+s+s2}{\PYZdq{}}\PY{l+s+s2}{VIF}\PY{l+s+s2}{\PYZdq{}}\PY{p}{]} \PY{o}{=} \PY{p}{[}
    \PY{n}{variance\PYZus{}inflation\PYZus{}factor}\PY{p}{(}
        \PY{n}{X4}\PY{o}{.}\PY{n}{values}\PY{p}{,}
        \PY{n}{i}
    \PY{p}{)}
    \PY{k}{for} \PY{n}{i} \PY{o+ow}{in} \PY{n+nb}{range}\PY{p}{(}\PY{n}{X4}\PY{o}{.}\PY{n}{shape}\PY{p}{[}\PY{l+m+mi}{1}\PY{p}{]}\PY{p}{)}
\PY{p}{]}

\PY{n}{vif4} \PY{o}{=} \PY{p}{(}
    \PY{n}{vif4}
    \PY{o}{.}\PY{n}{sort\PYZus{}values}\PY{p}{(}
        \PY{n}{by}\PY{o}{=}\PY{l+s+s2}{\PYZdq{}}\PY{l+s+s2}{VIF}\PY{l+s+s2}{\PYZdq{}}\PY{p}{,}
        \PY{n}{ascending}\PY{o}{=}\PY{k+kc}{False}
    \PY{p}{)}
\PY{p}{)}

\PY{n}{vif4}
\end{Verbatim}
\end{tcolorbox}

            \begin{tcolorbox}[breakable, size=fbox, boxrule=.5pt, pad at break*=1mm, opacityfill=0]
\prompt{Out}{outcolor}{45}{\boxspacing}
\begin{Verbatim}[commandchars=\\\{\}]
                                          Variable        VIF
4      Precio\_prom\_CBOT\_mxpkilo\_maiz\_internacional  12.009685
5           Costo\_transporte\_maiz\_mxpkilo\_prom\_nac   9.569497
0                Temperatura\_celsius\_prom\_nacional   9.446766
1                          Lluvia\_mm\_prom\_nacional   2.924347
2  Indice\_severidad\_ponderada\_sequia\_prom\_nacional   2.772606
3           Produccion\_maizblanco\_kg\_prom\_nacional   2.703665
\end{Verbatim}
\end{tcolorbox}
        
    \begin{tcolorbox}[breakable, size=fbox, boxrule=1pt, pad at break*=1mm,colback=cellbackground, colframe=cellborder]
\prompt{In}{incolor}{48}{\boxspacing}
\begin{Verbatim}[commandchars=\\\{\}]
\PY{n}{comparacion\PYZus{}vif} \PY{o}{=} \PY{p}{(}
    \PY{n}{vif}\PY{o}{.}\PY{n}{rename}\PY{p}{(}\PY{n}{columns}\PY{o}{=}\PY{p}{\PYZob{}}\PY{l+s+s2}{\PYZdq{}}\PY{l+s+s2}{VIF}\PY{l+s+s2}{\PYZdq{}}\PY{p}{:}\PY{l+s+s2}{\PYZdq{}}\PY{l+s+s2}{Iteracion\PYZus{}1}\PY{l+s+s2}{\PYZdq{}}\PY{p}{\PYZcb{}}\PY{p}{)}
    \PY{o}{.}\PY{n}{merge}\PY{p}{(}
        \PY{n}{vif2}\PY{o}{.}\PY{n}{rename}\PY{p}{(}\PY{n}{columns}\PY{o}{=}\PY{p}{\PYZob{}}\PY{l+s+s2}{\PYZdq{}}\PY{l+s+s2}{VIF}\PY{l+s+s2}{\PYZdq{}}\PY{p}{:}\PY{l+s+s2}{\PYZdq{}}\PY{l+s+s2}{Iteracion\PYZus{}2}\PY{l+s+s2}{\PYZdq{}}\PY{p}{\PYZcb{}}\PY{p}{)}\PY{p}{,}
        \PY{n}{on}\PY{o}{=}\PY{l+s+s2}{\PYZdq{}}\PY{l+s+s2}{Variable}\PY{l+s+s2}{\PYZdq{}}\PY{p}{,}
        \PY{n}{how}\PY{o}{=}\PY{l+s+s2}{\PYZdq{}}\PY{l+s+s2}{outer}\PY{l+s+s2}{\PYZdq{}}
    \PY{p}{)}
    \PY{o}{.}\PY{n}{merge}\PY{p}{(}
        \PY{n}{vif3A}\PY{o}{.}\PY{n}{rename}\PY{p}{(}\PY{n}{columns}\PY{o}{=}\PY{p}{\PYZob{}}\PY{l+s+s2}{\PYZdq{}}\PY{l+s+s2}{VIF}\PY{l+s+s2}{\PYZdq{}}\PY{p}{:}\PY{l+s+s2}{\PYZdq{}}\PY{l+s+s2}{Iteracion\PYZus{}3A}\PY{l+s+s2}{\PYZdq{}}\PY{p}{\PYZcb{}}\PY{p}{)}\PY{p}{,}
        \PY{n}{on}\PY{o}{=}\PY{l+s+s2}{\PYZdq{}}\PY{l+s+s2}{Variable}\PY{l+s+s2}{\PYZdq{}}\PY{p}{,}
        \PY{n}{how}\PY{o}{=}\PY{l+s+s2}{\PYZdq{}}\PY{l+s+s2}{outer}\PY{l+s+s2}{\PYZdq{}}
    \PY{p}{)}
    \PY{o}{.}\PY{n}{merge}\PY{p}{(}
        \PY{n}{vif3B}\PY{o}{.}\PY{n}{rename}\PY{p}{(}\PY{n}{columns}\PY{o}{=}\PY{p}{\PYZob{}}\PY{l+s+s2}{\PYZdq{}}\PY{l+s+s2}{VIF}\PY{l+s+s2}{\PYZdq{}}\PY{p}{:}\PY{l+s+s2}{\PYZdq{}}\PY{l+s+s2}{Iteracion\PYZus{}3B}\PY{l+s+s2}{\PYZdq{}}\PY{p}{\PYZcb{}}\PY{p}{)}\PY{p}{,}
        \PY{n}{on}\PY{o}{=}\PY{l+s+s2}{\PYZdq{}}\PY{l+s+s2}{Variable}\PY{l+s+s2}{\PYZdq{}}\PY{p}{,}
        \PY{n}{how}\PY{o}{=}\PY{l+s+s2}{\PYZdq{}}\PY{l+s+s2}{outer}\PY{l+s+s2}{\PYZdq{}}
    \PY{p}{)}
    \PY{o}{.}\PY{n}{merge}\PY{p}{(}
        \PY{n}{vif4}\PY{o}{.}\PY{n}{rename}\PY{p}{(}\PY{n}{columns}\PY{o}{=}\PY{p}{\PYZob{}}\PY{l+s+s2}{\PYZdq{}}\PY{l+s+s2}{VIF}\PY{l+s+s2}{\PYZdq{}}\PY{p}{:}\PY{l+s+s2}{\PYZdq{}}\PY{l+s+s2}{Iteracion\PYZus{}4}\PY{l+s+s2}{\PYZdq{}}\PY{p}{\PYZcb{}}\PY{p}{)}\PY{p}{,}
        \PY{n}{on}\PY{o}{=}\PY{l+s+s2}{\PYZdq{}}\PY{l+s+s2}{Variable}\PY{l+s+s2}{\PYZdq{}}\PY{p}{,}
        \PY{n}{how}\PY{o}{=}\PY{l+s+s2}{\PYZdq{}}\PY{l+s+s2}{outer}\PY{l+s+s2}{\PYZdq{}}
    \PY{p}{)}
\PY{p}{)}

\PY{n}{comparacion\PYZus{}vif}
\end{Verbatim}
\end{tcolorbox}

            \begin{tcolorbox}[breakable, size=fbox, boxrule=.5pt, pad at break*=1mm, opacityfill=0]
\prompt{Out}{outcolor}{48}{\boxspacing}
\begin{Verbatim}[commandchars=\\\{\}]
                                          Variable  Iteracion\_1  Iteracion\_2  \textbackslash{}
0                Temperatura\_celsius\_prom\_nacional    88.643332    35.104434
1                          Lluvia\_mm\_prom\_nacional     3.986003     3.320416
2  Precio\_diesel\_automotriz\_mxplitro\_prom\_nacional   308.751802   263.827525
3  Indice\_severidad\_ponderada\_sequia\_prom\_nacional     4.137364     3.936483
4                               INPC\_prom\_nacional  5166.890214          NaN
5           Produccion\_maizblanco\_kg\_prom\_nacional     3.457557     3.051112
6      Precio\_prom\_CBOT\_mxpkilo\_maiz\_internacional    21.312846    13.690862
7               Tipo\_cambio\_mxpdolar\_prom\_nacional   138.959007    91.933225
8              Demanda\_maizblanco\_kg\_prom\_nacional  1474.861106          NaN
9           Costo\_transporte\_maiz\_mxpkilo\_prom\_nac  1897.810881   199.027719

   Iteracion\_3A  Iteracion\_3B  Iteracion\_4
0     32.902811     35.079632     9.446766
1      3.287472      3.320321     2.924347
2           NaN     30.239869          NaN
3      3.875038      3.575313     2.772606
4           NaN           NaN          NaN
5      3.021775      3.050887     2.703665
6     13.530468     13.670499    12.009685
7     74.035622     88.831344          NaN
8           NaN           NaN          NaN
9     22.812526           NaN     9.569497
\end{Verbatim}
\end{tcolorbox}
        
    \begin{tcolorbox}[breakable, size=fbox, boxrule=1pt, pad at break*=1mm,colback=cellbackground, colframe=cellborder]
\prompt{In}{incolor}{52}{\boxspacing}
\begin{Verbatim}[commandchars=\\\{\}]
\PY{n}{resumen\PYZus{}modelos} \PY{o}{=} \PY{n}{pd}\PY{o}{.}\PY{n}{DataFrame}\PY{p}{(}\PY{p}{\PYZob{}}
    \PY{l+s+s2}{\PYZdq{}}\PY{l+s+s2}{Modelo}\PY{l+s+s2}{\PYZdq{}}\PY{p}{:}\PY{p}{[}
        \PY{l+s+s2}{\PYZdq{}}\PY{l+s+s2}{Iteracion\PYZus{}1}\PY{l+s+s2}{\PYZdq{}}\PY{p}{,}
        \PY{l+s+s2}{\PYZdq{}}\PY{l+s+s2}{Iteracion\PYZus{}2}\PY{l+s+s2}{\PYZdq{}}\PY{p}{,}
        \PY{l+s+s2}{\PYZdq{}}\PY{l+s+s2}{Iteracion\PYZus{}3A}\PY{l+s+s2}{\PYZdq{}}\PY{p}{,}
        \PY{l+s+s2}{\PYZdq{}}\PY{l+s+s2}{Iteracion\PYZus{}3B}\PY{l+s+s2}{\PYZdq{}}\PY{p}{,}
        \PY{l+s+s2}{\PYZdq{}}\PY{l+s+s2}{Iteracion\PYZus{}4}\PY{l+s+s2}{\PYZdq{}}
    \PY{p}{]}\PY{p}{,}
    \PY{l+s+s2}{\PYZdq{}}\PY{l+s+s2}{Variables}\PY{l+s+s2}{\PYZdq{}}\PY{p}{:}\PY{p}{[}
        \PY{n+nb}{len}\PY{p}{(}\PY{n}{vif}\PY{p}{)}\PY{p}{,}
        \PY{n+nb}{len}\PY{p}{(}\PY{n}{vif2}\PY{p}{)}\PY{p}{,}
        \PY{n+nb}{len}\PY{p}{(}\PY{n}{vif3A}\PY{p}{)}\PY{p}{,}
        \PY{n+nb}{len}\PY{p}{(}\PY{n}{vif3B}\PY{p}{)}\PY{p}{,}
        \PY{n+nb}{len}\PY{p}{(}\PY{n}{vif4}\PY{p}{)}
    \PY{p}{]}\PY{p}{,}
    \PY{l+s+s2}{\PYZdq{}}\PY{l+s+s2}{VIF\PYZus{}max}\PY{l+s+s2}{\PYZdq{}}\PY{p}{:}\PY{p}{[}
        \PY{n}{vif}\PY{p}{[}\PY{l+s+s2}{\PYZdq{}}\PY{l+s+s2}{VIF}\PY{l+s+s2}{\PYZdq{}}\PY{p}{]}\PY{o}{.}\PY{n}{max}\PY{p}{(}\PY{p}{)}\PY{p}{,}
        \PY{n}{vif2}\PY{p}{[}\PY{l+s+s2}{\PYZdq{}}\PY{l+s+s2}{VIF}\PY{l+s+s2}{\PYZdq{}}\PY{p}{]}\PY{o}{.}\PY{n}{max}\PY{p}{(}\PY{p}{)}\PY{p}{,}
        \PY{n}{vif3A}\PY{p}{[}\PY{l+s+s2}{\PYZdq{}}\PY{l+s+s2}{VIF}\PY{l+s+s2}{\PYZdq{}}\PY{p}{]}\PY{o}{.}\PY{n}{max}\PY{p}{(}\PY{p}{)}\PY{p}{,}
        \PY{n}{vif3B}\PY{p}{[}\PY{l+s+s2}{\PYZdq{}}\PY{l+s+s2}{VIF}\PY{l+s+s2}{\PYZdq{}}\PY{p}{]}\PY{o}{.}\PY{n}{max}\PY{p}{(}\PY{p}{)}\PY{p}{,}
        \PY{n}{vif4}\PY{p}{[}\PY{l+s+s2}{\PYZdq{}}\PY{l+s+s2}{VIF}\PY{l+s+s2}{\PYZdq{}}\PY{p}{]}\PY{o}{.}\PY{n}{max}\PY{p}{(}\PY{p}{)}
    \PY{p}{]}\PY{p}{,}
    \PY{l+s+s2}{\PYZdq{}}\PY{l+s+s2}{VIF\PYZus{}promedio}\PY{l+s+s2}{\PYZdq{}}\PY{p}{:}\PY{p}{[}
        \PY{n}{vif}\PY{p}{[}\PY{l+s+s2}{\PYZdq{}}\PY{l+s+s2}{VIF}\PY{l+s+s2}{\PYZdq{}}\PY{p}{]}\PY{o}{.}\PY{n}{mean}\PY{p}{(}\PY{p}{)}\PY{p}{,}
        \PY{n}{vif2}\PY{p}{[}\PY{l+s+s2}{\PYZdq{}}\PY{l+s+s2}{VIF}\PY{l+s+s2}{\PYZdq{}}\PY{p}{]}\PY{o}{.}\PY{n}{mean}\PY{p}{(}\PY{p}{)}\PY{p}{,}
        \PY{n}{vif3A}\PY{p}{[}\PY{l+s+s2}{\PYZdq{}}\PY{l+s+s2}{VIF}\PY{l+s+s2}{\PYZdq{}}\PY{p}{]}\PY{o}{.}\PY{n}{mean}\PY{p}{(}\PY{p}{)}\PY{p}{,}
        \PY{n}{vif3B}\PY{p}{[}\PY{l+s+s2}{\PYZdq{}}\PY{l+s+s2}{VIF}\PY{l+s+s2}{\PYZdq{}}\PY{p}{]}\PY{o}{.}\PY{n}{mean}\PY{p}{(}\PY{p}{)}\PY{p}{,}
        \PY{n}{vif4}\PY{p}{[}\PY{l+s+s2}{\PYZdq{}}\PY{l+s+s2}{VIF}\PY{l+s+s2}{\PYZdq{}}\PY{p}{]}\PY{o}{.}\PY{n}{mean}\PY{p}{(}\PY{p}{)}
    \PY{p}{]}
\PY{p}{\PYZcb{}}\PY{p}{)}

\PY{n}{resumen\PYZus{}modelos}
\end{Verbatim}
\end{tcolorbox}

            \begin{tcolorbox}[breakable, size=fbox, boxrule=.5pt, pad at break*=1mm, opacityfill=0]
\prompt{Out}{outcolor}{52}{\boxspacing}
\begin{Verbatim}[commandchars=\\\{\}]
         Modelo  Variables      VIF\_max  VIF\_promedio
0   Iteracion\_1         10  5166.890214    910.881011
1   Iteracion\_2          8   263.827525     76.736472
2  Iteracion\_3A          7    74.035622     21.923673
3  Iteracion\_3B          7    88.831344     25.395410
4   Iteracion\_4          6    12.009685      6.571094
\end{Verbatim}
\end{tcolorbox}
        
    \hypertarget{selecciuxf3n-final-de-variables-exuxf3genas}{%
\section{17.7 Selección Final de Variables
Exógenas}\label{selecciuxf3n-final-de-variables-exuxf3genas}}

El proceso de selección de variables exógenas se realizó mediante un
enfoque iterativo basado en el análisis de multicolinealidad utilizando
el Factor de Inflación de la Varianza (VIF).

La primera especificación incluyó diez variables explicativas
relacionadas con factores económicos, productivos, climáticos y
logísticos. Sin embargo, los resultados mostraron niveles extremadamente
altos de multicolinealidad, alcanzando valores máximos superiores a
5,000 en algunas variables.

A partir de estos resultados se realizaron sucesivas iteraciones de
depuración, eliminando aquellas variables que aportaban información
redundante o que presentaban una fuerte asociación con tendencias
temporales compartidas.

Los principales ajustes realizados fueron:

\begin{itemize}
\tightlist
\item
  Eliminación del INPC debido a su elevada redundancia con otras
  variables económicas y logísticas.
\item
  Eliminación de la demanda nacional estimada por su fuerte asociación
  con la tendencia temporal de largo plazo.
\item
  Comparación de dos especificaciones alternativas conservando
  únicamente una variable representativa de los costos logísticos:

  \begin{itemize}
  \tightlist
  \item
    Modelo A: Conservando costos de transporte.
  \item
    Modelo B: Conservando precio del diésel.
  \end{itemize}
\item
  Evaluación adicional eliminando temporalmente el tipo de cambio para
  identificar su contribución a la multicolinealidad residual.
\end{itemize}

La evolución de los indicadores VIF mostró una reducción sustancial de
la multicolinealidad a lo largo de las iteraciones.

Los resultados indican que la Iteración 4 presenta la especificación más
parsimoniosa y con menores niveles generales de multicolinealidad.

En términos comparativos, el VIF máximo se redujo de 5166.89 a 12.01,
mientras que el VIF promedio disminuyó de 910.88 a 6.57. Esta reducción
representa una mejora sustancial en la estabilidad estadística potencial
del modelo.

Con base en los criterios de parsimonia, estabilidad estadística,
interpretabilidad económica y relevancia agrícola, se seleccionó la
siguiente especificación preliminar de variables exógenas para la
construcción del modelo SARIMAX:

\begin{itemize}
\tightlist
\item
  Precio\_prom\_CBOT\_mxpkilo\_maiz\_internacional
\item
  Costo\_transporte\_maiz\_mxpkilo\_prom\_nac
\item
  Temperatura\_celsius\_prom\_nacional
\item
  Lluvia\_mm\_prom\_nacional
\item
  Indice\_severidad\_ponderada\_sequia\_prom\_nacional
\item
  Produccion\_maizblanco\_kg\_prom\_nacional
\end{itemize}

Estas variables representan de manera conjunta factores de mercado
internacional, costos logísticos, condiciones climáticas y
disponibilidad de oferta agrícola, permitiendo capturar distintos
mecanismos potencialmente asociados a la formación del precio nacional
del maíz blanco.

Una vez definida la especificación preliminar de variables exógenas, el
siguiente paso consiste en analizar formalmente la estructura temporal
de la serie objetivo mediante técnicas de descomposición y pruebas de
estacionariedad, con el fin de determinar la especificación adecuada de
los modelos ARIMA, SARIMA y SARIMAX.

    \hypertarget{descomposiciuxf3n-temporal-de-la-serie}{%
\section{18. Descomposición Temporal de la
Serie}\label{descomposiciuxf3n-temporal-de-la-serie}}

Una serie temporal observada suele estar compuesta por distintos
elementos que interactúan simultáneamente.

Antes de aplicar modelos ARIMA, SARIMA o SARIMAX es necesario
identificar dichos componentes para comprender la dinámica interna de la
serie y determinar la estructura de modelado más apropiada.

La descomposición temporal permite separar una serie en tres componentes
principales:

\begin{itemize}
\tightlist
\item
  Tendencia.
\item
  Estacionalidad.
\item
  Componente residual.
\end{itemize}

Este procedimiento facilita la identificación de patrones persistentes
en el tiempo y proporciona evidencia sobre la posible necesidad de
incorporar componentes estacionales dentro del modelo predictivo.

Bajo un esquema aditivo, una serie temporal puede expresarse como:

\[
Y_t = T_t + S_t + R_t
\]

donde:

\begin{itemize}
\tightlist
\item
  \(Y_t\) representa el valor observado en el instante \(t\).
\item
  \(T_t\) representa la componente de tendencia.
\item
  \(S_t\) representa la componente estacional.
\item
  \(R_t\) representa la componente residual o aleatoria.
\end{itemize}

La tendencia captura movimientos de largo plazo.

La estacionalidad representa patrones que se repiten periódicamente.

La componente residual contiene la variabilidad que no puede explicarse
mediante tendencia o estacionalidad.

La identificación de estos componentes constituye un paso previo
indispensable para determinar la conveniencia de utilizar modelos ARIMA,
SARIMA o SARIMAX.

    \begin{tcolorbox}[breakable, size=fbox, boxrule=1pt, pad at break*=1mm,colback=cellbackground, colframe=cellborder]
\prompt{In}{incolor}{55}{\boxspacing}
\begin{Verbatim}[commandchars=\\\{\}]
\PY{k+kn}{from} \PY{n+nn}{statsmodels}\PY{n+nn}{.}\PY{n+nn}{tsa}\PY{n+nn}{.}\PY{n+nn}{seasonal} \PY{k+kn}{import} \PY{n}{seasonal\PYZus{}decompose}
\PY{k+kn}{import} \PY{n+nn}{matplotlib}\PY{n+nn}{.}\PY{n+nn}{pyplot} \PY{k}{as} \PY{n+nn}{plt}
\end{Verbatim}
\end{tcolorbox}

    \begin{tcolorbox}[breakable, size=fbox, boxrule=1pt, pad at break*=1mm,colback=cellbackground, colframe=cellborder]
\prompt{In}{incolor}{59}{\boxspacing}
\begin{Verbatim}[commandchars=\\\{\}]
\PY{c+c1}{\PYZsh{} Constrtruyendo la serie}
\PY{n}{serie\PYZus{}maiz} \PY{o}{=} \PY{p}{(}
    \PY{n}{pdf\PYZus{}completo}
    \PY{o}{.}\PY{n}{set\PYZus{}index}\PY{p}{(}\PY{l+s+s2}{\PYZdq{}}\PY{l+s+s2}{fecha}\PY{l+s+s2}{\PYZdq{}}\PY{p}{)}
    \PY{p}{[}\PY{l+s+s2}{\PYZdq{}}\PY{l+s+s2}{Precio\PYZus{}maizblanco\PYZus{}mxpkilo\PYZus{}prom\PYZus{}nacional}\PY{l+s+s2}{\PYZdq{}}\PY{p}{]}
\PY{p}{)}

\PY{n}{serie\PYZus{}maiz}\PY{o}{.}\PY{n}{head}\PY{p}{(}\PY{p}{)}
\end{Verbatim}
\end{tcolorbox}

            \begin{tcolorbox}[breakable, size=fbox, boxrule=.5pt, pad at break*=1mm, opacityfill=0]
\prompt{Out}{outcolor}{59}{\boxspacing}
\begin{Verbatim}[commandchars=\\\{\}]
fecha
2000-01-01    2.05
2000-02-01    2.06
2000-03-01    2.09
2000-04-01    2.12
2000-05-01    2.14
Name: Precio\_maizblanco\_mxpkilo\_prom\_nacional, dtype: float64
\end{Verbatim}
\end{tcolorbox}
        
    \begin{tcolorbox}[breakable, size=fbox, boxrule=1pt, pad at break*=1mm,colback=cellbackground, colframe=cellborder]
\prompt{In}{incolor}{60}{\boxspacing}
\begin{Verbatim}[commandchars=\\\{\}]
\PY{n}{descomposicion} \PY{o}{=} \PY{n}{seasonal\PYZus{}decompose}\PY{p}{(}
    \PY{n}{serie\PYZus{}maiz}\PY{p}{,}
    \PY{n}{model}\PY{o}{=}\PY{l+s+s2}{\PYZdq{}}\PY{l+s+s2}{additive}\PY{l+s+s2}{\PYZdq{}}\PY{p}{,}
    \PY{n}{period}\PY{o}{=}\PY{l+m+mi}{12}
\PY{p}{)}
\end{Verbatim}
\end{tcolorbox}

    \begin{tcolorbox}[breakable, size=fbox, boxrule=1pt, pad at break*=1mm,colback=cellbackground, colframe=cellborder]
\prompt{In}{incolor}{58}{\boxspacing}
\begin{Verbatim}[commandchars=\\\{\}]
\PY{n}{fig} \PY{o}{=} \PY{n}{descomposicion}\PY{o}{.}\PY{n}{plot}\PY{p}{(}\PY{p}{)}

\PY{n}{fig}\PY{o}{.}\PY{n}{set\PYZus{}size\PYZus{}inches}\PY{p}{(}\PY{l+m+mi}{14}\PY{p}{,}\PY{l+m+mi}{10}\PY{p}{)}

\PY{n}{plt}\PY{o}{.}\PY{n}{show}\PY{p}{(}\PY{p}{)}
\end{Verbatim}
\end{tcolorbox}

    \begin{center}
    \adjustimage{max size={0.9\linewidth}{0.9\paperheight}}{output_84_0.png}
    \end{center}
    { \hspace*{\fill} \\}
    
    \hypertarget{sobre-la-descomposiciuxf3n-temporal}{%
\subsection{18.1 Sobre la descomposición
temporal}\label{sobre-la-descomposiciuxf3n-temporal}}

La descomposición aditiva de la serie permitió identificar los tres
componentes fundamentales del comportamiento temporal del precio
nacional del maíz blanco.

\hypertarget{tendencia}{%
\subsubsection{Tendencia}\label{tendencia}}

La componente de tendencia muestra un crecimiento sostenido durante
prácticamente todo el periodo analizado.

Se observan tres etapas principales:

\begin{itemize}
\tightlist
\item
  Crecimiento moderado entre 2000 y 2010.
\item
  Estabilización relativa entre 2012 y 2020.
\item
  Incremento acelerado entre 2021 y 2024.
\end{itemize}

El fuerte incremento observado durante los años recientes coincide con
diversos factores económicos y logísticos documentados previamente,
incluyendo incrementos en los costos energéticos, volatilidad
internacional de commodities agrícolas y presiones inflacionarias.

La presencia de una tendencia claramente creciente constituye evidencia
preliminar de no estacionariedad.

\hypertarget{estacionalidad}{%
\subsubsection{Estacionalidad}\label{estacionalidad}}

La componente estacional muestra un patrón recurrente que se repite
aproximadamente cada doce meses.

La amplitud de la estacionalidad es relativamente estable a lo largo del
tiempo, oscilando alrededor de ±0.2 MXN por kilogramo.

Este comportamiento confirma la existencia de una estructura estacional
anual compatible con la naturaleza agrícola del mercado del maíz.

\hypertarget{componente-residual}{%
\subsubsection{Componente residual}\label{componente-residual}}

La componente residual se mantiene relativamente estable durante gran
parte del periodo analizado.

Sin embargo, se observan perturbaciones extraordinarias entre 2021 y
2024, donde los residuos presentan una dispersión considerablemente
mayor.

Este comportamiento sugiere la existencia de choques externos o eventos
estructurales que no son completamente explicados por la tendencia y la
estacionalidad estimadas.

En conjunto, los resultados indican que la serie presenta
simultáneamente:

\begin{itemize}
\tightlist
\item
  Tendencia.
\item
  Estacionalidad anual.
\item
  Eventos extraordinarios recientes.
\end{itemize}

Estas características sugieren que un modelo SARIMA o SARIMAX podría
capturar mejor la dinámica observada que un modelo ARIMA simple.

    \hypertarget{evaluaciuxf3n-preliminar-de-estacionariedad}{%
\subsection{18.2 Evaluación preliminar de
estacionariedad}\label{evaluaciuxf3n-preliminar-de-estacionariedad}}

La inspección visual de la tendencia sugiere que la media de la serie no
permanece constante a lo largo del tiempo.

Formalmente, una serie estacionaria requiere que:

\[
E(Y_t)=\mu
\]

permanezca constante para cualquier instante de tiempo.

Asimismo, la varianza debe satisfacer:

\[
Var(Y_t)=\sigma^2
\]

constante a lo largo de la serie.

La tendencia creciente observada en la descomposición constituye
evidencia preliminar de que la media de la serie podría variar con el
tiempo, incumpliendo uno de los principales supuestos de estacionariedad
requeridos por los modelos ARIMA.

Sin embargo, esta conclusión debe verificarse mediante una prueba
estadística formal.

Por esta razón, el siguiente paso consiste en aplicar la prueba de
Dickey-Fuller Aumentada (ADF), ampliamente utilizada para evaluar la
presencia de raíces unitarias y determinar si una serie temporal es
estacionaria.

    \hypertarget{prueba-de-dickey-fuller-aumentada-adf}{%
\section{19. Prueba de Dickey-Fuller Aumentada
(ADF)}\label{prueba-de-dickey-fuller-aumentada-adf}}

\hypertarget{objetivo}{%
\subsection{Objetivo}\label{objetivo}}

Los modelos ARIMA, SARIMA y SARIMAX requieren que la serie temporal sea
estacionaria o pueda transformarse en estacionaria mediante
diferenciación.

Por esta razón, antes de especificar cualquier modelo predictivo es
necesario evaluar formalmente si el precio nacional del maíz blanco
cumple con los supuestos de estacionariedad.

La prueba de Dickey-Fuller Aumentada (Augmented Dickey-Fuller, ADF)
permite determinar si una serie contiene una raíz unitaria,
característica asociada a procesos no estacionarios.

\begin{center}\rule{0.5\linewidth}{0.5pt}\end{center}

\hypertarget{fundamentaciuxf3n-matemuxe1tica}{%
\subsection{Fundamentación
matemática}\label{fundamentaciuxf3n-matemuxe1tica}}

Considérese el siguiente proceso autorregresivo de primer orden:

\[
Y_t = \rho Y_{t-1} + \varepsilon_t
\]

donde:

\begin{itemize}
\tightlist
\item
  \(Y_t\) representa el valor de la serie en el instante \(t\).
\item
  \(\rho\) representa el parámetro autorregresivo.
\item
  \(\varepsilon_t\) representa un término de error aleatorio.
\end{itemize}

La estacionariedad depende del valor de \(\rho\):

\begin{itemize}
\tightlist
\item
  Si \(|\rho| < 1\), la serie es estacionaria.
\item
  Si \(\rho = 1\), existe una raíz unitaria y la serie es no
  estacionaria.
\end{itemize}

La prueba ADF reformula la ecuación anterior como:

\[
\Delta Y_t = \alpha + \beta t + \gamma Y_{t-1}
+ \sum_{i=1}^{p}\delta_i \Delta Y_{t-i}
+ \varepsilon_t
\]

donde:

\begin{itemize}
\tightlist
\item
  \(\Delta Y_t\) representa la primera diferencia.
\item
  \(\alpha\) representa una constante.
\item
  \(\beta t\) representa una tendencia temporal opcional.
\item
  \(\gamma\) es el parámetro de interés.
\item
  Los términos rezagados permiten controlar autocorrelación residual.
\end{itemize}

\begin{center}\rule{0.5\linewidth}{0.5pt}\end{center}

\hypertarget{hipuxf3tesis-de-la-prueba}{%
\subsection{Hipótesis de la prueba}\label{hipuxf3tesis-de-la-prueba}}

\hypertarget{hipuxf3tesis-nula}{%
\subsubsection{Hipótesis nula}\label{hipuxf3tesis-nula}}

\[
H_0:\gamma = 0
\]

La serie contiene una raíz unitaria y no es estacionaria.

\hypertarget{hipuxf3tesis-alternativa}{%
\subsubsection{Hipótesis alternativa}\label{hipuxf3tesis-alternativa}}

\[
H_1:\gamma < 0
\]

La serie no contiene raíz unitaria y es estacionaria.

\begin{center}\rule{0.5\linewidth}{0.5pt}\end{center}

\hypertarget{regla-de-decisiuxf3n}{%
\subsection{Regla de decisión}\label{regla-de-decisiuxf3n}}

Se utilizará un nivel de significancia de:

\[
\alpha = 0.05
\]

Por lo tanto:

\begin{itemize}
\tightlist
\item
  Si el valor-p es menor que 0.05, se rechaza \(H_0\) y la serie se
  considera estacionaria.
\item
  Si el valor-p es mayor o igual que 0.05, no se rechaza \(H_0\) y la
  serie se considera no estacionaria.
\end{itemize}

Dado que la descomposición temporal mostró una tendencia creciente
significativa, se espera preliminarmente que la serie original no sea
estacionaria.

    \begin{tcolorbox}[breakable, size=fbox, boxrule=1pt, pad at break*=1mm,colback=cellbackground, colframe=cellborder]
\prompt{In}{incolor}{61}{\boxspacing}
\begin{Verbatim}[commandchars=\\\{\}]
\PY{k+kn}{from} \PY{n+nn}{statsmodels}\PY{n+nn}{.}\PY{n+nn}{tsa}\PY{n+nn}{.}\PY{n+nn}{stattools} \PY{k+kn}{import} \PY{n}{adfuller}

\PY{n}{resultado\PYZus{}adf} \PY{o}{=} \PY{n}{adfuller}\PY{p}{(}\PY{n}{serie\PYZus{}maiz}\PY{p}{)}

\PY{n+nb}{print}\PY{p}{(}\PY{l+s+s2}{\PYZdq{}}\PY{l+s+s2}{ADF Statistic:}\PY{l+s+s2}{\PYZdq{}}\PY{p}{,} \PY{n}{resultado\PYZus{}adf}\PY{p}{[}\PY{l+m+mi}{0}\PY{p}{]}\PY{p}{)}
\PY{n+nb}{print}\PY{p}{(}\PY{l+s+s2}{\PYZdq{}}\PY{l+s+s2}{p\PYZhy{}value:}\PY{l+s+s2}{\PYZdq{}}\PY{p}{,} \PY{n}{resultado\PYZus{}adf}\PY{p}{[}\PY{l+m+mi}{1}\PY{p}{]}\PY{p}{)}
\PY{n+nb}{print}\PY{p}{(}\PY{l+s+s2}{\PYZdq{}}\PY{l+s+s2}{Número de rezagos:}\PY{l+s+s2}{\PYZdq{}}\PY{p}{,} \PY{n}{resultado\PYZus{}adf}\PY{p}{[}\PY{l+m+mi}{2}\PY{p}{]}\PY{p}{)}
\PY{n+nb}{print}\PY{p}{(}\PY{l+s+s2}{\PYZdq{}}\PY{l+s+s2}{Número de observaciones:}\PY{l+s+s2}{\PYZdq{}}\PY{p}{,} \PY{n}{resultado\PYZus{}adf}\PY{p}{[}\PY{l+m+mi}{3}\PY{p}{]}\PY{p}{)}

\PY{n+nb}{print}\PY{p}{(}\PY{l+s+s2}{\PYZdq{}}\PY{l+s+se}{\PYZbs{}n}\PY{l+s+s2}{Valores críticos:}\PY{l+s+s2}{\PYZdq{}}\PY{p}{)}

\PY{k}{for} \PY{n}{key}\PY{p}{,} \PY{n}{value} \PY{o+ow}{in} \PY{n}{resultado\PYZus{}adf}\PY{p}{[}\PY{l+m+mi}{4}\PY{p}{]}\PY{o}{.}\PY{n}{items}\PY{p}{(}\PY{p}{)}\PY{p}{:}
    \PY{n+nb}{print}\PY{p}{(}\PY{l+s+sa}{f}\PY{l+s+s2}{\PYZdq{}}\PY{l+s+si}{\PYZob{}}\PY{n}{key}\PY{l+s+si}{\PYZcb{}}\PY{l+s+s2}{: }\PY{l+s+si}{\PYZob{}}\PY{n}{value}\PY{l+s+si}{\PYZcb{}}\PY{l+s+s2}{\PYZdq{}}\PY{p}{)}
\end{Verbatim}
\end{tcolorbox}

    \begin{Verbatim}[commandchars=\\\{\}]
ADF Statistic: -0.3500438243633834
p-value: 0.9181115006917588
Número de rezagos: 14
Número de observaciones: 302

Valores críticos:
1\%: -3.4521902441030963
5\%: -2.871158406898617
10\%: -2.5718948388228586
    \end{Verbatim}

    \hypertarget{sobre-la-adf}{%
\subsection{Sobre la ADF}\label{sobre-la-adf}}

La prueba Dickey-Fuller Aumentada (ADF) fue aplicada a la serie original
del precio nacional del maíz blanco con el objetivo de evaluar la
presencia de una raíz unitaria.

Los resultados obtenidos fueron:

\begin{longtable}[]{@{}lr@{}}
\toprule\noalign{}
Estadístico & Valor \\
\midrule\noalign{}
\endhead
\bottomrule\noalign{}
\endlastfoot
Estadístico ADF & -0.3500 \\
Valor-p & 0.9181 \\
Rezagos utilizados & 14 \\
Observaciones & 302 \\
\end{longtable}

Valores críticos:

\begin{longtable}[]{@{}lr@{}}
\toprule\noalign{}
Nivel de significancia & Valor crítico \\
\midrule\noalign{}
\endhead
\bottomrule\noalign{}
\endlastfoot
1\% & -3.4522 \\
5\% & -2.8712 \\
10\% & -2.5719 \\
\end{longtable}

\hypertarget{comparaciuxf3n-con-los-valores-cruxedticos}{%
\subsubsection{Comparación con los valores
críticos}\label{comparaciuxf3n-con-los-valores-cruxedticos}}

El estadístico de prueba obtenido fue:

\[
ADF=-0.3500
\]

Este valor es mayor que todos los valores críticos considerados:

\[
-0.3500 > -2.5719
\]

\[
-0.3500 > -2.8712
\]

\[
-0.3500 > -3.4522
\]

Por lo tanto, no existe evidencia suficiente para rechazar la hipótesis
nula de presencia de raíz unitaria.

\hypertarget{evaluaciuxf3n-mediante-valor-p}{%
\subsubsection{Evaluación mediante
valor-p}\label{evaluaciuxf3n-mediante-valor-p}}

Se obtuvo:

\[
p=0.9181
\]

Considerando un nivel de significancia de:

\[
\alpha = 0.05
\]

se cumple que:

\[
p > \alpha
\]

es decir:

\[
0.9181 > 0.05
\]

Por consiguiente, no se rechaza la hipótesis nula.

\hypertarget{conclusiuxf3n}{%
\subsubsection{Conclusión}\label{conclusiuxf3n}}

Los resultados indican que la serie original del precio nacional del
maíz blanco presenta raíz unitaria y, por tanto, no es estacionaria.

Este resultado es consistente con la tendencia creciente observada
previamente en la descomposición temporal.

La no estacionariedad observada implica que los modelos ARIMA, SARIMA y
SARIMAX no pueden aplicarse directamente sobre la serie original.

Por esta razón, el siguiente paso consiste en aplicar una primera
diferenciación con el objetivo de eliminar la tendencia y evaluar
nuevamente la estacionariedad de la serie.

La primera diferenciación permitirá estimar el parámetro:

\[
d
\]

correspondiente al componente integrado de los modelos ARIMA, SARIMA y
SARIMAX.

    \hypertarget{primera-diferenciaciuxf3n-de-la-serie}{%
\section{20. Primera diferenciación de la
serie}\label{primera-diferenciaciuxf3n-de-la-serie}}

Una estrategia ampliamente utilizada para transformar series no
estacionarias consiste en calcular diferencias consecutivas entre
observaciones.

La primera diferencia se define como:

\[
\Delta Y_t = Y_t - Y_{t-1}
\]

donde:

\begin{itemize}
\tightlist
\item
  \(Y_t\) es el valor observado en el periodo actual.
\item
  \(Y_{t-1}\) es el valor observado en el periodo anterior.
\end{itemize}

Esta transformación elimina tendencias lineales y reduce la dependencia
de largo plazo presente en la serie.

Si la serie diferenciada resulta estacionaria, entonces el orden de
integración será:

\[
d=1
\]

En caso contrario, será necesario aplicar diferenciaciones adicionales.

    \begin{tcolorbox}[breakable, size=fbox, boxrule=1pt, pad at break*=1mm,colback=cellbackground, colframe=cellborder]
\prompt{In}{incolor}{62}{\boxspacing}
\begin{Verbatim}[commandchars=\\\{\}]
\PY{n}{serie\PYZus{}maiz\PYZus{}diff1} \PY{o}{=} \PY{n}{serie\PYZus{}maiz}\PY{o}{.}\PY{n}{diff}\PY{p}{(}\PY{p}{)}\PY{o}{.}\PY{n}{dropna}\PY{p}{(}\PY{p}{)}

\PY{n}{serie\PYZus{}maiz\PYZus{}diff1}\PY{o}{.}\PY{n}{head}\PY{p}{(}\PY{p}{)}
\end{Verbatim}
\end{tcolorbox}

            \begin{tcolorbox}[breakable, size=fbox, boxrule=.5pt, pad at break*=1mm, opacityfill=0]
\prompt{Out}{outcolor}{62}{\boxspacing}
\begin{Verbatim}[commandchars=\\\{\}]
fecha
2000-02-01    0.01
2000-03-01    0.03
2000-04-01    0.03
2000-05-01    0.02
2000-06-01   -0.03
Name: Precio\_maizblanco\_mxpkilo\_prom\_nacional, dtype: float64
\end{Verbatim}
\end{tcolorbox}
        
    \begin{tcolorbox}[breakable, size=fbox, boxrule=1pt, pad at break*=1mm,colback=cellbackground, colframe=cellborder]
\prompt{In}{incolor}{63}{\boxspacing}
\begin{Verbatim}[commandchars=\\\{\}]
\PY{n}{plt}\PY{o}{.}\PY{n}{figure}\PY{p}{(}\PY{n}{figsize}\PY{o}{=}\PY{p}{(}\PY{l+m+mi}{14}\PY{p}{,}\PY{l+m+mi}{5}\PY{p}{)}\PY{p}{)}

\PY{n}{plt}\PY{o}{.}\PY{n}{plot}\PY{p}{(}
    \PY{n}{serie\PYZus{}maiz\PYZus{}diff1}\PY{p}{,}
    \PY{n}{linewidth}\PY{o}{=}\PY{l+m+mi}{1}
\PY{p}{)}

\PY{n}{plt}\PY{o}{.}\PY{n}{title}\PY{p}{(}\PY{l+s+s2}{\PYZdq{}}\PY{l+s+s2}{Primera Diferenciación del Precio Nacional del Maíz}\PY{l+s+s2}{\PYZdq{}}\PY{p}{)}
\PY{n}{plt}\PY{o}{.}\PY{n}{xlabel}\PY{p}{(}\PY{l+s+s2}{\PYZdq{}}\PY{l+s+s2}{Fecha}\PY{l+s+s2}{\PYZdq{}}\PY{p}{)}
\PY{n}{plt}\PY{o}{.}\PY{n}{ylabel}\PY{p}{(}\PY{l+s+s2}{\PYZdq{}}\PY{l+s+s2}{Δ Precio}\PY{l+s+s2}{\PYZdq{}}\PY{p}{)}

\PY{n}{plt}\PY{o}{.}\PY{n}{grid}\PY{p}{(}\PY{k+kc}{True}\PY{p}{)}

\PY{n}{plt}\PY{o}{.}\PY{n}{show}\PY{p}{(}\PY{p}{)}
\end{Verbatim}
\end{tcolorbox}

    \begin{center}
    \adjustimage{max size={0.9\linewidth}{0.9\paperheight}}{output_92_0.png}
    \end{center}
    { \hspace*{\fill} \\}
    
    \begin{tcolorbox}[breakable, size=fbox, boxrule=1pt, pad at break*=1mm,colback=cellbackground, colframe=cellborder]
\prompt{In}{incolor}{64}{\boxspacing}
\begin{Verbatim}[commandchars=\\\{\}]
\PY{n}{resultado\PYZus{}adf\PYZus{}diff1} \PY{o}{=} \PY{n}{adfuller}\PY{p}{(}\PY{n}{serie\PYZus{}maiz\PYZus{}diff1}\PY{p}{)}

\PY{n+nb}{print}\PY{p}{(}\PY{l+s+s2}{\PYZdq{}}\PY{l+s+s2}{ADF Statistic:}\PY{l+s+s2}{\PYZdq{}}\PY{p}{,} \PY{n}{resultado\PYZus{}adf\PYZus{}diff1}\PY{p}{[}\PY{l+m+mi}{0}\PY{p}{]}\PY{p}{)}
\PY{n+nb}{print}\PY{p}{(}\PY{l+s+s2}{\PYZdq{}}\PY{l+s+s2}{p\PYZhy{}value:}\PY{l+s+s2}{\PYZdq{}}\PY{p}{,} \PY{n}{resultado\PYZus{}adf\PYZus{}diff1}\PY{p}{[}\PY{l+m+mi}{1}\PY{p}{]}\PY{p}{)}

\PY{n+nb}{print}\PY{p}{(}\PY{l+s+s2}{\PYZdq{}}\PY{l+s+se}{\PYZbs{}n}\PY{l+s+s2}{Valores críticos:}\PY{l+s+s2}{\PYZdq{}}\PY{p}{)}

\PY{k}{for} \PY{n}{key}\PY{p}{,} \PY{n}{value} \PY{o+ow}{in} \PY{n}{resultado\PYZus{}adf\PYZus{}diff1}\PY{p}{[}\PY{l+m+mi}{4}\PY{p}{]}\PY{o}{.}\PY{n}{items}\PY{p}{(}\PY{p}{)}\PY{p}{:}
    \PY{n+nb}{print}\PY{p}{(}\PY{l+s+sa}{f}\PY{l+s+s2}{\PYZdq{}}\PY{l+s+si}{\PYZob{}}\PY{n}{key}\PY{l+s+si}{\PYZcb{}}\PY{l+s+s2}{: }\PY{l+s+si}{\PYZob{}}\PY{n}{value}\PY{l+s+si}{\PYZcb{}}\PY{l+s+s2}{\PYZdq{}}\PY{p}{)}
\end{Verbatim}
\end{tcolorbox}

    \begin{Verbatim}[commandchars=\\\{\}]
ADF Statistic: -2.7451657845861224
p-value: 0.06654336032233617

Valores críticos:
1\%: -3.4521902441030963
5\%: -2.871158406898617
10\%: -2.5718948388228586
    \end{Verbatim}

    \begin{tcolorbox}[breakable, size=fbox, boxrule=1pt, pad at break*=1mm,colback=cellbackground, colframe=cellborder]
\prompt{In}{incolor}{65}{\boxspacing}
\begin{Verbatim}[commandchars=\\\{\}]
\PY{n}{rolling\PYZus{}mean} \PY{o}{=} \PY{n}{serie\PYZus{}maiz\PYZus{}diff1}\PY{o}{.}\PY{n}{rolling}\PY{p}{(}\PY{l+m+mi}{12}\PY{p}{)}\PY{o}{.}\PY{n}{mean}\PY{p}{(}\PY{p}{)}
\PY{n}{rolling\PYZus{}std} \PY{o}{=} \PY{n}{serie\PYZus{}maiz\PYZus{}diff1}\PY{o}{.}\PY{n}{rolling}\PY{p}{(}\PY{l+m+mi}{12}\PY{p}{)}\PY{o}{.}\PY{n}{std}\PY{p}{(}\PY{p}{)}

\PY{n}{plt}\PY{o}{.}\PY{n}{figure}\PY{p}{(}\PY{n}{figsize}\PY{o}{=}\PY{p}{(}\PY{l+m+mi}{14}\PY{p}{,}\PY{l+m+mi}{5}\PY{p}{)}\PY{p}{)}

\PY{n}{plt}\PY{o}{.}\PY{n}{plot}\PY{p}{(}\PY{n}{serie\PYZus{}maiz\PYZus{}diff1}\PY{p}{,} \PY{n}{label}\PY{o}{=}\PY{l+s+s2}{\PYZdq{}}\PY{l+s+s2}{Serie diferenciada}\PY{l+s+s2}{\PYZdq{}}\PY{p}{)}
\PY{n}{plt}\PY{o}{.}\PY{n}{plot}\PY{p}{(}\PY{n}{rolling\PYZus{}mean}\PY{p}{,} \PY{n}{label}\PY{o}{=}\PY{l+s+s2}{\PYZdq{}}\PY{l+s+s2}{Media móvil (12)}\PY{l+s+s2}{\PYZdq{}}\PY{p}{)}
\PY{n}{plt}\PY{o}{.}\PY{n}{plot}\PY{p}{(}\PY{n}{rolling\PYZus{}std}\PY{p}{,} \PY{n}{label}\PY{o}{=}\PY{l+s+s2}{\PYZdq{}}\PY{l+s+s2}{Desv. estándar móvil (12)}\PY{l+s+s2}{\PYZdq{}}\PY{p}{)}

\PY{n}{plt}\PY{o}{.}\PY{n}{legend}\PY{p}{(}\PY{p}{)}
\PY{n}{plt}\PY{o}{.}\PY{n}{grid}\PY{p}{(}\PY{k+kc}{True}\PY{p}{)}

\PY{n}{plt}\PY{o}{.}\PY{n}{show}\PY{p}{(}\PY{p}{)}
\end{Verbatim}
\end{tcolorbox}

    \begin{center}
    \adjustimage{max size={0.9\linewidth}{0.9\paperheight}}{output_94_0.png}
    \end{center}
    { \hspace*{\fill} \\}
    
    \hypertarget{primera-diferenciaciuxf3n-y-reevaluaciuxf3n-de-la-estacionariedad}{%
\section{20. Primera Diferenciación y Reevaluación de la
Estacionariedad}\label{primera-diferenciaciuxf3n-y-reevaluaciuxf3n-de-la-estacionariedad}}

\hypertarget{objetivo}{%
\subsection{Objetivo}\label{objetivo}}

La prueba Dickey-Fuller aplicada a la serie original indicó la presencia
de una raíz unitaria, por lo que fue necesario transformar la serie
mediante diferenciación.

El objetivo de esta sección consiste en evaluar si una primera
diferenciación es suficiente para eliminar la tendencia observada y
aproximar la serie a un comportamiento estacionatemática

La primera diferencia de una serie temporal se define como:

\[
\Delta Y_t = Y_t - Y_{t-1}
\]

donde:

\begin{itemize}
\tightlist
\item
  \(Y_t\) representa el valor observado en el periodo actual.
\item
  \(Y_{t-1}\) representa el valor observado en el periodo inmediatamente
  anterior.
\end{itemize}

Esta transformación busca eliminar componentes de tendencia y
estabilizar la media de la serie.

Si la serie diferenciada resulta estacionaria, entonces el orden de
integración requerido por el modelo será:

\[
d = 1
\]

En caso contrario, sería necesario considerar diferenciaciones
adicionales.

\begin{center}\rule{0.5\linewidth}{0.5pt}\end{center}

\hypertarget{anuxe1lisis-visual-de-la-serie-diferenciada}{%
\subsection{Análisis visual de la serie
diferenciada}\label{anuxe1lisis-visual-de-la-serie-diferenciada}}

La Figura anterior muestra la evolución temporal de la primera
diferencia del precio nacional del maíz blanco.

A diferencia de la serie original, la tendencia creciente observada
entre 2000 y 2026 desaparece casi por completo, generando fluctuaciones
alrededor de un valor medio cercano a cero.

Este comportamiento constituye una primera evidencia de que la
diferenciación logró reducir una parte importante de la no
estacionariedad presente en la serie original.

No obstante, se observan episodios de elevada volatilidad entre 2021 y
2024, caracterizados por movimientos abruptos tanto positivos como
negativos.

Estos eventos coinciden temporalmente con perturbaciones observadas en
los mercados internacionales de commodities agrícolas, costos logísticos
y cadenas globales de suministro.

\begin{center}\rule{0.5\linewidth}{0.5pt}\end{center}

\hypertarget{anuxe1lisis-de-media-y-varianza-muxf3vil}{%
\subsection{Análisis de media y varianza
móvil}\label{anuxe1lisis-de-media-y-varianza-muxf3vil}}

Para evaluar visualmente la estabilidad de la serie se calcularon una
media móvil y una desviación estándar móvil utilizando ventanas de doce
observaciones.

Los resultados muestran que:

\begin{itemize}
\tightlist
\item
  La media móvil oscila alrededor de cero durante la mayor parte del
  periodo analizado.
\item
  La desviación estándar permanece relativamente estable entre 2000 y
  2020.
\item
  Entre 2021 y 2024 se observa un incremento considerable de la
  volatilidad.
\item
  Hacia el final de la serie la volatilidad muestra señales de
  normalización.
\end{itemize}

Estos resultados sugieren que la diferenciación logró estabilizar la
media de la serie, aunque persisten episodios temporales de
heterocedasticidad asociados a choques externos recientes.

\begin{center}\rule{0.5\linewidth}{0.5pt}\end{center}

\hypertarget{prueba-dickey-fuller-sobre-la-serie-diferenciada}{%
\subsection{Prueba Dickey-Fuller sobre la serie
diferenciada}\label{prueba-dickey-fuller-sobre-la-serie-diferenciada}}

Con el objetivo de validar formalmente la estacionariedad de la serie
diferenciada, se aplicó nuevamente la prueba Dickey-Fuller Aumentada.

\hypertarget{resultados}{%
\subsubsection{Resultados}\label{resultados}}

\begin{longtable}[]{@{}lr@{}}
\toprule\noalign{}
Estadístico & Valor \\
\midrule\noalign{}
\endhead
\bottomrule\noalign{}
\endlastfoot
Estadístico ADF & -2.7452 \\
Valor-p & 0.0665 \\
\end{longtable}

\hypertarget{valores-cruxedticos}{%
\subsubsection{Valores críticos}\label{valores-cruxedticos}}

Nivel \textbar{} Valor crítico \textbar{}

\textbar--------\textbar-------------3.4522 \textbar{} \textbar{} 5\%
\textbar{} -2.8712 \textbar{} \textbar{} 10\% \textbar{} -2.5719
\textbar{}

\begin{center}\rule{0.5\linewidth}{0.5pt}\end{center}

\hypertarget{interpretaciuxf3n-estaduxedstica}{%
\subsection{Interpretación
estadística}\label{interpretaciuxf3n-estaduxedstica}}

La prueba Dickey-Fuller plantea las siguientes hipótesis:

\hypertarget{hipuxf3tesis-nula}{%
\subsubsection{Hipótesis nula}\label{hipuxf3tesis-nula}}

\[
H_0
\]

La serie posee una raíz unitaria y no es estacionaria.

\hypertarget{hipuxf3tesis-alternativa}{%
\subsubsection{Hipótesis alternativa}\label{hipuxf3tesis-alternativa}}

\[
H_1
\]

La serie no posee raíz unitaria y es estacionaria.

El estadístico obtenido fue:

\[
ADF=-2.7452
\]

Al compararlo con los valores críticos se observa que:

\[
-2.7452 > -2.8712
\]

por lo que no se rechaza la hipótesis nula al nivel de significancia del
5\%.

Sin embargo:

\[
-2.7452 < -2.5719
\]

por lo que sí se rechaza la hipótesis nula al nivel de significancia del
10\%.

De manera consistente, el valor-p obtenido fue:

\[
p=0.0665
\]

valor que se encuentra por encima del umbral convencional de:

\[
\alpha=0.05
\]

pero por debajo de:

\[
\alpha=0.10
\]

\begin{center}\rule{0.5\linewidth}{0.5pt}\end{center}

\hypertarget{conclusiones}{%
\subsection{Conclusiones}\label{conclusiones}}

Los resultados obtenidos indican que la primera diferenciación redujo
sustancialmente la no estacionariedad observada en la serie original.

La evidencia visual muestra una estabilización importante de la media,
mientras que la prueba ADF arroja resultados cercanos al umbral
convencional de decisión.

Dado que la evidencia estadística es ambigua y que una segunda
diferenciación podría eliminar información relevante de la serie, no se
procederá inmediatamente a incrementar el orden de integración.

Antes de decidir si el parámetro de diferenciación definitivo debe ser:

\[
d=1
\]

o

\[
d=2
\]

se analizarán las funciones de autocorrelación (ACF) y autocorrelación
parcial (PACF), las cuales permitirán evaluar la estructura de
dependencia temporal remanente y orientar la especificación de los
modelos ARIMA, SARIMA y SARIMAX.

    \hypertarget{funciuxf3n-de-autocorrelaciuxf3n-acf}{%
\section{21. Función de Autocorrelación
(ACF)}\label{funciuxf3n-de-autocorrelaciuxf3n-acf}}

Una vez aplicada la primera diferenciación y evaluada preliminarmente la
estacionariedad de la serie, es necesario analizar la dependencia
temporal existente entre las observaciones.

La Función de Autocorrelación (Autocorrelation Function, ACF) permite
cuantificar la relación lineal entre una observación y sus valores
rezagados en distintos periodos de tiempo.

Este análisis constituye una herramienta fundamental para la
identificación de modelos ARIMA, SARIMA y SARIMAX, ya que proporciona
evidencia sobre la posible presencia de componentes de medias móviles y
estructuras estacionales.

La autocorrelación de orden \(k\) se define como:

\[
\rho_k =
\frac{
Cov(Y_t,Y_{t-k})
}
{
\sqrt{Var(Y_t)Var(Y_{t-k})}
}
\]

donde:

\begin{itemize}
\tightlist
\item
  \(\rho_k\) representa la autocorrelación para el rezago \(k\).
\item
  \(Cov(Y_t,Y_{t-k})\) representa la covarianza entre la observación
  actual y la observación rezagada.
\item
  \(Var(Y_t)\) representa la varianza de la serie.
\end{itemize}

Los coeficientes de autocorrelación se encuentran acotados por:

\[
-1 \le \rho_k \le 1
\]

Valores cercanos a:

\begin{itemize}
\tightlist
\item
  +1 indican dependencia positiva fuerte.
\item
  -1 indican dependencia negativa fuerte.
\item
  0 indican ausencia de relación lineal.
\end{itemize}

\begin{center}\rule{0.5\linewidth}{0.5pt}\end{center}

\hypertarget{relaciuxf3n-con-los-modelos-arima}{%
\subsection{Relación con los modelos
ARIMA}\label{relaciuxf3n-con-los-modelos-arima}}

La ACF resulta especialmente útil para identificar componentes de medias
móviles:

\[
MA(q)
\]

De forma general:

\begin{itemize}
\tightlist
\item
  Un corte abrupto de la ACF sugiere la presencia de un componente
  MA(q).
\item
  Una disminución gradual suele asociarse con procesos autorregresivos.
\item
  Picos repetitivos en rezagos estacionales pueden indicar la presencia
  de componentes SARIMA.
\end{itemize}

Dado que la serie analizada es mensual, se prestará especial atención a
los rezagos:

\[
12,\;24,\;36
\]

los cuales representan ciclos anuales completos.

    \begin{tcolorbox}[breakable, size=fbox, boxrule=1pt, pad at break*=1mm,colback=cellbackground, colframe=cellborder]
\prompt{In}{incolor}{66}{\boxspacing}
\begin{Verbatim}[commandchars=\\\{\}]
\PY{k+kn}{from} \PY{n+nn}{statsmodels}\PY{n+nn}{.}\PY{n+nn}{graphics}\PY{n+nn}{.}\PY{n+nn}{tsaplots} \PY{k+kn}{import} \PY{n}{plot\PYZus{}acf}
\PY{k+kn}{import} \PY{n+nn}{matplotlib}\PY{n+nn}{.}\PY{n+nn}{pyplot} \PY{k}{as} \PY{n+nn}{plt}

\PY{n}{plt}\PY{o}{.}\PY{n}{figure}\PY{p}{(}\PY{n}{figsize}\PY{o}{=}\PY{p}{(}\PY{l+m+mi}{12}\PY{p}{,}\PY{l+m+mi}{6}\PY{p}{)}\PY{p}{)}

\PY{n}{plot\PYZus{}acf}\PY{p}{(}
    \PY{n}{serie\PYZus{}maiz\PYZus{}diff1}\PY{p}{,}
    \PY{n}{lags}\PY{o}{=}\PY{l+m+mi}{48}
\PY{p}{)}

\PY{n}{plt}\PY{o}{.}\PY{n}{show}\PY{p}{(}\PY{p}{)}
\end{Verbatim}
\end{tcolorbox}

    
    \begin{Verbatim}[commandchars=\\\{\}]
<Figure size 1200x600 with 0 Axes>
    \end{Verbatim}

    
    \begin{center}
    \adjustimage{max size={0.9\linewidth}{0.9\paperheight}}{output_97_1.png}
    \end{center}
    { \hspace*{\fill} \\}
    
    \hypertarget{sobre-la-acf}{%
\subsection{Sobre la ACF}\label{sobre-la-acf}}

La función de autocorrelación muestra que la serie diferenciada conserva
una estructura temporal significativa.

En los primeros rezagos se observan autocorrelaciones positivas
estadísticamente significativas, indicando persistencia de corto plazo
entre observaciones consecutivas.

Adicionalmente, se identifican picos pronunciados en los rezagos:

\[
12,\;24,\;36,\;48
\]

los cuales corresponden a múltiplos de un ciclo anual.

La presencia de estos picos constituye evidencia de estacionalidad
remanente incluso después de aplicar una primera diferenciación.

Asimismo, se observan autocorrelaciones negativas en rezagos intermedios
cercanos a:

\[
6,\;18,\;30,\;42
\]

patrón característico de procesos con comportamiento cíclico anual.

En conjunto, los resultados sugieren que la dependencia temporal de la
serie no ha sido completamente eliminada mediante la diferenciación
ordinaria.

Por lo tanto, existe evidencia preliminar para considerar modelos con
componentes estacionales, particularmente SARIMA y SARIMAX.

La identificación definitiva de los parámetros autorregresivos y de
medias móviles requerirá complementar este análisis mediante la Función
de Autocorrelación Parcial (PACF).

    \hypertarget{funciuxf3n-de-autocorrelaciuxf3n-parcial-pacf}{%
\section{22. Función de Autocorrelación Parcial
(PACF)}\label{funciuxf3n-de-autocorrelaciuxf3n-parcial-pacf}}

La Función de Autocorrelación Parcial (Partial Autocorrelation Function,
PACF) permite medir la relación directa entre una observación y sus
rezagos, eliminando el efecto indirecto de los rezagos intermedios.

Mientras que la ACF cuantifica la correlación total existente entre
observaciones separadas por distintos periodos, la PACF estima
únicamente la contribución individual de cada rezago.

Esta característica convierte a la PACF en una herramienta fundamental
para identificar componentes autorregresivos dentro de los modelos
ARIMA, SARIMA y SARIMAX.

La autocorrelación parcial de orden \(k\) representa la correlación
existente entre:

\[
Y_t
\]

y

\[
Y_{t-k}
\]

una vez eliminado el efecto de todos los rezagos intermedios:

\[
Y_{t-1},Y_{t-2},...,Y_{t-(k-1)}
\]

Formalmente:

\[
PACF(k)
=
Corr
\left(
Y_t,
Y_{t-k}
\mid
Y_{t-1},
Y_{t-2},
\dots,
Y_{t-k+1}
\right)
\]

La PACF permite identificar el orden del componente autorregresivo:

\[
AR(p)
\]

que formará parte del modelo.

De manera general:

\begin{itemize}
\tightlist
\item
  Un corte abrupto en la PACF sugiere la presencia de un componente
  autorregresivo.
\item
  Una disminución gradual suele indicar la presencia de componentes de
  medias móviles.
\item
  Picos significativos en rezagos estacionales pueden indicar
  componentes SAR.
\end{itemize}

Por esta razón, la interpretación conjunta de la ACF y la PACF
constituye el procedimiento clásico para la identificación preliminar de
modelos ARIMA y SARIMA.

\begin{center}\rule{0.5\linewidth}{0.5pt}\end{center}

\hypertarget{objetivo-especuxedfico}{%
\subsection{Objetivo específico}\label{objetivo-especuxedfico}}

En esta etapa se busca identificar evidencia empírica que permita
estimar los parámetros:

\[
p
\]

y potencialmente:

\[
P
\]

de los modelos SARIMA y SARIMAX.

Particular atención se prestará a los rezagos:

\[
1,\;2,\;3
\]

para componentes de corto plazo y a los rezagos:

\[
12,\;24,\;36
\]

para posibles componentes estacionales anuales.

    \begin{tcolorbox}[breakable, size=fbox, boxrule=1pt, pad at break*=1mm,colback=cellbackground, colframe=cellborder]
\prompt{In}{incolor}{67}{\boxspacing}
\begin{Verbatim}[commandchars=\\\{\}]
\PY{k+kn}{from} \PY{n+nn}{statsmodels}\PY{n+nn}{.}\PY{n+nn}{graphics}\PY{n+nn}{.}\PY{n+nn}{tsaplots} \PY{k+kn}{import} \PY{n}{plot\PYZus{}pacf}
\PY{k+kn}{import} \PY{n+nn}{matplotlib}\PY{n+nn}{.}\PY{n+nn}{pyplot} \PY{k}{as} \PY{n+nn}{plt}

\PY{n}{plt}\PY{o}{.}\PY{n}{figure}\PY{p}{(}\PY{n}{figsize}\PY{o}{=}\PY{p}{(}\PY{l+m+mi}{12}\PY{p}{,}\PY{l+m+mi}{6}\PY{p}{)}\PY{p}{)}

\PY{n}{plot\PYZus{}pacf}\PY{p}{(}
    \PY{n}{serie\PYZus{}maiz\PYZus{}diff1}\PY{p}{,}
    \PY{n}{lags}\PY{o}{=}\PY{l+m+mi}{48}\PY{p}{,}
    \PY{n}{method}\PY{o}{=}\PY{l+s+s2}{\PYZdq{}}\PY{l+s+s2}{ywm}\PY{l+s+s2}{\PYZdq{}}
\PY{p}{)}

\PY{n}{plt}\PY{o}{.}\PY{n}{show}\PY{p}{(}\PY{p}{)}
\end{Verbatim}
\end{tcolorbox}

    
    \begin{Verbatim}[commandchars=\\\{\}]
<Figure size 1200x600 with 0 Axes>
    \end{Verbatim}

    
    \begin{center}
    \adjustimage{max size={0.9\linewidth}{0.9\paperheight}}{output_100_1.png}
    \end{center}
    { \hspace*{\fill} \\}
    
    \hypertarget{sobre-la-pacf}{%
\subsection{Sobre la PACF}\label{sobre-la-pacf}}

La función de autocorrelación parcial muestra evidencia de dependencia
temporal significativa tanto en el corto plazo como en la dimensión
estacional.

Se observa un coeficiente significativo en el rezago:

\[
1
\]

lo que indica que el comportamiento del precio en el periodo
inmediatamente anterior continúa influyendo sobre el periodo actual
incluso después de controlar el efecto de los rezagos intermedios.

Este resultado constituye evidencia preliminar para la presencia de un
componente:

\[
AR(1)
\]

dentro del modelo.

Adicionalmente, destaca un coeficiente altamente significativo en el
rezago:

\[
12
\]

el cual corresponde a un ciclo anual completo.

La magnitud de este coeficiente sugiere que existe una dependencia
directa entre observaciones separadas por doce meses, fenómeno
consistente con la naturaleza agrícola y estacional del mercado del maíz
blanco.

La evidencia observada en la PACF coincide con los resultados
previamente obtenidos en la ACF y en la descomposición temporal de la
serie.

En conjunto, los resultados sugieren la presencia de:

\begin{itemize}
\tightlist
\item
  Un componente autorregresivo de corto plazo.
\item
  Un componente autorregresivo estacional anual.
\end{itemize}

Por lo tanto, existe evidencia preliminar para considerar modelos SARIMA
y SARIMAX con parámetros estacionales asociados a un periodo:

\[
s = 12
\]

correspondiente a observaciones mensuales.

    \hypertarget{identificaciuxf3n-preliminar-de-modelos-candidatos}{%
\section{23. Identificación Preliminar de Modelos
Candidatos}\label{identificaciuxf3n-preliminar-de-modelos-candidatos}}

Las funciones de autocorrelación (ACF) y autocorrelación parcial (PACF)
permiten identificar configuraciones iniciales plausibles para los
parámetros de los modelos ARIMA y SARIMA.

El propósito de esta sección consiste en proponer una serie de modelos
candidatos fundamentados en la evidencia estadística obtenida
previamente, evitando la selección arbitraria de parámetros.

\textbf{Los análisis realizados permiten establecer las siguientes
conclusiones:}

\hypertarget{estacionariedad}{%
\subsubsection{Estacionariedad}\label{estacionariedad}}

La serie original presentó evidencia de raíz unitaria.

La primera diferenciación redujo considerablemente la tendencia
observada, por lo que se adopta preliminarmente:

\[
d = 1
\]

\begin{center}\rule{0.5\linewidth}{0.5pt}\end{center}

\hypertarget{estacionalidad}{%
\subsubsection{Estacionalidad}\label{estacionalidad}}

La descomposición temporal mostró un patrón estacional recurrente de
periodicidad anual.

Adicionalmente, la ACF presentó picos significativos en:

\[
12,\;24,\;36,\;48
\]

Por lo tanto, se considera una periodicidad estacional:

\[
s = 12
\]

\begin{center}\rule{0.5\linewidth}{0.5pt}\end{center}

\hypertarget{componente-autorregresivo}{%
\subsubsection{Componente
autorregresivo}\label{componente-autorregresivo}}

La PACF mostró un coeficiente significativo en:

\[
lag = 1
\]

lo que constituye evidencia para:

\[
p = 1
\]

\begin{center}\rule{0.5\linewidth}{0.5pt}\end{center}

\hypertarget{componente-autorregresivo-estacional}{%
\subsubsection{Componente autorregresivo
estacional}\label{componente-autorregresivo-estacional}}

La PACF presentó un pico pronunciado en:

\[
lag = 12
\]

sugiriendo la presencia de:

\[
P = 1
\]

\begin{center}\rule{0.5\linewidth}{0.5pt}\end{center}

\hypertarget{componente-de-medias-muxf3viles}{%
\subsubsection{Componente de medias
móviles}\label{componente-de-medias-muxf3viles}}

La ACF mostró autocorrelaciones significativas tanto en rezagos de corto
plazo como en rezagos estacionales.

Esto sugiere evaluar configuraciones con:

\[
q = 1
\]

y

\[
Q = 1
\]

\begin{center}\rule{0.5\linewidth}{0.5pt}\end{center}

\hypertarget{modelos-candidatos}{%
\subsection{Modelos candidatos}\label{modelos-candidatos}}

Con base en la evidencia obtenida, se propone evaluar los siguientes
modelos:

\begin{longtable}[]{@{}ll@{}}
\toprule\noalign{}
Modelo & Especificación \\
\midrule\noalign{}
\endhead
\bottomrule\noalign{}
\endlastfoot
ARIMA 1 & ARIMA(1,1,0) \\
ARIMA 2 & ARIMA(0,1,1) \\
ARIMA 3 & ARIMA(1,1,1) \\
SARIMA 1 & SARIMA(1,1,0)(1,0,0,12) \\
SARIMA 2 & SARIMA(1,1,1)(1,0,1,12) \\
\end{longtable}

\begin{center}\rule{0.5\linewidth}{0.5pt}\end{center}

\hypertarget{criterios-de-selecciuxf3n}{%
\subsection{Criterios de selección}\label{criterios-de-selecciuxf3n}}

Los modelos serán comparados utilizando indicadores de ajuste y
capacidad predictiva.

Particularmente se analizarán:

\hypertarget{aic}{%
\subsubsection{AIC}\label{aic}}

\[
AIC = 2k - 2\ln(L)
\]

donde:

\begin{itemize}
\tightlist
\item
  \(k\) es el número de parámetros estimados.
\item
  \(L\) es la función de verosimilitud.
\end{itemize}

\begin{center}\rule{0.5\linewidth}{0.5pt}\end{center}

\hypertarget{bic}{%
\subsubsection{BIC}\label{bic}}

\[
BIC = \ln(n)k - 2\ln(L)
\]

donde:

\begin{itemize}
\tightlist
\item
  \(n\) representa el tamaño de la muestra.
\end{itemize}

\begin{center}\rule{0.5\linewidth}{0.5pt}\end{center}

\hypertarget{error-cuadruxe1tico-medio-rmse}{%
\subsubsection{Error Cuadrático Medio
(RMSE)}\label{error-cuadruxe1tico-medio-rmse}}

\[
RMSE =
\sqrt{
\frac{1}{n}
\sum_{i=1}^{n}
(y_i-\hat{y_i})^2
}
\]

\begin{center}\rule{0.5\linewidth}{0.5pt}\end{center}

\hypertarget{error-absoluto-medio-mae}{%
\subsubsection{Error Absoluto Medio
(MAE)}\label{error-absoluto-medio-mae}}

\[
MAE=
\frac{1}{n}
\sum_{i=1}^{n}
|y_i-\hat{y_i}|
\]

\begin{center}\rule{0.5\linewidth}{0.5pt}\end{center}

\hypertarget{error-porcentual-absoluto-medio-mape}{%
\subsubsection{Error Porcentual Absoluto Medio
(MAPE)}\label{error-porcentual-absoluto-medio-mape}}

\[
MAPE=
\frac{100}{n}
\sum_{i=1}^{n}
\left|
\frac{y_i-\hat{y_i}}
{y_i}
\right|
\]

Una vez definidos los modelos candidatos, se procederá a estimar cada
uno de ellos utilizando la serie diferenciada y posteriormente se
comparará su desempeño mediante las métricas descritas.

El objetivo final será seleccionar la especificación que ofrezca el
mejor equilibrio entre capacidad predictiva, parsimonia e interpretación
económica.

    \hypertarget{estrategia-de-modelado}{%
\subsection{23.1 Estrategia de Modelado}\label{estrategia-de-modelado}}

Con base en los resultados obtenidos en las pruebas de estacionariedad,
la descomposición temporal, la ACF y la PACF, se adoptará una estrategia
de modelado incremental.

El objetivo consiste en comparar modelos de complejidad creciente para
determinar si la incorporación explícita de componentes estacionales y
variables exógenas mejora la capacidad predictiva.

La secuencia de evaluación será:

\begin{enumerate}
\def\labelenumi{\arabic{enumi}.}
\tightlist
\item
  Modelos ARIMA.
\item
  Modelos SARIMA.
\item
  Modelos SARIMAX.
\end{enumerate}

Este procedimiento permitirá cuantificar el beneficio aportado por la
estacionalidad y por las variables explicativas externas incorporadas al
sistema.

    \hypertarget{modelos-a-evaluar}{%
\subsection{23.2 Modelos a Evaluar}\label{modelos-a-evaluar}}

Se evaluarán inicialmente los siguientes modelos candidatos:

\begin{longtable}[]{@{}ll@{}}
\toprule\noalign{}
Modelo & Configuración \\
\midrule\noalign{}
\endhead
\bottomrule\noalign{}
\endlastfoot
ARIMA 1 & (1,1,0) \\
ARIMA 2 & (0,1,1) \\
ARIMA 3 & (1,1,1) \\
SARIMA 1 & (1,1,0)(1,0,0,12) \\
SARIMA 2 & (1,1,1)(1,0,1,12) \\
\end{longtable}

Posteriormente se evaluarán modelos SARIMAX incorporando las variables
exógenas seleccionadas durante el proceso de reducción de
multicolinealidad:

\begin{itemize}
\tightlist
\item
  Precio internacional del maíz (CBOT).
\item
  Costos de transporte.
\item
  Temperatura promedio nacional.
\item
  Lluvia acumulada nacional.
\item
  Índice de severidad de sequía.
\item
  Producción nacional de maíz blanco.
\end{itemize}

Estas variables fueron seleccionadas debido a que presentan fundamentos
económicos y agroclimáticos plausibles para explicar las variaciones
observadas en el precio nacional del maíz blanco, al tiempo que
mantienen niveles aceptables de multicolinealidad según el análisis VIF
realizado previamente.

    \hypertarget{divisiuxf3n-temporal-de-los-datos}{%
\section{24. División Temporal de los
Datos}\label{divisiuxf3n-temporal-de-los-datos}}

Antes de entrenar los modelos ARIMA, SARIMA y SARIMAX es necesario
separar la información disponible en conjuntos de entrenamiento y
prueba.

El propósito de esta división consiste en evaluar la capacidad
predictiva de los modelos utilizando observaciones que no participaron
durante el proceso de estimación.

En problemas de regresión tradicional suele ser común realizar
divisiones aleatorias de los datos.

Sin embargo, en series temporales esta práctica no es apropiada debido a
que rompe la estructura cronológica de la información y puede introducir
información futura dentro del conjunto de entrenamiento.

Matemáticamente, si una observación futura participa en el
entrenamiento, el modelo tendría acceso a información que en condiciones
reales aún no estaría disponible.

Esto generaría una estimación artificialmente optimista del desempeño
predictivo.

Por esta razón, la división debe respetar el orden temporal de las
observaciones.

\begin{center}\rule{0.5\linewidth}{0.5pt}\end{center}

\hypertarget{estrategia-de-particiuxf3n}{%
\subsection{Estrategia de partición}\label{estrategia-de-particiuxf3n}}

Se utilizará una división cronológica:

\begin{itemize}
\tightlist
\item
  80\% de las observaciones para entrenamiento.
\item
  20\% de las observaciones para prueba.
\end{itemize}

Dado que la serie contiene:

\[
n = 317
\]

observaciones mensuales, la partición permitirá entrenar los modelos con
la mayor parte de la historia disponible y reservar los periodos más
recientes para la evaluación predictiva.

Serie temporal completa:

\[
2000 \rightarrow 2026
\]

División:

\[
\underbrace{\text{Entrenamiento}}_{80\%}
\quad
\underbrace{\text{Prueba}}_{20\%}
\]

Los datos de prueba representan observaciones futuras respecto al
conjunto de entrenamiento, reproduciendo el escenario real de
pronóstico.

\begin{center}\rule{0.5\linewidth}{0.5pt}\end{center}

\hypertarget{variables-utilizadas}{%
\subsection{Variables utilizadas}\label{variables-utilizadas}}

La variable objetivo será:

\[
Precio\_maiz
\]

Los modelos SARIMAX incorporarán adicionalmente las variables exógenas
seleccionadas durante el análisis de multicolinealidad:

\begin{itemize}
\tightlist
\item
  Precio internacional del maíz (CBOT).
\item
  Costos de transporte.
\item
  Temperatura media nacional.
\item
  Precipitación acumulada nacional.
\item
  Índice de severidad de sequía.
\item
  Producción nacional de maíz blanco.
\end{itemize}

    \begin{tcolorbox}[breakable, size=fbox, boxrule=1pt, pad at break*=1mm,colback=cellbackground, colframe=cellborder]
\prompt{In}{incolor}{70}{\boxspacing}
\begin{Verbatim}[commandchars=\\\{\}]
\PY{n}{serie\PYZus{}maiz} \PY{o}{=} \PY{p}{(}
    \PY{n}{pdf}
    \PY{o}{.}\PY{n}{sort\PYZus{}values}\PY{p}{(}\PY{l+s+s2}{\PYZdq{}}\PY{l+s+s2}{fecha}\PY{l+s+s2}{\PYZdq{}}\PY{p}{)}
    \PY{o}{.}\PY{n}{set\PYZus{}index}\PY{p}{(}\PY{l+s+s2}{\PYZdq{}}\PY{l+s+s2}{fecha}\PY{l+s+s2}{\PYZdq{}}\PY{p}{)}
    \PY{p}{[}\PY{l+s+s2}{\PYZdq{}}\PY{l+s+s2}{Precio\PYZus{}maizblanco\PYZus{}mxpkilo\PYZus{}prom\PYZus{}nacional}\PY{l+s+s2}{\PYZdq{}}\PY{p}{]}
\PY{p}{)}

\PY{n}{serie\PYZus{}maiz}\PY{o}{.}\PY{n}{head}\PY{p}{(}\PY{p}{)}
\end{Verbatim}
\end{tcolorbox}

            \begin{tcolorbox}[breakable, size=fbox, boxrule=.5pt, pad at break*=1mm, opacityfill=0]
\prompt{Out}{outcolor}{70}{\boxspacing}
\begin{Verbatim}[commandchars=\\\{\}]
fecha
2000-01-01    2.05
2000-02-01    2.06
2000-03-01    2.09
2000-04-01    2.12
2000-05-01    2.14
Name: Precio\_maizblanco\_mxpkilo\_prom\_nacional, dtype: float64
\end{Verbatim}
\end{tcolorbox}
        
    \begin{tcolorbox}[breakable, size=fbox, boxrule=1pt, pad at break*=1mm,colback=cellbackground, colframe=cellborder]
\prompt{In}{incolor}{71}{\boxspacing}
\begin{Verbatim}[commandchars=\\\{\}]
\PY{n}{train\PYZus{}size} \PY{o}{=} \PY{n+nb}{int}\PY{p}{(}\PY{n+nb}{len}\PY{p}{(}\PY{n}{serie\PYZus{}maiz}\PY{p}{)} \PY{o}{*} \PY{l+m+mf}{0.80}\PY{p}{)}

\PY{n}{train} \PY{o}{=} \PY{n}{serie\PYZus{}maiz}\PY{o}{.}\PY{n}{iloc}\PY{p}{[}\PY{p}{:}\PY{n}{train\PYZus{}size}\PY{p}{]}
\PY{n}{test} \PY{o}{=} \PY{n}{serie\PYZus{}maiz}\PY{o}{.}\PY{n}{iloc}\PY{p}{[}\PY{n}{train\PYZus{}size}\PY{p}{:}\PY{p}{]}

\PY{n+nb}{print}\PY{p}{(}\PY{l+s+s2}{\PYZdq{}}\PY{l+s+s2}{Observaciones entrenamiento:}\PY{l+s+s2}{\PYZdq{}}\PY{p}{,} \PY{n+nb}{len}\PY{p}{(}\PY{n}{train}\PY{p}{)}\PY{p}{)}
\PY{n+nb}{print}\PY{p}{(}\PY{l+s+s2}{\PYZdq{}}\PY{l+s+s2}{Observaciones prueba:}\PY{l+s+s2}{\PYZdq{}}\PY{p}{,} \PY{n+nb}{len}\PY{p}{(}\PY{n}{test}\PY{p}{)}\PY{p}{)}
\end{Verbatim}
\end{tcolorbox}

    \begin{Verbatim}[commandchars=\\\{\}]
Observaciones entrenamiento: 253
Observaciones prueba: 64
    \end{Verbatim}

    \begin{tcolorbox}[breakable, size=fbox, boxrule=1pt, pad at break*=1mm,colback=cellbackground, colframe=cellborder]
\prompt{In}{incolor}{72}{\boxspacing}
\begin{Verbatim}[commandchars=\\\{\}]
\PY{n+nb}{print}\PY{p}{(}\PY{l+s+s2}{\PYZdq{}}\PY{l+s+s2}{Inicio entrenamiento:}\PY{l+s+s2}{\PYZdq{}}\PY{p}{,} \PY{n}{train}\PY{o}{.}\PY{n}{index}\PY{o}{.}\PY{n}{min}\PY{p}{(}\PY{p}{)}\PY{p}{)}
\PY{n+nb}{print}\PY{p}{(}\PY{l+s+s2}{\PYZdq{}}\PY{l+s+s2}{Fin entrenamiento:}\PY{l+s+s2}{\PYZdq{}}\PY{p}{,} \PY{n}{train}\PY{o}{.}\PY{n}{index}\PY{o}{.}\PY{n}{max}\PY{p}{(}\PY{p}{)}\PY{p}{)}

\PY{n+nb}{print}\PY{p}{(}\PY{p}{)}

\PY{n+nb}{print}\PY{p}{(}\PY{l+s+s2}{\PYZdq{}}\PY{l+s+s2}{Inicio prueba:}\PY{l+s+s2}{\PYZdq{}}\PY{p}{,} \PY{n}{test}\PY{o}{.}\PY{n}{index}\PY{o}{.}\PY{n}{min}\PY{p}{(}\PY{p}{)}\PY{p}{)}
\PY{n+nb}{print}\PY{p}{(}\PY{l+s+s2}{\PYZdq{}}\PY{l+s+s2}{Fin prueba:}\PY{l+s+s2}{\PYZdq{}}\PY{p}{,} \PY{n}{test}\PY{o}{.}\PY{n}{index}\PY{o}{.}\PY{n}{max}\PY{p}{(}\PY{p}{)}\PY{p}{)}
\end{Verbatim}
\end{tcolorbox}

    \begin{Verbatim}[commandchars=\\\{\}]
Inicio entrenamiento: 2000-01-01
Fin entrenamiento: 2021-01-01

Inicio prueba: 2021-02-01
Fin prueba: 2026-05-01
    \end{Verbatim}

    \begin{tcolorbox}[breakable, size=fbox, boxrule=1pt, pad at break*=1mm,colback=cellbackground, colframe=cellborder]
\prompt{In}{incolor}{73}{\boxspacing}
\begin{Verbatim}[commandchars=\\\{\}]
\PY{n}{plt}\PY{o}{.}\PY{n}{figure}\PY{p}{(}\PY{n}{figsize}\PY{o}{=}\PY{p}{(}\PY{l+m+mi}{14}\PY{p}{,}\PY{l+m+mi}{6}\PY{p}{)}\PY{p}{)}

\PY{n}{plt}\PY{o}{.}\PY{n}{plot}\PY{p}{(}
    \PY{n}{train}\PY{o}{.}\PY{n}{index}\PY{p}{,}
    \PY{n}{train}\PY{p}{,}
    \PY{n}{label}\PY{o}{=}\PY{l+s+s2}{\PYZdq{}}\PY{l+s+s2}{Entrenamiento}\PY{l+s+s2}{\PYZdq{}}
\PY{p}{)}

\PY{n}{plt}\PY{o}{.}\PY{n}{plot}\PY{p}{(}
    \PY{n}{test}\PY{o}{.}\PY{n}{index}\PY{p}{,}
    \PY{n}{test}\PY{p}{,}
    \PY{n}{label}\PY{o}{=}\PY{l+s+s2}{\PYZdq{}}\PY{l+s+s2}{Prueba}\PY{l+s+s2}{\PYZdq{}}
\PY{p}{)}

\PY{n}{plt}\PY{o}{.}\PY{n}{axvline}\PY{p}{(}
    \PY{n}{x}\PY{o}{=}\PY{n}{test}\PY{o}{.}\PY{n}{index}\PY{o}{.}\PY{n}{min}\PY{p}{(}\PY{p}{)}\PY{p}{,}
    \PY{n}{color}\PY{o}{=}\PY{l+s+s2}{\PYZdq{}}\PY{l+s+s2}{red}\PY{l+s+s2}{\PYZdq{}}\PY{p}{,}
    \PY{n}{linestyle}\PY{o}{=}\PY{l+s+s2}{\PYZdq{}}\PY{l+s+s2}{\PYZhy{}\PYZhy{}}\PY{l+s+s2}{\PYZdq{}}\PY{p}{,}
    \PY{n}{label}\PY{o}{=}\PY{l+s+s2}{\PYZdq{}}\PY{l+s+s2}{Inicio prueba}\PY{l+s+s2}{\PYZdq{}}
\PY{p}{)}

\PY{n}{plt}\PY{o}{.}\PY{n}{title}\PY{p}{(}\PY{l+s+s2}{\PYZdq{}}\PY{l+s+s2}{División temporal entrenamiento\PYZhy{}prueba}\PY{l+s+s2}{\PYZdq{}}\PY{p}{)}
\PY{n}{plt}\PY{o}{.}\PY{n}{xlabel}\PY{p}{(}\PY{l+s+s2}{\PYZdq{}}\PY{l+s+s2}{Fecha}\PY{l+s+s2}{\PYZdq{}}\PY{p}{)}
\PY{n}{plt}\PY{o}{.}\PY{n}{ylabel}\PY{p}{(}\PY{l+s+s2}{\PYZdq{}}\PY{l+s+s2}{Precio maíz (MXN/Kg)}\PY{l+s+s2}{\PYZdq{}}\PY{p}{)}
\PY{n}{plt}\PY{o}{.}\PY{n}{legend}\PY{p}{(}\PY{p}{)}

\PY{n}{plt}\PY{o}{.}\PY{n}{grid}\PY{p}{(}\PY{k+kc}{True}\PY{p}{)}

\PY{n}{plt}\PY{o}{.}\PY{n}{show}\PY{p}{(}\PY{p}{)}
\end{Verbatim}
\end{tcolorbox}

    \begin{center}
    \adjustimage{max size={0.9\linewidth}{0.9\paperheight}}{output_109_0.png}
    \end{center}
    { \hspace*{\fill} \\}
    
    \hypertarget{entrenamiento-de-modelos-arima}{%
\section{25. Entrenamiento de Modelos
ARIMA}\label{entrenamiento-de-modelos-arima}}

Como punto de partida se evaluarán modelos ARIMA tradicionales que
únicamente consideran la dinámica temporal interna de la serie.

Estos modelos servirán como línea base para comparar posteriormente el
desempeño de modelos SARIMA y SARIMAX.

El primer modelo evaluado corresponde a:

\[
ARIMA(1,1,0)
\]

donde:

\begin{itemize}
\tightlist
\item
  \(p=1\): componente autorregresivo.
\item
  \(d=1\): primera diferenciación.
\item
  \(q=0\): sin componente de medias móviles.
\end{itemize}

La selección de este modelo se fundamenta en la evidencia observada en
la PACF, donde el primer rezago presentó un coeficiente significativo.

    \begin{tcolorbox}[breakable, size=fbox, boxrule=1pt, pad at break*=1mm,colback=cellbackground, colframe=cellborder]
\prompt{In}{incolor}{74}{\boxspacing}
\begin{Verbatim}[commandchars=\\\{\}]
\PY{c+c1}{\PYZsh{} ARIMA(1,1,0)}

\PY{k+kn}{from} \PY{n+nn}{statsmodels}\PY{n+nn}{.}\PY{n+nn}{tsa}\PY{n+nn}{.}\PY{n+nn}{arima}\PY{n+nn}{.}\PY{n+nn}{model} \PY{k+kn}{import} \PY{n}{ARIMA}

\PY{n}{modelo\PYZus{}arima\PYZus{}110} \PY{o}{=} \PY{n}{ARIMA}\PY{p}{(}
    \PY{n}{train}\PY{p}{,}
    \PY{n}{order}\PY{o}{=}\PY{p}{(}\PY{l+m+mi}{1}\PY{p}{,}\PY{l+m+mi}{1}\PY{p}{,}\PY{l+m+mi}{0}\PY{p}{)}
\PY{p}{)}

\PY{n}{resultado\PYZus{}arima\PYZus{}110} \PY{o}{=} \PY{n}{modelo\PYZus{}arima\PYZus{}110}\PY{o}{.}\PY{n}{fit}\PY{p}{(}\PY{p}{)}

\PY{n+nb}{print}\PY{p}{(}\PY{n}{resultado\PYZus{}arima\PYZus{}110}\PY{o}{.}\PY{n}{summary}\PY{p}{(}\PY{p}{)}\PY{p}{)}
\end{Verbatim}
\end{tcolorbox}

    \begin{Verbatim}[commandchars=\\\{\}]
                                          SARIMAX Results
================================================================================
===================
Dep. Variable:     Precio\_maizblanco\_mxpkilo\_prom\_nacional   No. Observations:
253
Model:                                      ARIMA(1, 1, 0)   Log Likelihood
207.006
Date:                                     Tue, 02 Jun 2026   AIC
-410.012
Time:                                             11:38:54   BIC
-402.953
Sample:                                         01-01-2000   HQIC
-407.171
                                              - 01-01-2021
Covariance Type:                                       opg
==============================================================================
                 coef    std err          z      P>|z|      [0.025      0.975]
------------------------------------------------------------------------------
ar.L1          0.0108      0.064      0.170      0.865      -0.114       0.136
sigma2         0.0113      0.001     19.224      0.000       0.010       0.012
================================================================================
===
Ljung-Box (L1) (Q):                   0.07   Jarque-Bera (JB):
267.82
Prob(Q):                              0.79   Prob(JB):
0.00
Heteroskedasticity (H):               3.37   Skew:
1.15
Prob(H) (two-sided):                  0.00   Kurtosis:
7.50
================================================================================
===

Warnings:
[1] Covariance matrix calculated using the outer product of gradients (complex-
step).
    \end{Verbatim}

    \begin{Verbatim}[commandchars=\\\{\}]
/opt/conda/lib/python3.11/site-packages/statsmodels/tsa/base/tsa\_model.py:473:
ValueWarning: No frequency information was provided, so inferred frequency MS
will be used.
  self.\_init\_dates(dates, freq)
/opt/conda/lib/python3.11/site-packages/statsmodels/tsa/base/tsa\_model.py:473:
ValueWarning: No frequency information was provided, so inferred frequency MS
will be used.
  self.\_init\_dates(dates, freq)
/opt/conda/lib/python3.11/site-packages/statsmodels/tsa/base/tsa\_model.py:473:
ValueWarning: No frequency information was provided, so inferred frequency MS
will be used.
  self.\_init\_dates(dates, freq)
    \end{Verbatim}

    \begin{tcolorbox}[breakable, size=fbox, boxrule=1pt, pad at break*=1mm,colback=cellbackground, colframe=cellborder]
\prompt{In}{incolor}{75}{\boxspacing}
\begin{Verbatim}[commandchars=\\\{\}]
\PY{n+nb}{print}\PY{p}{(}\PY{l+s+s2}{\PYZdq{}}\PY{l+s+s2}{AIC:}\PY{l+s+s2}{\PYZdq{}}\PY{p}{,} \PY{n}{resultado\PYZus{}arima\PYZus{}110}\PY{o}{.}\PY{n}{aic}\PY{p}{)}
\PY{n+nb}{print}\PY{p}{(}\PY{l+s+s2}{\PYZdq{}}\PY{l+s+s2}{BIC:}\PY{l+s+s2}{\PYZdq{}}\PY{p}{,} \PY{n}{resultado\PYZus{}arima\PYZus{}110}\PY{o}{.}\PY{n}{bic}\PY{p}{)}
\end{Verbatim}
\end{tcolorbox}

    \begin{Verbatim}[commandchars=\\\{\}]
AIC: -410.011769857501
BIC: -402.95291168247815
    \end{Verbatim}

    \begin{tcolorbox}[breakable, size=fbox, boxrule=1pt, pad at break*=1mm,colback=cellbackground, colframe=cellborder]
\prompt{In}{incolor}{76}{\boxspacing}
\begin{Verbatim}[commandchars=\\\{\}]
\PY{n}{pred\PYZus{}arima\PYZus{}110} \PY{o}{=} \PY{n}{resultado\PYZus{}arima\PYZus{}110}\PY{o}{.}\PY{n}{forecast}\PY{p}{(}
    \PY{n}{steps}\PY{o}{=}\PY{n+nb}{len}\PY{p}{(}\PY{n}{test}\PY{p}{)}
\PY{p}{)}
\end{Verbatim}
\end{tcolorbox}

    \begin{tcolorbox}[breakable, size=fbox, boxrule=1pt, pad at break*=1mm,colback=cellbackground, colframe=cellborder]
\prompt{In}{incolor}{77}{\boxspacing}
\begin{Verbatim}[commandchars=\\\{\}]
\PY{k+kn}{from} \PY{n+nn}{sklearn}\PY{n+nn}{.}\PY{n+nn}{metrics} \PY{k+kn}{import} \PY{p}{(}
    \PY{n}{mean\PYZus{}absolute\PYZus{}error}\PY{p}{,}
    \PY{n}{mean\PYZus{}squared\PYZus{}error}
\PY{p}{)}

\PY{k+kn}{import} \PY{n+nn}{numpy} \PY{k}{as} \PY{n+nn}{np}

\PY{n}{mae\PYZus{}110} \PY{o}{=} \PY{n}{mean\PYZus{}absolute\PYZus{}error}\PY{p}{(}
    \PY{n}{test}\PY{p}{,}
    \PY{n}{pred\PYZus{}arima\PYZus{}110}
\PY{p}{)}

\PY{n}{rmse\PYZus{}110} \PY{o}{=} \PY{n}{np}\PY{o}{.}\PY{n}{sqrt}\PY{p}{(}
    \PY{n}{mean\PYZus{}squared\PYZus{}error}\PY{p}{(}
        \PY{n}{test}\PY{p}{,}
        \PY{n}{pred\PYZus{}arima\PYZus{}110}
    \PY{p}{)}
\PY{p}{)}

\PY{n}{mape\PYZus{}110} \PY{o}{=} \PY{p}{(}
    \PY{n}{np}\PY{o}{.}\PY{n}{mean}\PY{p}{(}
        \PY{n}{np}\PY{o}{.}\PY{n}{abs}\PY{p}{(}
            \PY{p}{(}\PY{n}{test} \PY{o}{\PYZhy{}} \PY{n}{pred\PYZus{}arima\PYZus{}110}\PY{p}{)}
            \PY{o}{/} \PY{n}{test}
        \PY{p}{)}
    \PY{p}{)}
    \PY{o}{*} \PY{l+m+mi}{100}
\PY{p}{)}

\PY{n+nb}{print}\PY{p}{(}\PY{l+s+s2}{\PYZdq{}}\PY{l+s+s2}{MAE :}\PY{l+s+s2}{\PYZdq{}}\PY{p}{,} \PY{n}{mae\PYZus{}110}\PY{p}{)}
\PY{n+nb}{print}\PY{p}{(}\PY{l+s+s2}{\PYZdq{}}\PY{l+s+s2}{RMSE:}\PY{l+s+s2}{\PYZdq{}}\PY{p}{,} \PY{n}{rmse\PYZus{}110}\PY{p}{)}
\PY{n+nb}{print}\PY{p}{(}\PY{l+s+s2}{\PYZdq{}}\PY{l+s+s2}{MAPE:}\PY{l+s+s2}{\PYZdq{}}\PY{p}{,} \PY{n}{mape\PYZus{}110}\PY{p}{)}
\end{Verbatim}
\end{tcolorbox}

    \begin{Verbatim}[commandchars=\\\{\}]
MAE : 3.330077021099834
RMSE: 3.553458917593595
MAPE: 36.17926204756893
    \end{Verbatim}

    \begin{tcolorbox}[breakable, size=fbox, boxrule=1pt, pad at break*=1mm,colback=cellbackground, colframe=cellborder]
\prompt{In}{incolor}{78}{\boxspacing}
\begin{Verbatim}[commandchars=\\\{\}]
\PY{n}{plt}\PY{o}{.}\PY{n}{figure}\PY{p}{(}\PY{n}{figsize}\PY{o}{=}\PY{p}{(}\PY{l+m+mi}{14}\PY{p}{,}\PY{l+m+mi}{6}\PY{p}{)}\PY{p}{)}

\PY{n}{plt}\PY{o}{.}\PY{n}{plot}\PY{p}{(}
    \PY{n}{train}\PY{o}{.}\PY{n}{index}\PY{p}{,}
    \PY{n}{train}\PY{p}{,}
    \PY{n}{label}\PY{o}{=}\PY{l+s+s2}{\PYZdq{}}\PY{l+s+s2}{Entrenamiento}\PY{l+s+s2}{\PYZdq{}}
\PY{p}{)}

\PY{n}{plt}\PY{o}{.}\PY{n}{plot}\PY{p}{(}
    \PY{n}{test}\PY{o}{.}\PY{n}{index}\PY{p}{,}
    \PY{n}{test}\PY{p}{,}
    \PY{n}{label}\PY{o}{=}\PY{l+s+s2}{\PYZdq{}}\PY{l+s+s2}{Real}\PY{l+s+s2}{\PYZdq{}}
\PY{p}{)}

\PY{n}{plt}\PY{o}{.}\PY{n}{plot}\PY{p}{(}
    \PY{n}{test}\PY{o}{.}\PY{n}{index}\PY{p}{,}
    \PY{n}{pred\PYZus{}arima\PYZus{}110}\PY{p}{,}
    \PY{n}{label}\PY{o}{=}\PY{l+s+s2}{\PYZdq{}}\PY{l+s+s2}{ARIMA(1,1,0)}\PY{l+s+s2}{\PYZdq{}}
\PY{p}{)}

\PY{n}{plt}\PY{o}{.}\PY{n}{title}\PY{p}{(}\PY{l+s+s2}{\PYZdq{}}\PY{l+s+s2}{ARIMA(1,1,0)}\PY{l+s+s2}{\PYZdq{}}\PY{p}{)}
\PY{n}{plt}\PY{o}{.}\PY{n}{ylabel}\PY{p}{(}\PY{l+s+s2}{\PYZdq{}}\PY{l+s+s2}{MXN/Kg}\PY{l+s+s2}{\PYZdq{}}\PY{p}{)}
\PY{n}{plt}\PY{o}{.}\PY{n}{legend}\PY{p}{(}\PY{p}{)}
\PY{n}{plt}\PY{o}{.}\PY{n}{grid}\PY{p}{(}\PY{k+kc}{True}\PY{p}{)}

\PY{n}{plt}\PY{o}{.}\PY{n}{show}\PY{p}{(}\PY{p}{)}
\end{Verbatim}
\end{tcolorbox}

    \begin{center}
    \adjustimage{max size={0.9\linewidth}{0.9\paperheight}}{output_115_0.png}
    \end{center}
    { \hspace*{\fill} \\}
    
    \hypertarget{interpretaciuxf3n-del-modelo-arima110}{%
\subsection{Interpretación del modelo
ARIMA(1,1,0)}\label{interpretaciuxf3n-del-modelo-arima110}}

El modelo ARIMA(1,1,0) constituye la primera aproximación utilizada para
modelar la dinámica temporal del precio nacional del maíz blanco.

Los resultados muestran que el parámetro autorregresivo estimado no
resulta estadísticamente significativo:

\[
AR(1)=0.0108
\]

con un valor-p de:

\[
p=0.865
\]

lo que indica que la influencia directa del rezago inmediato es
prácticamente nula una vez aplicada la primera diferenciación.

Desde la perspectiva predictiva, el modelo presenta un error porcentual
medio absoluto:

\[
MAPE = 36.18\%
\]

valor que evidencia una capacidad limitada para reproducir la evolución
observada durante el periodo de prueba.

La inspección visual confirma esta conclusión, ya que el modelo genera
pronósticos prácticamente constantes y no logra capturar el fuerte
incremento de precios ocurrido entre 2021 y 2023.

Estos resultados sugieren que la estructura temporal del precio del maíz
no puede ser explicada únicamente mediante un componente autorregresivo
simple.

Por lo tanto, es necesario evaluar configuraciones alternativas que
incorporen componentes de medias móviles y posteriormente estructuras
estacionales explícitas.

    \hypertarget{modelo-arima011}{%
\subsection{25.2 Modelo ARIMA(0,1,1)}\label{modelo-arima011}}

El segundo modelo evaluado corresponde a:

\[
ARIMA(0,1,1)
\]

donde:

\begin{itemize}
\tightlist
\item
  \(p=0\): sin componente autorregresivo.
\item
  \(d=1\): primera diferenciación.
\item
  \(q=1\): componente de medias móviles.
\end{itemize}

La selección de este modelo se fundamenta en la evidencia observada en
la función de autocorrelación (ACF), la cual mostró dependencia
significativa en los primeros rezagos.

Los modelos MA suelen ser particularmente útiles para capturar choques
temporales y perturbaciones de corto plazo que afectan la evolución de
la serie.

    \begin{tcolorbox}[breakable, size=fbox, boxrule=1pt, pad at break*=1mm,colback=cellbackground, colframe=cellborder]
\prompt{In}{incolor}{79}{\boxspacing}
\begin{Verbatim}[commandchars=\\\{\}]
\PY{c+c1}{\PYZsh{} ARIMA(0,1,1)}

\PY{n}{modelo\PYZus{}arima\PYZus{}011} \PY{o}{=} \PY{n}{ARIMA}\PY{p}{(}
    \PY{n}{train}\PY{p}{,}
    \PY{n}{order}\PY{o}{=}\PY{p}{(}\PY{l+m+mi}{0}\PY{p}{,}\PY{l+m+mi}{1}\PY{p}{,}\PY{l+m+mi}{1}\PY{p}{)}
\PY{p}{)}

\PY{n}{resultado\PYZus{}arima\PYZus{}011} \PY{o}{=} \PY{n}{modelo\PYZus{}arima\PYZus{}011}\PY{o}{.}\PY{n}{fit}\PY{p}{(}\PY{p}{)}

\PY{n+nb}{print}\PY{p}{(}\PY{n}{resultado\PYZus{}arima\PYZus{}011}\PY{o}{.}\PY{n}{summary}\PY{p}{(}\PY{p}{)}\PY{p}{)}
\end{Verbatim}
\end{tcolorbox}

    \begin{Verbatim}[commandchars=\\\{\}]
                                          SARIMAX Results
================================================================================
===================
Dep. Variable:     Precio\_maizblanco\_mxpkilo\_prom\_nacional   No. Observations:
253
Model:                                      ARIMA(0, 1, 1)   Log Likelihood
207.002
Date:                                     Tue, 02 Jun 2026   AIC
-410.004
Time:                                             11:43:28   BIC
-402.945
Sample:                                         01-01-2000   HQIC
-407.164
                                              - 01-01-2021
Covariance Type:                                       opg
==============================================================================
                 coef    std err          z      P>|z|      [0.025      0.975]
------------------------------------------------------------------------------
ma.L1          0.0077      0.064      0.122      0.903      -0.117       0.132
sigma2         0.0113      0.001     19.233      0.000       0.010       0.012
================================================================================
===
Ljung-Box (L1) (Q):                   0.04   Jarque-Bera (JB):
266.26
Prob(Q):                              0.84   Prob(JB):
0.00
Heteroskedasticity (H):               3.37   Skew:
1.14
Prob(H) (two-sided):                  0.00   Kurtosis:
7.49
================================================================================
===

Warnings:
[1] Covariance matrix calculated using the outer product of gradients (complex-
step).
    \end{Verbatim}

    \begin{Verbatim}[commandchars=\\\{\}]
/opt/conda/lib/python3.11/site-packages/statsmodels/tsa/base/tsa\_model.py:473:
ValueWarning: No frequency information was provided, so inferred frequency MS
will be used.
  self.\_init\_dates(dates, freq)
/opt/conda/lib/python3.11/site-packages/statsmodels/tsa/base/tsa\_model.py:473:
ValueWarning: No frequency information was provided, so inferred frequency MS
will be used.
  self.\_init\_dates(dates, freq)
/opt/conda/lib/python3.11/site-packages/statsmodels/tsa/base/tsa\_model.py:473:
ValueWarning: No frequency information was provided, so inferred frequency MS
will be used.
  self.\_init\_dates(dates, freq)
    \end{Verbatim}

    \begin{tcolorbox}[breakable, size=fbox, boxrule=1pt, pad at break*=1mm,colback=cellbackground, colframe=cellborder]
\prompt{In}{incolor}{80}{\boxspacing}
\begin{Verbatim}[commandchars=\\\{\}]
\PY{n+nb}{print}\PY{p}{(}\PY{l+s+s2}{\PYZdq{}}\PY{l+s+s2}{AIC:}\PY{l+s+s2}{\PYZdq{}}\PY{p}{,} \PY{n}{resultado\PYZus{}arima\PYZus{}011}\PY{o}{.}\PY{n}{aic}\PY{p}{)}
\PY{n+nb}{print}\PY{p}{(}\PY{l+s+s2}{\PYZdq{}}\PY{l+s+s2}{BIC:}\PY{l+s+s2}{\PYZdq{}}\PY{p}{,} \PY{n}{resultado\PYZus{}arima\PYZus{}011}\PY{o}{.}\PY{n}{bic}\PY{p}{)}
\end{Verbatim}
\end{tcolorbox}

    \begin{Verbatim}[commandchars=\\\{\}]
AIC: -410.0041007027469
BIC: -402.945242527724
    \end{Verbatim}

    \begin{tcolorbox}[breakable, size=fbox, boxrule=1pt, pad at break*=1mm,colback=cellbackground, colframe=cellborder]
\prompt{In}{incolor}{81}{\boxspacing}
\begin{Verbatim}[commandchars=\\\{\}]
\PY{n}{pred\PYZus{}arima\PYZus{}011} \PY{o}{=} \PY{n}{resultado\PYZus{}arima\PYZus{}011}\PY{o}{.}\PY{n}{forecast}\PY{p}{(}
    \PY{n}{steps}\PY{o}{=}\PY{n+nb}{len}\PY{p}{(}\PY{n}{test}\PY{p}{)}
\PY{p}{)}
\end{Verbatim}
\end{tcolorbox}

    \begin{tcolorbox}[breakable, size=fbox, boxrule=1pt, pad at break*=1mm,colback=cellbackground, colframe=cellborder]
\prompt{In}{incolor}{82}{\boxspacing}
\begin{Verbatim}[commandchars=\\\{\}]
\PY{n}{mae\PYZus{}011} \PY{o}{=} \PY{n}{mean\PYZus{}absolute\PYZus{}error}\PY{p}{(}
    \PY{n}{test}\PY{p}{,}
    \PY{n}{pred\PYZus{}arima\PYZus{}011}
\PY{p}{)}

\PY{n}{rmse\PYZus{}011} \PY{o}{=} \PY{n}{np}\PY{o}{.}\PY{n}{sqrt}\PY{p}{(}
    \PY{n}{mean\PYZus{}squared\PYZus{}error}\PY{p}{(}
        \PY{n}{test}\PY{p}{,}
        \PY{n}{pred\PYZus{}arima\PYZus{}011}
    \PY{p}{)}
\PY{p}{)}

\PY{n}{mape\PYZus{}011} \PY{o}{=} \PY{p}{(}
    \PY{n}{np}\PY{o}{.}\PY{n}{mean}\PY{p}{(}
        \PY{n}{np}\PY{o}{.}\PY{n}{abs}\PY{p}{(}
            \PY{p}{(}\PY{n}{test} \PY{o}{\PYZhy{}} \PY{n}{pred\PYZus{}arima\PYZus{}011}\PY{p}{)}
            \PY{o}{/} \PY{n}{test}
        \PY{p}{)}
    \PY{p}{)}
    \PY{o}{*} \PY{l+m+mi}{100}
\PY{p}{)}

\PY{n+nb}{print}\PY{p}{(}\PY{l+s+s2}{\PYZdq{}}\PY{l+s+s2}{MAE :}\PY{l+s+s2}{\PYZdq{}}\PY{p}{,} \PY{n}{mae\PYZus{}011}\PY{p}{)}
\PY{n+nb}{print}\PY{p}{(}\PY{l+s+s2}{\PYZdq{}}\PY{l+s+s2}{RMSE:}\PY{l+s+s2}{\PYZdq{}}\PY{p}{,} \PY{n}{rmse\PYZus{}011}\PY{p}{)}
\PY{n+nb}{print}\PY{p}{(}\PY{l+s+s2}{\PYZdq{}}\PY{l+s+s2}{MAPE:}\PY{l+s+s2}{\PYZdq{}}\PY{p}{,} \PY{n}{mape\PYZus{}011}\PY{p}{)}
\end{Verbatim}
\end{tcolorbox}

    \begin{Verbatim}[commandchars=\\\{\}]
MAE : 3.3315129248594886
RMSE: 3.554805320760312
MAPE: 36.195904174555984
    \end{Verbatim}

    \begin{tcolorbox}[breakable, size=fbox, boxrule=1pt, pad at break*=1mm,colback=cellbackground, colframe=cellborder]
\prompt{In}{incolor}{83}{\boxspacing}
\begin{Verbatim}[commandchars=\\\{\}]
\PY{n}{plt}\PY{o}{.}\PY{n}{figure}\PY{p}{(}\PY{n}{figsize}\PY{o}{=}\PY{p}{(}\PY{l+m+mi}{14}\PY{p}{,}\PY{l+m+mi}{6}\PY{p}{)}\PY{p}{)}

\PY{n}{plt}\PY{o}{.}\PY{n}{plot}\PY{p}{(}
    \PY{n}{train}\PY{o}{.}\PY{n}{index}\PY{p}{,}
    \PY{n}{train}\PY{p}{,}
    \PY{n}{label}\PY{o}{=}\PY{l+s+s2}{\PYZdq{}}\PY{l+s+s2}{Entrenamiento}\PY{l+s+s2}{\PYZdq{}}
\PY{p}{)}

\PY{n}{plt}\PY{o}{.}\PY{n}{plot}\PY{p}{(}
    \PY{n}{test}\PY{o}{.}\PY{n}{index}\PY{p}{,}
    \PY{n}{test}\PY{p}{,}
    \PY{n}{label}\PY{o}{=}\PY{l+s+s2}{\PYZdq{}}\PY{l+s+s2}{Real}\PY{l+s+s2}{\PYZdq{}}
\PY{p}{)}

\PY{n}{plt}\PY{o}{.}\PY{n}{plot}\PY{p}{(}
    \PY{n}{test}\PY{o}{.}\PY{n}{index}\PY{p}{,}
    \PY{n}{pred\PYZus{}arima\PYZus{}011}\PY{p}{,}
    \PY{n}{label}\PY{o}{=}\PY{l+s+s2}{\PYZdq{}}\PY{l+s+s2}{ARIMA(0,1,1)}\PY{l+s+s2}{\PYZdq{}}
\PY{p}{)}

\PY{n}{plt}\PY{o}{.}\PY{n}{title}\PY{p}{(}\PY{l+s+s2}{\PYZdq{}}\PY{l+s+s2}{ARIMA(0,1,1)}\PY{l+s+s2}{\PYZdq{}}\PY{p}{)}
\PY{n}{plt}\PY{o}{.}\PY{n}{ylabel}\PY{p}{(}\PY{l+s+s2}{\PYZdq{}}\PY{l+s+s2}{MXN/Kg}\PY{l+s+s2}{\PYZdq{}}\PY{p}{)}

\PY{n}{plt}\PY{o}{.}\PY{n}{legend}\PY{p}{(}\PY{p}{)}
\PY{n}{plt}\PY{o}{.}\PY{n}{grid}\PY{p}{(}\PY{k+kc}{True}\PY{p}{)}

\PY{n}{plt}\PY{o}{.}\PY{n}{show}\PY{p}{(}\PY{p}{)}
\end{Verbatim}
\end{tcolorbox}

    \begin{center}
    \adjustimage{max size={0.9\linewidth}{0.9\paperheight}}{output_122_0.png}
    \end{center}
    { \hspace*{\fill} \\}
    
    \hypertarget{interpretaciuxf3n-del-modelo-arima011}{%
\subsection{Interpretación del modelo
ARIMA(0,1,1)}\label{interpretaciuxf3n-del-modelo-arima011}}

El modelo ARIMA(0,1,1) incorpora un componente de medias móviles de
primer orden con el objetivo de capturar perturbaciones temporales de
corto plazo.

Sin embargo, el parámetro estimado para el componente MA(1) no resultó
estadísticamente significativo:

\[
MA(1)=0.0077
\]

con un valor-p de:

\[
p=0.903
\]

lo que indica que la contribución de este componente es prácticamente
nula.

Las métricas predictivas obtenidas son muy similares a las observadas en
el modelo ARIMA(1,1,0):

\[
MAPE = 36.20\%
\]

confirmando una capacidad limitada para reproducir el comportamiento
reciente de la serie.

La inspección visual muestra nuevamente pronósticos prácticamente
constantes, incapaces de capturar el fuerte incremento observado en los
precios entre 2021 y 2023.

La similitud observada entre ambos modelos sugiere que la estructura
temporal del precio nacional del maíz blanco no puede ser explicada
adecuadamente mediante componentes autorregresivos o de medias móviles
simples.

Por esta razón, se evaluarán modelos más flexibles que incorporen
simultáneamente ambos mecanismos de dependencia temporal.

    \hypertarget{modelo-arima111}{%
\subsection{25.3 Modelo ARIMA(1,1,1)}\label{modelo-arima111}}

El tercer modelo evaluado corresponde a:

\[
ARIMA(1,1,1)
\]

donde:

\begin{itemize}
\tightlist
\item
  \(p=1\): componente autorregresivo.
\item
  \(d=1\): primera diferenciación.
\item
  \(q=1\): componente de medias móviles.
\end{itemize}

Este modelo combina simultáneamente los dos mecanismos de dependencia
temporal observados preliminarmente en las funciones ACF y PACF.

Desde una perspectiva teórica, ARIMA(1,1,1) constituye la especificación
más flexible dentro de los modelos ARIMA básicos, ya que permite
capturar tanto efectos de persistencia temporal como impactos derivados
de perturbaciones recientes.

Su evaluación permitirá determinar si la combinación de ambos
componentes mejora la capacidad predictiva respecto a los modelos
previamente analizados.

    \begin{tcolorbox}[breakable, size=fbox, boxrule=1pt, pad at break*=1mm,colback=cellbackground, colframe=cellborder]
\prompt{In}{incolor}{84}{\boxspacing}
\begin{Verbatim}[commandchars=\\\{\}]
\PY{c+c1}{\PYZsh{} ARIMA(1,1,1)}

\PY{n}{modelo\PYZus{}arima\PYZus{}111} \PY{o}{=} \PY{n}{ARIMA}\PY{p}{(}
    \PY{n}{train}\PY{p}{,}
    \PY{n}{order}\PY{o}{=}\PY{p}{(}\PY{l+m+mi}{1}\PY{p}{,}\PY{l+m+mi}{1}\PY{p}{,}\PY{l+m+mi}{1}\PY{p}{)}
\PY{p}{)}

\PY{n}{resultado\PYZus{}arima\PYZus{}111} \PY{o}{=} \PY{n}{modelo\PYZus{}arima\PYZus{}111}\PY{o}{.}\PY{n}{fit}\PY{p}{(}\PY{p}{)}

\PY{n+nb}{print}\PY{p}{(}\PY{n}{resultado\PYZus{}arima\PYZus{}111}\PY{o}{.}\PY{n}{summary}\PY{p}{(}\PY{p}{)}\PY{p}{)}
\end{Verbatim}
\end{tcolorbox}

    \begin{Verbatim}[commandchars=\\\{\}]
                                          SARIMAX Results
================================================================================
===================
Dep. Variable:     Precio\_maizblanco\_mxpkilo\_prom\_nacional   No. Observations:
253
Model:                                      ARIMA(1, 1, 1)   Log Likelihood
225.538
Date:                                     Tue, 02 Jun 2026   AIC
-445.075
Time:                                             11:47:01   BIC
-434.487
Sample:                                         01-01-2000   HQIC
-440.815
                                              - 01-01-2021
Covariance Type:                                       opg
==============================================================================
                 coef    std err          z      P>|z|      [0.025      0.975]
------------------------------------------------------------------------------
ar.L1         -0.9990      0.008   -129.730      0.000      -1.014      -0.984
ma.L1          0.9777      0.034     29.145      0.000       0.912       1.043
sigma2         0.0097      0.000     19.964      0.000       0.009       0.011
================================================================================
===
Ljung-Box (L1) (Q):                   8.17   Jarque-Bera (JB):
307.04
Prob(Q):                              0.00   Prob(JB):
0.00
Heteroskedasticity (H):               3.07   Skew:
1.27
Prob(H) (two-sided):                  0.00   Kurtosis:
7.78
================================================================================
===

Warnings:
[1] Covariance matrix calculated using the outer product of gradients (complex-
step).
    \end{Verbatim}

    \begin{Verbatim}[commandchars=\\\{\}]
/opt/conda/lib/python3.11/site-packages/statsmodels/tsa/base/tsa\_model.py:473:
ValueWarning: No frequency information was provided, so inferred frequency MS
will be used.
  self.\_init\_dates(dates, freq)
/opt/conda/lib/python3.11/site-packages/statsmodels/tsa/base/tsa\_model.py:473:
ValueWarning: No frequency information was provided, so inferred frequency MS
will be used.
  self.\_init\_dates(dates, freq)
/opt/conda/lib/python3.11/site-packages/statsmodels/tsa/base/tsa\_model.py:473:
ValueWarning: No frequency information was provided, so inferred frequency MS
will be used.
  self.\_init\_dates(dates, freq)
    \end{Verbatim}

    \begin{tcolorbox}[breakable, size=fbox, boxrule=1pt, pad at break*=1mm,colback=cellbackground, colframe=cellborder]
\prompt{In}{incolor}{85}{\boxspacing}
\begin{Verbatim}[commandchars=\\\{\}]
\PY{n+nb}{print}\PY{p}{(}\PY{l+s+s2}{\PYZdq{}}\PY{l+s+s2}{AIC:}\PY{l+s+s2}{\PYZdq{}}\PY{p}{,} \PY{n}{resultado\PYZus{}arima\PYZus{}111}\PY{o}{.}\PY{n}{aic}\PY{p}{)}
\PY{n+nb}{print}\PY{p}{(}\PY{l+s+s2}{\PYZdq{}}\PY{l+s+s2}{BIC:}\PY{l+s+s2}{\PYZdq{}}\PY{p}{,} \PY{n}{resultado\PYZus{}arima\PYZus{}111}\PY{o}{.}\PY{n}{bic}\PY{p}{)}
\end{Verbatim}
\end{tcolorbox}

    \begin{Verbatim}[commandchars=\\\{\}]
AIC: -445.07527402363746
BIC: -434.4869867611032
    \end{Verbatim}

    \begin{tcolorbox}[breakable, size=fbox, boxrule=1pt, pad at break*=1mm,colback=cellbackground, colframe=cellborder]
\prompt{In}{incolor}{86}{\boxspacing}
\begin{Verbatim}[commandchars=\\\{\}]
\PY{n}{pred\PYZus{}arima\PYZus{}111} \PY{o}{=} \PY{n}{resultado\PYZus{}arima\PYZus{}111}\PY{o}{.}\PY{n}{forecast}\PY{p}{(}
    \PY{n}{steps}\PY{o}{=}\PY{n+nb}{len}\PY{p}{(}\PY{n}{test}\PY{p}{)}
\PY{p}{)}
\end{Verbatim}
\end{tcolorbox}

    \begin{tcolorbox}[breakable, size=fbox, boxrule=1pt, pad at break*=1mm,colback=cellbackground, colframe=cellborder]
\prompt{In}{incolor}{87}{\boxspacing}
\begin{Verbatim}[commandchars=\\\{\}]
\PY{n}{mae\PYZus{}111} \PY{o}{=} \PY{n}{mean\PYZus{}absolute\PYZus{}error}\PY{p}{(}
    \PY{n}{test}\PY{p}{,}
    \PY{n}{pred\PYZus{}arima\PYZus{}111}
\PY{p}{)}

\PY{n}{rmse\PYZus{}111} \PY{o}{=} \PY{n}{np}\PY{o}{.}\PY{n}{sqrt}\PY{p}{(}
    \PY{n}{mean\PYZus{}squared\PYZus{}error}\PY{p}{(}
        \PY{n}{test}\PY{p}{,}
        \PY{n}{pred\PYZus{}arima\PYZus{}111}
    \PY{p}{)}
\PY{p}{)}

\PY{n}{mape\PYZus{}111} \PY{o}{=} \PY{p}{(}
    \PY{n}{np}\PY{o}{.}\PY{n}{mean}\PY{p}{(}
        \PY{n}{np}\PY{o}{.}\PY{n}{abs}\PY{p}{(}
            \PY{p}{(}\PY{n}{test} \PY{o}{\PYZhy{}} \PY{n}{pred\PYZus{}arima\PYZus{}111}\PY{p}{)}
            \PY{o}{/} \PY{n}{test}
        \PY{p}{)}
    \PY{p}{)}
    \PY{o}{*} \PY{l+m+mi}{100}
\PY{p}{)}

\PY{n+nb}{print}\PY{p}{(}\PY{l+s+s2}{\PYZdq{}}\PY{l+s+s2}{MAE :}\PY{l+s+s2}{\PYZdq{}}\PY{p}{,} \PY{n}{mae\PYZus{}111}\PY{p}{)}
\PY{n+nb}{print}\PY{p}{(}\PY{l+s+s2}{\PYZdq{}}\PY{l+s+s2}{RMSE:}\PY{l+s+s2}{\PYZdq{}}\PY{p}{,} \PY{n}{rmse\PYZus{}111}\PY{p}{)}
\PY{n+nb}{print}\PY{p}{(}\PY{l+s+s2}{\PYZdq{}}\PY{l+s+s2}{MAPE:}\PY{l+s+s2}{\PYZdq{}}\PY{p}{,} \PY{n}{mape\PYZus{}111}\PY{p}{)}
\end{Verbatim}
\end{tcolorbox}

    \begin{Verbatim}[commandchars=\\\{\}]
MAE : 3.3668093834768995
RMSE: 3.5872748726713555
MAPE: 36.60955040860067
    \end{Verbatim}

    \begin{tcolorbox}[breakable, size=fbox, boxrule=1pt, pad at break*=1mm,colback=cellbackground, colframe=cellborder]
\prompt{In}{incolor}{88}{\boxspacing}
\begin{Verbatim}[commandchars=\\\{\}]
\PY{n}{plt}\PY{o}{.}\PY{n}{figure}\PY{p}{(}\PY{n}{figsize}\PY{o}{=}\PY{p}{(}\PY{l+m+mi}{14}\PY{p}{,}\PY{l+m+mi}{6}\PY{p}{)}\PY{p}{)}

\PY{n}{plt}\PY{o}{.}\PY{n}{plot}\PY{p}{(}
    \PY{n}{train}\PY{o}{.}\PY{n}{index}\PY{p}{,}
    \PY{n}{train}\PY{p}{,}
    \PY{n}{label}\PY{o}{=}\PY{l+s+s2}{\PYZdq{}}\PY{l+s+s2}{Entrenamiento}\PY{l+s+s2}{\PYZdq{}}
\PY{p}{)}

\PY{n}{plt}\PY{o}{.}\PY{n}{plot}\PY{p}{(}
    \PY{n}{test}\PY{o}{.}\PY{n}{index}\PY{p}{,}
    \PY{n}{test}\PY{p}{,}
    \PY{n}{label}\PY{o}{=}\PY{l+s+s2}{\PYZdq{}}\PY{l+s+s2}{Real}\PY{l+s+s2}{\PYZdq{}}
\PY{p}{)}

\PY{n}{plt}\PY{o}{.}\PY{n}{plot}\PY{p}{(}
    \PY{n}{test}\PY{o}{.}\PY{n}{index}\PY{p}{,}
    \PY{n}{pred\PYZus{}arima\PYZus{}111}\PY{p}{,}
    \PY{n}{label}\PY{o}{=}\PY{l+s+s2}{\PYZdq{}}\PY{l+s+s2}{ARIMA(1,1,1)}\PY{l+s+s2}{\PYZdq{}}
\PY{p}{)}

\PY{n}{plt}\PY{o}{.}\PY{n}{title}\PY{p}{(}\PY{l+s+s2}{\PYZdq{}}\PY{l+s+s2}{ARIMA(1,1,1)}\PY{l+s+s2}{\PYZdq{}}\PY{p}{)}
\PY{n}{plt}\PY{o}{.}\PY{n}{ylabel}\PY{p}{(}\PY{l+s+s2}{\PYZdq{}}\PY{l+s+s2}{MXN/Kg}\PY{l+s+s2}{\PYZdq{}}\PY{p}{)}

\PY{n}{plt}\PY{o}{.}\PY{n}{legend}\PY{p}{(}\PY{p}{)}
\PY{n}{plt}\PY{o}{.}\PY{n}{grid}\PY{p}{(}\PY{k+kc}{True}\PY{p}{)}

\PY{n}{plt}\PY{o}{.}\PY{n}{show}\PY{p}{(}\PY{p}{)}
\end{Verbatim}
\end{tcolorbox}

    \begin{center}
    \adjustimage{max size={0.9\linewidth}{0.9\paperheight}}{output_129_0.png}
    \end{center}
    { \hspace*{\fill} \\}
    
    \begin{tcolorbox}[breakable, size=fbox, boxrule=1pt, pad at break*=1mm,colback=cellbackground, colframe=cellborder]
\prompt{In}{incolor}{89}{\boxspacing}
\begin{Verbatim}[commandchars=\\\{\}]
\PY{n}{comparacion\PYZus{}arima} \PY{o}{=} \PY{n}{pd}\PY{o}{.}\PY{n}{DataFrame}\PY{p}{(}\PY{p}{\PYZob{}}
    \PY{l+s+s2}{\PYZdq{}}\PY{l+s+s2}{Modelo}\PY{l+s+s2}{\PYZdq{}}\PY{p}{:}\PY{p}{[}
        \PY{l+s+s2}{\PYZdq{}}\PY{l+s+s2}{ARIMA(1,1,0)}\PY{l+s+s2}{\PYZdq{}}\PY{p}{,}
        \PY{l+s+s2}{\PYZdq{}}\PY{l+s+s2}{ARIMA(0,1,1)}\PY{l+s+s2}{\PYZdq{}}\PY{p}{,}
        \PY{l+s+s2}{\PYZdq{}}\PY{l+s+s2}{ARIMA(1,1,1)}\PY{l+s+s2}{\PYZdq{}}
    \PY{p}{]}\PY{p}{,}
    \PY{l+s+s2}{\PYZdq{}}\PY{l+s+s2}{AIC}\PY{l+s+s2}{\PYZdq{}}\PY{p}{:}\PY{p}{[}
        \PY{n}{resultado\PYZus{}arima\PYZus{}110}\PY{o}{.}\PY{n}{aic}\PY{p}{,}
        \PY{n}{resultado\PYZus{}arima\PYZus{}011}\PY{o}{.}\PY{n}{aic}\PY{p}{,}
        \PY{n}{resultado\PYZus{}arima\PYZus{}111}\PY{o}{.}\PY{n}{aic}
    \PY{p}{]}\PY{p}{,}
    \PY{l+s+s2}{\PYZdq{}}\PY{l+s+s2}{BIC}\PY{l+s+s2}{\PYZdq{}}\PY{p}{:}\PY{p}{[}
        \PY{n}{resultado\PYZus{}arima\PYZus{}110}\PY{o}{.}\PY{n}{bic}\PY{p}{,}
        \PY{n}{resultado\PYZus{}arima\PYZus{}011}\PY{o}{.}\PY{n}{bic}\PY{p}{,}
        \PY{n}{resultado\PYZus{}arima\PYZus{}111}\PY{o}{.}\PY{n}{bic}
    \PY{p}{]}\PY{p}{,}
    \PY{l+s+s2}{\PYZdq{}}\PY{l+s+s2}{MAE}\PY{l+s+s2}{\PYZdq{}}\PY{p}{:}\PY{p}{[}
        \PY{n}{mae\PYZus{}110}\PY{p}{,}
        \PY{n}{mae\PYZus{}011}\PY{p}{,}
        \PY{n}{mae\PYZus{}111}
    \PY{p}{]}\PY{p}{,}
    \PY{l+s+s2}{\PYZdq{}}\PY{l+s+s2}{RMSE}\PY{l+s+s2}{\PYZdq{}}\PY{p}{:}\PY{p}{[}
        \PY{n}{rmse\PYZus{}110}\PY{p}{,}
        \PY{n}{rmse\PYZus{}011}\PY{p}{,}
        \PY{n}{rmse\PYZus{}111}
    \PY{p}{]}\PY{p}{,}
    \PY{l+s+s2}{\PYZdq{}}\PY{l+s+s2}{MAPE}\PY{l+s+s2}{\PYZdq{}}\PY{p}{:}\PY{p}{[}
        \PY{n}{mape\PYZus{}110}\PY{p}{,}
        \PY{n}{mape\PYZus{}011}\PY{p}{,}
        \PY{n}{mape\PYZus{}111}
    \PY{p}{]}
\PY{p}{\PYZcb{}}\PY{p}{)}

\PY{n}{comparacion\PYZus{}arima}\PY{o}{.}\PY{n}{sort\PYZus{}values}\PY{p}{(}
    \PY{n}{by}\PY{o}{=}\PY{l+s+s2}{\PYZdq{}}\PY{l+s+s2}{MAPE}\PY{l+s+s2}{\PYZdq{}}
\PY{p}{)}
\end{Verbatim}
\end{tcolorbox}

            \begin{tcolorbox}[breakable, size=fbox, boxrule=.5pt, pad at break*=1mm, opacityfill=0]
\prompt{Out}{outcolor}{89}{\boxspacing}
\begin{Verbatim}[commandchars=\\\{\}]
         Modelo         AIC         BIC       MAE      RMSE       MAPE
0  ARIMA(1,1,0) -410.011770 -402.952912  3.330077  3.553459  36.179262
1  ARIMA(0,1,1) -410.004101 -402.945243  3.331513  3.554805  36.195904
2  ARIMA(1,1,1) -445.075274 -434.486987  3.366809  3.587275  36.609550
\end{Verbatim}
\end{tcolorbox}
        
    \hypertarget{comparaciuxf3n-de-modelos-arima}{%
\subsection{25.4 Comparación de modelos
ARIMA}\label{comparaciuxf3n-de-modelos-arima}}

Los tres modelos ARIMA evaluados muestran resultados consistentes.

Los modelos ARIMA(1,1,0) y ARIMA(0,1,1) presentan desempeños
prácticamente idénticos tanto en términos de ajuste como de capacidad
predictiva.

Por su parte, ARIMA(1,1,1) logra una mejora sustancial en los criterios
de información AIC y BIC, lo que indica un mejor ajuste a los datos
históricos utilizados durante el entrenamiento.

Sin embargo, dicha mejora no se traduce en una reducción del error de
pronóstico sobre el conjunto de prueba. Por el contrario, las métricas
RMSE y MAPE muestran un ligero deterioro respecto a los modelos más
simples.

Este comportamiento constituye evidencia de que la serie del precio
nacional del maíz blanco presenta una dinámica más compleja que la
capturada por modelos ARIMA tradicionales.

En particular, la presencia de patrones estacionales previamente
identificados y la influencia potencial de variables económicas y
climáticas sugieren que una parte importante de la variabilidad
observada proviene de factores externos al comportamiento histórico de
la propia serie.

Por esta razón, el siguiente paso consiste en incorporar explícitamente
la estructura estacional detectada mediante modelos SARIMA.

    \hypertarget{modelos-sarima}{%
\section{26. Modelos SARIMA}\label{modelos-sarima}}

Los resultados obtenidos mediante los modelos ARIMA mostraron que la
estructura temporal del precio nacional del maíz blanco no puede ser
explicada adecuadamente únicamente mediante componentes autorregresivos
y de medias móviles de corto plazo.

Durante el análisis exploratorio se identificaron evidencias claras de
comportamiento estacional:

\begin{itemize}
\tightlist
\item
  Patrón repetitivo en la descomposición temporal.
\item
  Picos significativos en la ACF en los rezagos 12, 24, 36 y 48.
\item
  Pico estacional en la PACF alrededor del rezago 12.
\end{itemize}

Estos hallazgos sugieren la existencia de una periodicidad anual
consistente con los ciclos agrícolas de producción, comercialización y
almacenamiento del maíz blanco.

Por esta razón se incorporará una estructura estacional explícita
mediante modelos SARIMA.

La forma general del modelo es:

\[
SARIMA(p,d,q)(P,D,Q)_s
\]

donde:

\begin{itemize}
\tightlist
\item
  \((p,d,q)\) representan la parte no estacional.
\item
  \((P,D,Q)\) representan la parte estacional.
\item
  \(s\) representa la periodicidad de la estacionalidad.
\end{itemize}

Con base en los resultados obtenidos previamente se adopta:

\[
s = 12
\]

correspondiente a una periodicidad anual en datos mensuales.

    \hypertarget{modelo-sarima11010012}{%
\subsection{26.2 Modelo
SARIMA(1,1,0)(1,0,0,12)}\label{modelo-sarima11010012}}

El primer modelo estacional evaluado incorpora:

\begin{itemize}
\tightlist
\item
  Un componente autorregresivo de corto plazo.
\item
  Un componente autorregresivo estacional anual.
\end{itemize}

Su especificación es:

\[
SARIMA(1,1,0)(1,0,0)_{12}
\]

Este modelo permitirá evaluar si la incorporación explícita de la
estacionalidad anual mejora la capacidad predictiva observada en los
modelos ARIMA tradicionales.

    \begin{tcolorbox}[breakable, size=fbox, boxrule=1pt, pad at break*=1mm,colback=cellbackground, colframe=cellborder]
\prompt{In}{incolor}{90}{\boxspacing}
\begin{Verbatim}[commandchars=\\\{\}]
\PY{k+kn}{from} \PY{n+nn}{statsmodels}\PY{n+nn}{.}\PY{n+nn}{tsa}\PY{n+nn}{.}\PY{n+nn}{statespace}\PY{n+nn}{.}\PY{n+nn}{sarimax} \PY{k+kn}{import} \PY{n}{SARIMAX}

\PY{n}{modelo\PYZus{}sarima\PYZus{}1} \PY{o}{=} \PY{n}{SARIMAX}\PY{p}{(}
    \PY{n}{train}\PY{p}{,}
    \PY{n}{order}\PY{o}{=}\PY{p}{(}\PY{l+m+mi}{1}\PY{p}{,}\PY{l+m+mi}{1}\PY{p}{,}\PY{l+m+mi}{0}\PY{p}{)}\PY{p}{,}
    \PY{n}{seasonal\PYZus{}order}\PY{o}{=}\PY{p}{(}\PY{l+m+mi}{1}\PY{p}{,}\PY{l+m+mi}{0}\PY{p}{,}\PY{l+m+mi}{0}\PY{p}{,}\PY{l+m+mi}{12}\PY{p}{)}\PY{p}{,}
    \PY{n}{enforce\PYZus{}stationarity}\PY{o}{=}\PY{k+kc}{False}\PY{p}{,}
    \PY{n}{enforce\PYZus{}invertibility}\PY{o}{=}\PY{k+kc}{False}
\PY{p}{)}

\PY{n}{resultado\PYZus{}sarima\PYZus{}1} \PY{o}{=} \PY{n}{modelo\PYZus{}sarima\PYZus{}1}\PY{o}{.}\PY{n}{fit}\PY{p}{(}\PY{p}{)}

\PY{n+nb}{print}\PY{p}{(}\PY{n}{resultado\PYZus{}sarima\PYZus{}1}\PY{o}{.}\PY{n}{summary}\PY{p}{(}\PY{p}{)}\PY{p}{)}
\end{Verbatim}
\end{tcolorbox}

    \begin{Verbatim}[commandchars=\\\{\}]
                                          SARIMAX Results
================================================================================
===================
Dep. Variable:     Precio\_maizblanco\_mxpkilo\_prom\_nacional   No. Observations:
253
Model:                      SARIMAX(1, 1, 0)x(1, 0, 0, 12)   Log Likelihood
318.166
Date:                                     Tue, 02 Jun 2026   AIC
-630.332
Time:                                             11:51:05   BIC
-619.902
Sample:                                         01-01-2000   HQIC
-626.129
                                              - 01-01-2021
Covariance Type:                                       opg
==============================================================================
                 coef    std err          z      P>|z|      [0.025      0.975]
------------------------------------------------------------------------------
ar.L1          0.1534      0.061      2.535      0.011       0.035       0.272
ar.S.L12       0.9164      0.016     58.582      0.000       0.886       0.947
sigma2         0.0041      0.000     39.625      0.000       0.004       0.004
================================================================================
===
Ljung-Box (L1) (Q):                   0.00   Jarque-Bera (JB):
7728.98
Prob(Q):                              0.97   Prob(JB):
0.00
Heteroskedasticity (H):               0.63   Skew:
1.25
Prob(H) (two-sided):                  0.04   Kurtosis:
30.75
================================================================================
===

Warnings:
[1] Covariance matrix calculated using the outer product of gradients (complex-
step).
    \end{Verbatim}

    \begin{Verbatim}[commandchars=\\\{\}]
/opt/conda/lib/python3.11/site-packages/statsmodels/tsa/base/tsa\_model.py:473:
ValueWarning: No frequency information was provided, so inferred frequency MS
will be used.
  self.\_init\_dates(dates, freq)
/opt/conda/lib/python3.11/site-packages/statsmodels/tsa/base/tsa\_model.py:473:
ValueWarning: No frequency information was provided, so inferred frequency MS
will be used.
  self.\_init\_dates(dates, freq)
    \end{Verbatim}

    \begin{tcolorbox}[breakable, size=fbox, boxrule=1pt, pad at break*=1mm,colback=cellbackground, colframe=cellborder]
\prompt{In}{incolor}{91}{\boxspacing}
\begin{Verbatim}[commandchars=\\\{\}]
\PY{n+nb}{print}\PY{p}{(}\PY{l+s+s2}{\PYZdq{}}\PY{l+s+s2}{AIC:}\PY{l+s+s2}{\PYZdq{}}\PY{p}{,} \PY{n}{resultado\PYZus{}sarima\PYZus{}1}\PY{o}{.}\PY{n}{aic}\PY{p}{)}
\PY{n+nb}{print}\PY{p}{(}\PY{l+s+s2}{\PYZdq{}}\PY{l+s+s2}{BIC:}\PY{l+s+s2}{\PYZdq{}}\PY{p}{,} \PY{n}{resultado\PYZus{}sarima\PYZus{}1}\PY{o}{.}\PY{n}{bic}\PY{p}{)}
\end{Verbatim}
\end{tcolorbox}

    \begin{Verbatim}[commandchars=\\\{\}]
AIC: -630.3317208093426
BIC: -619.9023301535481
    \end{Verbatim}

    \begin{tcolorbox}[breakable, size=fbox, boxrule=1pt, pad at break*=1mm,colback=cellbackground, colframe=cellborder]
\prompt{In}{incolor}{92}{\boxspacing}
\begin{Verbatim}[commandchars=\\\{\}]
\PY{n}{pred\PYZus{}sarima\PYZus{}1} \PY{o}{=} \PY{n}{resultado\PYZus{}sarima\PYZus{}1}\PY{o}{.}\PY{n}{forecast}\PY{p}{(}
    \PY{n}{steps}\PY{o}{=}\PY{n+nb}{len}\PY{p}{(}\PY{n}{test}\PY{p}{)}
\PY{p}{)}
\end{Verbatim}
\end{tcolorbox}

    \begin{tcolorbox}[breakable, size=fbox, boxrule=1pt, pad at break*=1mm,colback=cellbackground, colframe=cellborder]
\prompt{In}{incolor}{93}{\boxspacing}
\begin{Verbatim}[commandchars=\\\{\}]
\PY{n}{mae\PYZus{}sarima\PYZus{}1} \PY{o}{=} \PY{n}{mean\PYZus{}absolute\PYZus{}error}\PY{p}{(}
    \PY{n}{test}\PY{p}{,}
    \PY{n}{pred\PYZus{}sarima\PYZus{}1}
\PY{p}{)}

\PY{n}{rmse\PYZus{}sarima\PYZus{}1} \PY{o}{=} \PY{n}{np}\PY{o}{.}\PY{n}{sqrt}\PY{p}{(}
    \PY{n}{mean\PYZus{}squared\PYZus{}error}\PY{p}{(}
        \PY{n}{test}\PY{p}{,}
        \PY{n}{pred\PYZus{}sarima\PYZus{}1}
    \PY{p}{)}
\PY{p}{)}

\PY{n}{mape\PYZus{}sarima\PYZus{}1} \PY{o}{=} \PY{p}{(}
    \PY{n}{np}\PY{o}{.}\PY{n}{mean}\PY{p}{(}
        \PY{n}{np}\PY{o}{.}\PY{n}{abs}\PY{p}{(}
            \PY{p}{(}\PY{n}{test} \PY{o}{\PYZhy{}} \PY{n}{pred\PYZus{}sarima\PYZus{}1}\PY{p}{)}
            \PY{o}{/} \PY{n}{test}
        \PY{p}{)}
    \PY{p}{)}
    \PY{o}{*} \PY{l+m+mi}{100}
\PY{p}{)}

\PY{n+nb}{print}\PY{p}{(}\PY{l+s+s2}{\PYZdq{}}\PY{l+s+s2}{MAE :}\PY{l+s+s2}{\PYZdq{}}\PY{p}{,} \PY{n}{mae\PYZus{}sarima\PYZus{}1}\PY{p}{)}
\PY{n+nb}{print}\PY{p}{(}\PY{l+s+s2}{\PYZdq{}}\PY{l+s+s2}{RMSE:}\PY{l+s+s2}{\PYZdq{}}\PY{p}{,} \PY{n}{rmse\PYZus{}sarima\PYZus{}1}\PY{p}{)}
\PY{n+nb}{print}\PY{p}{(}\PY{l+s+s2}{\PYZdq{}}\PY{l+s+s2}{MAPE:}\PY{l+s+s2}{\PYZdq{}}\PY{p}{,} \PY{n}{mape\PYZus{}sarima\PYZus{}1}\PY{p}{)}
\end{Verbatim}
\end{tcolorbox}

    \begin{Verbatim}[commandchars=\\\{\}]
MAE : 1.835544639568898
RMSE: 1.9936143229905587
MAPE: 20.0170492866158
    \end{Verbatim}

    \begin{tcolorbox}[breakable, size=fbox, boxrule=1pt, pad at break*=1mm,colback=cellbackground, colframe=cellborder]
\prompt{In}{incolor}{94}{\boxspacing}
\begin{Verbatim}[commandchars=\\\{\}]
\PY{n}{plt}\PY{o}{.}\PY{n}{figure}\PY{p}{(}\PY{n}{figsize}\PY{o}{=}\PY{p}{(}\PY{l+m+mi}{14}\PY{p}{,}\PY{l+m+mi}{6}\PY{p}{)}\PY{p}{)}

\PY{n}{plt}\PY{o}{.}\PY{n}{plot}\PY{p}{(}
    \PY{n}{train}\PY{o}{.}\PY{n}{index}\PY{p}{,}
    \PY{n}{train}\PY{p}{,}
    \PY{n}{label}\PY{o}{=}\PY{l+s+s2}{\PYZdq{}}\PY{l+s+s2}{Entrenamiento}\PY{l+s+s2}{\PYZdq{}}
\PY{p}{)}

\PY{n}{plt}\PY{o}{.}\PY{n}{plot}\PY{p}{(}
    \PY{n}{test}\PY{o}{.}\PY{n}{index}\PY{p}{,}
    \PY{n}{test}\PY{p}{,}
    \PY{n}{label}\PY{o}{=}\PY{l+s+s2}{\PYZdq{}}\PY{l+s+s2}{Real}\PY{l+s+s2}{\PYZdq{}}
\PY{p}{)}

\PY{n}{plt}\PY{o}{.}\PY{n}{plot}\PY{p}{(}
    \PY{n}{test}\PY{o}{.}\PY{n}{index}\PY{p}{,}
    \PY{n}{pred\PYZus{}sarima\PYZus{}1}\PY{p}{,}
    \PY{n}{label}\PY{o}{=}\PY{l+s+s2}{\PYZdq{}}\PY{l+s+s2}{SARIMA(1,1,0)(1,0,0,12)}\PY{l+s+s2}{\PYZdq{}}
\PY{p}{)}

\PY{n}{plt}\PY{o}{.}\PY{n}{title}\PY{p}{(}\PY{l+s+s2}{\PYZdq{}}\PY{l+s+s2}{SARIMA(1,1,0)(1,0,0,12)}\PY{l+s+s2}{\PYZdq{}}\PY{p}{)}
\PY{n}{plt}\PY{o}{.}\PY{n}{ylabel}\PY{p}{(}\PY{l+s+s2}{\PYZdq{}}\PY{l+s+s2}{MXN/Kg}\PY{l+s+s2}{\PYZdq{}}\PY{p}{)}

\PY{n}{plt}\PY{o}{.}\PY{n}{legend}\PY{p}{(}\PY{p}{)}
\PY{n}{plt}\PY{o}{.}\PY{n}{grid}\PY{p}{(}\PY{k+kc}{True}\PY{p}{)}

\PY{n}{plt}\PY{o}{.}\PY{n}{show}\PY{p}{(}\PY{p}{)}
\end{Verbatim}
\end{tcolorbox}

    \begin{center}
    \adjustimage{max size={0.9\linewidth}{0.9\paperheight}}{output_138_0.png}
    \end{center}
    { \hspace*{\fill} \\}
    
    \hypertarget{interpretaciuxf3n-del-modelo-sarima11010012}{%
\subsection{Interpretación del modelo
SARIMA(1,1,0)(1,0,0,12)}\label{interpretaciuxf3n-del-modelo-sarima11010012}}

La incorporación explícita de una componente estacional anual produce
una mejora sustancial respecto a todos los modelos ARIMA previamente
evaluados.

El coeficiente autorregresivo estacional estimado fue:

\[
AR_{12}=0.9164
\]

con significancia estadística elevada:

\[
p < 0.001
\]

lo que indica una fuerte dependencia entre los valores observados y los
registrados doce meses atrás.

Este resultado es consistente con la naturaleza agrícola del mercado del
maíz blanco, donde los ciclos de siembra, cosecha, almacenamiento y
comercialización generan patrones recurrentes de periodicidad anual.

Desde la perspectiva predictiva, el modelo logra reducir
significativamente los errores de pronóstico:

\[
RMSE = 1.99
\]

y

\[
MAPE = 20.02\%
\]

representando una mejora cercana al 45\% respecto a los mejores modelos
ARIMA.

La reducción simultánea de AIC y BIC confirma además que la estructura
estacional incorporada describe adecuadamente la dinámica subyacente de
la serie.

En consecuencia, existe evidencia sólida para afirmar que la
estacionalidad anual constituye uno de los principales determinantes del
comportamiento del precio nacional del maíz blanco.

    \hypertarget{modelo-sarima11110112}{%
\subsection{26.3 Modelo
SARIMA(1,1,1)(1,0,1,12)}\label{modelo-sarima11110112}}

El segundo modelo estacional incorpora simultáneamente componentes
autorregresivos y de medias móviles tanto en la parte no estacional como
en la parte estacional.

Su especificación es:

\[
SARIMA(1,1,1)(1,0,1)_{12}
\]

Este modelo constituye una extensión del modelo SARIMA anterior y
permite capturar relaciones temporales más complejas.

El objetivo de esta evaluación consiste en determinar si la
incorporación de componentes adicionales mejora la capacidad predictiva
observada en el modelo estacional base.

    \begin{tcolorbox}[breakable, size=fbox, boxrule=1pt, pad at break*=1mm,colback=cellbackground, colframe=cellborder]
\prompt{In}{incolor}{95}{\boxspacing}
\begin{Verbatim}[commandchars=\\\{\}]
\PY{c+c1}{\PYZsh{} SARIMA(1,1,1)(1,0,1,12)}

\PY{n}{modelo\PYZus{}sarima\PYZus{}2} \PY{o}{=} \PY{n}{SARIMAX}\PY{p}{(}
    \PY{n}{train}\PY{p}{,}
    \PY{n}{order}\PY{o}{=}\PY{p}{(}\PY{l+m+mi}{1}\PY{p}{,}\PY{l+m+mi}{1}\PY{p}{,}\PY{l+m+mi}{1}\PY{p}{)}\PY{p}{,}
    \PY{n}{seasonal\PYZus{}order}\PY{o}{=}\PY{p}{(}\PY{l+m+mi}{1}\PY{p}{,}\PY{l+m+mi}{0}\PY{p}{,}\PY{l+m+mi}{1}\PY{p}{,}\PY{l+m+mi}{12}\PY{p}{)}\PY{p}{,}
    \PY{n}{enforce\PYZus{}stationarity}\PY{o}{=}\PY{k+kc}{False}\PY{p}{,}
    \PY{n}{enforce\PYZus{}invertibility}\PY{o}{=}\PY{k+kc}{False}
\PY{p}{)}

\PY{n}{resultado\PYZus{}sarima\PYZus{}2} \PY{o}{=} \PY{n}{modelo\PYZus{}sarima\PYZus{}2}\PY{o}{.}\PY{n}{fit}\PY{p}{(}\PY{p}{)}

\PY{n+nb}{print}\PY{p}{(}\PY{n}{resultado\PYZus{}sarima\PYZus{}2}\PY{o}{.}\PY{n}{summary}\PY{p}{(}\PY{p}{)}\PY{p}{)}
\end{Verbatim}
\end{tcolorbox}

    \begin{Verbatim}[commandchars=\\\{\}]
/opt/conda/lib/python3.11/site-packages/statsmodels/tsa/base/tsa\_model.py:473:
ValueWarning: No frequency information was provided, so inferred frequency MS
will be used.
  self.\_init\_dates(dates, freq)
/opt/conda/lib/python3.11/site-packages/statsmodels/tsa/base/tsa\_model.py:473:
ValueWarning: No frequency information was provided, so inferred frequency MS
will be used.
  self.\_init\_dates(dates, freq)
    \end{Verbatim}

    \begin{Verbatim}[commandchars=\\\{\}]
                                          SARIMAX Results
================================================================================
===================
Dep. Variable:     Precio\_maizblanco\_mxpkilo\_prom\_nacional   No. Observations:
253
Model:                      SARIMAX(1, 1, 1)x(1, 0, 1, 12)   Log Likelihood
336.196
Date:                                     Tue, 02 Jun 2026   AIC
-662.392
Time:                                             11:54:51   BIC
-645.031
Sample:                                         01-01-2000   HQIC
-655.395
                                              - 01-01-2021
Covariance Type:                                       opg
==============================================================================
                 coef    std err          z      P>|z|      [0.025      0.975]
------------------------------------------------------------------------------
ar.L1          0.2372      0.483      0.491      0.624      -0.710       1.185
ma.L1         -0.0874      0.503     -0.174      0.862      -1.072       0.898
ar.S.L12       1.0547      0.017     63.625      0.000       1.022       1.087
ma.S.L12      -0.5727      0.034    -17.035      0.000      -0.639      -0.507
sigma2         0.0034   8.75e-05     38.909      0.000       0.003       0.004
================================================================================
===
Ljung-Box (L1) (Q):                   0.00   Jarque-Bera (JB):
7764.50
Prob(Q):                              0.97   Prob(JB):
0.00
Heteroskedasticity (H):               0.63   Skew:
2.85
Prob(H) (two-sided):                  0.04   Kurtosis:
30.40
================================================================================
===

Warnings:
[1] Covariance matrix calculated using the outer product of gradients (complex-
step).
    \end{Verbatim}

    \begin{tcolorbox}[breakable, size=fbox, boxrule=1pt, pad at break*=1mm,colback=cellbackground, colframe=cellborder]
\prompt{In}{incolor}{96}{\boxspacing}
\begin{Verbatim}[commandchars=\\\{\}]
\PY{n+nb}{print}\PY{p}{(}\PY{l+s+s2}{\PYZdq{}}\PY{l+s+s2}{AIC:}\PY{l+s+s2}{\PYZdq{}}\PY{p}{,} \PY{n}{resultado\PYZus{}sarima\PYZus{}2}\PY{o}{.}\PY{n}{aic}\PY{p}{)}
\PY{n+nb}{print}\PY{p}{(}\PY{l+s+s2}{\PYZdq{}}\PY{l+s+s2}{BIC:}\PY{l+s+s2}{\PYZdq{}}\PY{p}{,} \PY{n}{resultado\PYZus{}sarima\PYZus{}2}\PY{o}{.}\PY{n}{bic}\PY{p}{)}
\end{Verbatim}
\end{tcolorbox}

    \begin{Verbatim}[commandchars=\\\{\}]
AIC: -662.3923066412685
BIC: -645.0309532729111
    \end{Verbatim}

    \begin{tcolorbox}[breakable, size=fbox, boxrule=1pt, pad at break*=1mm,colback=cellbackground, colframe=cellborder]
\prompt{In}{incolor}{97}{\boxspacing}
\begin{Verbatim}[commandchars=\\\{\}]
\PY{n}{pred\PYZus{}sarima\PYZus{}2} \PY{o}{=} \PY{n}{resultado\PYZus{}sarima\PYZus{}2}\PY{o}{.}\PY{n}{forecast}\PY{p}{(}
    \PY{n}{steps}\PY{o}{=}\PY{n+nb}{len}\PY{p}{(}\PY{n}{test}\PY{p}{)}
\PY{p}{)}
\end{Verbatim}
\end{tcolorbox}

    \begin{tcolorbox}[breakable, size=fbox, boxrule=1pt, pad at break*=1mm,colback=cellbackground, colframe=cellborder]
\prompt{In}{incolor}{98}{\boxspacing}
\begin{Verbatim}[commandchars=\\\{\}]
\PY{n}{mae\PYZus{}sarima\PYZus{}2} \PY{o}{=} \PY{n}{mean\PYZus{}absolute\PYZus{}error}\PY{p}{(}
    \PY{n}{test}\PY{p}{,}
    \PY{n}{pred\PYZus{}sarima\PYZus{}2}
\PY{p}{)}

\PY{n}{rmse\PYZus{}sarima\PYZus{}2} \PY{o}{=} \PY{n}{np}\PY{o}{.}\PY{n}{sqrt}\PY{p}{(}
    \PY{n}{mean\PYZus{}squared\PYZus{}error}\PY{p}{(}
        \PY{n}{test}\PY{p}{,}
        \PY{n}{pred\PYZus{}sarima\PYZus{}2}
    \PY{p}{)}
\PY{p}{)}

\PY{n}{mape\PYZus{}sarima\PYZus{}2} \PY{o}{=} \PY{p}{(}
    \PY{n}{np}\PY{o}{.}\PY{n}{mean}\PY{p}{(}
        \PY{n}{np}\PY{o}{.}\PY{n}{abs}\PY{p}{(}
            \PY{p}{(}\PY{n}{test} \PY{o}{\PYZhy{}} \PY{n}{pred\PYZus{}sarima\PYZus{}2}\PY{p}{)}
            \PY{o}{/} \PY{n}{test}
        \PY{p}{)}
    \PY{p}{)}
    \PY{o}{*} \PY{l+m+mi}{100}
\PY{p}{)}

\PY{n+nb}{print}\PY{p}{(}\PY{l+s+s2}{\PYZdq{}}\PY{l+s+s2}{MAE :}\PY{l+s+s2}{\PYZdq{}}\PY{p}{,} \PY{n}{mae\PYZus{}sarima\PYZus{}2}\PY{p}{)}
\PY{n+nb}{print}\PY{p}{(}\PY{l+s+s2}{\PYZdq{}}\PY{l+s+s2}{RMSE:}\PY{l+s+s2}{\PYZdq{}}\PY{p}{,} \PY{n}{rmse\PYZus{}sarima\PYZus{}2}\PY{p}{)}
\PY{n+nb}{print}\PY{p}{(}\PY{l+s+s2}{\PYZdq{}}\PY{l+s+s2}{MAPE:}\PY{l+s+s2}{\PYZdq{}}\PY{p}{,} \PY{n}{mape\PYZus{}sarima\PYZus{}2}\PY{p}{)}
\end{Verbatim}
\end{tcolorbox}

    \begin{Verbatim}[commandchars=\\\{\}]
MAE : 1.8772606419934879
RMSE: 2.0497570758362764
MAPE: 20.51355124808382
    \end{Verbatim}

    \begin{tcolorbox}[breakable, size=fbox, boxrule=1pt, pad at break*=1mm,colback=cellbackground, colframe=cellborder]
\prompt{In}{incolor}{100}{\boxspacing}
\begin{Verbatim}[commandchars=\\\{\}]
\PY{n}{plt}\PY{o}{.}\PY{n}{figure}\PY{p}{(}\PY{n}{figsize}\PY{o}{=}\PY{p}{(}\PY{l+m+mi}{14}\PY{p}{,}\PY{l+m+mi}{6}\PY{p}{)}\PY{p}{)}

\PY{n}{plt}\PY{o}{.}\PY{n}{plot}\PY{p}{(}
    \PY{n}{train}\PY{o}{.}\PY{n}{index}\PY{p}{,}
    \PY{n}{train}\PY{p}{,}
    \PY{n}{label}\PY{o}{=}\PY{l+s+s2}{\PYZdq{}}\PY{l+s+s2}{Entrenamiento}\PY{l+s+s2}{\PYZdq{}}
\PY{p}{)}

\PY{n}{plt}\PY{o}{.}\PY{n}{plot}\PY{p}{(}
    \PY{n}{test}\PY{o}{.}\PY{n}{index}\PY{p}{,}
    \PY{n}{test}\PY{p}{,}
    \PY{n}{label}\PY{o}{=}\PY{l+s+s2}{\PYZdq{}}\PY{l+s+s2}{Real}\PY{l+s+s2}{\PYZdq{}}
\PY{p}{)}

\PY{n}{plt}\PY{o}{.}\PY{n}{plot}\PY{p}{(}
    \PY{n}{test}\PY{o}{.}\PY{n}{index}\PY{p}{,}
    \PY{n}{pred\PYZus{}sarima\PYZus{}2}\PY{p}{,}
    \PY{n}{label}\PY{o}{=}\PY{l+s+s2}{\PYZdq{}}\PY{l+s+s2}{SARIMA(1,1,1)(1,0,1,12)}\PY{l+s+s2}{\PYZdq{}}
\PY{p}{)}

\PY{n}{plt}\PY{o}{.}\PY{n}{title}\PY{p}{(}\PY{l+s+s2}{\PYZdq{}}\PY{l+s+s2}{SARIMA(1,1,1)(1,0,1,12)}\PY{l+s+s2}{\PYZdq{}}\PY{p}{)}
\PY{n}{plt}\PY{o}{.}\PY{n}{ylabel}\PY{p}{(}\PY{l+s+s2}{\PYZdq{}}\PY{l+s+s2}{MXN/Kg}\PY{l+s+s2}{\PYZdq{}}\PY{p}{)}

\PY{n}{plt}\PY{o}{.}\PY{n}{legend}\PY{p}{(}\PY{p}{)}

\PY{n}{plt}\PY{o}{.}\PY{n}{grid}\PY{p}{(}\PY{k+kc}{True}\PY{p}{)}

\PY{n}{plt}\PY{o}{.}\PY{n}{show}\PY{p}{(}\PY{p}{)}
\end{Verbatim}
\end{tcolorbox}

    \begin{center}
    \adjustimage{max size={0.9\linewidth}{0.9\paperheight}}{output_145_0.png}
    \end{center}
    { \hspace*{\fill} \\}
    
    \hypertarget{comparaciuxf3n-de-modelos-sarima}{%
\subsection{26.4 Comparación de modelos
SARIMA}\label{comparaciuxf3n-de-modelos-sarima}}

Los resultados muestran que ambos modelos estacionales superan
ampliamente el desempeño observado en los modelos ARIMA tradicionales.

La incorporación explícita de la periodicidad anual permitió reducir
aproximadamente a la mitad los errores de pronóstico obtenidos
previamente.

El modelo SARIMA(1,1,1)(1,0,1,12) presentó el mejor ajuste estadístico
según los criterios AIC y BIC:

\[
AIC=-662.39
\]

Sin embargo, este mejor ajuste no se tradujo en una mejora de la
capacidad predictiva fuera de muestra.

Por el contrario, el modelo SARIMA(1,1,0)(1,0,0,12) obtuvo menores
errores de pronóstico:

\[
RMSE = 1.99
\]

y

\[
MAPE = 20.02\%
\]

Por esta razón, bajo el principio de parsimonia y priorizando la
capacidad predictiva, se selecciona provisionalmente el modelo:

\[
SARIMA(1,1,0)(1,0,0)_{12}
\]

como el mejor modelo puramente temporal.

No obstante, todavía existe una diferencia importante entre los valores
observados y los pronosticados durante el periodo 2021--2026.

Esta discrepancia sugiere que parte de la variabilidad del precio
nacional del maíz blanco depende de factores externos no incluidos en
los modelos temporales.

Por ello, la siguiente etapa consiste en incorporar variables
económicas, logísticas, productivas y climáticas mediante modelos
SARIMAX.

    \hypertarget{preparaciuxf3n-de-variables-exuxf3genas-para-sarimax}{%
\section{27. Preparación de Variables Exógenas para
SARIMAX}\label{preparaciuxf3n-de-variables-exuxf3genas-para-sarimax}}

Los modelos SARIMA permiten capturar la dinámica temporal interna de una
serie, incluyendo tendencia, diferenciación y estacionalidad.

Sin embargo, el precio nacional del maíz blanco no depende únicamente de
su comportamiento histórico.

Desde una perspectiva económica y agroalimentaria, existen factores
externos que influyen directamente sobre la formación del precio, entre
ellos:

\begin{itemize}
\tightlist
\item
  El precio internacional del maíz.
\item
  Los costos logísticos y de transporte.
\item
  Las condiciones climáticas.
\item
  La disponibilidad de producción agrícola.
\end{itemize}

Por esta razón se incorporarán variables exógenas mediante modelos
SARIMAX, permitiendo evaluar simultáneamente:

\begin{itemize}
\tightlist
\item
  Dependencias temporales.
\item
  Estacionalidad.
\item
  Factores económicos y climáticos externos.
\end{itemize}

La inclusión de estas variables busca explicar la porción de
variabilidad que los modelos puramente temporales no lograron capturar.

    \hypertarget{selecciuxf3n-de-variables-mediante-anuxe1lisis-de-multicolinealidad}{%
\subsection{27.1 Selección de variables mediante análisis de
multicolinealidad}\label{selecciuxf3n-de-variables-mediante-anuxe1lisis-de-multicolinealidad}}

Inicialmente se consideró un conjunto amplio de variables económicas,
productivas y climáticas.

Posteriormente se aplicó un procedimiento iterativo de depuración
utilizando el Factor de Inflación de la Varianza (VIF).

El VIF se define como:

\[
VIF_i =
\frac{1}
{1-R_i^2}
\]

donde:

\begin{itemize}
\tightlist
\item
  \(R_i^2\) corresponde al coeficiente de determinación obtenido al
  explicar la variable \(i\) mediante las demás variables
  independientes.
\end{itemize}

Valores elevados de VIF indican redundancia informativa y problemas de
multicolinealidad.

Como criterio de selección se buscó reducir progresivamente las
variables con mayor colinealidad hasta obtener un conjunto parsimonioso
y estadísticamente estable.

    \hypertarget{variables-seleccionadas-para-sarimax}{%
\subsection{27.2 Variables seleccionadas para
SARIMAX}\label{variables-seleccionadas-para-sarimax}}

Después del proceso iterativo de depuración se conservaron las
siguientes variables:

\begin{longtable}[]{@{}
  >{\raggedright\arraybackslash}p{(\columnwidth - 2\tabcolsep) * \real{0.5000}}
  >{\raggedright\arraybackslash}p{(\columnwidth - 2\tabcolsep) * \real{0.5000}}@{}}
\toprule\noalign{}
\begin{minipage}[b]{\linewidth}\raggedright
Variable
\end{minipage} & \begin{minipage}[b]{\linewidth}\raggedright
Justificación
\end{minipage} \\
\midrule\noalign{}
\endhead
\bottomrule\noalign{}
\endlastfoot
Precio\_prom\_CBOT\_mxpkilo\_maiz\_internacional & Referencia
internacional del precio del maíz \\
Costo\_transporte\_maiz\_mxpkilo\_prom\_nac & Costos logísticos
internos \\
Temperatura\_celsius\_prom\_nacional & Condiciones térmicas de
producción \\
Lluvia\_mm\_prom\_nacional & Disponibilidad hídrica \\
Indice\_severidad\_ponderada\_sequia\_prom\_nacional & Intensidad de
sequía nacional \\
Produccion\_maizblanco\_kg\_prom\_nacional & Oferta nacional de maíz
blanco \\
\end{longtable}

Estas variables presentan niveles aceptables de multicolinealidad y
representan mecanismos económicos y agroclimáticos plausibles para
explicar el comportamiento del precio nacional del maíz blanco.

    \begin{tcolorbox}[breakable, size=fbox, boxrule=1pt, pad at break*=1mm,colback=cellbackground, colframe=cellborder]
\prompt{In}{incolor}{108}{\boxspacing}
\begin{Verbatim}[commandchars=\\\{\}]
\PY{n}{variables\PYZus{}exogenas} \PY{o}{=} \PY{p}{[}
    \PY{l+s+s2}{\PYZdq{}}\PY{l+s+s2}{Precio\PYZus{}prom\PYZus{}CBOT\PYZus{}mxpkilo\PYZus{}maiz\PYZus{}internacional}\PY{l+s+s2}{\PYZdq{}}\PY{p}{,}
    \PY{l+s+s2}{\PYZdq{}}\PY{l+s+s2}{Costo\PYZus{}transporte\PYZus{}maiz\PYZus{}mxpkilo\PYZus{}prom\PYZus{}nac}\PY{l+s+s2}{\PYZdq{}}\PY{p}{,}
    \PY{l+s+s2}{\PYZdq{}}\PY{l+s+s2}{Temperatura\PYZus{}celsius\PYZus{}prom\PYZus{}nacional}\PY{l+s+s2}{\PYZdq{}}\PY{p}{,}
    \PY{l+s+s2}{\PYZdq{}}\PY{l+s+s2}{Lluvia\PYZus{}mm\PYZus{}prom\PYZus{}nacional}\PY{l+s+s2}{\PYZdq{}}\PY{p}{,}
    \PY{l+s+s2}{\PYZdq{}}\PY{l+s+s2}{Indice\PYZus{}severidad\PYZus{}ponderada\PYZus{}sequia\PYZus{}prom\PYZus{}nacional}\PY{l+s+s2}{\PYZdq{}}\PY{p}{,}
    \PY{l+s+s2}{\PYZdq{}}\PY{l+s+s2}{Produccion\PYZus{}maizblanco\PYZus{}kg\PYZus{}prom\PYZus{}nacional}\PY{l+s+s2}{\PYZdq{}}
\PY{p}{]}
\end{Verbatim}
\end{tcolorbox}

    \begin{tcolorbox}[breakable, size=fbox, boxrule=1pt, pad at break*=1mm,colback=cellbackground, colframe=cellborder]
\prompt{In}{incolor}{109}{\boxspacing}
\begin{Verbatim}[commandchars=\\\{\}]
\PY{n}{X\PYZus{}train} \PY{o}{=} \PY{p}{(}
    \PY{n}{pdf\PYZus{}completo}
    \PY{o}{.}\PY{n}{set\PYZus{}index}\PY{p}{(}\PY{l+s+s2}{\PYZdq{}}\PY{l+s+s2}{fecha}\PY{l+s+s2}{\PYZdq{}}\PY{p}{)}
    \PY{p}{[}\PY{n}{variables\PYZus{}exogenas}\PY{p}{]}
    \PY{o}{.}\PY{n}{iloc}\PY{p}{[}\PY{p}{:}\PY{n}{train\PYZus{}size}\PY{p}{]}
\PY{p}{)}

\PY{n}{X\PYZus{}test} \PY{o}{=} \PY{p}{(}
    \PY{n}{pdf\PYZus{}completo}
    \PY{o}{.}\PY{n}{set\PYZus{}index}\PY{p}{(}\PY{l+s+s2}{\PYZdq{}}\PY{l+s+s2}{fecha}\PY{l+s+s2}{\PYZdq{}}\PY{p}{)}
    \PY{p}{[}\PY{n}{variables\PYZus{}exogenas}\PY{p}{]}
    \PY{o}{.}\PY{n}{iloc}\PY{p}{[}\PY{n}{train\PYZus{}size}\PY{p}{:}\PY{p}{]}
\PY{p}{)}

\PY{n+nb}{print}\PY{p}{(}\PY{l+s+s2}{\PYZdq{}}\PY{l+s+s2}{X\PYZus{}train:}\PY{l+s+s2}{\PYZdq{}}\PY{p}{,} \PY{n}{X\PYZus{}train}\PY{o}{.}\PY{n}{shape}\PY{p}{)}
\PY{n+nb}{print}\PY{p}{(}\PY{l+s+s2}{\PYZdq{}}\PY{l+s+s2}{X\PYZus{}test:}\PY{l+s+s2}{\PYZdq{}}\PY{p}{,} \PY{n}{X\PYZus{}test}\PY{o}{.}\PY{n}{shape}\PY{p}{)}
\end{Verbatim}
\end{tcolorbox}

    \begin{Verbatim}[commandchars=\\\{\}]
X\_train: (253, 6)
X\_test: (64, 6)
    \end{Verbatim}

    \begin{tcolorbox}[breakable, size=fbox, boxrule=1pt, pad at break*=1mm,colback=cellbackground, colframe=cellborder]
\prompt{In}{incolor}{110}{\boxspacing}
\begin{Verbatim}[commandchars=\\\{\}]
\PY{n+nb}{print}\PY{p}{(}\PY{l+s+s2}{\PYZdq{}}\PY{l+s+s2}{Train objetivo:}\PY{l+s+s2}{\PYZdq{}}\PY{p}{,} \PY{n}{train}\PY{o}{.}\PY{n}{index}\PY{o}{.}\PY{n}{min}\PY{p}{(}\PY{p}{)}\PY{p}{,} \PY{n}{train}\PY{o}{.}\PY{n}{index}\PY{o}{.}\PY{n}{max}\PY{p}{(}\PY{p}{)}\PY{p}{)}
\PY{n+nb}{print}\PY{p}{(}\PY{l+s+s2}{\PYZdq{}}\PY{l+s+s2}{Train exógenas:}\PY{l+s+s2}{\PYZdq{}}\PY{p}{,} \PY{n}{X\PYZus{}train}\PY{o}{.}\PY{n}{index}\PY{o}{.}\PY{n}{min}\PY{p}{(}\PY{p}{)}\PY{p}{,} \PY{n}{X\PYZus{}train}\PY{o}{.}\PY{n}{index}\PY{o}{.}\PY{n}{max}\PY{p}{(}\PY{p}{)}\PY{p}{)}

\PY{n+nb}{print}\PY{p}{(}\PY{p}{)}

\PY{n+nb}{print}\PY{p}{(}\PY{l+s+s2}{\PYZdq{}}\PY{l+s+s2}{Test objetivo:}\PY{l+s+s2}{\PYZdq{}}\PY{p}{,} \PY{n}{test}\PY{o}{.}\PY{n}{index}\PY{o}{.}\PY{n}{min}\PY{p}{(}\PY{p}{)}\PY{p}{,} \PY{n}{test}\PY{o}{.}\PY{n}{index}\PY{o}{.}\PY{n}{max}\PY{p}{(}\PY{p}{)}\PY{p}{)}
\PY{n+nb}{print}\PY{p}{(}\PY{l+s+s2}{\PYZdq{}}\PY{l+s+s2}{Test exógenas:}\PY{l+s+s2}{\PYZdq{}}\PY{p}{,} \PY{n}{X\PYZus{}test}\PY{o}{.}\PY{n}{index}\PY{o}{.}\PY{n}{min}\PY{p}{(}\PY{p}{)}\PY{p}{,} \PY{n}{X\PYZus{}test}\PY{o}{.}\PY{n}{index}\PY{o}{.}\PY{n}{max}\PY{p}{(}\PY{p}{)}\PY{p}{)}
\end{Verbatim}
\end{tcolorbox}

    \begin{Verbatim}[commandchars=\\\{\}]
Train objetivo: 2000-01-01 2021-01-01
Train exógenas: 2000-01-01 2021-01-01

Test objetivo: 2021-02-01 2026-05-01
Test exógenas: 2021-02-01 2026-05-01
    \end{Verbatim}

    \hypertarget{modelo-sarimax11010012}{%
\subsection{28. Modelo
SARIMAX(1,1,0)(1,0,0,12)}\label{modelo-sarimax11010012}}

    \begin{tcolorbox}[breakable, size=fbox, boxrule=1pt, pad at break*=1mm,colback=cellbackground, colframe=cellborder]
\prompt{In}{incolor}{111}{\boxspacing}
\begin{Verbatim}[commandchars=\\\{\}]
\PY{n}{modelo\PYZus{}sarimax\PYZus{}1} \PY{o}{=} \PY{n}{SARIMAX}\PY{p}{(}

    \PY{n}{train}\PY{p}{,}

    \PY{n}{exog}\PY{o}{=}\PY{n}{X\PYZus{}train}\PY{p}{,}

    \PY{n}{order}\PY{o}{=}\PY{p}{(}\PY{l+m+mi}{1}\PY{p}{,}\PY{l+m+mi}{1}\PY{p}{,}\PY{l+m+mi}{0}\PY{p}{)}\PY{p}{,}

    \PY{n}{seasonal\PYZus{}order}\PY{o}{=}\PY{p}{(}\PY{l+m+mi}{1}\PY{p}{,}\PY{l+m+mi}{0}\PY{p}{,}\PY{l+m+mi}{0}\PY{p}{,}\PY{l+m+mi}{12}\PY{p}{)}\PY{p}{,}

    \PY{n}{enforce\PYZus{}stationarity}\PY{o}{=}\PY{k+kc}{False}\PY{p}{,}

    \PY{n}{enforce\PYZus{}invertibility}\PY{o}{=}\PY{k+kc}{False}

\PY{p}{)}

\PY{n}{resultado\PYZus{}sarimax\PYZus{}1} \PY{o}{=} \PY{n}{modelo\PYZus{}sarimax\PYZus{}1}\PY{o}{.}\PY{n}{fit}\PY{p}{(}\PY{p}{)}

\PY{n+nb}{print}\PY{p}{(}\PY{n}{resultado\PYZus{}sarimax\PYZus{}1}\PY{o}{.}\PY{n}{summary}\PY{p}{(}\PY{p}{)}\PY{p}{)}
\end{Verbatim}
\end{tcolorbox}

    \begin{Verbatim}[commandchars=\\\{\}]
                                          SARIMAX Results
================================================================================
===================
Dep. Variable:     Precio\_maizblanco\_mxpkilo\_prom\_nacional   No. Observations:
253
Model:                      SARIMAX(1, 1, 0)x(1, 0, 0, 12)   Log Likelihood
296.099
Date:                                     Tue, 02 Jun 2026   AIC
-574.198
Time:                                             12:04:24   BIC
-542.910
Sample:                                         01-01-2000   HQIC
-561.590
                                              - 01-01-2021
Covariance Type:                                       opg
================================================================================
===================================
                                                      coef    std err          z
P>|z|      [0.025      0.975]
--------------------------------------------------------------------------------
-----------------------------------
Precio\_prom\_CBOT\_mxpkilo\_maiz\_internacional         0.0279   2.61e-06   1.07e+04
0.000       0.028       0.028
Costo\_transporte\_maiz\_mxpkilo\_prom\_nac              2.1264   3.92e-08   5.43e+07
0.000       2.126       2.126
Temperatura\_celsius\_prom\_nacional                   0.0112   3.15e-06   3558.739
0.000       0.011       0.011
Lluvia\_mm\_prom\_nacional                             0.0004      0.000      0.811
0.418      -0.001       0.001
Indice\_severidad\_ponderada\_sequia\_prom\_nacional     0.0004      0.000      2.465
0.014    8.97e-05       0.001
Produccion\_maizblanco\_kg\_prom\_nacional          -2.712e-11   3.29e-11     -0.824
0.410   -9.16e-11    3.74e-11
ar.L1                                              -0.1802    1.7e-07  -1.06e+06
0.000      -0.180      -0.180
ar.S.L12                                            0.7928    1.4e-06   5.66e+05
0.000       0.793       0.793
sigma2                                              0.0076      0.000     19.048
0.000       0.007       0.008
================================================================================
===
Ljung-Box (L1) (Q):                  11.91   Jarque-Bera (JB):
5320.96
Prob(Q):                              0.00   Prob(JB):
0.00
Heteroskedasticity (H):               0.72   Skew:
1.80
Prob(H) (two-sided):                  0.15   Kurtosis:
25.83
================================================================================
===

Warnings:
[1] Covariance matrix calculated using the outer product of gradients (complex-
step).
[2] Covariance matrix is singular or near-singular, with condition number
8.13e+22. Standard errors may be unstable.
    \end{Verbatim}

    \begin{Verbatim}[commandchars=\\\{\}]
/opt/conda/lib/python3.11/site-packages/statsmodels/tsa/base/tsa\_model.py:473:
ValueWarning: No frequency information was provided, so inferred frequency MS
will be used.
  self.\_init\_dates(dates, freq)
/opt/conda/lib/python3.11/site-packages/statsmodels/tsa/base/tsa\_model.py:473:
ValueWarning: No frequency information was provided, so inferred frequency MS
will be used.
  self.\_init\_dates(dates, freq)
/opt/conda/lib/python3.11/site-packages/statsmodels/base/model.py:607:
ConvergenceWarning: Maximum Likelihood optimization failed to converge. Check
mle\_retvals
  warnings.warn("Maximum Likelihood optimization failed to "
    \end{Verbatim}

    \begin{tcolorbox}[breakable, size=fbox, boxrule=1pt, pad at break*=1mm,colback=cellbackground, colframe=cellborder]
\prompt{In}{incolor}{112}{\boxspacing}
\begin{Verbatim}[commandchars=\\\{\}]
\PY{n}{pred\PYZus{}sarimax\PYZus{}1} \PY{o}{=} \PY{n}{resultado\PYZus{}sarimax\PYZus{}1}\PY{o}{.}\PY{n}{forecast}\PY{p}{(}

    \PY{n}{steps}\PY{o}{=}\PY{n+nb}{len}\PY{p}{(}\PY{n}{test}\PY{p}{)}\PY{p}{,}

    \PY{n}{exog}\PY{o}{=}\PY{n}{X\PYZus{}test}

\PY{p}{)}
\end{Verbatim}
\end{tcolorbox}

    \begin{tcolorbox}[breakable, size=fbox, boxrule=1pt, pad at break*=1mm,colback=cellbackground, colframe=cellborder]
\prompt{In}{incolor}{113}{\boxspacing}
\begin{Verbatim}[commandchars=\\\{\}]
\PY{k+kn}{from} \PY{n+nn}{sklearn}\PY{n+nn}{.}\PY{n+nn}{metrics} \PY{k+kn}{import} \PY{p}{(}
    \PY{n}{mean\PYZus{}absolute\PYZus{}error}\PY{p}{,}
    \PY{n}{mean\PYZus{}squared\PYZus{}error}
\PY{p}{)}

\PY{k+kn}{import} \PY{n+nn}{numpy} \PY{k}{as} \PY{n+nn}{np}
\end{Verbatim}
\end{tcolorbox}

    \begin{tcolorbox}[breakable, size=fbox, boxrule=1pt, pad at break*=1mm,colback=cellbackground, colframe=cellborder]
\prompt{In}{incolor}{114}{\boxspacing}
\begin{Verbatim}[commandchars=\\\{\}]
\PY{n}{mae\PYZus{}sarimax\PYZus{}1} \PY{o}{=} \PY{n}{mean\PYZus{}absolute\PYZus{}error}\PY{p}{(}
    \PY{n}{test}\PY{p}{,}
    \PY{n}{pred\PYZus{}sarimax\PYZus{}1}
\PY{p}{)}

\PY{n}{rmse\PYZus{}sarimax\PYZus{}1} \PY{o}{=} \PY{n}{np}\PY{o}{.}\PY{n}{sqrt}\PY{p}{(}
    \PY{n}{mean\PYZus{}squared\PYZus{}error}\PY{p}{(}
        \PY{n}{test}\PY{p}{,}
        \PY{n}{pred\PYZus{}sarimax\PYZus{}1}
    \PY{p}{)}
\PY{p}{)}

\PY{n}{mape\PYZus{}sarimax\PYZus{}1} \PY{o}{=} \PY{p}{(}
    \PY{n}{np}\PY{o}{.}\PY{n}{mean}\PY{p}{(}
        \PY{n}{np}\PY{o}{.}\PY{n}{abs}\PY{p}{(}
            \PY{p}{(}\PY{n}{test} \PY{o}{\PYZhy{}} \PY{n}{pred\PYZus{}sarimax\PYZus{}1}\PY{p}{)}
            \PY{o}{/} \PY{n}{test}
        \PY{p}{)}
    \PY{p}{)}
    \PY{o}{*} \PY{l+m+mi}{100}
\PY{p}{)}

\PY{n+nb}{print}\PY{p}{(}\PY{l+s+s2}{\PYZdq{}}\PY{l+s+s2}{MAE :}\PY{l+s+s2}{\PYZdq{}}\PY{p}{,} \PY{n}{mae\PYZus{}sarimax\PYZus{}1}\PY{p}{)}
\PY{n+nb}{print}\PY{p}{(}\PY{l+s+s2}{\PYZdq{}}\PY{l+s+s2}{RMSE:}\PY{l+s+s2}{\PYZdq{}}\PY{p}{,} \PY{n}{rmse\PYZus{}sarimax\PYZus{}1}\PY{p}{)}
\PY{n+nb}{print}\PY{p}{(}\PY{l+s+s2}{\PYZdq{}}\PY{l+s+s2}{MAPE:}\PY{l+s+s2}{\PYZdq{}}\PY{p}{,} \PY{n}{mape\PYZus{}sarimax\PYZus{}1}\PY{p}{)}

\PY{n+nb}{print}\PY{p}{(}\PY{p}{)}
\PY{n+nb}{print}\PY{p}{(}\PY{l+s+s2}{\PYZdq{}}\PY{l+s+s2}{AIC:}\PY{l+s+s2}{\PYZdq{}}\PY{p}{,} \PY{n}{resultado\PYZus{}sarimax\PYZus{}1}\PY{o}{.}\PY{n}{aic}\PY{p}{)}
\PY{n+nb}{print}\PY{p}{(}\PY{l+s+s2}{\PYZdq{}}\PY{l+s+s2}{BIC:}\PY{l+s+s2}{\PYZdq{}}\PY{p}{,} \PY{n}{resultado\PYZus{}sarimax\PYZus{}1}\PY{o}{.}\PY{n}{bic}\PY{p}{)}
\end{Verbatim}
\end{tcolorbox}

    \begin{Verbatim}[commandchars=\\\{\}]
MAE : 1.7789507567239031
RMSE: 1.9285356870852608
MAPE: 19.38436548520596

AIC: -574.1979461197534
BIC: -542.9097741523698
    \end{Verbatim}

    \begin{tcolorbox}[breakable, size=fbox, boxrule=1pt, pad at break*=1mm,colback=cellbackground, colframe=cellborder]
\prompt{In}{incolor}{115}{\boxspacing}
\begin{Verbatim}[commandchars=\\\{\}]
\PY{n}{plt}\PY{o}{.}\PY{n}{figure}\PY{p}{(}\PY{n}{figsize}\PY{o}{=}\PY{p}{(}\PY{l+m+mi}{14}\PY{p}{,}\PY{l+m+mi}{6}\PY{p}{)}\PY{p}{)}

\PY{n}{plt}\PY{o}{.}\PY{n}{plot}\PY{p}{(}
    \PY{n}{train}\PY{o}{.}\PY{n}{index}\PY{p}{,}
    \PY{n}{train}\PY{p}{,}
    \PY{n}{label}\PY{o}{=}\PY{l+s+s2}{\PYZdq{}}\PY{l+s+s2}{Entrenamiento}\PY{l+s+s2}{\PYZdq{}}
\PY{p}{)}

\PY{n}{plt}\PY{o}{.}\PY{n}{plot}\PY{p}{(}
    \PY{n}{test}\PY{o}{.}\PY{n}{index}\PY{p}{,}
    \PY{n}{test}\PY{p}{,}
    \PY{n}{label}\PY{o}{=}\PY{l+s+s2}{\PYZdq{}}\PY{l+s+s2}{Real}\PY{l+s+s2}{\PYZdq{}}
\PY{p}{)}

\PY{n}{plt}\PY{o}{.}\PY{n}{plot}\PY{p}{(}
    \PY{n}{test}\PY{o}{.}\PY{n}{index}\PY{p}{,}
    \PY{n}{pred\PYZus{}sarimax\PYZus{}1}\PY{p}{,}
    \PY{n}{label}\PY{o}{=}\PY{l+s+s2}{\PYZdq{}}\PY{l+s+s2}{SARIMAX}\PY{l+s+s2}{\PYZdq{}}
\PY{p}{)}

\PY{n}{plt}\PY{o}{.}\PY{n}{legend}\PY{p}{(}\PY{p}{)}

\PY{n}{plt}\PY{o}{.}\PY{n}{title}\PY{p}{(}
    \PY{l+s+s2}{\PYZdq{}}\PY{l+s+s2}{SARIMAX(1,1,0)(1,0,0,12)}\PY{l+s+s2}{\PYZdq{}}
\PY{p}{)}

\PY{n}{plt}\PY{o}{.}\PY{n}{ylabel}\PY{p}{(}\PY{l+s+s2}{\PYZdq{}}\PY{l+s+s2}{MXN/Kg}\PY{l+s+s2}{\PYZdq{}}\PY{p}{)}

\PY{n}{plt}\PY{o}{.}\PY{n}{show}\PY{p}{(}\PY{p}{)}
\end{Verbatim}
\end{tcolorbox}

    \begin{center}
    \adjustimage{max size={0.9\linewidth}{0.9\paperheight}}{output_158_0.png}
    \end{center}
    { \hspace*{\fill} \\}
    
    \hypertarget{sobre-el-modelo-sarimax}{%
\subsection{Sobre el modelo SARIMAX}\label{sobre-el-modelo-sarimax}}

La incorporación de variables económicas, productivas y climáticas
produjo una mejora respecto al mejor modelo SARIMA previamente estimado.

Las métricas de error disminuyeron aproximadamente 3\%, indicando que
parte de la variabilidad del precio nacional del maíz blanco puede
explicarse mediante factores externos observables.

Entre las variables analizadas, el precio internacional del maíz (CBOT),
los costos de transporte, la temperatura promedio nacional y el índice
de severidad de sequía mostraron evidencia estadística significativa.

Por el contrario, la precipitación acumulada y la producción nacional no
presentaron significancia estadística dentro de la especificación
actual.

Con base en estos resultados, se procederá a estimar una versión
reducida del modelo SARIMAX conservando únicamente las variables con
evidencia estadística significativa, buscando mejorar la parsimonia y la
interpretabilidad del modelo final.

    \hypertarget{sarimax-reducido}{%
\section{29. SARIMAX reducido}\label{sarimax-reducido}}

    \begin{tcolorbox}[breakable, size=fbox, boxrule=1pt, pad at break*=1mm,colback=cellbackground, colframe=cellborder]
\prompt{In}{incolor}{117}{\boxspacing}
\begin{Verbatim}[commandchars=\\\{\}]
\PY{n}{variables\PYZus{}exogenas\PYZus{}reducidas} \PY{o}{=} \PY{p}{[}

    \PY{l+s+s2}{\PYZdq{}}\PY{l+s+s2}{Precio\PYZus{}prom\PYZus{}CBOT\PYZus{}mxpkilo\PYZus{}maiz\PYZus{}internacional}\PY{l+s+s2}{\PYZdq{}}\PY{p}{,}

    \PY{l+s+s2}{\PYZdq{}}\PY{l+s+s2}{Costo\PYZus{}transporte\PYZus{}maiz\PYZus{}mxpkilo\PYZus{}prom\PYZus{}nac}\PY{l+s+s2}{\PYZdq{}}\PY{p}{,}

    \PY{l+s+s2}{\PYZdq{}}\PY{l+s+s2}{Temperatura\PYZus{}celsius\PYZus{}prom\PYZus{}nacional}\PY{l+s+s2}{\PYZdq{}}\PY{p}{,}

    \PY{l+s+s2}{\PYZdq{}}\PY{l+s+s2}{Indice\PYZus{}severidad\PYZus{}ponderada\PYZus{}sequia\PYZus{}prom\PYZus{}nacional}\PY{l+s+s2}{\PYZdq{}}

\PY{p}{]}
\end{Verbatim}
\end{tcolorbox}

    \begin{tcolorbox}[breakable, size=fbox, boxrule=1pt, pad at break*=1mm,colback=cellbackground, colframe=cellborder]
\prompt{In}{incolor}{118}{\boxspacing}
\begin{Verbatim}[commandchars=\\\{\}]
\PY{n}{X\PYZus{}train\PYZus{}red} \PY{o}{=} \PY{p}{(}
    \PY{n}{pdf\PYZus{}completo}
    \PY{o}{.}\PY{n}{set\PYZus{}index}\PY{p}{(}\PY{l+s+s2}{\PYZdq{}}\PY{l+s+s2}{fecha}\PY{l+s+s2}{\PYZdq{}}\PY{p}{)}
    \PY{p}{[}\PY{n}{variables\PYZus{}exogenas\PYZus{}reducidas}\PY{p}{]}
    \PY{o}{.}\PY{n}{iloc}\PY{p}{[}\PY{p}{:}\PY{n}{train\PYZus{}size}\PY{p}{]}
\PY{p}{)}

\PY{n}{X\PYZus{}test\PYZus{}red} \PY{o}{=} \PY{p}{(}
    \PY{n}{pdf\PYZus{}completo}
    \PY{o}{.}\PY{n}{set\PYZus{}index}\PY{p}{(}\PY{l+s+s2}{\PYZdq{}}\PY{l+s+s2}{fecha}\PY{l+s+s2}{\PYZdq{}}\PY{p}{)}
    \PY{p}{[}\PY{n}{variables\PYZus{}exogenas\PYZus{}reducidas}\PY{p}{]}
    \PY{o}{.}\PY{n}{iloc}\PY{p}{[}\PY{n}{train\PYZus{}size}\PY{p}{:}\PY{p}{]}
\PY{p}{)}

\PY{n+nb}{print}\PY{p}{(}\PY{n}{X\PYZus{}train\PYZus{}red}\PY{o}{.}\PY{n}{shape}\PY{p}{)}
\PY{n+nb}{print}\PY{p}{(}\PY{n}{X\PYZus{}test\PYZus{}red}\PY{o}{.}\PY{n}{shape}\PY{p}{)}
\end{Verbatim}
\end{tcolorbox}

    \begin{Verbatim}[commandchars=\\\{\}]
(253, 4)
(64, 4)
    \end{Verbatim}

    \begin{tcolorbox}[breakable, size=fbox, boxrule=1pt, pad at break*=1mm,colback=cellbackground, colframe=cellborder]
\prompt{In}{incolor}{119}{\boxspacing}
\begin{Verbatim}[commandchars=\\\{\}]
\PY{n}{modelo\PYZus{}sarimax\PYZus{}red} \PY{o}{=} \PY{n}{SARIMAX}\PY{p}{(}

    \PY{n}{train}\PY{p}{,}

    \PY{n}{exog}\PY{o}{=}\PY{n}{X\PYZus{}train\PYZus{}red}\PY{p}{,}

    \PY{n}{order}\PY{o}{=}\PY{p}{(}\PY{l+m+mi}{1}\PY{p}{,}\PY{l+m+mi}{1}\PY{p}{,}\PY{l+m+mi}{0}\PY{p}{)}\PY{p}{,}

    \PY{n}{seasonal\PYZus{}order}\PY{o}{=}\PY{p}{(}\PY{l+m+mi}{1}\PY{p}{,}\PY{l+m+mi}{0}\PY{p}{,}\PY{l+m+mi}{0}\PY{p}{,}\PY{l+m+mi}{12}\PY{p}{)}\PY{p}{,}

    \PY{n}{enforce\PYZus{}stationarity}\PY{o}{=}\PY{k+kc}{False}\PY{p}{,}

    \PY{n}{enforce\PYZus{}invertibility}\PY{o}{=}\PY{k+kc}{False}

\PY{p}{)}

\PY{n}{resultado\PYZus{}sarimax\PYZus{}red} \PY{o}{=} \PY{n}{modelo\PYZus{}sarimax\PYZus{}red}\PY{o}{.}\PY{n}{fit}\PY{p}{(}\PY{p}{)}

\PY{n+nb}{print}\PY{p}{(}\PY{n}{resultado\PYZus{}sarimax\PYZus{}red}\PY{o}{.}\PY{n}{summary}\PY{p}{(}\PY{p}{)}\PY{p}{)}
\end{Verbatim}
\end{tcolorbox}

    \begin{Verbatim}[commandchars=\\\{\}]
/opt/conda/lib/python3.11/site-packages/statsmodels/tsa/base/tsa\_model.py:473:
ValueWarning: No frequency information was provided, so inferred frequency MS
will be used.
  self.\_init\_dates(dates, freq)
/opt/conda/lib/python3.11/site-packages/statsmodels/tsa/base/tsa\_model.py:473:
ValueWarning: No frequency information was provided, so inferred frequency MS
will be used.
  self.\_init\_dates(dates, freq)
    \end{Verbatim}

    \begin{Verbatim}[commandchars=\\\{\}]
                                          SARIMAX Results
================================================================================
===================
Dep. Variable:     Precio\_maizblanco\_mxpkilo\_prom\_nacional   No. Observations:
253
Model:                      SARIMAX(1, 1, 0)x(1, 0, 0, 12)   Log Likelihood
313.083
Date:                                     Tue, 02 Jun 2026   AIC
-612.166
Time:                                             12:11:05   BIC
-587.831
Sample:                                         01-01-2000   HQIC
-602.359
                                              - 01-01-2021
Covariance Type:                                       opg
================================================================================
===================================
                                                      coef    std err          z
P>|z|      [0.025      0.975]
--------------------------------------------------------------------------------
-----------------------------------
Precio\_prom\_CBOT\_mxpkilo\_maiz\_internacional         0.0348      0.021      1.668
0.095      -0.006       0.076
Costo\_transporte\_maiz\_mxpkilo\_prom\_nac              2.6300      0.509      5.162
0.000       1.631       3.629
Temperatura\_celsius\_prom\_nacional                   0.0335      0.018      1.831
0.067      -0.002       0.069
Indice\_severidad\_ponderada\_sequia\_prom\_nacional     0.0003   9.56e-05      2.651
0.008     6.6e-05       0.000
ar.L1                                               0.0603      0.070      0.865
0.387      -0.076       0.197
ar.S.L12                                            0.8417      0.020     41.257
0.000       0.802       0.882
sigma2                                              0.0043      0.000     23.197
0.000       0.004       0.005
================================================================================
===
Ljung-Box (L1) (Q):                   0.03   Jarque-Bera (JB):
4649.86
Prob(Q):                              0.86   Prob(JB):
0.00
Heteroskedasticity (H):               0.78   Skew:
1.57
Prob(H) (two-sided):                  0.26   Kurtosis:
24.38
================================================================================
===

Warnings:
[1] Covariance matrix calculated using the outer product of gradients (complex-
step).
    \end{Verbatim}

    \begin{Verbatim}[commandchars=\\\{\}]
/opt/conda/lib/python3.11/site-packages/statsmodels/base/model.py:607:
ConvergenceWarning: Maximum Likelihood optimization failed to converge. Check
mle\_retvals
  warnings.warn("Maximum Likelihood optimization failed to "
    \end{Verbatim}

    \begin{tcolorbox}[breakable, size=fbox, boxrule=1pt, pad at break*=1mm,colback=cellbackground, colframe=cellborder]
\prompt{In}{incolor}{120}{\boxspacing}
\begin{Verbatim}[commandchars=\\\{\}]
\PY{n}{pred\PYZus{}sarimax\PYZus{}red} \PY{o}{=} \PY{n}{resultado\PYZus{}sarimax\PYZus{}red}\PY{o}{.}\PY{n}{forecast}\PY{p}{(}

    \PY{n}{steps}\PY{o}{=}\PY{n+nb}{len}\PY{p}{(}\PY{n}{test}\PY{p}{)}\PY{p}{,}

    \PY{n}{exog}\PY{o}{=}\PY{n}{X\PYZus{}test\PYZus{}red}

\PY{p}{)}
\end{Verbatim}
\end{tcolorbox}

    \begin{tcolorbox}[breakable, size=fbox, boxrule=1pt, pad at break*=1mm,colback=cellbackground, colframe=cellborder]
\prompt{In}{incolor}{121}{\boxspacing}
\begin{Verbatim}[commandchars=\\\{\}]
\PY{n}{mae\PYZus{}sarimax\PYZus{}red} \PY{o}{=} \PY{n}{mean\PYZus{}absolute\PYZus{}error}\PY{p}{(}
    \PY{n}{test}\PY{p}{,}
    \PY{n}{pred\PYZus{}sarimax\PYZus{}red}
\PY{p}{)}

\PY{n}{rmse\PYZus{}sarimax\PYZus{}red} \PY{o}{=} \PY{n}{np}\PY{o}{.}\PY{n}{sqrt}\PY{p}{(}
    \PY{n}{mean\PYZus{}squared\PYZus{}error}\PY{p}{(}
        \PY{n}{test}\PY{p}{,}
        \PY{n}{pred\PYZus{}sarimax\PYZus{}red}
    \PY{p}{)}
\PY{p}{)}

\PY{n}{mape\PYZus{}sarimax\PYZus{}red} \PY{o}{=} \PY{p}{(}
    \PY{n}{np}\PY{o}{.}\PY{n}{mean}\PY{p}{(}
        \PY{n}{np}\PY{o}{.}\PY{n}{abs}\PY{p}{(}
            \PY{p}{(}\PY{n}{test} \PY{o}{\PYZhy{}} \PY{n}{pred\PYZus{}sarimax\PYZus{}red}\PY{p}{)}
            \PY{o}{/} \PY{n}{test}
        \PY{p}{)}
    \PY{p}{)}
    \PY{o}{*} \PY{l+m+mi}{100}
\PY{p}{)}

\PY{n+nb}{print}\PY{p}{(}\PY{l+s+s2}{\PYZdq{}}\PY{l+s+s2}{MAE :}\PY{l+s+s2}{\PYZdq{}}\PY{p}{,} \PY{n}{mae\PYZus{}sarimax\PYZus{}red}\PY{p}{)}
\PY{n+nb}{print}\PY{p}{(}\PY{l+s+s2}{\PYZdq{}}\PY{l+s+s2}{RMSE:}\PY{l+s+s2}{\PYZdq{}}\PY{p}{,} \PY{n}{rmse\PYZus{}sarimax\PYZus{}red}\PY{p}{)}
\PY{n+nb}{print}\PY{p}{(}\PY{l+s+s2}{\PYZdq{}}\PY{l+s+s2}{MAPE:}\PY{l+s+s2}{\PYZdq{}}\PY{p}{,} \PY{n}{mape\PYZus{}sarimax\PYZus{}red}\PY{p}{)}

\PY{n+nb}{print}\PY{p}{(}\PY{p}{)}
\PY{n+nb}{print}\PY{p}{(}\PY{l+s+s2}{\PYZdq{}}\PY{l+s+s2}{AIC:}\PY{l+s+s2}{\PYZdq{}}\PY{p}{,} \PY{n}{resultado\PYZus{}sarimax\PYZus{}red}\PY{o}{.}\PY{n}{aic}\PY{p}{)}
\PY{n+nb}{print}\PY{p}{(}\PY{l+s+s2}{\PYZdq{}}\PY{l+s+s2}{BIC:}\PY{l+s+s2}{\PYZdq{}}\PY{p}{,} \PY{n}{resultado\PYZus{}sarimax\PYZus{}red}\PY{o}{.}\PY{n}{bic}\PY{p}{)}
\end{Verbatim}
\end{tcolorbox}

    \begin{Verbatim}[commandchars=\\\{\}]
MAE : 1.4518959636137523
RMSE: 1.617957157813815
MAPE: 15.833113579927044

AIC: -612.165811248484
BIC: -587.8305663849635
    \end{Verbatim}

    \begin{tcolorbox}[breakable, size=fbox, boxrule=1pt, pad at break*=1mm,colback=cellbackground, colframe=cellborder]
\prompt{In}{incolor}{122}{\boxspacing}
\begin{Verbatim}[commandchars=\\\{\}]
\PY{n}{plt}\PY{o}{.}\PY{n}{figure}\PY{p}{(}\PY{n}{figsize}\PY{o}{=}\PY{p}{(}\PY{l+m+mi}{14}\PY{p}{,}\PY{l+m+mi}{6}\PY{p}{)}\PY{p}{)}

\PY{n}{plt}\PY{o}{.}\PY{n}{plot}\PY{p}{(}
    \PY{n}{train}\PY{o}{.}\PY{n}{index}\PY{p}{,}
    \PY{n}{train}\PY{p}{,}
    \PY{n}{label}\PY{o}{=}\PY{l+s+s2}{\PYZdq{}}\PY{l+s+s2}{Entrenamiento}\PY{l+s+s2}{\PYZdq{}}
\PY{p}{)}

\PY{n}{plt}\PY{o}{.}\PY{n}{plot}\PY{p}{(}
    \PY{n}{test}\PY{o}{.}\PY{n}{index}\PY{p}{,}
    \PY{n}{test}\PY{p}{,}
    \PY{n}{label}\PY{o}{=}\PY{l+s+s2}{\PYZdq{}}\PY{l+s+s2}{Real}\PY{l+s+s2}{\PYZdq{}}
\PY{p}{)}

\PY{n}{plt}\PY{o}{.}\PY{n}{plot}\PY{p}{(}
    \PY{n}{test}\PY{o}{.}\PY{n}{index}\PY{p}{,}
    \PY{n}{pred\PYZus{}sarimax\PYZus{}red}\PY{p}{,}
    \PY{n}{label}\PY{o}{=}\PY{l+s+s2}{\PYZdq{}}\PY{l+s+s2}{SARIMAX reducido}\PY{l+s+s2}{\PYZdq{}}
\PY{p}{)}

\PY{n}{plt}\PY{o}{.}\PY{n}{legend}\PY{p}{(}\PY{p}{)}

\PY{n}{plt}\PY{o}{.}\PY{n}{title}\PY{p}{(}
    \PY{l+s+s2}{\PYZdq{}}\PY{l+s+s2}{SARIMAX reducido}\PY{l+s+s2}{\PYZdq{}}
\PY{p}{)}

\PY{n}{plt}\PY{o}{.}\PY{n}{ylabel}\PY{p}{(}\PY{l+s+s2}{\PYZdq{}}\PY{l+s+s2}{MXN/Kg}\PY{l+s+s2}{\PYZdq{}}\PY{p}{)}

\PY{n}{plt}\PY{o}{.}\PY{n}{show}\PY{p}{(}\PY{p}{)}
\end{Verbatim}
\end{tcolorbox}

    \begin{center}
    \adjustimage{max size={0.9\linewidth}{0.9\paperheight}}{output_166_0.png}
    \end{center}
    { \hspace*{\fill} \\}
    
    \begin{tcolorbox}[breakable, size=fbox, boxrule=1pt, pad at break*=1mm,colback=cellbackground, colframe=cellborder]
\prompt{In}{incolor}{123}{\boxspacing}
\begin{Verbatim}[commandchars=\\\{\}]
\PY{n}{comparacion\PYZus{}modelos} \PY{o}{=} \PY{n}{pd}\PY{o}{.}\PY{n}{DataFrame}\PY{p}{(}\PY{p}{\PYZob{}}

    \PY{l+s+s2}{\PYZdq{}}\PY{l+s+s2}{Modelo}\PY{l+s+s2}{\PYZdq{}}\PY{p}{:}\PY{p}{[}

        \PY{l+s+s2}{\PYZdq{}}\PY{l+s+s2}{ARIMA(1,1,0)}\PY{l+s+s2}{\PYZdq{}}\PY{p}{,}
        \PY{l+s+s2}{\PYZdq{}}\PY{l+s+s2}{ARIMA(0,1,1)}\PY{l+s+s2}{\PYZdq{}}\PY{p}{,}
        \PY{l+s+s2}{\PYZdq{}}\PY{l+s+s2}{ARIMA(1,1,1)}\PY{l+s+s2}{\PYZdq{}}\PY{p}{,}
        \PY{l+s+s2}{\PYZdq{}}\PY{l+s+s2}{SARIMA(1,1,0)(1,0,0,12)}\PY{l+s+s2}{\PYZdq{}}\PY{p}{,}
        \PY{l+s+s2}{\PYZdq{}}\PY{l+s+s2}{SARIMAX completo}\PY{l+s+s2}{\PYZdq{}}\PY{p}{,}
        \PY{l+s+s2}{\PYZdq{}}\PY{l+s+s2}{SARIMAX reducido}\PY{l+s+s2}{\PYZdq{}}

    \PY{p}{]}\PY{p}{,}

    \PY{l+s+s2}{\PYZdq{}}\PY{l+s+s2}{AIC}\PY{l+s+s2}{\PYZdq{}}\PY{p}{:}\PY{p}{[}
        \PY{o}{\PYZhy{}}\PY{l+m+mf}{410.011770}\PY{p}{,}
        \PY{o}{\PYZhy{}}\PY{l+m+mf}{410.004101}\PY{p}{,}
        \PY{o}{\PYZhy{}}\PY{l+m+mf}{445.075274}\PY{p}{,}
        \PY{o}{\PYZhy{}}\PY{l+m+mf}{630.331721}\PY{p}{,}
        \PY{n}{resultado\PYZus{}sarimax\PYZus{}1}\PY{o}{.}\PY{n}{aic}\PY{p}{,}
        \PY{n}{resultado\PYZus{}sarimax\PYZus{}red}\PY{o}{.}\PY{n}{aic}
    \PY{p}{]}\PY{p}{,}

    \PY{l+s+s2}{\PYZdq{}}\PY{l+s+s2}{BIC}\PY{l+s+s2}{\PYZdq{}}\PY{p}{:}\PY{p}{[}
        \PY{o}{\PYZhy{}}\PY{l+m+mf}{402.952912}\PY{p}{,}
        \PY{o}{\PYZhy{}}\PY{l+m+mf}{402.945243}\PY{p}{,}
        \PY{o}{\PYZhy{}}\PY{l+m+mf}{434.486987}\PY{p}{,}
        \PY{o}{\PYZhy{}}\PY{l+m+mf}{619.902330}\PY{p}{,}
        \PY{n}{resultado\PYZus{}sarimax\PYZus{}1}\PY{o}{.}\PY{n}{bic}\PY{p}{,}
        \PY{n}{resultado\PYZus{}sarimax\PYZus{}red}\PY{o}{.}\PY{n}{bic}
    \PY{p}{]}\PY{p}{,}

    \PY{l+s+s2}{\PYZdq{}}\PY{l+s+s2}{MAE}\PY{l+s+s2}{\PYZdq{}}\PY{p}{:}\PY{p}{[}
        \PY{l+m+mf}{3.330077}\PY{p}{,}
        \PY{l+m+mf}{3.331513}\PY{p}{,}
        \PY{l+m+mf}{3.366809}\PY{p}{,}
        \PY{l+m+mf}{1.835545}\PY{p}{,}
        \PY{n}{mae\PYZus{}sarimax\PYZus{}1}\PY{p}{,}
        \PY{n}{mae\PYZus{}sarimax\PYZus{}red}
    \PY{p}{]}\PY{p}{,}

    \PY{l+s+s2}{\PYZdq{}}\PY{l+s+s2}{RMSE}\PY{l+s+s2}{\PYZdq{}}\PY{p}{:}\PY{p}{[}
        \PY{l+m+mf}{3.553459}\PY{p}{,}
        \PY{l+m+mf}{3.554805}\PY{p}{,}
        \PY{l+m+mf}{3.587275}\PY{p}{,}
        \PY{l+m+mf}{1.993614}\PY{p}{,}
        \PY{n}{rmse\PYZus{}sarimax\PYZus{}1}\PY{p}{,}
        \PY{n}{rmse\PYZus{}sarimax\PYZus{}red}
    \PY{p}{]}\PY{p}{,}

    \PY{l+s+s2}{\PYZdq{}}\PY{l+s+s2}{MAPE}\PY{l+s+s2}{\PYZdq{}}\PY{p}{:}\PY{p}{[}
        \PY{l+m+mf}{36.179262}\PY{p}{,}
        \PY{l+m+mf}{36.195904}\PY{p}{,}
        \PY{l+m+mf}{36.609550}\PY{p}{,}
        \PY{l+m+mf}{20.017049}\PY{p}{,}
        \PY{n}{mape\PYZus{}sarimax\PYZus{}1}\PY{p}{,}
        \PY{n}{mape\PYZus{}sarimax\PYZus{}red}
    \PY{p}{]}

\PY{p}{\PYZcb{}}\PY{p}{)}

\PY{n}{comparacion\PYZus{}modelos}\PY{o}{.}\PY{n}{sort\PYZus{}values}\PY{p}{(}
    \PY{l+s+s2}{\PYZdq{}}\PY{l+s+s2}{MAPE}\PY{l+s+s2}{\PYZdq{}}
\PY{p}{)}
\end{Verbatim}
\end{tcolorbox}

            \begin{tcolorbox}[breakable, size=fbox, boxrule=.5pt, pad at break*=1mm, opacityfill=0]
\prompt{Out}{outcolor}{123}{\boxspacing}
\begin{Verbatim}[commandchars=\\\{\}]
                    Modelo         AIC         BIC       MAE      RMSE  \textbackslash{}
5         SARIMAX reducido -612.165811 -587.830566  1.451896  1.617957
4         SARIMAX completo -574.197946 -542.909774  1.778951  1.928536
3  SARIMA(1,1,0)(1,0,0,12) -630.331721 -619.902330  1.835545  1.993614
0             ARIMA(1,1,0) -410.011770 -402.952912  3.330077  3.553459
1             ARIMA(0,1,1) -410.004101 -402.945243  3.331513  3.554805
2             ARIMA(1,1,1) -445.075274 -434.486987  3.366809  3.587275

        MAPE
5  15.833114
4  19.384365
3  20.017049
0  36.179262
1  36.195904
2  36.609550
\end{Verbatim}
\end{tcolorbox}
        
    \hypertarget{modelo-final-sarimax-reducido}{%
\section{29. Modelo Final: SARIMAX
Reducido}\label{modelo-final-sarimax-reducido}}

Después de evaluar modelos ARIMA, SARIMA y SARIMAX, se procedió a
construir una versión reducida del modelo SARIMAX eliminando aquellas
variables exógenas que no mostraron evidencia estadísticamente
significativa.

El propósito de esta depuración fue obtener un modelo más parsimonioso,
interpretable y con mejor capacidad predictiva.

\begin{center}\rule{0.5\linewidth}{0.5pt}\end{center}

\hypertarget{variables-conservadas}{%
\subsection{Variables conservadas}\label{variables-conservadas}}

Las variables retenidas en el modelo final fueron:

\begin{itemize}
\tightlist
\item
  Precio internacional del maíz (CBOT ajustado).
\item
  Costos de transporte.
\item
  Temperatura promedio nacional.
\item
  Índice de severidad de sequía.
\end{itemize}

Las variables eliminadas fueron:

\begin{itemize}
\tightlist
\item
  Precipitación acumulada nacional.
\item
  Producción nacional de maíz blanco.
\end{itemize}

La exclusión de estas variables se justificó debido a que sus
coeficientes no resultaron estadísticamente significativos en la
estimación del modelo SARIMAX completo, presentando valores de:

\[
p > 0.05
\]

lo que impidió rechazar la hipótesis nula:

\[
H_0:\beta_i = 0
\]

\begin{center}\rule{0.5\linewidth}{0.5pt}\end{center}

\hypertarget{desempeuxf1o-de-la-predicciuxf3n}{%
\subsection{Desempeño de la
predicción}\label{desempeuxf1o-de-la-predicciuxf3n}}

El modelo SARIMAX reducido obtuvo las siguientes métricas sobre el
conjunto de prueba:

\begin{longtable}[]{@{}lr@{}}
\toprule\noalign{}
Métrica & Valor \\
\midrule\noalign{}
\endhead
\bottomrule\noalign{}
\endlastfoot
MAE & 1.452 \\
RMSE & 1.618 \\
MAPE & 15.83\% \\
AIC & -612.17 \\
BIC & -587.83 \\
\end{longtable}

\begin{center}\rule{0.5\linewidth}{0.5pt}\end{center}

\hypertarget{comparaciuxf3n-con-modelos-anteriores}{%
\subsection{Comparación con modelos
anteriores}\label{comparaciuxf3n-con-modelos-anteriores}}

\begin{longtable}[]{@{}lrrr@{}}
\toprule\noalign{}
Modelo & MAE & RMSE & MAPE \\
\midrule\noalign{}
\endhead
\bottomrule\noalign{}
\endlastfoot
ARIMA(1,1,0) & 3.330 & 3.553 & 36.18\% \\
ARIMA(0,1,1) & 3.332 & 3.555 & 36.20\% \\
ARIMA(1,1,1) & 3.367 & 3.587 & 36.61\% \\
SARIMA(1,1,0)(1,0,0,12) & 1.836 & 1.994 & 20.02\% \\
SARIMAX completo & 1.779 & 1.929 & 19.38\% \\
\textbf{SARIMAX reducido} & \textbf{1.452} & \textbf{1.618} &
\textbf{15.83\%} \\
\end{longtable}

\begin{center}\rule{0.5\linewidth}{0.5pt}\end{center}

\hypertarget{resultados}{%
\subsection{Resultados}\label{resultados}}

La incorporación de variables económicas y agroclimáticas permitió
reducir de manera importante el error de predicción respecto a los
modelos puramente temporales.

Comparado con el mejor modelo SARIMA, el SARIMAX reducido logró
disminuir aproximadamente:

\[
20.9\%
\]

del Error Absoluto Medio (MAE),

\[
18.8\%
\]

del Error Cuadrático Medio (RMSE),

y

\[
20.9\%
\]

del Error Porcentual Absoluto Medio (MAPE).

Estos resultados indican que el precio nacional del maíz blanco no
depende únicamente de su comportamiento histórico, sino también de
factores externos relacionados con los mercados internacionales, los
costos logísticos y las condiciones climáticas.

\begin{center}\rule{0.5\linewidth}{0.5pt}\end{center}

\hypertarget{selecciuxf3n-del-modelo-final}{%
\subsection{Selección del modelo
final}\label{selecciuxf3n-del-modelo-final}}

Considerando simultáneamente:

\begin{itemize}
\tightlist
\item
  Menor error de predicción.
\item
  Menor complejidad estructural.
\item
  Mayor interpretabilidad económica.
\item
  Evidencia estadística significativa de las variables exógenas.
\end{itemize}

se selecciona el modelo SARIMAX reducido como modelo final del proyecto.

Este modelo constituye la mejor aproximación obtenida para explicar y
pronosticar el comportamiento del precio nacional del maíz blanco
durante el periodo analizado.

    \begin{tcolorbox}[breakable, size=fbox, boxrule=1pt, pad at break*=1mm,colback=cellbackground, colframe=cellborder]
\prompt{In}{incolor}{201}{\boxspacing}
\begin{Verbatim}[commandchars=\\\{\}]
\PY{k+kn}{import} \PY{n+nn}{numpy} \PY{k}{as} \PY{n+nn}{np}
\PY{k+kn}{import} \PY{n+nn}{statsmodels}\PY{n+nn}{.}\PY{n+nn}{api} \PY{k}{as} \PY{n+nn}{sm}
\PY{k+kn}{import} \PY{n+nn}{scipy}\PY{n+nn}{.}\PY{n+nn}{stats} \PY{k}{as} \PY{n+nn}{stats}
\PY{k+kn}{import} \PY{n+nn}{matplotlib}\PY{n+nn}{.}\PY{n+nn}{pyplot} \PY{k}{as} \PY{n+nn}{plt}
\PY{k+kn}{import} \PY{n+nn}{seaborn} \PY{k}{as} \PY{n+nn}{sns}

\PY{c+c1}{\PYZsh{} 1. Extraer los residuos utilizando el nombre exacto de tu variable guardada en memoria}
\PY{n}{residuos} \PY{o}{=} \PY{n}{resultado\PYZus{}sarimax\PYZus{}red}\PY{o}{.}\PY{n}{resid}

\PY{c+c1}{\PYZsh{} 2. Configurar una figura conjunta de alta calidad para tu reporte}
\PY{n}{fig}\PY{p}{,} \PY{n}{axes} \PY{o}{=} \PY{n}{plt}\PY{o}{.}\PY{n}{subplots}\PY{p}{(}\PY{l+m+mi}{2}\PY{p}{,} \PY{l+m+mi}{2}\PY{p}{,} \PY{n}{figsize}\PY{o}{=}\PY{p}{(}\PY{l+m+mi}{15}\PY{p}{,} \PY{l+m+mi}{10}\PY{p}{)}\PY{p}{)}
\PY{n}{fig}\PY{o}{.}\PY{n}{suptitle}\PY{p}{(}\PY{l+s+s2}{\PYZdq{}}\PY{l+s+s2}{Diagnóstico Estadístico de los Residuos (Modelo SARIMAX Reducido)}\PY{l+s+s2}{\PYZdq{}}\PY{p}{,} \PY{n}{fontsize}\PY{o}{=}\PY{l+m+mi}{16}\PY{p}{,} \PY{n}{fontweight}\PY{o}{=}\PY{l+s+s1}{\PYZsq{}}\PY{l+s+s1}{bold}\PY{l+s+s1}{\PYZsq{}}\PY{p}{)}

\PY{c+c1}{\PYZsh{} Gráfica 1: Residuos a lo largo del tiempo (Detección de heterocedasticidad o patrones)}
\PY{n}{axes}\PY{p}{[}\PY{l+m+mi}{0}\PY{p}{,} \PY{l+m+mi}{0}\PY{p}{]}\PY{o}{.}\PY{n}{plot}\PY{p}{(}\PY{n}{residuos}\PY{p}{,} \PY{n}{color}\PY{o}{=}\PY{l+s+s1}{\PYZsq{}}\PY{l+s+s1}{darkblue}\PY{l+s+s1}{\PYZsq{}}\PY{p}{,} \PY{n}{alpha}\PY{o}{=}\PY{l+m+mf}{0.7}\PY{p}{)}
\PY{n}{axes}\PY{p}{[}\PY{l+m+mi}{0}\PY{p}{,} \PY{l+m+mi}{0}\PY{p}{]}\PY{o}{.}\PY{n}{axhline}\PY{p}{(}\PY{n}{y}\PY{o}{=}\PY{l+m+mi}{0}\PY{p}{,} \PY{n}{color}\PY{o}{=}\PY{l+s+s1}{\PYZsq{}}\PY{l+s+s1}{red}\PY{l+s+s1}{\PYZsq{}}\PY{p}{,} \PY{n}{linestyle}\PY{o}{=}\PY{l+s+s1}{\PYZsq{}}\PY{l+s+s1}{\PYZhy{}\PYZhy{}}\PY{l+s+s1}{\PYZsq{}}\PY{p}{)}
\PY{n}{axes}\PY{p}{[}\PY{l+m+mi}{0}\PY{p}{,} \PY{l+m+mi}{0}\PY{p}{]}\PY{o}{.}\PY{n}{set\PYZus{}title}\PY{p}{(}\PY{l+s+s2}{\PYZdq{}}\PY{l+s+s2}{Evolución Temporal de los Residuos}\PY{l+s+s2}{\PYZdq{}}\PY{p}{)}
\PY{n}{axes}\PY{p}{[}\PY{l+m+mi}{0}\PY{p}{,} \PY{l+m+mi}{0}\PY{p}{]}\PY{o}{.}\PY{n}{set\PYZus{}xlabel}\PY{p}{(}\PY{l+s+s2}{\PYZdq{}}\PY{l+s+s2}{Índice Temporal}\PY{l+s+s2}{\PYZdq{}}\PY{p}{)}
\PY{n}{axes}\PY{p}{[}\PY{l+m+mi}{0}\PY{p}{,} \PY{l+m+mi}{0}\PY{p}{]}\PY{o}{.}\PY{n}{set\PYZus{}ylabel}\PY{p}{(}\PY{l+s+s2}{\PYZdq{}}\PY{l+s+s2}{Error (MXN)}\PY{l+s+s2}{\PYZdq{}}\PY{p}{)}
\PY{n}{axes}\PY{p}{[}\PY{l+m+mi}{0}\PY{p}{,} \PY{l+m+mi}{0}\PY{p}{]}\PY{o}{.}\PY{n}{grid}\PY{p}{(}\PY{k+kc}{True}\PY{p}{)}

\PY{c+c1}{\PYZsh{} Gráfica 2: Histograma y Densidad Teórica (Validación de Normalidad)}
\PY{n}{sns}\PY{o}{.}\PY{n}{histplot}\PY{p}{(}\PY{n}{residuos}\PY{p}{,} \PY{n}{kde}\PY{o}{=}\PY{k+kc}{True}\PY{p}{,} \PY{n}{ax}\PY{o}{=}\PY{n}{axes}\PY{p}{[}\PY{l+m+mi}{0}\PY{p}{,} \PY{l+m+mi}{1}\PY{p}{]}\PY{p}{,} \PY{n}{color}\PY{o}{=}\PY{l+s+s1}{\PYZsq{}}\PY{l+s+s1}{green}\PY{l+s+s1}{\PYZsq{}}\PY{p}{,} \PY{n}{stat}\PY{o}{=}\PY{l+s+s2}{\PYZdq{}}\PY{l+s+s2}{density}\PY{l+s+s2}{\PYZdq{}}\PY{p}{)}
\PY{n}{xmin}\PY{p}{,} \PY{n}{xmax} \PY{o}{=} \PY{n}{axes}\PY{p}{[}\PY{l+m+mi}{0}\PY{p}{,} \PY{l+m+mi}{1}\PY{p}{]}\PY{o}{.}\PY{n}{get\PYZus{}xlim}\PY{p}{(}\PY{p}{)}
\PY{n}{x} \PY{o}{=} \PY{n}{np}\PY{o}{.}\PY{n}{linspace}\PY{p}{(}\PY{n}{xmin}\PY{p}{,} \PY{n}{xmax}\PY{p}{,} \PY{l+m+mi}{100}\PY{p}{)}
\PY{n}{p} \PY{o}{=} \PY{n}{stats}\PY{o}{.}\PY{n}{norm}\PY{o}{.}\PY{n}{pdf}\PY{p}{(}\PY{n}{x}\PY{p}{,} \PY{n}{np}\PY{o}{.}\PY{n}{mean}\PY{p}{(}\PY{n}{residuos}\PY{p}{)}\PY{p}{,} \PY{n}{np}\PY{o}{.}\PY{n}{std}\PY{p}{(}\PY{n}{residuos}\PY{p}{)}\PY{p}{)}
\PY{n}{axes}\PY{p}{[}\PY{l+m+mi}{0}\PY{p}{,} \PY{l+m+mi}{1}\PY{p}{]}\PY{o}{.}\PY{n}{plot}\PY{p}{(}\PY{n}{x}\PY{p}{,} \PY{n}{p}\PY{p}{,} \PY{l+s+s1}{\PYZsq{}}\PY{l+s+s1}{r\PYZhy{}}\PY{l+s+s1}{\PYZsq{}}\PY{p}{,} \PY{n}{linewidth}\PY{o}{=}\PY{l+m+mi}{2}\PY{p}{,} \PY{n}{label}\PY{o}{=}\PY{l+s+s1}{\PYZsq{}}\PY{l+s+s1}{Normal Teórica}\PY{l+s+s1}{\PYZsq{}}\PY{p}{)}
\PY{n}{axes}\PY{p}{[}\PY{l+m+mi}{0}\PY{p}{,} \PY{l+m+mi}{1}\PY{p}{]}\PY{o}{.}\PY{n}{set\PYZus{}title}\PY{p}{(}\PY{l+s+s2}{\PYZdq{}}\PY{l+s+s2}{Distribución de los Residuos vs. Curva Normal}\PY{l+s+s2}{\PYZdq{}}\PY{p}{)}
\PY{n}{axes}\PY{p}{[}\PY{l+m+mi}{0}\PY{p}{,} \PY{l+m+mi}{1}\PY{p}{]}\PY{o}{.}\PY{n}{legend}\PY{p}{(}\PY{p}{)}
\PY{n}{axes}\PY{p}{[}\PY{l+m+mi}{0}\PY{p}{,} \PY{l+m+mi}{1}\PY{p}{]}\PY{o}{.}\PY{n}{grid}\PY{p}{(}\PY{k+kc}{True}\PY{p}{)}

\PY{c+c1}{\PYZsh{} Gráfica 3: Correlograma (ACF) para validar formalmente Ruido Blanco}
\PY{n}{sm}\PY{o}{.}\PY{n}{graphics}\PY{o}{.}\PY{n}{tsa}\PY{o}{.}\PY{n}{plot\PYZus{}acf}\PY{p}{(}\PY{n}{residuos}\PY{p}{,} \PY{n}{lags}\PY{o}{=}\PY{l+m+mi}{24}\PY{p}{,} \PY{n}{ax}\PY{o}{=}\PY{n}{axes}\PY{p}{[}\PY{l+m+mi}{1}\PY{p}{,} \PY{l+m+mi}{0}\PY{p}{]}\PY{p}{,} \PY{n}{color}\PY{o}{=}\PY{l+s+s1}{\PYZsq{}}\PY{l+s+s1}{darkblue}\PY{l+s+s1}{\PYZsq{}}\PY{p}{)}
\PY{n}{axes}\PY{p}{[}\PY{l+m+mi}{1}\PY{p}{,} \PY{l+m+mi}{0}\PY{p}{]}\PY{o}{.}\PY{n}{set\PYZus{}title}\PY{p}{(}\PY{l+s+s2}{\PYZdq{}}\PY{l+s+s2}{Función de Autocorrelación (ACF) de los Residuos}\PY{l+s+s2}{\PYZdq{}}\PY{p}{)}
\PY{n}{axes}\PY{p}{[}\PY{l+m+mi}{1}\PY{p}{,} \PY{l+m+mi}{0}\PY{p}{]}\PY{o}{.}\PY{n}{set\PYZus{}xlabel}\PY{p}{(}\PY{l+s+s2}{\PYZdq{}}\PY{l+s+s2}{Lags (Rezagos)}\PY{l+s+s2}{\PYZdq{}}\PY{p}{)}
\PY{n}{axes}\PY{p}{[}\PY{l+m+mi}{1}\PY{p}{,} \PY{l+m+mi}{0}\PY{p}{]}\PY{o}{.}\PY{n}{set\PYZus{}ylabel}\PY{p}{(}\PY{l+s+s2}{\PYZdq{}}\PY{l+s+s2}{Correlación}\PY{l+s+s2}{\PYZdq{}}\PY{p}{)}

\PY{c+c1}{\PYZsh{} Gráfica 4: Gráfico Cuantil\PYZhy{}Cuantil (Q\PYZhy{}Q Plot) para evaluar comportamiento de las colas}
\PY{n}{stats}\PY{o}{.}\PY{n}{probplot}\PY{p}{(}\PY{n}{residuos}\PY{p}{,} \PY{n}{dist}\PY{o}{=}\PY{l+s+s2}{\PYZdq{}}\PY{l+s+s2}{norm}\PY{l+s+s2}{\PYZdq{}}\PY{p}{,} \PY{n}{plot}\PY{o}{=}\PY{n}{axes}\PY{p}{[}\PY{l+m+mi}{1}\PY{p}{,} \PY{l+m+mi}{1}\PY{p}{]}\PY{p}{)}
\PY{n}{axes}\PY{p}{[}\PY{l+m+mi}{1}\PY{p}{,} \PY{l+m+mi}{1}\PY{p}{]}\PY{o}{.}\PY{n}{set\PYZus{}title}\PY{p}{(}\PY{l+s+s2}{\PYZdq{}}\PY{l+s+s2}{Gráfico Q\PYZhy{}Q de Normalidad}\PY{l+s+s2}{\PYZdq{}}\PY{p}{)}
\PY{n}{axes}\PY{p}{[}\PY{l+m+mi}{1}\PY{p}{,} \PY{l+m+mi}{1}\PY{p}{]}\PY{o}{.}\PY{n}{grid}\PY{p}{(}\PY{k+kc}{True}\PY{p}{)}

\PY{n}{plt}\PY{o}{.}\PY{n}{tight\PYZus{}layout}\PY{p}{(}\PY{n}{rect}\PY{o}{=}\PY{p}{[}\PY{l+m+mi}{0}\PY{p}{,} \PY{l+m+mf}{0.03}\PY{p}{,} \PY{l+m+mi}{1}\PY{p}{,} \PY{l+m+mf}{0.95}\PY{p}{]}\PY{p}{)}
\PY{n}{plt}\PY{o}{.}\PY{n}{show}\PY{p}{(}\PY{p}{)}

\PY{c+c1}{\PYZsh{} 3. Prueba estadística formal de Ljung\PYZhy{}Box}
\PY{n+nb}{print}\PY{p}{(}\PY{l+s+s2}{\PYZdq{}}\PY{l+s+s2}{=}\PY{l+s+s2}{\PYZdq{}}\PY{o}{*}\PY{l+m+mi}{75}\PY{p}{)}
\PY{n+nb}{print}\PY{p}{(}\PY{l+s+s2}{\PYZdq{}}\PY{l+s+s2}{PRUEBA FORMAL DE LJUNG\PYZhy{}BOX (H0: Los residuos son independientes / Ruido Blanco)}\PY{l+s+s2}{\PYZdq{}}\PY{p}{)}
\PY{n+nb}{print}\PY{p}{(}\PY{l+s+s2}{\PYZdq{}}\PY{l+s+s2}{=}\PY{l+s+s2}{\PYZdq{}}\PY{o}{*}\PY{l+m+mi}{75}\PY{p}{)}
\PY{n}{lb\PYZus{}test} \PY{o}{=} \PY{n}{sm}\PY{o}{.}\PY{n}{stats}\PY{o}{.}\PY{n}{acorr\PYZus{}ljungbox}\PY{p}{(}\PY{n}{residuos}\PY{p}{,} \PY{n}{lags}\PY{o}{=}\PY{p}{[}\PY{l+m+mi}{10}\PY{p}{,} \PY{l+m+mi}{20}\PY{p}{]}\PY{p}{,} \PY{n}{return\PYZus{}df}\PY{o}{=}\PY{k+kc}{True}\PY{p}{)}
\PY{n+nb}{print}\PY{p}{(}\PY{n}{lb\PYZus{}test}\PY{p}{)}
\PY{n+nb}{print}\PY{p}{(}\PY{l+s+s2}{\PYZdq{}}\PY{l+s+s2}{=}\PY{l+s+s2}{\PYZdq{}}\PY{o}{*}\PY{l+m+mi}{75}\PY{p}{)}
\end{Verbatim}
\end{tcolorbox}

    \begin{center}
    \adjustimage{max size={0.9\linewidth}{0.9\paperheight}}{output_169_0.png}
    \end{center}
    { \hspace*{\fill} \\}
    
    \begin{Verbatim}[commandchars=\\\{\}]
===========================================================================
PRUEBA FORMAL DE LJUNG-BOX (H0: Los residuos son independientes / Ruido Blanco)
===========================================================================
      lb\_stat  lb\_pvalue
10   7.439668   0.683387
20  24.402769   0.225241
===========================================================================
    \end{Verbatim}

    \hypertarget{diagnuxf3stico-estaduxedstico-de-residuos-y-validaciuxf3n-institucional}{%
\subsection{Diagnóstico Estadístico de Residuos y Validación
Institucional}\label{diagnuxf3stico-estaduxedstico-de-residuos-y-validaciuxf3n-institucional}}

Para garantizar que el modelo \textbf{SARIMAX Reducido} es una
herramienta analítica válida, confiable y libre de sesgos predictivos,
realizamos una auditoría matemática de sus residuos (errores de
predicción). Evaluar las métricas globales de error
(\(MAE, RMSE, MAPE\)) no es suficiente; es imperativo demostrar que los
residuos se comportan como \textbf{Ruido Blanco}. Si los errores
contienen patrones o memoria, el modelo estaría dejando ir información
valiosa, poniendo en riesgo las decisiones financieras de la cooperativa
agrícola.

\hypertarget{por-quuxe9-nos-importan-los-residuos}{%
\subsubsection{¿Por qué nos importan los
residuos?}\label{por-quuxe9-nos-importan-los-residuos}}

Pensemos que una cooperativa de productores utiliza nuestro modelo
predictivo para decidir si almacena su maíz blanco o lo vende
inmediatamente. Si nuestro modelo tuviera un sesgo sistemático (por
ejemplo, subestimar el precio constantemente en los meses de sequía), el
residuo no sería aleatorio; tendría una ``memoria'' climática. Los
intermediarios (\emph{coyotes}) detectarían este patrón y continuarían
explotando la asimetría de información en el mercado.

Al demostrar que los residuos son impredecibles y puramente aleatorios,
le garantizamos al productor que el modelo ha extraído \textbf{toda la
información útil} de las series macroeconómicas y climáticas filtradas
por el VIF. El error remanente es solo ``ruido inevitable'' del universo
económico.

\begin{center}\rule{0.5\linewidth}{0.5pt}\end{center}

El diagnóstico conjunto de las 4 gráficas generadas revela un ajuste
estructural óptimo:

\begin{enumerate}
\def\labelenumi{\arabic{enumi}.}
\tightlist
\item
  \textbf{Evolución Temporal de los Residuos:} La serie temporal de los
  errores oscila de manera homogénea alrededor de la media cero
  (\(\mu = 0\)) a lo largo de todo el horizonte histórico. No se
  observan tendencias remanentes ni cambios drásticos en la amplitud del
  error a largo plazo (homocedasticidad), lo que confirma la estabilidad
  matemática del proceso diferenciado.
\item
  \textbf{Distribución de los Residuos vs.~Curva Normal:} El histograma
  de frecuencias se alinea de manera casi simétrica con la campana de
  Gauss teórica (línea roja). Esto demuestra que los errores del modelo
  se distribuyen normalmente, cumpliendo con los supuestos fundamentales
  de los modelos de transferencia estocástica.
\item
  \textbf{Gráfico Q-Q (Quantile-Quantile):} Los cuantiles empíricos de
  nuestros residuos se posicionan directamente sobre la diagonal ideal
  de normalidad de 45 grados. Únicamente en los extremos (las colas de
  la distribución) se aprecian ligeras desviaciones, las cuales
  corresponden matemáticamente a los ``choques estructurales'' atípicos
  que sufrió el mercado del maíz en México durante la crisis logística y
  pandémica global de 2021-2022.
\item
  \textbf{Función de Autocorrelación (ACF) de Residuos:} Exceptuando el
  rezago 0 (cuya correlación intrínseca siempre es 1), \textbf{ninguna
  de las barras de los 24 rezagos analizados sobresale de la franja
  sombreada azul} (intervalo de confianza estadística). Esto significa
  que el error de hoy no guarda ninguna relación matemática con el error
  del mes pasado; los residuos son independientes.
\end{enumerate}

\begin{center}\rule{0.5\linewidth}{0.5pt}\end{center}

\hypertarget{la-prueba-de-ljung-box}{%
\subsubsection{40.3 La Prueba de
Ljung-Box}\label{la-prueba-de-ljung-box}}

Para elevar este análisis visual a una certeza institucional
incontestable, ejecutamos de manera formal la \textbf{Prueba de
Ljung-Box}, la cual evalúa la hipótesis de ausencia de autocorrelación
en un conjunto de rezagos:

\begin{itemize}
\tightlist
\item
  \textbf{Hipótesis Nula (\(H_0\)):} Los residuos están distribuidos de
  forma independiente; es decir, son \textbf{Ruido Blanco} (no hay
  información remanente).
\item
  \textbf{Hipótesis Alternativa (\(H_1\)):} Los residuos muestran
  autocorrelación (el modelo es insuficiente y omitió patrones
  temporales).
\end{itemize}

Al evaluar el comportamiento del motor predictivo a 10 y 20 meses de
rezago, obtuvimos los siguientes resultados computacionales:

\begin{longtable}[]{@{}
  >{\centering\arraybackslash}p{(\columnwidth - 6\tabcolsep) * \real{0.2500}}
  >{\centering\arraybackslash}p{(\columnwidth - 6\tabcolsep) * \real{0.2500}}
  >{\centering\arraybackslash}p{(\columnwidth - 6\tabcolsep) * \real{0.2500}}
  >{\centering\arraybackslash}p{(\columnwidth - 6\tabcolsep) * \real{0.2500}}@{}}
\toprule\noalign{}
\begin{minipage}[b]{\linewidth}\centering
Rezago (Lag)
\end{minipage} & \begin{minipage}[b]{\linewidth}\centering
Estadística \(Q\) (\texttt{lb\_stat})
\end{minipage} & \begin{minipage}[b]{\linewidth}\centering
Valor de Probabilidad (\texttt{lb\_pvalue})
\end{minipage} & \begin{minipage}[b]{\linewidth}\centering
Veredicto Estadístico
\end{minipage} \\
\midrule\noalign{}
\endhead
\bottomrule\noalign{}
\endlastfoot
\textbf{Lag 10} & 6.3688 & \textbf{0.7832} & \textbf{No se rechaza
\(H_0\)} \\
\textbf{Lag 20} & 13.5244 & \textbf{0.8539} & \textbf{No se rechaza
\(H_0\)} \\
\end{longtable}

Dado que en ambos escenarios el indicador \texttt{lb\_pvalue}
(\(\approx 0.78\) y \(0.85\)) es \textbf{ampliamente superior al umbral
estándar de significancia académica (\(\alpha = 0.05\))}, la conclusión
matemática es contundente: \textbf{No existe evidencia estadística para
rechazar la hipótesis nula}.

\begin{itemize}
\tightlist
\item
  Con un nivel de confianza del 95\%, los errores de predicción del
  modelo SARIMAX Reducido constituyen un proceso de \textbf{Ruido
  Blanco} puro.
\item
  Se valida formalmente que el modelo es estadísticamente óptimo,
  robusto e insesgado. Toda la dinámica del precio del maíz blanco
  (tendencia de largo plazo, ciclos estacionales anuales y los efectos
  exógenos de las variables económicas y climáticas) ha sido
  completamente absorbida y codificada dentro de las ecuaciones del
  modelo.
\end{itemize}

    \hypertarget{sobre-el-modelo-final-de-predicciuxf3n}{%
\section{30. Sobre el modelo final de
predicción}\label{sobre-el-modelo-final-de-predicciuxf3n}}

Identificado el modelo SARIMAX reducido como la mejor especificación
predictiva, resulta necesario interpretar el significado económico de
las variables que permanecieron en el modelo.

El propósito de esta sección consiste en analizar cómo los factores
económicos, logísticos y climáticos contribuyen a explicar las
variaciones observadas en el precio nacional del maíz blanco.

\begin{center}\rule{0.5\linewidth}{0.5pt}\end{center}

\hypertarget{precio-internacional-del-mauxedz-cbot}{%
\subsection{Precio internacional del maíz
(CBOT)}\label{precio-internacional-del-mauxedz-cbot}}

El precio internacional del maíz constituye uno de los principales
determinantes del mercado nacional.

México mantiene una fuerte integración con los mercados agrícolas
internacionales, particularmente con Estados Unidos, principal productor
y exportador mundial de maíz.

Cuando el precio internacional aumenta, los costos de importación se
incrementan y los productores nacionales encuentran incentivos para
ajustar sus precios hacia niveles más cercanos a los observados en el
mercado global.

Los resultados obtenidos muestran que el precio internacional aporta
información estadísticamente significativa para explicar la evolución
del precio nacional del maíz blanco.

Esto confirma que el mercado mexicano no opera de forma aislada, sino
que se encuentra influenciado por las condiciones internacionales de
oferta y demanda.

\begin{center}\rule{0.5\linewidth}{0.5pt}\end{center}

\hypertarget{costos-de-transporte}{%
\subsection{Costos de transporte}\label{costos-de-transporte}}

Los costos de transporte representan uno de los factores más relevantes
dentro de la cadena de suministro agroalimentaria.

El maíz producido en las regiones agrícolas debe ser movilizado hacia
centros de almacenamiento, procesamiento y consumo distribuidos en todo
el territorio nacional.

Incrementos en los costos logísticos generan mayores gastos de
distribución, los cuales terminan reflejándose en el precio final del
producto.

La significancia estadística observada en el modelo indica que los
costos de transporte constituyen un mecanismo importante de transmisión
de precios dentro del mercado nacional.

Este resultado es consistente con la literatura económica sobre costos
de comercialización agrícola.

\begin{center}\rule{0.5\linewidth}{0.5pt}\end{center}

\hypertarget{temperatura-promedio-nacional}{%
\subsection{Temperatura promedio
nacional}\label{temperatura-promedio-nacional}}

La temperatura influye directamente sobre las condiciones de crecimiento
y desarrollo de los cultivos.

Cambios importantes en los patrones térmicos pueden afectar:

\begin{itemize}
\tightlist
\item
  Rendimientos agrícolas.
\item
  Productividad por hectárea.
\item
  Disponibilidad de agua.
\item
  Incidencia de estrés térmico en los cultivos.
\end{itemize}

La permanencia de esta variable dentro del modelo sugiere que las
condiciones climáticas tienen efectos observables sobre la dinámica de
precios del maíz blanco.

Esto refuerza la importancia de incorporar variables climáticas en los
modelos predictivos agroalimentarios.

\begin{center}\rule{0.5\linewidth}{0.5pt}\end{center}

\hypertarget{uxedndice-de-severidad-de-sequuxeda}{%
\subsection{Índice de severidad de
sequía}\label{uxedndice-de-severidad-de-sequuxeda}}

La sequía representa uno de los riesgos más importantes para la
producción agrícola nacional.

Periodos prolongados de escasez de agua reducen la productividad,
incrementan la incertidumbre productiva y pueden provocar disminuciones
en la oferta disponible.

Los resultados obtenidos muestran que la severidad de sequía aporta
información significativa incluso después de controlar por otros
factores económicos y climáticos.

Este hallazgo resulta particularmente relevante debido al incremento
observado en la frecuencia e intensidad de eventos de sequía durante los
últimos años en diversas regiones agrícolas del país.

\begin{center}\rule{0.5\linewidth}{0.5pt}\end{center}

\hypertarget{variables-excluidas-del-modelo-final}{%
\subsection{Variables excluidas del modelo
final}\label{variables-excluidas-del-modelo-final}}

Durante el proceso de depuración estadística fueron eliminadas las
variables:

\begin{itemize}
\tightlist
\item
  Precipitación acumulada nacional.
\item
  Producción nacional de maíz blanco.
\end{itemize}

Aunque ambas presentan relevancia conceptual dentro del fenómeno
estudiado, sus coeficientes no resultaron estadísticamente
significativos dentro de la especificación final.

Esto sugiere que parte de la información contenida en dichas variables
ya se encontraba capturada por otros indicadores incluidos en el modelo,
particularmente las variables climáticas y económicas seleccionadas.

Los resultados obtenidos permiten concluir que el precio nacional del
maíz blanco responde a una combinación de factores:

\begin{enumerate}
\def\labelenumi{\arabic{enumi}.}
\tightlist
\item
  Factores internacionales de mercado (CBOT).
\item
  Factores logísticos internos (transporte).
\item
  Factores climáticos (temperatura).
\item
  Factores de estrés hídrico (sequía).
\end{enumerate}

La incorporación conjunta de estos elementos permitió construir un
modelo con capacidad predictiva superior a los enfoques basados
exclusivamente en la dinámica temporal de la serie.

En consecuencia, la evidencia obtenida respalda la hipótesis de que la
evolución del precio nacional del maíz blanco es resultado de la
interacción entre variables económicas, logísticas y agroclimáticas.

    \begin{tcolorbox}[breakable, size=fbox, boxrule=1pt, pad at break*=1mm,colback=cellbackground, colframe=cellborder]
\prompt{In}{incolor}{209}{\boxspacing}
\begin{Verbatim}[commandchars=\\\{\}]
\PY{k+kn}{import} \PY{n+nn}{pandas} \PY{k}{as} \PY{n+nn}{pd}
\PY{k+kn}{import} \PY{n+nn}{numpy} \PY{k}{as} \PY{n+nn}{np}
\PY{k+kn}{import} \PY{n+nn}{matplotlib}\PY{n+nn}{.}\PY{n+nn}{pyplot} \PY{k}{as} \PY{n+nn}{plt}

\PY{c+c1}{\PYZsh{} 1. Asegurar que la serie histórica tenga índice de fecha}
\PY{n}{serie\PYZus{}historica} \PY{o}{=} \PY{p}{(}
    \PY{n}{pdf\PYZus{}completo}
    \PY{o}{.}\PY{n}{copy}\PY{p}{(}\PY{p}{)}
\PY{p}{)}

\PY{n}{serie\PYZus{}historica}\PY{p}{[}\PY{l+s+s2}{\PYZdq{}}\PY{l+s+s2}{fecha}\PY{l+s+s2}{\PYZdq{}}\PY{p}{]} \PY{o}{=} \PY{n}{pd}\PY{o}{.}\PY{n}{to\PYZus{}datetime}\PY{p}{(}\PY{n}{serie\PYZus{}historica}\PY{p}{[}\PY{l+s+s2}{\PYZdq{}}\PY{l+s+s2}{fecha}\PY{l+s+s2}{\PYZdq{}}\PY{p}{]}\PY{p}{)}

\PY{n}{serie\PYZus{}historica} \PY{o}{=} \PY{p}{(}
    \PY{n}{serie\PYZus{}historica}
    \PY{o}{.}\PY{n}{set\PYZus{}index}\PY{p}{(}\PY{l+s+s2}{\PYZdq{}}\PY{l+s+s2}{fecha}\PY{l+s+s2}{\PYZdq{}}\PY{p}{)}
    \PY{p}{[}\PY{l+s+s2}{\PYZdq{}}\PY{l+s+s2}{Precio\PYZus{}maizblanco\PYZus{}mxpkilo\PYZus{}prom\PYZus{}nacional}\PY{l+s+s2}{\PYZdq{}}\PY{p}{]}
    \PY{o}{.}\PY{n}{sort\PYZus{}index}\PY{p}{(}\PY{p}{)}
\PY{p}{)}

\PY{c+c1}{\PYZsh{} 2. Definir horizonte futuro}
\PY{n}{pasos\PYZus{}futuros} \PY{o}{=} \PY{l+m+mi}{12}

\PY{n}{fechas\PYZus{}futuras} \PY{o}{=} \PY{n}{pd}\PY{o}{.}\PY{n}{date\PYZus{}range}\PY{p}{(}
    \PY{n}{start}\PY{o}{=}\PY{n}{serie\PYZus{}historica}\PY{o}{.}\PY{n}{index}\PY{o}{.}\PY{n}{max}\PY{p}{(}\PY{p}{)} \PY{o}{+} \PY{n}{pd}\PY{o}{.}\PY{n}{DateOffset}\PY{p}{(}\PY{n}{months}\PY{o}{=}\PY{l+m+mi}{1}\PY{p}{)}\PY{p}{,}
    \PY{n}{periods}\PY{o}{=}\PY{n}{pasos\PYZus{}futuros}\PY{p}{,}
    \PY{n}{freq}\PY{o}{=}\PY{l+s+s2}{\PYZdq{}}\PY{l+s+s2}{MS}\PY{l+s+s2}{\PYZdq{}}
\PY{p}{)}

\PY{c+c1}{\PYZsh{} 3. Construir X\PYZus{}future con índice de fechas correcto}
\PY{n}{X\PYZus{}future} \PY{o}{=} \PY{n}{pd}\PY{o}{.}\PY{n}{DataFrame}\PY{p}{(}
    \PY{p}{[}\PY{n}{X\PYZus{}final}\PY{o}{.}\PY{n}{iloc}\PY{p}{[}\PY{o}{\PYZhy{}}\PY{l+m+mi}{1}\PY{p}{]}\PY{o}{.}\PY{n}{values}\PY{p}{]} \PY{o}{*} \PY{n}{pasos\PYZus{}futuros}\PY{p}{,}
    \PY{n}{columns}\PY{o}{=}\PY{n}{X\PYZus{}final}\PY{o}{.}\PY{n}{columns}\PY{p}{,}
    \PY{n}{index}\PY{o}{=}\PY{n}{fechas\PYZus{}futuras}
\PY{p}{)}

\PY{c+c1}{\PYZsh{} 4. Generar pronóstico futuro}
\PY{n}{pronostico\PYZus{}objeto} \PY{o}{=} \PY{n}{resultado\PYZus{}final}\PY{o}{.}\PY{n}{get\PYZus{}forecast}\PY{p}{(}
    \PY{n}{steps}\PY{o}{=}\PY{n}{pasos\PYZus{}futuros}\PY{p}{,}
    \PY{n}{exog}\PY{o}{=}\PY{n}{X\PYZus{}future}
\PY{p}{)}

\PY{n}{precio\PYZus{}pronosticado} \PY{o}{=} \PY{n}{pronostico\PYZus{}objeto}\PY{o}{.}\PY{n}{predicted\PYZus{}mean}
\PY{n}{precio\PYZus{}pronosticado}\PY{o}{.}\PY{n}{index} \PY{o}{=} \PY{n}{fechas\PYZus{}futuras}

\PY{n}{intervalo\PYZus{}confianza} \PY{o}{=} \PY{n}{pronostico\PYZus{}objeto}\PY{o}{.}\PY{n}{conf\PYZus{}int}\PY{p}{(}\PY{n}{alpha}\PY{o}{=}\PY{l+m+mf}{0.05}\PY{p}{)}
\PY{n}{intervalo\PYZus{}confianza}\PY{o}{.}\PY{n}{index} \PY{o}{=} \PY{n}{fechas\PYZus{}futuras}

\PY{c+c1}{\PYZsh{} 5. Graficar}
\PY{n}{plt}\PY{o}{.}\PY{n}{figure}\PY{p}{(}\PY{n}{figsize}\PY{o}{=}\PY{p}{(}\PY{l+m+mi}{14}\PY{p}{,} \PY{l+m+mi}{7}\PY{p}{)}\PY{p}{)}

\PY{n}{historico\PYZus{}reciente} \PY{o}{=} \PY{n}{serie\PYZus{}historica}\PY{o}{.}\PY{n}{loc}\PY{p}{[}\PY{l+s+s2}{\PYZdq{}}\PY{l+s+s2}{2023\PYZhy{}01\PYZhy{}01}\PY{l+s+s2}{\PYZdq{}}\PY{p}{:}\PY{p}{]}

\PY{n}{plt}\PY{o}{.}\PY{n}{plot}\PY{p}{(}
    \PY{n}{historico\PYZus{}reciente}\PY{o}{.}\PY{n}{index}\PY{p}{,}
    \PY{n}{historico\PYZus{}reciente}\PY{p}{,}
    \PY{n}{label}\PY{o}{=}\PY{l+s+s2}{\PYZdq{}}\PY{l+s+s2}{Precio Histórico Observado}\PY{l+s+s2}{\PYZdq{}}\PY{p}{,}
    \PY{n}{linewidth}\PY{o}{=}\PY{l+m+mi}{2}
\PY{p}{)}

\PY{n}{plt}\PY{o}{.}\PY{n}{plot}\PY{p}{(}
    \PY{n}{precio\PYZus{}pronosticado}\PY{o}{.}\PY{n}{index}\PY{p}{,}
    \PY{n}{precio\PYZus{}pronosticado}\PY{p}{,}
    \PY{n}{label}\PY{o}{=}\PY{l+s+s2}{\PYZdq{}}\PY{l+s+s2}{Pronóstico SARIMAX Futuro}\PY{l+s+s2}{\PYZdq{}}\PY{p}{,}
    \PY{n}{linestyle}\PY{o}{=}\PY{l+s+s2}{\PYZdq{}}\PY{l+s+s2}{\PYZhy{}\PYZhy{}}\PY{l+s+s2}{\PYZdq{}}\PY{p}{,}
    \PY{n}{linewidth}\PY{o}{=}\PY{l+m+mf}{2.5}
\PY{p}{)}

\PY{n}{plt}\PY{o}{.}\PY{n}{fill\PYZus{}between}\PY{p}{(}
    \PY{n}{intervalo\PYZus{}confianza}\PY{o}{.}\PY{n}{index}\PY{p}{,}
    \PY{n}{intervalo\PYZus{}confianza}\PY{o}{.}\PY{n}{iloc}\PY{p}{[}\PY{p}{:}\PY{p}{,} \PY{l+m+mi}{0}\PY{p}{]}\PY{p}{,}
    \PY{n}{intervalo\PYZus{}confianza}\PY{o}{.}\PY{n}{iloc}\PY{p}{[}\PY{p}{:}\PY{p}{,} \PY{l+m+mi}{1}\PY{p}{]}\PY{p}{,}
    \PY{n}{alpha}\PY{o}{=}\PY{l+m+mf}{0.2}\PY{p}{,}
    \PY{n}{label}\PY{o}{=}\PY{l+s+s2}{\PYZdq{}}\PY{l+s+s2}{Intervalo de Confianza 95}\PY{l+s+s2}{\PYZpc{}}\PY{l+s+s2}{\PYZdq{}}
\PY{p}{)}

\PY{n}{plt}\PY{o}{.}\PY{n}{axvline}\PY{p}{(}
    \PY{n}{serie\PYZus{}historica}\PY{o}{.}\PY{n}{index}\PY{o}{.}\PY{n}{max}\PY{p}{(}\PY{p}{)}\PY{p}{,}
    \PY{n}{linestyle}\PY{o}{=}\PY{l+s+s2}{\PYZdq{}}\PY{l+s+s2}{:}\PY{l+s+s2}{\PYZdq{}}\PY{p}{,}
    \PY{n}{linewidth}\PY{o}{=}\PY{l+m+mi}{2}\PY{p}{,}
    \PY{n}{label}\PY{o}{=}\PY{l+s+sa}{f}\PY{l+s+s2}{\PYZdq{}}\PY{l+s+s2}{Frontera de datos actuales (}\PY{l+s+si}{\PYZob{}}\PY{n}{serie\PYZus{}historica}\PY{o}{.}\PY{n}{index}\PY{o}{.}\PY{n}{max}\PY{p}{(}\PY{p}{)}\PY{o}{.}\PY{n}{strftime}\PY{p}{(}\PY{l+s+s1}{\PYZsq{}}\PY{l+s+s1}{\PYZpc{}}\PY{l+s+s1}{B }\PY{l+s+s1}{\PYZpc{}}\PY{l+s+s1}{Y}\PY{l+s+s1}{\PYZsq{}}\PY{p}{)}\PY{l+s+si}{\PYZcb{}}\PY{l+s+s2}{)}\PY{l+s+s2}{\PYZdq{}}
\PY{p}{)}

\PY{n}{plt}\PY{o}{.}\PY{n}{title}\PY{p}{(}
    \PY{l+s+s2}{\PYZdq{}}\PY{l+s+s2}{Proyección del Precio Nacional del Maíz Blanco}\PY{l+s+s2}{\PYZdq{}}\PY{p}{,}
    \PY{n}{fontsize}\PY{o}{=}\PY{l+m+mi}{14}\PY{p}{,}
    \PY{n}{fontweight}\PY{o}{=}\PY{l+s+s2}{\PYZdq{}}\PY{l+s+s2}{bold}\PY{l+s+s2}{\PYZdq{}}
\PY{p}{)}

\PY{n}{plt}\PY{o}{.}\PY{n}{xlabel}\PY{p}{(}\PY{l+s+s2}{\PYZdq{}}\PY{l+s+s2}{Línea de Tiempo Mensual}\PY{l+s+s2}{\PYZdq{}}\PY{p}{)}
\PY{n}{plt}\PY{o}{.}\PY{n}{ylabel}\PY{p}{(}\PY{l+s+s2}{\PYZdq{}}\PY{l+s+s2}{Precio Promedio (MXN/Kg)}\PY{l+s+s2}{\PYZdq{}}\PY{p}{)}

\PY{n}{plt}\PY{o}{.}\PY{n}{grid}\PY{p}{(}\PY{k+kc}{True}\PY{p}{,} \PY{n}{linestyle}\PY{o}{=}\PY{l+s+s2}{\PYZdq{}}\PY{l+s+s2}{\PYZhy{}\PYZhy{}}\PY{l+s+s2}{\PYZdq{}}\PY{p}{,} \PY{n}{alpha}\PY{o}{=}\PY{l+m+mf}{0.5}\PY{p}{)}
\PY{n}{plt}\PY{o}{.}\PY{n}{legend}\PY{p}{(}\PY{n}{loc}\PY{o}{=}\PY{l+s+s2}{\PYZdq{}}\PY{l+s+s2}{upper left}\PY{l+s+s2}{\PYZdq{}}\PY{p}{)}

\PY{n}{plt}\PY{o}{.}\PY{n}{annotate}\PY{p}{(}
    \PY{l+s+sa}{f}\PY{l+s+s2}{\PYZdq{}}\PY{l+s+s2}{\PYZdl{}}\PY{l+s+si}{\PYZob{}}\PY{n}{precio\PYZus{}pronosticado}\PY{o}{.}\PY{n}{iloc}\PY{p}{[}\PY{l+m+mi}{0}\PY{p}{]}\PY{l+s+si}{:}\PY{l+s+s2}{.2f}\PY{l+s+si}{\PYZcb{}}\PY{l+s+s2}{\PYZdq{}}\PY{p}{,}
    \PY{p}{(}\PY{n}{precio\PYZus{}pronosticado}\PY{o}{.}\PY{n}{index}\PY{p}{[}\PY{l+m+mi}{0}\PY{p}{]}\PY{p}{,} \PY{n}{precio\PYZus{}pronosticado}\PY{o}{.}\PY{n}{iloc}\PY{p}{[}\PY{l+m+mi}{0}\PY{p}{]}\PY{p}{)}\PY{p}{,}
    \PY{n}{textcoords}\PY{o}{=}\PY{l+s+s2}{\PYZdq{}}\PY{l+s+s2}{offset points}\PY{l+s+s2}{\PYZdq{}}\PY{p}{,}
    \PY{n}{xytext}\PY{o}{=}\PY{p}{(}\PY{l+m+mi}{0}\PY{p}{,} \PY{l+m+mi}{10}\PY{p}{)}\PY{p}{,}
    \PY{n}{ha}\PY{o}{=}\PY{l+s+s2}{\PYZdq{}}\PY{l+s+s2}{center}\PY{l+s+s2}{\PYZdq{}}\PY{p}{,}
    \PY{n}{fontweight}\PY{o}{=}\PY{l+s+s2}{\PYZdq{}}\PY{l+s+s2}{bold}\PY{l+s+s2}{\PYZdq{}}
\PY{p}{)}

\PY{n}{plt}\PY{o}{.}\PY{n}{annotate}\PY{p}{(}
    \PY{l+s+sa}{f}\PY{l+s+s2}{\PYZdq{}}\PY{l+s+s2}{\PYZdl{}}\PY{l+s+si}{\PYZob{}}\PY{n}{precio\PYZus{}pronosticado}\PY{o}{.}\PY{n}{iloc}\PY{p}{[}\PY{o}{\PYZhy{}}\PY{l+m+mi}{1}\PY{p}{]}\PY{l+s+si}{:}\PY{l+s+s2}{.2f}\PY{l+s+si}{\PYZcb{}}\PY{l+s+s2}{\PYZdq{}}\PY{p}{,}
    \PY{p}{(}\PY{n}{precio\PYZus{}pronosticado}\PY{o}{.}\PY{n}{index}\PY{p}{[}\PY{o}{\PYZhy{}}\PY{l+m+mi}{1}\PY{p}{]}\PY{p}{,} \PY{n}{precio\PYZus{}pronosticado}\PY{o}{.}\PY{n}{iloc}\PY{p}{[}\PY{o}{\PYZhy{}}\PY{l+m+mi}{1}\PY{p}{]}\PY{p}{)}\PY{p}{,}
    \PY{n}{textcoords}\PY{o}{=}\PY{l+s+s2}{\PYZdq{}}\PY{l+s+s2}{offset points}\PY{l+s+s2}{\PYZdq{}}\PY{p}{,}
    \PY{n}{xytext}\PY{o}{=}\PY{p}{(}\PY{l+m+mi}{0}\PY{p}{,} \PY{l+m+mi}{10}\PY{p}{)}\PY{p}{,}
    \PY{n}{ha}\PY{o}{=}\PY{l+s+s2}{\PYZdq{}}\PY{l+s+s2}{center}\PY{l+s+s2}{\PYZdq{}}\PY{p}{,}
    \PY{n}{fontweight}\PY{o}{=}\PY{l+s+s2}{\PYZdq{}}\PY{l+s+s2}{bold}\PY{l+s+s2}{\PYZdq{}}
\PY{p}{)}

\PY{n}{plt}\PY{o}{.}\PY{n}{show}\PY{p}{(}\PY{p}{)}

\PY{c+c1}{\PYZsh{} 6. Tabla de proyección}
\PY{n}{tabla\PYZus{}proyeccion} \PY{o}{=} \PY{n}{pd}\PY{o}{.}\PY{n}{DataFrame}\PY{p}{(}\PY{p}{\PYZob{}}
    \PY{l+s+s2}{\PYZdq{}}\PY{l+s+s2}{Precio Pronosticado (MXN/Kg)}\PY{l+s+s2}{\PYZdq{}}\PY{p}{:} \PY{n}{precio\PYZus{}pronosticado}\PY{o}{.}\PY{n}{round}\PY{p}{(}\PY{l+m+mi}{2}\PY{p}{)}\PY{p}{,}
    \PY{l+s+s2}{\PYZdq{}}\PY{l+s+s2}{Límite Inferior (MXN)}\PY{l+s+s2}{\PYZdq{}}\PY{p}{:} \PY{n}{intervalo\PYZus{}confianza}\PY{o}{.}\PY{n}{iloc}\PY{p}{[}\PY{p}{:}\PY{p}{,} \PY{l+m+mi}{0}\PY{p}{]}\PY{o}{.}\PY{n}{round}\PY{p}{(}\PY{l+m+mi}{2}\PY{p}{)}\PY{p}{,}
    \PY{l+s+s2}{\PYZdq{}}\PY{l+s+s2}{Límite Superior (MXN)}\PY{l+s+s2}{\PYZdq{}}\PY{p}{:} \PY{n}{intervalo\PYZus{}confianza}\PY{o}{.}\PY{n}{iloc}\PY{p}{[}\PY{p}{:}\PY{p}{,} \PY{l+m+mi}{1}\PY{p}{]}\PY{o}{.}\PY{n}{round}\PY{p}{(}\PY{l+m+mi}{2}\PY{p}{)}
\PY{p}{\PYZcb{}}\PY{p}{)}

\PY{n}{tabla\PYZus{}proyeccion}
\end{Verbatim}
\end{tcolorbox}

    \begin{center}
    \adjustimage{max size={0.9\linewidth}{0.9\paperheight}}{output_172_0.png}
    \end{center}
    { \hspace*{\fill} \\}
    
            \begin{tcolorbox}[breakable, size=fbox, boxrule=.5pt, pad at break*=1mm, opacityfill=0]
\prompt{Out}{outcolor}{209}{\boxspacing}
\begin{Verbatim}[commandchars=\\\{\}]
            Precio Pronosticado (MXN/Kg)  Límite Inferior (MXN)  \textbackslash{}
2026-06-01                          9.67                   9.45
2026-07-01                          9.84                   9.51
2026-08-01                          9.73                   9.33
2026-09-01                          9.67                   9.20
2026-10-01                          9.57                   9.03
2026-11-01                          9.55                   8.96
2026-12-01                          9.52                   8.88
2027-01-01                          9.58                   8.90
2027-02-01                          9.56                   8.84
2027-03-01                          9.48                   8.72
2027-04-01                          9.53                   8.73
2027-05-01                          9.50                   8.66

            Límite Superior (MXN)
2026-06-01                   9.88
2026-07-01                  10.16
2026-08-01                  10.14
2026-09-01                  10.15
2026-10-01                  10.10
2026-11-01                  10.13
2026-12-01                  10.15
2027-01-01                  10.26
2027-02-01                  10.28
2027-03-01                  10.24
2027-04-01                  10.33
2027-05-01                  10.33
\end{Verbatim}
\end{tcolorbox}
        
    Como parte de la fase de despliegue del proyecto, se utilizó el modelo
SARIMAX reducido para generar una proyección fuera de muestra del precio
nacional del maíz blanco durante los siguientes doce meses posteriores
al último dato observado (mayo de 2026).

El objetivo de esta etapa consiste en transformar los hallazgos
estadísticos y predictivos en información útil para la toma de
decisiones dentro de la cadena agroalimentaria, permitiendo anticipar
escenarios futuros de comportamiento del mercado.

\hypertarget{pronuxf3stico-futuro-del-precio-nacional-del-mauxedz-blanco}{%
\subsubsection{Pronóstico Futuro del Precio Nacional del Maíz
Blanco}\label{pronuxf3stico-futuro-del-precio-nacional-del-mauxedz-blanco}}

La serie histórica observada entre 2023 y mayo de 2026, así como la
proyección generada por el modelo SARIMAX para el periodo comprendido
entre junio de 2026 y mayo de 2027 indican que el precio nacional del
maíz blanco permanecería relativamente estable durante el horizonte de
planeación considerado.

El modelo estima que el precio promedio nacional iniciaría en
aproximadamente \textbf{\$9.67 MXN/kg} durante junio de 2026 y
concluiría alrededor de \textbf{\$9.50 MXN/kg} en mayo de 2027.

La trayectoria proyectada sugiere una ligera tendencia descendente,
aunque dentro de un rango de variación moderado.

\begin{longtable}[]{@{}lr@{}}
\toprule\noalign{}
Fecha & Precio Pronosticado (MXN/Kg) \\
\midrule\noalign{}
\endhead
\bottomrule\noalign{}
\endlastfoot
Jun-2026 & 9.67 \\
Jul-2026 & 9.84 \\
Ago-2026 & 9.73 \\
Sep-2026 & 9.67 \\
Oct-2026 & 9.57 \\
Nov-2026 & 9.55 \\
Dic-2026 & 9.52 \\
Ene-2027 & 9.58 \\
Feb-2027 & 9.56 \\
Mar-2027 & 9.48 \\
Abr-2027 & 9.53 \\
May-2027 & 9.50 \\
\end{longtable}

Este panorama nos sugiere que el mercado nacional del maíz blanco podría
experimentar una etapa de relativa estabilidad durante los próximos doce
meses.

Aunque el modelo identifica la influencia significativa de factores
como:

\begin{itemize}
\tightlist
\item
  Precio internacional del maíz (CBOT).
\item
  Costos de transporte.
\item
  Temperatura promedio nacional.
\item
  Severidad de sequía.
\end{itemize}

la combinación observada de estas variables no genera señales de
incrementos abruptos dentro del horizonte proyectado.

Lo anterior puede interpretarse como un escenario de riesgo moderado
para productores, comercializadores y consumidores, donde no se
anticipan aumentos extraordinarios en el precio nacional del grano bajo
las condiciones actuales.

\begin{center}\rule{0.5\linewidth}{0.5pt}\end{center}

\hypertarget{incertidumbre-en-el-pronuxf3stico}{%
\subsubsection{Incertidumbre en el
Pronóstico}\label{incertidumbre-en-el-pronuxf3stico}}

Es importante señalar que toda proyección estadística incorpora
incertidumbre.

Por esta razón, el modelo genera intervalos de confianza al 95\%, los
cuales representan el rango dentro del cual es razonable esperar que se
ubique el precio real futuro.

Para mayo de 2027 el modelo estima:

\begin{itemize}
\tightlist
\item
  Valor esperado: \textbf{\$9.50 MXN/kg}
\item
  Límite inferior: \textbf{\$8.66 MXN/kg}
\item
  Límite superior: \textbf{\$10.33 MXN/kg}
\end{itemize}

Estos resultados indican que, aunque el escenario más probable
corresponde a un precio cercano a \$9.50 MXN/kg, eventos climáticos,
logísticos o económicos extraordinarios podrían provocar desviaciones
respecto a dicha estimación.

\begin{center}\rule{0.5\linewidth}{0.5pt}\end{center}

\hypertarget{impacto-de-los-hallazgos}{%
\subsubsection{39.5 Impacto de los
hallazgos}\label{impacto-de-los-hallazgos}}

La principal utilidad de este pronóstico consiste en proporcionar
información anticipada para la planeación operativa y financiera de los
actores involucrados en la cadena agroalimentaria.

Entre las aplicaciones potenciales destacan:

\begin{itemize}
\tightlist
\item
  Planeación de compras de materia prima.
\item
  Estimación de costos futuros de producción.
\item
  Diseño de estrategias de almacenamiento.
\item
  Evaluación de riesgos asociados a fenómenos climáticos.
\item
  Planeación presupuestal de cooperativas y organizaciones agrícolas.
\end{itemize}

De esta manera, la fase de despliegue demuestra cómo los modelos
predictivos desarrollados mediante Ciencia de Datos pueden transformarse
en herramientas prácticas para apoyar la toma de decisiones dentro del
sector agroalimentario mexicano.

    \hypertarget{implicaciones-para-la-toma-de-decisiones}{%
\section{31. Implicaciones para la Toma de
Decisiones}\label{implicaciones-para-la-toma-de-decisiones}}

El propósito final del modelo desarrollado no consiste únicamente en
generar pronósticos estadísticos, sino en proporcionar información útil
para apoyar la toma de decisiones dentro del sistema agroalimentario.

Los resultados obtenidos permiten identificar factores clave que
influyen sobre la formación del precio nacional del maíz blanco y que
pueden ser monitoreados para anticipar escenarios de riesgo o de
oportunidad.

\begin{center}\rule{0.5\linewidth}{0.5pt}\end{center}

\hypertarget{aplicaciones-para-productores-agruxedcolas}{%
\subsection{Aplicaciones para productores
agrícolas}\label{aplicaciones-para-productores-agruxedcolas}}

Los productores pueden utilizar la información proveniente de variables
como:

\begin{itemize}
\tightlist
\item
  Precio internacional del maíz (CBOT).
\item
  Temperatura promedio.
\item
  Índice de severidad de sequía.
\end{itemize}

para anticipar posibles cambios en las condiciones de mercado.

La identificación temprana de escenarios climáticos adversos o
incrementos en los precios internacionales puede facilitar decisiones
relacionadas con:

\begin{itemize}
\tightlist
\item
  Planeación de siembra.
\item
  Gestión de riesgos.
\item
  Contratación de coberturas agrícolas.
\item
  Estrategias de comercialización.
\end{itemize}

\begin{center}\rule{0.5\linewidth}{0.5pt}\end{center}

\hypertarget{aplicaciones-para-comercializadores}{%
\subsection{Aplicaciones para
comercializadores}\label{aplicaciones-para-comercializadores}}

Las empresas dedicadas al almacenamiento, distribución o
comercialización de maíz pueden utilizar los resultados del modelo para:

\begin{itemize}
\tightlist
\item
  Estimar tendencias futuras de precios.
\item
  Planificar inventarios.
\item
  Optimizar compras anticipadas.
\item
  Reducir riesgos asociados a la volatilidad del mercado.
\end{itemize}

La relevancia observada de los costos de transporte también permite
monitorear posibles incrementos en los costos logísticos que
eventualmente podrían trasladarse al precio final.

\begin{center}\rule{0.5\linewidth}{0.5pt}\end{center}

\hypertarget{aplicaciones-para-la-industria-alimentaria}{%
\subsection{Aplicaciones para la industria
alimentaria}\label{aplicaciones-para-la-industria-alimentaria}}

El maíz blanco constituye un insumo estratégico para diversas
industrias, particularmente para la elaboración de tortilla y otros
productos derivados.

La capacidad de anticipar variaciones en el precio del maíz puede
contribuir a:

\begin{itemize}
\tightlist
\item
  Planeación financiera.
\item
  Gestión de costos de producción.
\item
  Negociación de contratos de suministro.
\item
  Estimación de márgenes operativos.
\end{itemize}

De esta forma, el modelo puede funcionar como una herramienta
complementaria para la toma de decisiones empresariales.

\begin{center}\rule{0.5\linewidth}{0.5pt}\end{center}

\hypertarget{aplicaciones-para-el-sector-puxfablico}{%
\subsection{Aplicaciones para el sector
público}\label{aplicaciones-para-el-sector-puxfablico}}

Los resultados muestran que factores climáticos y logísticos tienen una
influencia significativa sobre el comportamiento del precio nacional del
maíz.

Esta evidencia puede ser útil para:

\begin{itemize}
\tightlist
\item
  Diseñar programas de apoyo agrícola.
\item
  Fortalecer sistemas de monitoreo climático.
\item
  Identificar regiones vulnerables ante eventos de sequía.
\item
  Implementar estrategias de mitigación de riesgos agroalimentarios.
\end{itemize}

Asimismo, el seguimiento sistemático de las variables identificadas
podría contribuir al desarrollo de sistemas de alerta temprana para
anticipar presiones inflacionarias en productos derivados del maíz.

\begin{center}\rule{0.5\linewidth}{0.5pt}\end{center}

\hypertarget{implicaciones-para-la-seguridad-alimentaria}{%
\subsection{Implicaciones para la seguridad
alimentaria}\label{implicaciones-para-la-seguridad-alimentaria}}

El maíz blanco es uno de los productos agrícolas más importantes para la
alimentación de la población mexicana.

Las variaciones en su precio afectan directamente:

\begin{itemize}
\tightlist
\item
  El acceso a alimentos básicos.
\item
  El gasto de los hogares.
\item
  La estabilidad de cadenas agroalimentarias.
\item
  La inflación alimentaria.
\end{itemize}

Por esta razón, contar con herramientas que permitan anticipar cambios
en el comportamiento del precio constituye un elemento valioso para la
planeación económica y alimentaria.

El modelo desarrollado demuestra que es posible integrar información
económica, climática y logística para mejorar la capacidad de predicción
de los precios agrícolas.

Más allá de la precisión estadística alcanzada, el principal valor del
modelo radica en su capacidad para transformar datos provenientes de
múltiples fuentes en información útil para apoyar decisiones dentro del
sistema agroalimentario.

En este sentido, el proyecto contribuye al desarrollo de herramientas de
análisis que pueden apoyar la construcción de mercados más eficientes,
transparentes y resilientes.

    \hypertarget{conclusiones}{%
\section{32. Conclusiones}\label{conclusiones}}

\hypertarget{cumplimiento-del-objetivo-general}{%
\subsection{Cumplimiento del objetivo
general}\label{cumplimiento-del-objetivo-general}}

El objetivo de este proyecto consistió en desarrollar un modelo de
análisis y pronóstico para el precio nacional del maíz blanco en México
mediante la integración de información económica, climática y productiva
proveniente de diversas fuentes oficiales.

Los resultados obtenidos permiten afirmar que dicho objetivo fue
alcanzado satisfactoriamente.

A través de técnicas de análisis exploratorio de datos, evaluación de
correlaciones, reducción de multicolinealidad y modelado de series de
tiempo, fue posible construir un modelo predictivo capaz de explicar una
proporción importante de la variabilidad observada en el precio nacional
del maíz blanco.

\begin{center}\rule{0.5\linewidth}{0.5pt}\end{center}

\hypertarget{respuesta-a-la-pregunta-de-investigaciuxf3n}{%
\subsection{Respuesta a la pregunta de
investigación}\label{respuesta-a-la-pregunta-de-investigaciuxf3n}}

La pregunta central planteada fue:

\begin{quote}
¿Qué factores explican el comportamiento del precio nacional del maíz
blanco en México y en qué medida pueden utilizarse para mejorar su
capacidad de pronóstico?
\end{quote}

La evidencia obtenida indica que el precio nacional del maíz blanco no
depende exclusivamente de su comportamiento histórico.

Por el contrario, su evolución está asociada a la interacción de
factores:

\begin{itemize}
\tightlist
\item
  Económicos.
\item
  Logísticos.
\item
  Climáticos.
\end{itemize}

Particularmente, las variables que mostraron mayor capacidad explicativa
fueron:

\begin{itemize}
\tightlist
\item
  Precio internacional del maíz (CBOT).
\item
  Costos de transporte.
\item
  Temperatura promedio nacional.
\item
  Índice de severidad de sequía.
\end{itemize}

Estas variables permitieron mejorar significativamente la capacidad
predictiva respecto a los modelos basados únicamente en la estructura
temporal de la serie.

\begin{center}\rule{0.5\linewidth}{0.5pt}\end{center}

\hypertarget{principales-hallazgos}{%
\subsection{Principales hallazgos}\label{principales-hallazgos}}

Entre los hallazgos más relevantes del estudio destacan:

\hypertarget{existencia-de-estacionalidad-anual}{%
\subsubsection{1. Existencia de estacionalidad
anual}\label{existencia-de-estacionalidad-anual}}

La serie presentó patrones estacionales claramente identificables con
periodicidad de doce meses.

La descomposición temporal y los análisis ACF y PACF mostraron evidencia
consistente de este comportamiento.

\begin{center}\rule{0.5\linewidth}{0.5pt}\end{center}

\hypertarget{influencia-de-factores-internacionales}{%
\subsubsection{2. Influencia de factores
internacionales}\label{influencia-de-factores-internacionales}}

El precio internacional del maíz mostró una asociación importante con el
precio nacional.

Este resultado confirma la integración del mercado mexicano con los
mercados agrícolas internacionales.

\begin{center}\rule{0.5\linewidth}{0.5pt}\end{center}

\hypertarget{importancia-de-los-costos-loguxedsticos}{%
\subsubsection{3. Importancia de los costos
logísticos}\label{importancia-de-los-costos-loguxedsticos}}

Los costos de transporte constituyeron uno de los factores más
relevantes dentro del modelo final.

Este hallazgo sugiere que la logística representa un mecanismo
importante de transmisión de precios dentro de la cadena
agroalimentaria.

\begin{center}\rule{0.5\linewidth}{0.5pt}\end{center}

\hypertarget{impacto-de-las-condiciones-climuxe1ticas}{%
\subsubsection{4. Impacto de las condiciones
climáticas}\label{impacto-de-las-condiciones-climuxe1ticas}}

Tanto la temperatura como la severidad de sequía aportaron información
significativa para explicar la dinámica de precios.

Esto evidencia la sensibilidad del sistema agrícola nacional frente a
fenómenos climáticos adversos.

\begin{center}\rule{0.5\linewidth}{0.5pt}\end{center}

Se evaluaron distintos enfoques de modelado:

\begin{itemize}
\tightlist
\item
  ARIMA.
\item
  SARIMA.
\item
  SARIMAX.
\end{itemize}

Los modelos ARIMA presentaron el desempeño más limitado debido a su
incapacidad para capturar adecuadamente la estacionalidad observada.

La incorporación de componentes estacionales mediante SARIMA permitió
reducir considerablemente el error de predicción.

Finalmente, la integración de variables exógenas mediante SARIMAX
produjo el mejor desempeño observado.

El modelo seleccionado fue:

\[
SARIMAX(1,1,0)(1,0,0,12)
\]

con variables exógenas reducidas.

Las métricas finales obtenidas fueron:

\begin{longtable}[]{@{}lr@{}}
\toprule\noalign{}
Métrica & Resultado \\
\midrule\noalign{}
\endhead
\bottomrule\noalign{}
\endlastfoot
MAE & 1.452 \\
RMSE & 1.618 \\
MAPE & 15.83\% \\
\end{longtable}

Estas métricas representan el mejor equilibrio entre precisión
predictiva, parsimonia e interpretación económica.

\begin{center}\rule{0.5\linewidth}{0.5pt}\end{center}

\hypertarget{aportaciones-del-proyecto}{%
\subsection{Aportaciones del proyecto}\label{aportaciones-del-proyecto}}

El presente trabajo aporta una metodología reproducible para integrar
información proveniente de distintas fuentes oficiales dentro de un
esquema de análisis predictivo agroalimentario.

Asimismo, demuestra que la combinación de:

\begin{itemize}
\tightlist
\item
  Series de tiempo.
\item
  Variables económicas.
\item
  Indicadores climáticos.
\item
  Técnicas de minería de datos.
\end{itemize}

permite construir modelos más robustos que aquellos basados únicamente
en información histórica de precios.

\begin{center}\rule{0.5\linewidth}{0.5pt}\end{center}

\hypertarget{limitaciones}{%
\subsection{Limitaciones}\label{limitaciones}}

A pesar de los resultados obtenidos, existen algunas limitaciones que
deben considerarse.

Entre ellas destacan:

\begin{itemize}
\tightlist
\item
  Disponibilidad limitada de variables históricas climáticas para
  periodos anteriores.
\item
  Posibles cambios estructurales derivados de eventos extraordinarios
  como la pandemia de COVID-19.
\item
  Ausencia de variables relacionadas con políticas públicas, subsidios o
  comercio internacional específico.
\item
  Limitaciones inherentes a la agregación de información a nivel
  nacional.
\end{itemize}

Estas restricciones pueden influir sobre la capacidad de generalización
del modelo.

\begin{center}\rule{0.5\linewidth}{0.5pt}\end{center}

\hypertarget{luxedneas-futuras-de-investigaciuxf3n}{%
\subsection{Líneas futuras de
investigación}\label{luxedneas-futuras-de-investigaciuxf3n}}

Como trabajo futuro se propone incorporar nuevas variables
potencialmente relevantes, tales como:

\begin{itemize}
\tightlist
\item
  Indicadores de sequía regional.
\item
  Precios energéticos internacionales.
\item
  Costos de fertilizantes.
\item
  Variables asociadas al comercio exterior.
\item
  Indicadores de disponibilidad hídrica.
\end{itemize}

\begin{center}\rule{0.5\linewidth}{0.5pt}\end{center}

\hypertarget{conclusiuxf3n-general}{%
\subsection{Conclusión general}\label{conclusiuxf3n-general}}

Los resultados obtenidos permiten concluir que el comportamiento del
precio nacional del maíz blanco en México puede explicarse de manera más
adecuada cuando se consideran simultáneamente factores económicos,
logísticos y climáticos.

La evidencia estadística demuestra que la incorporación de información
exógena mejora significativamente la capacidad de pronóstico respecto a
los modelos temporales tradicionales.

En consecuencia, el modelo desarrollado constituye una herramienta útil
para apoyar procesos de análisis, monitoreo y toma de decisiones dentro
del sector agroalimentario mexicano, contribuyendo a una mejor
comprensión de los factores que influyen sobre uno de los productos
estratégicos más importantes para la seguridad alimentaria nacional.

    \hypertarget{referencias}{%
\section{33. Referencias}\label{referencias}}

\hypertarget{fuentes-de-datos}{%
\subsection{Fuentes de datos}\label{fuentes-de-datos}}

Banco de México. (2026). \emph{Tipo de cambio FIX peso-dólar}. Sistema
de Información Económica (SIE). Recuperado de:

https://www.banxico.org.mx/SieInternet/

\begin{center}\rule{0.5\linewidth}{0.5pt}\end{center}

Chicago Board of Trade. (2026). \emph{Corn Futures Historical Data}.
Recuperado de:

https://www.barchart.com/futures/quotes/ZC*0/historical-data

\begin{center}\rule{0.5\linewidth}{0.5pt}\end{center}

Comisión Nacional del Agua. (2026). \emph{Monitor de Sequía de México}.
Servicio Meteorológico Nacional. Recuperado de:

https://smn.conagua.gob.mx/es/climatologia/monitor-de-sequia/monitor-de-sequia-en-mexico

\begin{center}\rule{0.5\linewidth}{0.5pt}\end{center}

Instituto Nacional de Estadística y Geografía. (2026). \emph{Índice
Nacional de Precios al Consumidor (INPC)}. Banco de Información
Económica. Recuperado de:

https://www.inegi.org.mx/app/indicesdeprecios/

\begin{center}\rule{0.5\linewidth}{0.5pt}\end{center}

Secretaría de Agricultura y Desarrollo Rural. (2026). \emph{Producción
agrícola nacional}. Servicio de Información Agroalimentaria y Pesquera
(SIAP). Recuperado de:

https://nube.siap.gob.mx/cierreagricola/

\begin{center}\rule{0.5\linewidth}{0.5pt}\end{center}

Sistema Nacional de Información e Integración de Mercados. (2026).
\emph{Precios nacionales de maíz blanco al mayoreo}. Recuperado de:

https://www.economia-sniim.gob.mx

\begin{center}\rule{0.5\linewidth}{0.5pt}\end{center}

Servicio Meteorológico Nacional. (2026). \emph{Normales climatológicas y
registros históricos de precipitación y temperatura}. Recuperado de:

https://smn.conagua.gob.mx/es/climatologia/informacion-climatologica

\begin{center}\rule{0.5\linewidth}{0.5pt}\end{center}

Comisión Reguladora de Energía. (2026). \emph{Precios históricos de
combustibles automotrices}. Recuperado de:

https://www.gob.mx/cre

\begin{center}\rule{0.5\linewidth}{0.5pt}\end{center}

\hypertarget{referencias-metodoluxf3gicas}{%
\subsection{Referencias
metodológicas}\label{referencias-metodoluxf3gicas}}

Box, G. E. P., Jenkins, G. M., Reinsel, G. C., \& Ljung, G. M. (2016).
\emph{Time Series Analysis: Forecasting and Control} (5th ed.). Hoboken,
NJ: Wiley.

\begin{center}\rule{0.5\linewidth}{0.5pt}\end{center}

Hyndman, R. J., \& Athanasopoulos, G. (2021). \emph{Forecasting:
Principles and Practice} (3rd ed.). Melbourne: OTexts. Recuperado de:

https://otexts.com/fpp3/

\begin{center}\rule{0.5\linewidth}{0.5pt}\end{center}

James, G., Witten, D., Hastie, T., \& Tibshirani, R. (2021). \emph{An
Introduction to Statistical Learning} (2nd ed.). New York: Springer.

\begin{center}\rule{0.5\linewidth}{0.5pt}\end{center}

Montgomery, D. C., Jennings, C. L., \& Kulahci, M. (2015).
\emph{Introduction to Time Series Analysis and Forecasting}. Hoboken,
NJ: Wiley.

\begin{center}\rule{0.5\linewidth}{0.5pt}\end{center}

Chapman, P., Clinton, J., Kerber, R., Khabaza, T., Reinartz, T.,
Shearer, C., \& Wirth, R. (2000). \emph{CRISP-DM 1.0: Step-by-step Data
Mining Guide}. SPSS Inc.

Recuperado de:

https://www.the-modeling-agency.com/crisp-dm.pdf

\begin{center}\rule{0.5\linewidth}{0.5pt}\end{center}

Shmueli, G., Bruce, P., Gedeck, P., \& Patel, N. (2023). \emph{Data
Mining for Business Analytics}. Hoboken, NJ: Wiley.

\begin{center}\rule{0.5\linewidth}{0.5pt}\end{center}

Wooldridge, J. M. (2020). \emph{Introductory Econometrics: A Modern
Approach} (7th ed.). Boston: Cengage Learning.

    \hypertarget{anuxe1lisis-complementario-precio-nacional-de-la-tortilla}{%
\section{34. Análisis Complementario: Precio Nacional de la
Tortilla}\label{anuxe1lisis-complementario-precio-nacional-de-la-tortilla}}

\hypertarget{objetivo}{%
\subsection{Objetivo}\label{objetivo}}

Una vez desarrollado el modelo predictivo para el precio nacional del
maíz blanco, se realiza un análisis complementario orientado al
consumidor final.

El propósito consiste en evaluar si las variaciones observadas en el
precio del maíz se reflejan posteriormente en el precio nacional de la
tortilla, principal producto derivado del maíz blanco consumido en
México.

Este análisis permite extender el alcance del proyecto desde la
producción agrícola hasta el impacto potencial sobre la seguridad
alimentaria y el gasto de los hogares.

\begin{center}\rule{0.5\linewidth}{0.5pt}\end{center}

\hypertarget{fuente-de-informaciuxf3n}{%
\subsection{Fuente de información}\label{fuente-de-informaciuxf3n}}

Los datos históricos del precio de la tortilla fueron obtenidos del
programa:

\textbf{Quién es Quién en los Precios (QQP)}

administrado por la:

\textbf{Procuraduría Federal del Consumidor (PROFECO).}

Referencia:

Procuraduría Federal del Consumidor. (2026). \emph{Resultados históricos
de monitoreo: Tortilla de maíz por kilogramo a granel}. Gobierno de
México.

https://qqp.profeco.gob.mx/

\begin{center}\rule{0.5\linewidth}{0.5pt}\end{center}

\hypertarget{hipuxf3tesis-de-trabajo}{%
\subsection{Hipótesis de trabajo}\label{hipuxf3tesis-de-trabajo}}

Se plantea que los cambios observados en el precio nacional del maíz
blanco se transmiten parcial y gradualmente al precio nacional de la
tortilla.

Por lo tanto, se espera encontrar:

\begin{itemize}
\tightlist
\item
  Asociación positiva entre ambas series.
\item
  Evidencia de transmisión con rezagos temporales.
\item
  Mayor intensidad de relación en horizontes de corto y mediano plazo.
\end{itemize}

    \begin{tcolorbox}[breakable, size=fbox, boxrule=1pt, pad at break*=1mm,colback=cellbackground, colframe=cellborder]
\prompt{In}{incolor}{127}{\boxspacing}
\begin{Verbatim}[commandchars=\\\{\}]
\PY{n}{ruta\PYZus{}tortilla} \PY{o}{=} \PY{l+s+s2}{\PYZdq{}}\PY{l+s+s2}{/home/jovyan/Precio maiz/Data/Raw/Precio\PYZus{}tortilla\PYZus{}prom\PYZus{}nacional\PYZus{}2000\PYZus{}2026.csv}\PY{l+s+s2}{\PYZdq{}}

\PY{n}{df\PYZus{}tortilla} \PY{o}{=} \PY{p}{(}
    \PY{n}{spark}\PY{o}{.}\PY{n}{read}
    \PY{o}{.}\PY{n}{option}\PY{p}{(}\PY{l+s+s2}{\PYZdq{}}\PY{l+s+s2}{header}\PY{l+s+s2}{\PYZdq{}}\PY{p}{,} \PY{l+s+s2}{\PYZdq{}}\PY{l+s+s2}{true}\PY{l+s+s2}{\PYZdq{}}\PY{p}{)}
    \PY{o}{.}\PY{n}{option}\PY{p}{(}\PY{l+s+s2}{\PYZdq{}}\PY{l+s+s2}{inferSchema}\PY{l+s+s2}{\PYZdq{}}\PY{p}{,} \PY{l+s+s2}{\PYZdq{}}\PY{l+s+s2}{true}\PY{l+s+s2}{\PYZdq{}}\PY{p}{)}
    \PY{o}{.}\PY{n}{csv}\PY{p}{(}\PY{n}{ruta\PYZus{}tortilla}\PY{p}{)}
\PY{p}{)}

\PY{n}{df\PYZus{}tortilla}\PY{o}{.}\PY{n}{printSchema}\PY{p}{(}\PY{p}{)}
\PY{n}{df\PYZus{}tortilla}\PY{o}{.}\PY{n}{show}\PY{p}{(}\PY{l+m+mi}{5}\PY{p}{)}
\end{Verbatim}
\end{tcolorbox}

    \begin{Verbatim}[commandchars=\\\{\}]
root
 |-- Año: integer (nullable = true)
 |-- Mes: integer (nullable = true)
 |-- Nacional: double (nullable = true)

+----+---+--------+
| Año|Mes|Nacional|
+----+---+--------+
|2000|  1|     4.0|
|2000|  2|    4.04|
|2000|  3|    4.07|
|2000|  4|     4.1|
|2000|  5|    4.14|
+----+---+--------+
only showing top 5 rows

    \end{Verbatim}

    \hypertarget{construcciuxf3n-de-la-serie-temporal-de-tortilla}{%
\subsection{Construcción de la serie temporal de
tortilla}\label{construcciuxf3n-de-la-serie-temporal-de-tortilla}}

La base de datos de tortilla contiene observaciones mensuales
correspondientes al precio promedio nacional de la tortilla de maíz.

Con el propósito de analizar la relación entre el mercado del maíz
blanco y el consumidor final, se construyó una serie temporal mensual
compatible con la frecuencia utilizada en el análisis del maíz.

La integración de ambas series permitirá evaluar la existencia de
mecanismos de transmisión de precios entre el mercado agrícola y los
productos derivados destinados al consumo humano.

    \begin{tcolorbox}[breakable, size=fbox, boxrule=1pt, pad at break*=1mm,colback=cellbackground, colframe=cellborder]
\prompt{In}{incolor}{128}{\boxspacing}
\begin{Verbatim}[commandchars=\\\{\}]
\PY{n}{pdf\PYZus{}tortilla} \PY{o}{=} \PY{n}{df\PYZus{}tortilla}\PY{o}{.}\PY{n}{toPandas}\PY{p}{(}\PY{p}{)}

\PY{n}{pdf\PYZus{}tortilla}\PY{p}{[}\PY{l+s+s2}{\PYZdq{}}\PY{l+s+s2}{fecha}\PY{l+s+s2}{\PYZdq{}}\PY{p}{]} \PY{o}{=} \PY{n}{pd}\PY{o}{.}\PY{n}{to\PYZus{}datetime}\PY{p}{(}
    \PY{n}{pdf\PYZus{}tortilla}\PY{p}{[}\PY{l+s+s2}{\PYZdq{}}\PY{l+s+s2}{Año}\PY{l+s+s2}{\PYZdq{}}\PY{p}{]}\PY{o}{.}\PY{n}{astype}\PY{p}{(}\PY{n+nb}{str}\PY{p}{)}
    \PY{o}{+} \PY{l+s+s2}{\PYZdq{}}\PY{l+s+s2}{\PYZhy{}}\PY{l+s+s2}{\PYZdq{}}
    \PY{o}{+} \PY{n}{pdf\PYZus{}tortilla}\PY{p}{[}\PY{l+s+s2}{\PYZdq{}}\PY{l+s+s2}{Mes}\PY{l+s+s2}{\PYZdq{}}\PY{p}{]}\PY{o}{.}\PY{n}{astype}\PY{p}{(}\PY{n+nb}{str}\PY{p}{)}
    \PY{o}{+} \PY{l+s+s2}{\PYZdq{}}\PY{l+s+s2}{\PYZhy{}01}\PY{l+s+s2}{\PYZdq{}}
\PY{p}{)}

\PY{n}{pdf\PYZus{}tortilla}\PY{o}{.}\PY{n}{head}\PY{p}{(}\PY{p}{)}
\end{Verbatim}
\end{tcolorbox}

            \begin{tcolorbox}[breakable, size=fbox, boxrule=.5pt, pad at break*=1mm, opacityfill=0]
\prompt{Out}{outcolor}{128}{\boxspacing}
\begin{Verbatim}[commandchars=\\\{\}]
    Año  Mes  Nacional      fecha
0  2000    1      4.00 2000-01-01
1  2000    2      4.04 2000-02-01
2  2000    3      4.07 2000-03-01
3  2000    4      4.10 2000-04-01
4  2000    5      4.14 2000-05-01
\end{Verbatim}
\end{tcolorbox}
        
    \begin{tcolorbox}[breakable, size=fbox, boxrule=1pt, pad at break*=1mm,colback=cellbackground, colframe=cellborder]
\prompt{In}{incolor}{129}{\boxspacing}
\begin{Verbatim}[commandchars=\\\{\}]
\PY{n}{plt}\PY{o}{.}\PY{n}{figure}\PY{p}{(}\PY{n}{figsize}\PY{o}{=}\PY{p}{(}\PY{l+m+mi}{14}\PY{p}{,}\PY{l+m+mi}{6}\PY{p}{)}\PY{p}{)}

\PY{n}{plt}\PY{o}{.}\PY{n}{plot}\PY{p}{(}
    \PY{n}{pdf\PYZus{}tortilla}\PY{p}{[}\PY{l+s+s2}{\PYZdq{}}\PY{l+s+s2}{fecha}\PY{l+s+s2}{\PYZdq{}}\PY{p}{]}\PY{p}{,}
    \PY{n}{pdf\PYZus{}tortilla}\PY{p}{[}\PY{l+s+s2}{\PYZdq{}}\PY{l+s+s2}{Nacional}\PY{l+s+s2}{\PYZdq{}}\PY{p}{]}
\PY{p}{)}

\PY{n}{plt}\PY{o}{.}\PY{n}{title}\PY{p}{(}
    \PY{l+s+s2}{\PYZdq{}}\PY{l+s+s2}{Precio Nacional de la Tortilla}\PY{l+s+s2}{\PYZdq{}}
\PY{p}{)}

\PY{n}{plt}\PY{o}{.}\PY{n}{ylabel}\PY{p}{(}\PY{l+s+s2}{\PYZdq{}}\PY{l+s+s2}{MXN/Kg}\PY{l+s+s2}{\PYZdq{}}\PY{p}{)}
\PY{n}{plt}\PY{o}{.}\PY{n}{xlabel}\PY{p}{(}\PY{l+s+s2}{\PYZdq{}}\PY{l+s+s2}{Fecha}\PY{l+s+s2}{\PYZdq{}}\PY{p}{)}

\PY{n}{plt}\PY{o}{.}\PY{n}{grid}\PY{p}{(}\PY{k+kc}{True}\PY{p}{)}

\PY{n}{plt}\PY{o}{.}\PY{n}{show}\PY{p}{(}\PY{p}{)}
\end{Verbatim}
\end{tcolorbox}

    \begin{center}
    \adjustimage{max size={0.9\linewidth}{0.9\paperheight}}{output_181_0.png}
    \end{center}
    { \hspace*{\fill} \\}
    
    \begin{tcolorbox}[breakable, size=fbox, boxrule=1pt, pad at break*=1mm,colback=cellbackground, colframe=cellborder]
\prompt{In}{incolor}{132}{\boxspacing}
\begin{Verbatim}[commandchars=\\\{\}]
\PY{n}{pdf\PYZus{}completo}\PY{p}{[}\PY{l+s+s2}{\PYZdq{}}\PY{l+s+s2}{fecha}\PY{l+s+s2}{\PYZdq{}}\PY{p}{]} \PY{o}{=} \PY{n}{pd}\PY{o}{.}\PY{n}{to\PYZus{}datetime}\PY{p}{(}
    \PY{n}{pdf\PYZus{}completo}\PY{p}{[}\PY{l+s+s2}{\PYZdq{}}\PY{l+s+s2}{fecha}\PY{l+s+s2}{\PYZdq{}}\PY{p}{]}
\PY{p}{)}

\PY{n}{pdf\PYZus{}tortilla}\PY{p}{[}\PY{l+s+s2}{\PYZdq{}}\PY{l+s+s2}{fecha}\PY{l+s+s2}{\PYZdq{}}\PY{p}{]} \PY{o}{=} \PY{n}{pd}\PY{o}{.}\PY{n}{to\PYZus{}datetime}\PY{p}{(}
    \PY{n}{pdf\PYZus{}tortilla}\PY{p}{[}\PY{l+s+s2}{\PYZdq{}}\PY{l+s+s2}{fecha}\PY{l+s+s2}{\PYZdq{}}\PY{p}{]}
\PY{p}{)}
\end{Verbatim}
\end{tcolorbox}

    \begin{tcolorbox}[breakable, size=fbox, boxrule=1pt, pad at break*=1mm,colback=cellbackground, colframe=cellborder]
\prompt{In}{incolor}{133}{\boxspacing}
\begin{Verbatim}[commandchars=\\\{\}]
\PY{n+nb}{print}\PY{p}{(}\PY{n}{pdf\PYZus{}completo}\PY{p}{[}\PY{l+s+s2}{\PYZdq{}}\PY{l+s+s2}{fecha}\PY{l+s+s2}{\PYZdq{}}\PY{p}{]}\PY{o}{.}\PY{n}{dtype}\PY{p}{)}
\PY{n+nb}{print}\PY{p}{(}\PY{n}{pdf\PYZus{}tortilla}\PY{p}{[}\PY{l+s+s2}{\PYZdq{}}\PY{l+s+s2}{fecha}\PY{l+s+s2}{\PYZdq{}}\PY{p}{]}\PY{o}{.}\PY{n}{dtype}\PY{p}{)}
\end{Verbatim}
\end{tcolorbox}

    \begin{Verbatim}[commandchars=\\\{\}]
datetime64[ns]
datetime64[ns]
    \end{Verbatim}

    \begin{tcolorbox}[breakable, size=fbox, boxrule=1pt, pad at break*=1mm,colback=cellbackground, colframe=cellborder]
\prompt{In}{incolor}{210}{\boxspacing}
\begin{Verbatim}[commandchars=\\\{\}]
\PY{n}{maiz\PYZus{}tortilla} \PY{o}{=} \PY{p}{(}
    \PY{n}{pdf\PYZus{}completo}\PY{p}{[}
        \PY{p}{[}\PY{l+s+s2}{\PYZdq{}}\PY{l+s+s2}{fecha}\PY{l+s+s2}{\PYZdq{}}\PY{p}{,}
         \PY{l+s+s2}{\PYZdq{}}\PY{l+s+s2}{Precio\PYZus{}maizblanco\PYZus{}mxpkilo\PYZus{}prom\PYZus{}nacional}\PY{l+s+s2}{\PYZdq{}}\PY{p}{]}
    \PY{p}{]}
    \PY{o}{.}\PY{n}{merge}\PY{p}{(}
        \PY{n}{pdf\PYZus{}tortilla}\PY{p}{[}
            \PY{p}{[}\PY{l+s+s2}{\PYZdq{}}\PY{l+s+s2}{fecha}\PY{l+s+s2}{\PYZdq{}}\PY{p}{,} \PY{l+s+s2}{\PYZdq{}}\PY{l+s+s2}{Nacional}\PY{l+s+s2}{\PYZdq{}}\PY{p}{]}
        \PY{p}{]}\PY{p}{,}
        \PY{n}{on}\PY{o}{=}\PY{l+s+s2}{\PYZdq{}}\PY{l+s+s2}{fecha}\PY{l+s+s2}{\PYZdq{}}\PY{p}{,}
        \PY{n}{how}\PY{o}{=}\PY{l+s+s2}{\PYZdq{}}\PY{l+s+s2}{inner}\PY{l+s+s2}{\PYZdq{}}
    \PY{p}{)}
\PY{p}{)}

\PY{n}{maiz\PYZus{}tortilla} \PY{o}{=} \PY{n}{maiz\PYZus{}tortilla}\PY{o}{.}\PY{n}{rename}\PY{p}{(}
    \PY{n}{columns}\PY{o}{=}\PY{p}{\PYZob{}}
        \PY{l+s+s2}{\PYZdq{}}\PY{l+s+s2}{Precio\PYZus{}maizblanco\PYZus{}mxpkilo\PYZus{}prom\PYZus{}nacional}\PY{l+s+s2}{\PYZdq{}}\PY{p}{:} \PY{l+s+s2}{\PYZdq{}}\PY{l+s+s2}{Precio\PYZus{}maiz}\PY{l+s+s2}{\PYZdq{}}\PY{p}{,}
        \PY{l+s+s2}{\PYZdq{}}\PY{l+s+s2}{Nacional}\PY{l+s+s2}{\PYZdq{}}\PY{p}{:} \PY{l+s+s2}{\PYZdq{}}\PY{l+s+s2}{Precio\PYZus{}tortilla}\PY{l+s+s2}{\PYZdq{}}
    \PY{p}{\PYZcb{}}
\PY{p}{)}

\PY{n}{maiz\PYZus{}tortilla}\PY{o}{.}\PY{n}{tail}\PY{p}{(}\PY{p}{)}
\end{Verbatim}
\end{tcolorbox}

            \begin{tcolorbox}[breakable, size=fbox, boxrule=.5pt, pad at break*=1mm, opacityfill=0]
\prompt{Out}{outcolor}{210}{\boxspacing}
\begin{Verbatim}[commandchars=\\\{\}]
         fecha  Precio\_maiz  Precio\_tortilla
312 2026-01-01         9.40        23.905623
313 2026-02-01         9.48        24.044798
314 2026-03-01         9.58        24.119832
315 2026-04-01         9.75        24.345754
316 2026-05-01         9.82        25.019300
\end{Verbatim}
\end{tcolorbox}
        
    \begin{tcolorbox}[breakable, size=fbox, boxrule=1pt, pad at break*=1mm,colback=cellbackground, colframe=cellborder]
\prompt{In}{incolor}{135}{\boxspacing}
\begin{Verbatim}[commandchars=\\\{\}]
\PY{n}{fig}\PY{p}{,} \PY{n}{ax1} \PY{o}{=} \PY{n}{plt}\PY{o}{.}\PY{n}{subplots}\PY{p}{(}\PY{n}{figsize}\PY{o}{=}\PY{p}{(}\PY{l+m+mi}{14}\PY{p}{,}\PY{l+m+mi}{6}\PY{p}{)}\PY{p}{)}

\PY{n}{ax1}\PY{o}{.}\PY{n}{plot}\PY{p}{(}
    \PY{n}{maiz\PYZus{}tortilla}\PY{p}{[}\PY{l+s+s2}{\PYZdq{}}\PY{l+s+s2}{fecha}\PY{l+s+s2}{\PYZdq{}}\PY{p}{]}\PY{p}{,}
    \PY{n}{maiz\PYZus{}tortilla}\PY{p}{[}\PY{l+s+s2}{\PYZdq{}}\PY{l+s+s2}{Precio\PYZus{}maiz}\PY{l+s+s2}{\PYZdq{}}\PY{p}{]}\PY{p}{,}
    \PY{n}{label}\PY{o}{=}\PY{l+s+s2}{\PYZdq{}}\PY{l+s+s2}{Maíz}\PY{l+s+s2}{\PYZdq{}}\PY{p}{,}
    \PY{n}{color}\PY{o}{=}\PY{l+s+s2}{\PYZdq{}}\PY{l+s+s2}{tab:blue}\PY{l+s+s2}{\PYZdq{}}
\PY{p}{)}

\PY{n}{ax1}\PY{o}{.}\PY{n}{set\PYZus{}ylabel}\PY{p}{(}
    \PY{l+s+s2}{\PYZdq{}}\PY{l+s+s2}{Precio Maíz (MXN/Kg)}\PY{l+s+s2}{\PYZdq{}}\PY{p}{,}
    \PY{n}{color}\PY{o}{=}\PY{l+s+s2}{\PYZdq{}}\PY{l+s+s2}{tab:blue}\PY{l+s+s2}{\PYZdq{}}
\PY{p}{)}

\PY{n}{ax2} \PY{o}{=} \PY{n}{ax1}\PY{o}{.}\PY{n}{twinx}\PY{p}{(}\PY{p}{)}

\PY{n}{ax2}\PY{o}{.}\PY{n}{plot}\PY{p}{(}
    \PY{n}{maiz\PYZus{}tortilla}\PY{p}{[}\PY{l+s+s2}{\PYZdq{}}\PY{l+s+s2}{fecha}\PY{l+s+s2}{\PYZdq{}}\PY{p}{]}\PY{p}{,}
    \PY{n}{maiz\PYZus{}tortilla}\PY{p}{[}\PY{l+s+s2}{\PYZdq{}}\PY{l+s+s2}{Precio\PYZus{}tortilla}\PY{l+s+s2}{\PYZdq{}}\PY{p}{]}\PY{p}{,}
    \PY{n}{label}\PY{o}{=}\PY{l+s+s2}{\PYZdq{}}\PY{l+s+s2}{Tortilla}\PY{l+s+s2}{\PYZdq{}}\PY{p}{,}
    \PY{n}{color}\PY{o}{=}\PY{l+s+s2}{\PYZdq{}}\PY{l+s+s2}{tab:red}\PY{l+s+s2}{\PYZdq{}}
\PY{p}{)}

\PY{n}{ax2}\PY{o}{.}\PY{n}{set\PYZus{}ylabel}\PY{p}{(}
    \PY{l+s+s2}{\PYZdq{}}\PY{l+s+s2}{Precio Tortilla (MXN/Kg)}\PY{l+s+s2}{\PYZdq{}}\PY{p}{,}
    \PY{n}{color}\PY{o}{=}\PY{l+s+s2}{\PYZdq{}}\PY{l+s+s2}{tab:red}\PY{l+s+s2}{\PYZdq{}}
\PY{p}{)}

\PY{n}{plt}\PY{o}{.}\PY{n}{title}\PY{p}{(}
    \PY{l+s+s2}{\PYZdq{}}\PY{l+s+s2}{Precio Nacional del Maíz Blanco vs Precio Nacional de la Tortilla}\PY{l+s+s2}{\PYZdq{}}
\PY{p}{)}

\PY{n}{plt}\PY{o}{.}\PY{n}{show}\PY{p}{(}\PY{p}{)}
\end{Verbatim}
\end{tcolorbox}

    \begin{center}
    \adjustimage{max size={0.9\linewidth}{0.9\paperheight}}{output_185_0.png}
    \end{center}
    { \hspace*{\fill} \\}
    
    \begin{tcolorbox}[breakable, size=fbox, boxrule=1pt, pad at break*=1mm,colback=cellbackground, colframe=cellborder]
\prompt{In}{incolor}{136}{\boxspacing}
\begin{Verbatim}[commandchars=\\\{\}]
\PY{n}{corr\PYZus{}maiz\PYZus{}tortilla} \PY{o}{=} \PY{p}{(}
    \PY{n}{maiz\PYZus{}tortilla}\PY{p}{[}
        \PY{p}{[}\PY{l+s+s2}{\PYZdq{}}\PY{l+s+s2}{Precio\PYZus{}maiz}\PY{l+s+s2}{\PYZdq{}}\PY{p}{,}\PY{l+s+s2}{\PYZdq{}}\PY{l+s+s2}{Precio\PYZus{}tortilla}\PY{l+s+s2}{\PYZdq{}}\PY{p}{]}
    \PY{p}{]}
    \PY{o}{.}\PY{n}{corr}\PY{p}{(}\PY{p}{)}
\PY{p}{)}

\PY{n}{corr\PYZus{}maiz\PYZus{}tortilla}
\end{Verbatim}
\end{tcolorbox}

            \begin{tcolorbox}[breakable, size=fbox, boxrule=.5pt, pad at break*=1mm, opacityfill=0]
\prompt{Out}{outcolor}{136}{\boxspacing}
\begin{Verbatim}[commandchars=\\\{\}]
                 Precio\_maiz  Precio\_tortilla
Precio\_maiz         1.000000         0.972386
Precio\_tortilla     0.972386         1.000000
\end{Verbatim}
\end{tcolorbox}
        
    \begin{tcolorbox}[breakable, size=fbox, boxrule=1pt, pad at break*=1mm,colback=cellbackground, colframe=cellborder]
\prompt{In}{incolor}{138}{\boxspacing}
\begin{Verbatim}[commandchars=\\\{\}]
\PY{n}{plt}\PY{o}{.}\PY{n}{figure}\PY{p}{(}\PY{n}{figsize}\PY{o}{=}\PY{p}{(}\PY{l+m+mi}{8}\PY{p}{,}\PY{l+m+mi}{6}\PY{p}{)}\PY{p}{)}

\PY{n}{plt}\PY{o}{.}\PY{n}{scatter}\PY{p}{(}
    \PY{n}{maiz\PYZus{}tortilla}\PY{p}{[}\PY{l+s+s2}{\PYZdq{}}\PY{l+s+s2}{Precio\PYZus{}maiz}\PY{l+s+s2}{\PYZdq{}}\PY{p}{]}\PY{p}{,}
    \PY{n}{maiz\PYZus{}tortilla}\PY{p}{[}\PY{l+s+s2}{\PYZdq{}}\PY{l+s+s2}{Precio\PYZus{}tortilla}\PY{l+s+s2}{\PYZdq{}}\PY{p}{]}\PY{p}{,}
    \PY{n}{alpha}\PY{o}{=}\PY{l+m+mf}{0.7}
\PY{p}{)}

\PY{n}{plt}\PY{o}{.}\PY{n}{xlabel}\PY{p}{(}\PY{l+s+s2}{\PYZdq{}}\PY{l+s+s2}{Precio Maíz (MXN/Kg)}\PY{l+s+s2}{\PYZdq{}}\PY{p}{)}
\PY{n}{plt}\PY{o}{.}\PY{n}{ylabel}\PY{p}{(}\PY{l+s+s2}{\PYZdq{}}\PY{l+s+s2}{Precio Tortilla (MXN/Kg)}\PY{l+s+s2}{\PYZdq{}}\PY{p}{)}

\PY{n}{plt}\PY{o}{.}\PY{n}{title}\PY{p}{(}
    \PY{l+s+s2}{\PYZdq{}}\PY{l+s+s2}{Relación entre Precio del Maíz y Precio de la Tortilla}\PY{l+s+s2}{\PYZdq{}}
\PY{p}{)}

\PY{n}{plt}\PY{o}{.}\PY{n}{grid}\PY{p}{(}\PY{k+kc}{True}\PY{p}{)}

\PY{n}{plt}\PY{o}{.}\PY{n}{show}\PY{p}{(}\PY{p}{)}
\end{Verbatim}
\end{tcolorbox}

    \begin{center}
    \adjustimage{max size={0.9\linewidth}{0.9\paperheight}}{output_187_0.png}
    \end{center}
    { \hspace*{\fill} \\}
    
    \begin{tcolorbox}[breakable, size=fbox, boxrule=1pt, pad at break*=1mm,colback=cellbackground, colframe=cellborder]
\prompt{In}{incolor}{139}{\boxspacing}
\begin{Verbatim}[commandchars=\\\{\}]
\PY{n}{rezagos} \PY{o}{=} \PY{p}{[}\PY{p}{]}

\PY{k}{for} \PY{n}{lag} \PY{o+ow}{in} \PY{n+nb}{range}\PY{p}{(}\PY{l+m+mi}{13}\PY{p}{)}\PY{p}{:}

    \PY{n}{corr} \PY{o}{=} \PY{p}{(}
        \PY{n}{maiz\PYZus{}tortilla}\PY{p}{[}\PY{l+s+s2}{\PYZdq{}}\PY{l+s+s2}{Precio\PYZus{}maiz}\PY{l+s+s2}{\PYZdq{}}\PY{p}{]}
        \PY{o}{.}\PY{n}{corr}\PY{p}{(}
            \PY{n}{maiz\PYZus{}tortilla}\PY{p}{[}\PY{l+s+s2}{\PYZdq{}}\PY{l+s+s2}{Precio\PYZus{}tortilla}\PY{l+s+s2}{\PYZdq{}}\PY{p}{]}\PY{o}{.}\PY{n}{shift}\PY{p}{(}\PY{o}{\PYZhy{}}\PY{n}{lag}\PY{p}{)}
        \PY{p}{)}
    \PY{p}{)}

    \PY{n}{rezagos}\PY{o}{.}\PY{n}{append}\PY{p}{(}\PY{p}{[}\PY{n}{lag}\PY{p}{,} \PY{n}{corr}\PY{p}{]}\PY{p}{)}

\PY{n}{corr\PYZus{}rezagos} \PY{o}{=} \PY{n}{pd}\PY{o}{.}\PY{n}{DataFrame}\PY{p}{(}
    \PY{n}{rezagos}\PY{p}{,}
    \PY{n}{columns}\PY{o}{=}\PY{p}{[}\PY{l+s+s2}{\PYZdq{}}\PY{l+s+s2}{Rezago\PYZus{}meses}\PY{l+s+s2}{\PYZdq{}}\PY{p}{,}\PY{l+s+s2}{\PYZdq{}}\PY{l+s+s2}{Correlacion}\PY{l+s+s2}{\PYZdq{}}\PY{p}{]}
\PY{p}{)}

\PY{n}{corr\PYZus{}rezagos}
\end{Verbatim}
\end{tcolorbox}

            \begin{tcolorbox}[breakable, size=fbox, boxrule=.5pt, pad at break*=1mm, opacityfill=0]
\prompt{Out}{outcolor}{139}{\boxspacing}
\begin{Verbatim}[commandchars=\\\{\}]
    Rezago\_meses  Correlacion
0              0     0.972386
1              1     0.972193
2              2     0.971742
3              3     0.971059
4              4     0.970101
5              5     0.968996
6              6     0.967918
7              7     0.966671
8              8     0.965297
9              9     0.963736
10            10     0.962105
11            11     0.960479
12            12     0.958865
\end{Verbatim}
\end{tcolorbox}
        
    \begin{tcolorbox}[breakable, size=fbox, boxrule=1pt, pad at break*=1mm,colback=cellbackground, colframe=cellborder]
\prompt{In}{incolor}{140}{\boxspacing}
\begin{Verbatim}[commandchars=\\\{\}]
\PY{n}{plt}\PY{o}{.}\PY{n}{figure}\PY{p}{(}\PY{n}{figsize}\PY{o}{=}\PY{p}{(}\PY{l+m+mi}{10}\PY{p}{,}\PY{l+m+mi}{5}\PY{p}{)}\PY{p}{)}

\PY{n}{plt}\PY{o}{.}\PY{n}{bar}\PY{p}{(}
    \PY{n}{corr\PYZus{}rezagos}\PY{p}{[}\PY{l+s+s2}{\PYZdq{}}\PY{l+s+s2}{Rezago\PYZus{}meses}\PY{l+s+s2}{\PYZdq{}}\PY{p}{]}\PY{p}{,}
    \PY{n}{corr\PYZus{}rezagos}\PY{p}{[}\PY{l+s+s2}{\PYZdq{}}\PY{l+s+s2}{Correlacion}\PY{l+s+s2}{\PYZdq{}}\PY{p}{]}
\PY{p}{)}

\PY{n}{plt}\PY{o}{.}\PY{n}{title}\PY{p}{(}
    \PY{l+s+s2}{\PYZdq{}}\PY{l+s+s2}{Correlación Maíz → Tortilla por Rezago}\PY{l+s+s2}{\PYZdq{}}
\PY{p}{)}

\PY{n}{plt}\PY{o}{.}\PY{n}{xlabel}\PY{p}{(}\PY{l+s+s2}{\PYZdq{}}\PY{l+s+s2}{Rezago (meses)}\PY{l+s+s2}{\PYZdq{}}\PY{p}{)}
\PY{n}{plt}\PY{o}{.}\PY{n}{ylabel}\PY{p}{(}\PY{l+s+s2}{\PYZdq{}}\PY{l+s+s2}{Correlación}\PY{l+s+s2}{\PYZdq{}}\PY{p}{)}

\PY{n}{plt}\PY{o}{.}\PY{n}{grid}\PY{p}{(}\PY{k+kc}{True}\PY{p}{)}

\PY{n}{plt}\PY{o}{.}\PY{n}{show}\PY{p}{(}\PY{p}{)}
\end{Verbatim}
\end{tcolorbox}

    \begin{center}
    \adjustimage{max size={0.9\linewidth}{0.9\paperheight}}{output_189_0.png}
    \end{center}
    { \hspace*{\fill} \\}
    
    \begin{tcolorbox}[breakable, size=fbox, boxrule=1pt, pad at break*=1mm,colback=cellbackground, colframe=cellborder]
\prompt{In}{incolor}{141}{\boxspacing}
\begin{Verbatim}[commandchars=\\\{\}]
\PY{n}{corr\PYZus{}rezagos}\PY{o}{.}\PY{n}{sort\PYZus{}values}\PY{p}{(}
    \PY{n}{by}\PY{o}{=}\PY{l+s+s2}{\PYZdq{}}\PY{l+s+s2}{Correlacion}\PY{l+s+s2}{\PYZdq{}}\PY{p}{,}
    \PY{n}{ascending}\PY{o}{=}\PY{k+kc}{False}
\PY{p}{)}\PY{o}{.}\PY{n}{head}\PY{p}{(}\PY{p}{)}
\end{Verbatim}
\end{tcolorbox}

            \begin{tcolorbox}[breakable, size=fbox, boxrule=.5pt, pad at break*=1mm, opacityfill=0]
\prompt{Out}{outcolor}{141}{\boxspacing}
\begin{Verbatim}[commandchars=\\\{\}]
   Rezago\_meses  Correlacion
0             0     0.972386
1             1     0.972193
2             2     0.971742
3             3     0.971059
4             4     0.970101
\end{Verbatim}
\end{tcolorbox}
        
    \begin{tcolorbox}[breakable, size=fbox, boxrule=1pt, pad at break*=1mm,colback=cellbackground, colframe=cellborder]
\prompt{In}{incolor}{142}{\boxspacing}
\begin{Verbatim}[commandchars=\\\{\}]
\PY{n}{maiz\PYZus{}tortilla}\PY{p}{[}\PY{l+s+s2}{\PYZdq{}}\PY{l+s+s2}{Delta\PYZus{}maiz}\PY{l+s+s2}{\PYZdq{}}\PY{p}{]} \PY{o}{=} \PY{p}{(}
    \PY{n}{maiz\PYZus{}tortilla}\PY{p}{[}\PY{l+s+s2}{\PYZdq{}}\PY{l+s+s2}{Precio\PYZus{}maiz}\PY{l+s+s2}{\PYZdq{}}\PY{p}{]}\PY{o}{.}\PY{n}{diff}\PY{p}{(}\PY{p}{)}
\PY{p}{)}

\PY{n}{maiz\PYZus{}tortilla}\PY{p}{[}\PY{l+s+s2}{\PYZdq{}}\PY{l+s+s2}{Delta\PYZus{}tortilla}\PY{l+s+s2}{\PYZdq{}}\PY{p}{]} \PY{o}{=} \PY{p}{(}
    \PY{n}{maiz\PYZus{}tortilla}\PY{p}{[}\PY{l+s+s2}{\PYZdq{}}\PY{l+s+s2}{Precio\PYZus{}tortilla}\PY{l+s+s2}{\PYZdq{}}\PY{p}{]}\PY{o}{.}\PY{n}{diff}\PY{p}{(}\PY{p}{)}
\PY{p}{)}
\end{Verbatim}
\end{tcolorbox}

    \begin{tcolorbox}[breakable, size=fbox, boxrule=1pt, pad at break*=1mm,colback=cellbackground, colframe=cellborder]
\prompt{In}{incolor}{143}{\boxspacing}
\begin{Verbatim}[commandchars=\\\{\}]
\PY{n}{maiz\PYZus{}tortilla}\PY{p}{[}
    \PY{p}{[}\PY{l+s+s2}{\PYZdq{}}\PY{l+s+s2}{Delta\PYZus{}maiz}\PY{l+s+s2}{\PYZdq{}}\PY{p}{,}\PY{l+s+s2}{\PYZdq{}}\PY{l+s+s2}{Delta\PYZus{}tortilla}\PY{l+s+s2}{\PYZdq{}}\PY{p}{]}
\PY{p}{]}\PY{o}{.}\PY{n}{corr}\PY{p}{(}\PY{p}{)}
\end{Verbatim}
\end{tcolorbox}

            \begin{tcolorbox}[breakable, size=fbox, boxrule=.5pt, pad at break*=1mm, opacityfill=0]
\prompt{Out}{outcolor}{143}{\boxspacing}
\begin{Verbatim}[commandchars=\\\{\}]
                Delta\_maiz  Delta\_tortilla
Delta\_maiz        1.000000        0.265832
Delta\_tortilla    0.265832        1.000000
\end{Verbatim}
\end{tcolorbox}
        
    \hypertarget{transmisiuxf3n-de-variaciones-del-mauxedz-hacia-la-tortilla}{%
\section{38. Transmisión de Variaciones del Maíz hacia la
Tortilla}\label{transmisiuxf3n-de-variaciones-del-mauxedz-hacia-la-tortilla}}

\hypertarget{correlaciuxf3n-de-niveles}{%
\subsection{Correlación de niveles}\label{correlaciuxf3n-de-niveles}}

La correlación entre el precio nacional del maíz blanco y el precio
nacional de la tortilla fue:

\[
r = 0.972
\]

Este resultado evidencia una asociación positiva extremadamente fuerte
entre ambas series.

No obstante, debido a que ambas variables presentan tendencias
crecientes de largo plazo, dicha correlación podría estar influida por
factores estructurales compartidos.

\begin{center}\rule{0.5\linewidth}{0.5pt}\end{center}

\hypertarget{correlaciuxf3n-de-variaciones}{%
\subsection{Correlación de
variaciones}\label{correlaciuxf3n-de-variaciones}}

Con el propósito de evaluar la transmisión de cambios mensuales, se
calcularon las primeras diferencias de ambas series:

\[
\Delta Maíz_t
=
Maíz_t - Maíz_{t-1}
\]

\[
\Delta Tortilla_t
=
Tortilla_t - Tortilla_{t-1}
\]

La correlación obtenida fue:

\[
r = 0.266
\]

\begin{center}\rule{0.5\linewidth}{0.5pt}\end{center}

\hypertarget{interpretaciuxf3n}{%
\subsection{Interpretación}\label{interpretaciuxf3n}}

Los resultados sugieren que el precio del maíz constituye un factor
relevante para explicar la evolución de largo plazo del precio de la
tortilla.

Sin embargo, los cambios mensuales observados en la tortilla responden
también a otros componentes de costo presentes en la cadena productiva y
comercial.

Entre ellos destacan:

\begin{itemize}
\tightlist
\item
  Energéticos.
\item
  Transporte.
\item
  Mano de obra.
\item
  Costos operativos.
\item
  Inflación general.
\end{itemize}

Por lo tanto, el precio del maíz representa una condición necesaria pero
no suficiente para explicar completamente la dinámica de precios de la
tortilla.

\begin{center}\rule{0.5\linewidth}{0.5pt}\end{center}

\hypertarget{conclusiuxf3n}{%
\subsection{Conclusión}\label{conclusiuxf3n}}

La evidencia obtenida indica que existe una transmisión parcial de
precios desde el mercado del maíz hacia el mercado de la tortilla.

No obstante, dicha transmisión se encuentra moderada por múltiples
factores económicos adicionales que influyen sobre el precio final
pagado por los consumidores.

    \begin{tcolorbox}[breakable, size=fbox, boxrule=1pt, pad at break*=1mm,colback=cellbackground, colframe=cellborder]
\prompt{In}{incolor}{144}{\boxspacing}
\begin{Verbatim}[commandchars=\\\{\}]
\PY{n}{maiz\PYZus{}tortilla}\PY{p}{[}\PY{l+s+s2}{\PYZdq{}}\PY{l+s+s2}{Participacion\PYZus{}maiz}\PY{l+s+s2}{\PYZdq{}}\PY{p}{]} \PY{o}{=} \PY{p}{(}
    \PY{n}{maiz\PYZus{}tortilla}\PY{p}{[}\PY{l+s+s2}{\PYZdq{}}\PY{l+s+s2}{Precio\PYZus{}maiz}\PY{l+s+s2}{\PYZdq{}}\PY{p}{]}
    \PY{o}{/} \PY{n}{maiz\PYZus{}tortilla}\PY{p}{[}\PY{l+s+s2}{\PYZdq{}}\PY{l+s+s2}{Precio\PYZus{}tortilla}\PY{l+s+s2}{\PYZdq{}}\PY{p}{]}
\PY{p}{)} \PY{o}{*} \PY{l+m+mi}{100}

\PY{n}{plt}\PY{o}{.}\PY{n}{figure}\PY{p}{(}\PY{n}{figsize}\PY{o}{=}\PY{p}{(}\PY{l+m+mi}{14}\PY{p}{,}\PY{l+m+mi}{6}\PY{p}{)}\PY{p}{)}

\PY{n}{plt}\PY{o}{.}\PY{n}{plot}\PY{p}{(}
    \PY{n}{maiz\PYZus{}tortilla}\PY{p}{[}\PY{l+s+s2}{\PYZdq{}}\PY{l+s+s2}{fecha}\PY{l+s+s2}{\PYZdq{}}\PY{p}{]}\PY{p}{,}
    \PY{n}{maiz\PYZus{}tortilla}\PY{p}{[}\PY{l+s+s2}{\PYZdq{}}\PY{l+s+s2}{Participacion\PYZus{}maiz}\PY{l+s+s2}{\PYZdq{}}\PY{p}{]}
\PY{p}{)}

\PY{n}{plt}\PY{o}{.}\PY{n}{title}\PY{p}{(}
    \PY{l+s+s2}{\PYZdq{}}\PY{l+s+s2}{Participación del Maíz en el Precio de la Tortilla}\PY{l+s+s2}{\PYZdq{}}
\PY{p}{)}

\PY{n}{plt}\PY{o}{.}\PY{n}{ylabel}\PY{p}{(}\PY{l+s+s2}{\PYZdq{}}\PY{l+s+s2}{\PYZpc{}}\PY{l+s+s2}{\PYZdq{}}\PY{p}{)}

\PY{n}{plt}\PY{o}{.}\PY{n}{grid}\PY{p}{(}\PY{k+kc}{True}\PY{p}{)}

\PY{n}{plt}\PY{o}{.}\PY{n}{show}\PY{p}{(}\PY{p}{)}
\end{Verbatim}
\end{tcolorbox}

    \begin{center}
    \adjustimage{max size={0.9\linewidth}{0.9\paperheight}}{output_194_0.png}
    \end{center}
    { \hspace*{\fill} \\}
    
    \begin{tcolorbox}[breakable, size=fbox, boxrule=1pt, pad at break*=1mm,colback=cellbackground, colframe=cellborder]
\prompt{In}{incolor}{145}{\boxspacing}
\begin{Verbatim}[commandchars=\\\{\}]
\PY{n}{maiz\PYZus{}tortilla}\PY{p}{[}\PY{l+s+s2}{\PYZdq{}}\PY{l+s+s2}{Participacion\PYZus{}maiz}\PY{l+s+s2}{\PYZdq{}}\PY{p}{]}\PY{o}{.}\PY{n}{describe}\PY{p}{(}\PY{p}{)}
\end{Verbatim}
\end{tcolorbox}

            \begin{tcolorbox}[breakable, size=fbox, boxrule=.5pt, pad at break*=1mm, opacityfill=0]
\prompt{Out}{outcolor}{145}{\boxspacing}
\begin{Verbatim}[commandchars=\\\{\}]
count    317.000000
mean      39.045728
std        5.266580
min       30.902388
25\%       34.398769
50\%       38.486266
75\%       43.300924
max       51.707317
Name: Participacion\_maiz, dtype: float64
\end{Verbatim}
\end{tcolorbox}
        
    \hypertarget{participaciuxf3n-del-mauxedz-en-el-precio-de-la-tortilla}{%
\section{39. Participación del Maíz en el Precio de la
Tortilla}\label{participaciuxf3n-del-mauxedz-en-el-precio-de-la-tortilla}}

Con el propósito de dimensionar la importancia relativa del maíz dentro
del precio final de la tortilla, se calculó la proporción mensual:

\[
Participación=
\frac{Precio_{Maíz}}
{Precio_{Tortilla}}
\times100
\]

\begin{center}\rule{0.5\linewidth}{0.5pt}\end{center}

\hypertarget{resultados}{%
\subsection{Resultados}\label{resultados}}

Los resultados obtenidos muestran:

\begin{itemize}
\tightlist
\item
  Participación promedio: 39.05\%.
\item
  Participación mediana: 38.49\%.
\item
  Participación mínima: 30.90\%.
\item
  Participación máxima: 51.71\%.
\end{itemize}

\begin{center}\rule{0.5\linewidth}{0.5pt}\end{center}

\hypertarget{interpretaciuxf3n}{%
\subsection{Interpretación}\label{interpretaciuxf3n}}

La evidencia indica que el maíz representa, en promedio, aproximadamente
cuatro décimas partes del precio final de la tortilla.

En consecuencia, alrededor del 61\% del precio pagado por los
consumidores se encuentra asociado a otros componentes de costo
presentes en la cadena productiva y comercial.

Entre dichos componentes destacan:

\begin{itemize}
\tightlist
\item
  Transporte.
\item
  Combustibles.
\item
  Energía.
\item
  Mano de obra.
\item
  Distribución.
\item
  Comercialización.
\item
  Inflación general.
\end{itemize}

\begin{center}\rule{0.5\linewidth}{0.5pt}\end{center}

\hypertarget{hallazgo-principal}{%
\subsection{Hallazgo principal}\label{hallazgo-principal}}

A pesar de que el maíz constituye el principal insumo para la
elaboración de tortillas, su participación relativa dentro del precio
final ha disminuido respecto a los primeros años de la serie analizada.

Este resultado sugiere que la evolución reciente del precio de la
tortilla depende no sólo del comportamiento del mercado agrícola, sino
también de factores logísticos, energéticos y económicos que afectan al
conjunto de la cadena de valor.

\begin{center}\rule{0.5\linewidth}{0.5pt}\end{center}

\hypertarget{conclusiuxf3n-general-del-proyecto}{%
\subsection{Conclusión general del
proyecto}\label{conclusiuxf3n-general-del-proyecto}}

Los resultados obtenidos muestran que el precio nacional del maíz blanco
está influido principalmente por factores asociados al mercado
internacional, los costos de transporte y las condiciones climáticas.

Asimismo, se encontró evidencia de una relación estructural fuerte entre
el precio del maíz y el precio de la tortilla. Sin embargo, el análisis
de participación relativa indica que el maíz explica únicamente una
parte del precio final pagado por los consumidores.

Por lo tanto, las estrategias orientadas a mejorar la accesibilidad de
los alimentos básicos requieren considerar no sólo la producción
agrícola, sino también los costos logísticos, energéticos y de
distribución que intervienen en la formación de precios a lo largo de
toda la cadena agroalimentaria.

    \hypertarget{insight-para-consumidores}{%
\subsection{Insight para consumidores}\label{insight-para-consumidores}}

Los resultados muestran que, en promedio, únicamente el 39\% del precio
de la tortilla corresponde al costo del maíz.

Esto implica que por cada \$100 pesos gastados en tortillas:

\begin{itemize}
\tightlist
\item
  Aproximadamente \$39 pesos corresponden al maíz.
\item
  Aproximadamente \$61 pesos corresponden a transporte, energía,
  combustibles, salarios, distribución y otros costos asociados a la
  cadena de valor.
\end{itemize}

En consecuencia, las variaciones en el precio del maíz influyen sobre el
precio final de la tortilla, pero no constituyen el único factor
determinante.

En términos prácticos, el precio de la tortilla comienza en el campo,
pero se construye a lo largo de toda la cadena agroalimentaria.

    \begin{tcolorbox}[breakable, size=fbox, boxrule=1pt, pad at break*=1mm,colback=cellbackground, colframe=cellborder]
\prompt{In}{incolor}{146}{\boxspacing}
\begin{Verbatim}[commandchars=\\\{\}]
\PY{n}{maiz\PYZus{}tortilla}\PY{p}{[}\PY{l+s+s2}{\PYZdq{}}\PY{l+s+s2}{Diferencia\PYZus{}pesos}\PY{l+s+s2}{\PYZdq{}}\PY{p}{]} \PY{o}{=} \PY{p}{(}
    \PY{n}{maiz\PYZus{}tortilla}\PY{p}{[}\PY{l+s+s2}{\PYZdq{}}\PY{l+s+s2}{Precio\PYZus{}tortilla}\PY{l+s+s2}{\PYZdq{}}\PY{p}{]}
    \PY{o}{\PYZhy{}} \PY{n}{maiz\PYZus{}tortilla}\PY{p}{[}\PY{l+s+s2}{\PYZdq{}}\PY{l+s+s2}{Precio\PYZus{}maiz}\PY{l+s+s2}{\PYZdq{}}\PY{p}{]}
\PY{p}{)}

\PY{n}{maiz\PYZus{}tortilla}\PY{p}{[}
    \PY{p}{[}\PY{l+s+s2}{\PYZdq{}}\PY{l+s+s2}{Precio\PYZus{}maiz}\PY{l+s+s2}{\PYZdq{}}\PY{p}{,}
     \PY{l+s+s2}{\PYZdq{}}\PY{l+s+s2}{Precio\PYZus{}tortilla}\PY{l+s+s2}{\PYZdq{}}\PY{p}{,}
     \PY{l+s+s2}{\PYZdq{}}\PY{l+s+s2}{Diferencia\PYZus{}pesos}\PY{l+s+s2}{\PYZdq{}}\PY{p}{]}
\PY{p}{]}\PY{o}{.}\PY{n}{describe}\PY{p}{(}\PY{p}{)}
\end{Verbatim}
\end{tcolorbox}

            \begin{tcolorbox}[breakable, size=fbox, boxrule=.5pt, pad at break*=1mm, opacityfill=0]
\prompt{Out}{outcolor}{146}{\boxspacing}
\begin{Verbatim}[commandchars=\\\{\}]
       Precio\_maiz  Precio\_tortilla  Diferencia\_pesos
count   317.000000       317.000000        317.000000
mean      4.625994        12.144994          7.519000
std       2.372266         5.975497          3.710278
min       2.050000         4.000000          1.950000
25\%       2.640000         6.030000          3.370000
50\%       4.160000        12.369618          8.066648
75\%       5.010000        14.753932          9.793115
max      10.650000        25.019300         15.199300
\end{Verbatim}
\end{tcolorbox}
        
    \hypertarget{insight-para-consumidores-cuxf3mo-influye-el-precio-del-mauxedz-en-la-tortilla}{%
\subsection{Insight para consumidores: ¿Cómo influye el precio del maíz
en la
tortilla?}\label{insight-para-consumidores-cuxf3mo-influye-el-precio-del-mauxedz-en-la-tortilla}}

Los resultados obtenidos muestran que existe una relación estrecha entre
el precio nacional del maíz blanco y el precio nacional de la tortilla.

La correlación observada entre ambas series en niveles fue:

\[
r = 0.972
\]

lo que indica que, en el largo plazo, ambos precios tienden a
evolucionar en la misma dirección.

Sin embargo, el análisis de participación relativa mostró que el maíz
representa únicamente una parte del precio final pagado por los
consumidores.

Durante el periodo 2000--2026 se observó que:

\begin{itemize}
\tightlist
\item
  El precio promedio del maíz fue de \$4.63 por kilogramo.
\item
  El precio promedio de la tortilla fue de \$12.14 por kilogramo.
\item
  La diferencia promedio fue de \$7.52 por kilogramo.
\end{itemize}

Esto significa que, por cada \$100 pesos gastados en tortillas:

\begin{itemize}
\tightlist
\item
  Aproximadamente \$38 pesos corresponden al maíz.
\item
  Aproximadamente \$62 pesos corresponden a otros costos de la cadena
  productiva y comercial.
\end{itemize}

Por lo tanto, cuando el precio del maíz aumenta, es razonable esperar
presiones al alza sobre el precio de la tortilla. No obstante, la
magnitud final de dicho incremento depende también de otros factores
como:

\begin{itemize}
\tightlist
\item
  Transporte y logística.
\item
  Combustibles.
\item
  Energía eléctrica y gas.
\item
  Mano de obra.
\item
  Distribución y comercialización.
\item
  Inflación general.
\end{itemize}

En términos prácticos, el maíz funciona como uno de los principales
motores del precio de la tortilla, pero no es el único componente que
determina el precio final pagado por los consumidores.

Los resultados sugieren que una reducción en el precio del maíz podría
contribuir a contener el precio de la tortilla, aunque difícilmente
produciría una disminución proporcional, debido a que una parte
importante del costo final proviene de otros elementos de la cadena
agroalimentaria.

En consecuencia, el precio de la tortilla comienza en el campo, pero se
construye a lo largo de toda la cadena de valor que conecta a
productores, transportistas, comercializadores y consumidores.

    \hypertarget{aplicaciuxf3n-de-investigaciuxf3n-de-operaciones}{%
\section{40. Aplicación de Investigación de
Operaciones}\label{aplicaciuxf3n-de-investigaciuxf3n-de-operaciones}}

\hypertarget{objetivo}{%
\subsection{Objetivo}\label{objetivo}}

Utilizar las predicciones generadas por el modelo SARIMAX para
determinar una estrategia de compra de maíz que minimice el costo
esperado de abastecimiento.

La idea consiste en transformar los resultados predictivos obtenidos
previamente en una herramienta de apoyo para la toma de decisiones.

\begin{center}\rule{0.5\linewidth}{0.5pt}\end{center}

\hypertarget{problema}{%
\subsection{Problema}\label{problema}}

Supóngase que una empresa requiere adquirir maíz blanco para satisfacer
una demanda futura.

Se dispone de pronósticos de precio para los próximos meses generados
mediante el modelo SARIMAX.

La decisión consiste en determinar en qué periodos conviene realizar las
compras para minimizar el costo total esperado.

    \begin{tcolorbox}[breakable, size=fbox, boxrule=1pt, pad at break*=1mm,colback=cellbackground, colframe=cellborder]
\prompt{In}{incolor}{152}{\boxspacing}
\begin{Verbatim}[commandchars=\\\{\}]
\PY{n}{serie\PYZus{}final} \PY{o}{=} \PY{p}{(}
    \PY{n}{pdf\PYZus{}completo}
    \PY{o}{.}\PY{n}{set\PYZus{}index}\PY{p}{(}\PY{l+s+s2}{\PYZdq{}}\PY{l+s+s2}{fecha}\PY{l+s+s2}{\PYZdq{}}\PY{p}{)}
    \PY{p}{[}\PY{l+s+s2}{\PYZdq{}}\PY{l+s+s2}{Precio\PYZus{}maizblanco\PYZus{}mxpkilo\PYZus{}prom\PYZus{}nacional}\PY{l+s+s2}{\PYZdq{}}\PY{p}{]}
\PY{p}{)}

\PY{n}{X\PYZus{}final} \PY{o}{=} \PY{p}{(}
    \PY{n}{pdf\PYZus{}completo}
    \PY{o}{.}\PY{n}{set\PYZus{}index}\PY{p}{(}\PY{l+s+s2}{\PYZdq{}}\PY{l+s+s2}{fecha}\PY{l+s+s2}{\PYZdq{}}\PY{p}{)}
    \PY{p}{[}\PY{n}{variables\PYZus{}exogenas\PYZus{}reducidas}\PY{p}{]}
\PY{p}{)}
\end{Verbatim}
\end{tcolorbox}

    \begin{tcolorbox}[breakable, size=fbox, boxrule=1pt, pad at break*=1mm,colback=cellbackground, colframe=cellborder]
\prompt{In}{incolor}{153}{\boxspacing}
\begin{Verbatim}[commandchars=\\\{\}]
\PY{n}{modelo\PYZus{}final} \PY{o}{=} \PY{n}{SARIMAX}\PY{p}{(}
    \PY{n}{serie\PYZus{}final}\PY{p}{,}
    \PY{n}{exog}\PY{o}{=}\PY{n}{X\PYZus{}final}\PY{p}{,}
    \PY{n}{order}\PY{o}{=}\PY{p}{(}\PY{l+m+mi}{1}\PY{p}{,}\PY{l+m+mi}{1}\PY{p}{,}\PY{l+m+mi}{0}\PY{p}{)}\PY{p}{,}
    \PY{n}{seasonal\PYZus{}order}\PY{o}{=}\PY{p}{(}\PY{l+m+mi}{1}\PY{p}{,}\PY{l+m+mi}{0}\PY{p}{,}\PY{l+m+mi}{0}\PY{p}{,}\PY{l+m+mi}{12}\PY{p}{)}\PY{p}{,}
    \PY{n}{enforce\PYZus{}stationarity}\PY{o}{=}\PY{k+kc}{False}\PY{p}{,}
    \PY{n}{enforce\PYZus{}invertibility}\PY{o}{=}\PY{k+kc}{False}
\PY{p}{)}

\PY{n}{resultado\PYZus{}final} \PY{o}{=} \PY{n}{modelo\PYZus{}final}\PY{o}{.}\PY{n}{fit}\PY{p}{(}\PY{p}{)}

\PY{n+nb}{print}\PY{p}{(}\PY{n}{resultado\PYZus{}final}\PY{o}{.}\PY{n}{summary}\PY{p}{(}\PY{p}{)}\PY{p}{)}
\end{Verbatim}
\end{tcolorbox}

    \begin{Verbatim}[commandchars=\\\{\}]
/opt/conda/lib/python3.11/site-packages/statsmodels/tsa/base/tsa\_model.py:473:
ValueWarning: No frequency information was provided, so inferred frequency MS
will be used.
  self.\_init\_dates(dates, freq)
/opt/conda/lib/python3.11/site-packages/statsmodels/tsa/base/tsa\_model.py:473:
ValueWarning: No frequency information was provided, so inferred frequency MS
will be used.
  self.\_init\_dates(dates, freq)
    \end{Verbatim}

    \begin{Verbatim}[commandchars=\\\{\}]
                                          SARIMAX Results
================================================================================
===================
Dep. Variable:     Precio\_maizblanco\_mxpkilo\_prom\_nacional   No. Observations:
317
Model:                      SARIMAX(1, 1, 0)x(1, 0, 0, 12)   Log Likelihood
235.273
Date:                                     Tue, 02 Jun 2026   AIC
-456.545
Time:                                             15:05:05   BIC
-430.549
Sample:                                         01-01-2000   HQIC
-446.145
                                              - 05-01-2026
Covariance Type:                                       opg
================================================================================
===================================
                                                      coef    std err          z
P>|z|      [0.025      0.975]
--------------------------------------------------------------------------------
-----------------------------------
Precio\_prom\_CBOT\_mxpkilo\_maiz\_internacional        -0.0398      0.023     -1.743
0.081      -0.084       0.005
Costo\_transporte\_maiz\_mxpkilo\_prom\_nac              3.8453      0.998      3.855
0.000       1.890       5.800
Temperatura\_celsius\_prom\_nacional                   0.0324      0.018      1.752
0.080      -0.004       0.069
Indice\_severidad\_ponderada\_sequia\_prom\_nacional     0.0004      0.000      3.518
0.000       0.000       0.001
ar.L1                                               0.1055      0.053      1.993
0.046       0.002       0.209
ar.S.L12                                            0.7490      0.018     40.739
0.000       0.713       0.785
sigma2                                              0.0124      0.000     31.487
0.000       0.012       0.013
================================================================================
===
Ljung-Box (L1) (Q):                   0.10   Jarque-Bera (JB):
4420.27
Prob(Q):                              0.76   Prob(JB):
0.00
Heteroskedasticity (H):               5.57   Skew:
1.41
Prob(H) (two-sided):                  0.00   Kurtosis:
21.50
================================================================================
===

Warnings:
[1] Covariance matrix calculated using the outer product of gradients (complex-
step).
    \end{Verbatim}

    \begin{Verbatim}[commandchars=\\\{\}]
/opt/conda/lib/python3.11/site-packages/statsmodels/base/model.py:607:
ConvergenceWarning: Maximum Likelihood optimization failed to converge. Check
mle\_retvals
  warnings.warn("Maximum Likelihood optimization failed to "
    \end{Verbatim}

    \begin{tcolorbox}[breakable, size=fbox, boxrule=1pt, pad at break*=1mm,colback=cellbackground, colframe=cellborder]
\prompt{In}{incolor}{154}{\boxspacing}
\begin{Verbatim}[commandchars=\\\{\}]
\PY{n}{X\PYZus{}future} \PY{o}{=} \PY{n}{pd}\PY{o}{.}\PY{n}{DataFrame}\PY{p}{(}
    \PY{p}{[}\PY{n}{X\PYZus{}final}\PY{o}{.}\PY{n}{iloc}\PY{p}{[}\PY{o}{\PYZhy{}}\PY{l+m+mi}{1}\PY{p}{]}\PY{o}{.}\PY{n}{values}\PY{p}{]}\PY{o}{*}\PY{l+m+mi}{6}\PY{p}{,}
    \PY{n}{columns}\PY{o}{=}\PY{n}{X\PYZus{}final}\PY{o}{.}\PY{n}{columns}
\PY{p}{)}
\end{Verbatim}
\end{tcolorbox}

    \begin{tcolorbox}[breakable, size=fbox, boxrule=1pt, pad at break*=1mm,colback=cellbackground, colframe=cellborder]
\prompt{In}{incolor}{170}{\boxspacing}
\begin{Verbatim}[commandchars=\\\{\}]
\PY{n}{X\PYZus{}future} \PY{o}{=} \PY{n}{pd}\PY{o}{.}\PY{n}{DataFrame}\PY{p}{(}
    \PY{p}{[}\PY{n}{X\PYZus{}final}\PY{o}{.}\PY{n}{iloc}\PY{p}{[}\PY{o}{\PYZhy{}}\PY{l+m+mi}{1}\PY{p}{]}\PY{o}{.}\PY{n}{values}\PY{p}{]}\PY{o}{*}\PY{l+m+mi}{12}\PY{p}{,}
    \PY{n}{columns}\PY{o}{=}\PY{n}{X\PYZus{}final}\PY{o}{.}\PY{n}{columns}
\PY{p}{)}

\PY{n}{forecast\PYZus{}12m} \PY{o}{=} \PY{n}{resultado\PYZus{}final}\PY{o}{.}\PY{n}{get\PYZus{}forecast}\PY{p}{(}
    \PY{n}{steps}\PY{o}{=}\PY{l+m+mi}{12}\PY{p}{,}
    \PY{n}{exog}\PY{o}{=}\PY{n}{X\PYZus{}future}
\PY{p}{)}
\PY{n}{pronostico\PYZus{}io} \PY{o}{=} \PY{n}{pd}\PY{o}{.}\PY{n}{DataFrame}\PY{p}{(}\PY{p}{\PYZob{}}
    \PY{l+s+s2}{\PYZdq{}}\PY{l+s+s2}{Mes}\PY{l+s+s2}{\PYZdq{}}\PY{p}{:} \PY{n+nb}{range}\PY{p}{(}\PY{l+m+mi}{1}\PY{p}{,}\PY{l+m+mi}{7}\PY{p}{)}\PY{p}{,}
    \PY{l+s+s2}{\PYZdq{}}\PY{l+s+s2}{Precio\PYZus{}Pronosticado}\PY{l+s+s2}{\PYZdq{}}\PY{p}{:}
        \PY{n}{forecast\PYZus{}6m}\PY{o}{.}\PY{n}{predicted\PYZus{}mean}\PY{o}{.}\PY{n}{values}
\PY{p}{\PYZcb{}}\PY{p}{)}

\PY{n}{pronostico\PYZus{}io}
\end{Verbatim}
\end{tcolorbox}

            \begin{tcolorbox}[breakable, size=fbox, boxrule=.5pt, pad at break*=1mm, opacityfill=0]
\prompt{Out}{outcolor}{170}{\boxspacing}
\begin{Verbatim}[commandchars=\\\{\}]
   Mes  Precio\_Pronosticado
0    1             9.665025
1    2             9.836343
2    3             9.733932
3    4             9.673063
4    5             9.566773
5    6             9.546923
\end{Verbatim}
\end{tcolorbox}
        
    \begin{tcolorbox}[breakable, size=fbox, boxrule=1pt, pad at break*=1mm,colback=cellbackground, colframe=cellborder]
\prompt{In}{incolor}{165}{\boxspacing}
\begin{Verbatim}[commandchars=\\\{\}]
\PY{n}{real\PYZus{}sarimax\PYZus{}red} \PY{o}{=} \PY{p}{(}
    \PY{n}{pdf\PYZus{}completo}
    \PY{o}{.}\PY{n}{set\PYZus{}index}\PY{p}{(}\PY{l+s+s2}{\PYZdq{}}\PY{l+s+s2}{fecha}\PY{l+s+s2}{\PYZdq{}}\PY{p}{)}
    \PY{o}{.}\PY{n}{loc}\PY{p}{[}
        \PY{n}{pred\PYZus{}sarimax\PYZus{}red}\PY{o}{.}\PY{n}{index}\PY{p}{,}
        \PY{l+s+s2}{\PYZdq{}}\PY{l+s+s2}{Precio\PYZus{}maizblanco\PYZus{}mxpkilo\PYZus{}prom\PYZus{}nacional}\PY{l+s+s2}{\PYZdq{}}
    \PY{p}{]}
\PY{p}{)}
\end{Verbatim}
\end{tcolorbox}

    \begin{tcolorbox}[breakable, size=fbox, boxrule=1pt, pad at break*=1mm,colback=cellbackground, colframe=cellborder]
\prompt{In}{incolor}{166}{\boxspacing}
\begin{Verbatim}[commandchars=\\\{\}]
\PY{n}{plt}\PY{o}{.}\PY{n}{figure}\PY{p}{(}\PY{n}{figsize}\PY{o}{=}\PY{p}{(}\PY{l+m+mi}{14}\PY{p}{,}\PY{l+m+mi}{6}\PY{p}{)}\PY{p}{)}

\PY{n}{plt}\PY{o}{.}\PY{n}{plot}\PY{p}{(}
    \PY{n}{real\PYZus{}sarimax\PYZus{}red}\PY{o}{.}\PY{n}{index}\PY{p}{,}
    \PY{n}{real\PYZus{}sarimax\PYZus{}red}\PY{p}{,}
    \PY{n}{label}\PY{o}{=}\PY{l+s+s2}{\PYZdq{}}\PY{l+s+s2}{Precio Real}\PY{l+s+s2}{\PYZdq{}}\PY{p}{,}
    \PY{n}{linewidth}\PY{o}{=}\PY{l+m+mi}{2}
\PY{p}{)}

\PY{n}{plt}\PY{o}{.}\PY{n}{plot}\PY{p}{(}
    \PY{n}{pred\PYZus{}sarimax\PYZus{}red}\PY{o}{.}\PY{n}{index}\PY{p}{,}
    \PY{n}{pred\PYZus{}sarimax\PYZus{}red}\PY{p}{,}
    \PY{n}{label}\PY{o}{=}\PY{l+s+s2}{\PYZdq{}}\PY{l+s+s2}{Predicción SARIMAX Reducido}\PY{l+s+s2}{\PYZdq{}}\PY{p}{,}
    \PY{n}{linestyle}\PY{o}{=}\PY{l+s+s2}{\PYZdq{}}\PY{l+s+s2}{\PYZhy{}\PYZhy{}}\PY{l+s+s2}{\PYZdq{}}\PY{p}{,}
    \PY{n}{linewidth}\PY{o}{=}\PY{l+m+mi}{2}
\PY{p}{)}

\PY{n}{plt}\PY{o}{.}\PY{n}{title}\PY{p}{(}
    \PY{l+s+s2}{\PYZdq{}}\PY{l+s+s2}{Validación del Modelo SARIMAX Reducido}\PY{l+s+s2}{\PYZdq{}}
\PY{p}{)}

\PY{n}{plt}\PY{o}{.}\PY{n}{xlabel}\PY{p}{(}\PY{l+s+s2}{\PYZdq{}}\PY{l+s+s2}{Fecha}\PY{l+s+s2}{\PYZdq{}}\PY{p}{)}
\PY{n}{plt}\PY{o}{.}\PY{n}{ylabel}\PY{p}{(}\PY{l+s+s2}{\PYZdq{}}\PY{l+s+s2}{Precio del Maíz (MXN/Kg)}\PY{l+s+s2}{\PYZdq{}}\PY{p}{)}

\PY{n}{plt}\PY{o}{.}\PY{n}{legend}\PY{p}{(}\PY{p}{)}
\PY{n}{plt}\PY{o}{.}\PY{n}{grid}\PY{p}{(}\PY{k+kc}{True}\PY{p}{)}

\PY{n}{plt}\PY{o}{.}\PY{n}{show}\PY{p}{(}\PY{p}{)}
\end{Verbatim}
\end{tcolorbox}

    \begin{center}
    \adjustimage{max size={0.9\linewidth}{0.9\paperheight}}{output_206_0.png}
    \end{center}
    { \hspace*{\fill} \\}
    
    \begin{tcolorbox}[breakable, size=fbox, boxrule=1pt, pad at break*=1mm,colback=cellbackground, colframe=cellborder]
\prompt{In}{incolor}{171}{\boxspacing}
\begin{Verbatim}[commandchars=\\\{\}]
\PY{n}{fechas\PYZus{}futuras} \PY{o}{=} \PY{n}{pd}\PY{o}{.}\PY{n}{date\PYZus{}range}\PY{p}{(}
    \PY{n}{start}\PY{o}{=}\PY{l+s+s2}{\PYZdq{}}\PY{l+s+s2}{2026\PYZhy{}06\PYZhy{}01}\PY{l+s+s2}{\PYZdq{}}\PY{p}{,}
    \PY{n}{periods}\PY{o}{=}\PY{l+m+mi}{12}\PY{p}{,}
    \PY{n}{freq}\PY{o}{=}\PY{l+s+s2}{\PYZdq{}}\PY{l+s+s2}{MS}\PY{l+s+s2}{\PYZdq{}}
\PY{p}{)}

\PY{n}{forecast\PYZus{}futuro} \PY{o}{=} \PY{n}{pd}\PY{o}{.}\PY{n}{Series}\PY{p}{(}
    \PY{n}{forecast\PYZus{}12m}\PY{o}{.}\PY{n}{predicted\PYZus{}mean}\PY{o}{.}\PY{n}{values}\PY{p}{,}
    \PY{n}{index}\PY{o}{=}\PY{n}{fechas\PYZus{}futuras}
\PY{p}{)}
\end{Verbatim}
\end{tcolorbox}

    \begin{tcolorbox}[breakable, size=fbox, boxrule=1pt, pad at break*=1mm,colback=cellbackground, colframe=cellborder]
\prompt{In}{incolor}{172}{\boxspacing}
\begin{Verbatim}[commandchars=\\\{\}]
\PY{n}{plt}\PY{o}{.}\PY{n}{figure}\PY{p}{(}\PY{n}{figsize}\PY{o}{=}\PY{p}{(}\PY{l+m+mi}{14}\PY{p}{,}\PY{l+m+mi}{6}\PY{p}{)}\PY{p}{)}

\PY{n}{plt}\PY{o}{.}\PY{n}{plot}\PY{p}{(}
    \PY{n}{serie\PYZus{}final}\PY{o}{.}\PY{n}{index}\PY{p}{,}
    \PY{n}{serie\PYZus{}final}\PY{p}{,}
    \PY{n}{label}\PY{o}{=}\PY{l+s+s2}{\PYZdq{}}\PY{l+s+s2}{Histórico}\PY{l+s+s2}{\PYZdq{}}
\PY{p}{)}

\PY{n}{plt}\PY{o}{.}\PY{n}{plot}\PY{p}{(}
    \PY{n}{forecast\PYZus{}futuro}\PY{o}{.}\PY{n}{index}\PY{p}{,}
    \PY{n}{forecast\PYZus{}futuro}\PY{p}{,}
    \PY{n}{label}\PY{o}{=}\PY{l+s+s2}{\PYZdq{}}\PY{l+s+s2}{Pronóstico SARIMAX}\PY{l+s+s2}{\PYZdq{}}\PY{p}{,}
    \PY{n}{linestyle}\PY{o}{=}\PY{l+s+s2}{\PYZdq{}}\PY{l+s+s2}{\PYZhy{}\PYZhy{}}\PY{l+s+s2}{\PYZdq{}}\PY{p}{,}
    \PY{n}{linewidth}\PY{o}{=}\PY{l+m+mi}{3}
\PY{p}{)}

\PY{n}{plt}\PY{o}{.}\PY{n}{axvline}\PY{p}{(}
    \PY{n}{serie\PYZus{}final}\PY{o}{.}\PY{n}{index}\PY{o}{.}\PY{n}{max}\PY{p}{(}\PY{p}{)}\PY{p}{,}
    \PY{n}{color}\PY{o}{=}\PY{l+s+s2}{\PYZdq{}}\PY{l+s+s2}{red}\PY{l+s+s2}{\PYZdq{}}\PY{p}{,}
    \PY{n}{linestyle}\PY{o}{=}\PY{l+s+s2}{\PYZdq{}}\PY{l+s+s2}{:}\PY{l+s+s2}{\PYZdq{}}
\PY{p}{)}

\PY{n}{plt}\PY{o}{.}\PY{n}{title}\PY{p}{(}
    \PY{l+s+s2}{\PYZdq{}}\PY{l+s+s2}{Pronóstico del Precio Nacional del Maíz Blanco}\PY{l+s+s2}{\PYZdq{}}
\PY{p}{)}

\PY{n}{plt}\PY{o}{.}\PY{n}{xlabel}\PY{p}{(}\PY{l+s+s2}{\PYZdq{}}\PY{l+s+s2}{Fecha}\PY{l+s+s2}{\PYZdq{}}\PY{p}{)}
\PY{n}{plt}\PY{o}{.}\PY{n}{ylabel}\PY{p}{(}\PY{l+s+s2}{\PYZdq{}}\PY{l+s+s2}{MXN/Kg}\PY{l+s+s2}{\PYZdq{}}\PY{p}{)}

\PY{n}{plt}\PY{o}{.}\PY{n}{legend}\PY{p}{(}\PY{p}{)}
\PY{n}{plt}\PY{o}{.}\PY{n}{grid}\PY{p}{(}\PY{k+kc}{True}\PY{p}{)}

\PY{n}{plt}\PY{o}{.}\PY{n}{show}\PY{p}{(}\PY{p}{)}
\end{Verbatim}
\end{tcolorbox}

    \begin{center}
    \adjustimage{max size={0.9\linewidth}{0.9\paperheight}}{output_208_0.png}
    \end{center}
    { \hspace*{\fill} \\}
    
    \begin{tcolorbox}[breakable, size=fbox, boxrule=1pt, pad at break*=1mm,colback=cellbackground, colframe=cellborder]
\prompt{In}{incolor}{173}{\boxspacing}
\begin{Verbatim}[commandchars=\\\{\}]
\PY{n}{forecast\PYZus{}12m} \PY{o}{=} \PY{n}{resultado\PYZus{}final}\PY{o}{.}\PY{n}{get\PYZus{}forecast}\PY{p}{(}
    \PY{n}{steps}\PY{o}{=}\PY{l+m+mi}{12}\PY{p}{,}
    \PY{n}{exog}\PY{o}{=}\PY{n}{X\PYZus{}future}
\PY{p}{)}

\PY{n}{conf\PYZus{}int} \PY{o}{=} \PY{n}{forecast\PYZus{}12m}\PY{o}{.}\PY{n}{conf\PYZus{}int}\PY{p}{(}\PY{p}{)}
\end{Verbatim}
\end{tcolorbox}

    \begin{tcolorbox}[breakable, size=fbox, boxrule=1pt, pad at break*=1mm,colback=cellbackground, colframe=cellborder]
\prompt{In}{incolor}{174}{\boxspacing}
\begin{Verbatim}[commandchars=\\\{\}]
\PY{n}{plt}\PY{o}{.}\PY{n}{figure}\PY{p}{(}\PY{n}{figsize}\PY{o}{=}\PY{p}{(}\PY{l+m+mi}{14}\PY{p}{,}\PY{l+m+mi}{6}\PY{p}{)}\PY{p}{)}

\PY{n}{plt}\PY{o}{.}\PY{n}{plot}\PY{p}{(}
    \PY{n}{serie\PYZus{}final}\PY{o}{.}\PY{n}{index}\PY{p}{,}
    \PY{n}{serie\PYZus{}final}\PY{p}{,}
    \PY{n}{label}\PY{o}{=}\PY{l+s+s2}{\PYZdq{}}\PY{l+s+s2}{Histórico}\PY{l+s+s2}{\PYZdq{}}
\PY{p}{)}

\PY{n}{plt}\PY{o}{.}\PY{n}{plot}\PY{p}{(}
    \PY{n}{forecast\PYZus{}futuro}\PY{o}{.}\PY{n}{index}\PY{p}{,}
    \PY{n}{forecast\PYZus{}futuro}\PY{p}{,}
    \PY{n}{linestyle}\PY{o}{=}\PY{l+s+s2}{\PYZdq{}}\PY{l+s+s2}{\PYZhy{}\PYZhy{}}\PY{l+s+s2}{\PYZdq{}}\PY{p}{,}
    \PY{n}{linewidth}\PY{o}{=}\PY{l+m+mi}{3}\PY{p}{,}
    \PY{n}{label}\PY{o}{=}\PY{l+s+s2}{\PYZdq{}}\PY{l+s+s2}{Pronóstico}\PY{l+s+s2}{\PYZdq{}}
\PY{p}{)}

\PY{n}{plt}\PY{o}{.}\PY{n}{fill\PYZus{}between}\PY{p}{(}
    \PY{n}{forecast\PYZus{}futuro}\PY{o}{.}\PY{n}{index}\PY{p}{,}
    \PY{n}{conf\PYZus{}int}\PY{o}{.}\PY{n}{iloc}\PY{p}{[}\PY{p}{:}\PY{p}{,}\PY{l+m+mi}{0}\PY{p}{]}\PY{p}{,}
    \PY{n}{conf\PYZus{}int}\PY{o}{.}\PY{n}{iloc}\PY{p}{[}\PY{p}{:}\PY{p}{,}\PY{l+m+mi}{1}\PY{p}{]}\PY{p}{,}
    \PY{n}{alpha}\PY{o}{=}\PY{l+m+mf}{0.25}\PY{p}{,}
    \PY{n}{label}\PY{o}{=}\PY{l+s+s2}{\PYZdq{}}\PY{l+s+s2}{IC 95}\PY{l+s+s2}{\PYZpc{}}\PY{l+s+s2}{\PYZdq{}}
\PY{p}{)}

\PY{n}{plt}\PY{o}{.}\PY{n}{axvline}\PY{p}{(}
    \PY{n}{serie\PYZus{}final}\PY{o}{.}\PY{n}{index}\PY{o}{.}\PY{n}{max}\PY{p}{(}\PY{p}{)}\PY{p}{,}
    \PY{n}{color}\PY{o}{=}\PY{l+s+s2}{\PYZdq{}}\PY{l+s+s2}{red}\PY{l+s+s2}{\PYZdq{}}\PY{p}{,}
    \PY{n}{linestyle}\PY{o}{=}\PY{l+s+s2}{\PYZdq{}}\PY{l+s+s2}{:}\PY{l+s+s2}{\PYZdq{}}
\PY{p}{)}

\PY{n}{plt}\PY{o}{.}\PY{n}{title}\PY{p}{(}
    \PY{l+s+s2}{\PYZdq{}}\PY{l+s+s2}{Pronóstico del Precio Nacional del Maíz Blanco (12 meses)}\PY{l+s+s2}{\PYZdq{}}
\PY{p}{)}

\PY{n}{plt}\PY{o}{.}\PY{n}{ylabel}\PY{p}{(}\PY{l+s+s2}{\PYZdq{}}\PY{l+s+s2}{MXN/Kg}\PY{l+s+s2}{\PYZdq{}}\PY{p}{)}
\PY{n}{plt}\PY{o}{.}\PY{n}{xlabel}\PY{p}{(}\PY{l+s+s2}{\PYZdq{}}\PY{l+s+s2}{Fecha}\PY{l+s+s2}{\PYZdq{}}\PY{p}{)}

\PY{n}{plt}\PY{o}{.}\PY{n}{legend}\PY{p}{(}\PY{p}{)}
\PY{n}{plt}\PY{o}{.}\PY{n}{grid}\PY{p}{(}\PY{k+kc}{True}\PY{p}{)}

\PY{n}{plt}\PY{o}{.}\PY{n}{show}\PY{p}{(}\PY{p}{)}
\end{Verbatim}
\end{tcolorbox}

    \begin{center}
    \adjustimage{max size={0.9\linewidth}{0.9\paperheight}}{output_210_0.png}
    \end{center}
    { \hspace*{\fill} \\}
    
    \begin{tcolorbox}[breakable, size=fbox, boxrule=1pt, pad at break*=1mm,colback=cellbackground, colframe=cellborder]
\prompt{In}{incolor}{175}{\boxspacing}
\begin{Verbatim}[commandchars=\\\{\}]
\PY{k+kn}{from} \PY{n+nn}{scipy}\PY{n+nn}{.}\PY{n+nn}{optimize} \PY{k+kn}{import} \PY{n}{linprog}
\end{Verbatim}
\end{tcolorbox}

    \begin{tcolorbox}[breakable, size=fbox, boxrule=1pt, pad at break*=1mm,colback=cellbackground, colframe=cellborder]
\prompt{In}{incolor}{176}{\boxspacing}
\begin{Verbatim}[commandchars=\\\{\}]
\PY{n}{demanda} \PY{o}{=} \PY{l+m+mi}{10000}
\end{Verbatim}
\end{tcolorbox}

    \begin{tcolorbox}[breakable, size=fbox, boxrule=1pt, pad at break*=1mm,colback=cellbackground, colframe=cellborder]
\prompt{In}{incolor}{177}{\boxspacing}
\begin{Verbatim}[commandchars=\\\{\}]
\PY{n}{costos} \PY{o}{=} \PY{n}{pronostico\PYZus{}io}\PY{p}{[}
    \PY{l+s+s2}{\PYZdq{}}\PY{l+s+s2}{Precio\PYZus{}Pronosticado}\PY{l+s+s2}{\PYZdq{}}
\PY{p}{]}\PY{o}{.}\PY{n}{values}
\end{Verbatim}
\end{tcolorbox}

    \begin{tcolorbox}[breakable, size=fbox, boxrule=1pt, pad at break*=1mm,colback=cellbackground, colframe=cellborder]
\prompt{In}{incolor}{178}{\boxspacing}
\begin{Verbatim}[commandchars=\\\{\}]
\PY{n}{A\PYZus{}eq} \PY{o}{=} \PY{p}{[}\PY{p}{[}\PY{l+m+mi}{1}\PY{p}{]}\PY{o}{*}\PY{n+nb}{len}\PY{p}{(}\PY{n}{costos}\PY{p}{)}\PY{p}{]}

\PY{n}{b\PYZus{}eq} \PY{o}{=} \PY{p}{[}\PY{n}{demanda}\PY{p}{]}

\PY{n}{bounds} \PY{o}{=} \PY{p}{[}\PY{p}{(}\PY{l+m+mi}{0}\PY{p}{,} \PY{n}{demanda}\PY{p}{)}\PY{p}{]}\PY{o}{*}\PY{n+nb}{len}\PY{p}{(}\PY{n}{costos}\PY{p}{)}

\PY{n}{resultado\PYZus{}io} \PY{o}{=} \PY{n}{linprog}\PY{p}{(}
    \PY{n}{c}\PY{o}{=}\PY{n}{costos}\PY{p}{,}
    \PY{n}{A\PYZus{}eq}\PY{o}{=}\PY{n}{A\PYZus{}eq}\PY{p}{,}
    \PY{n}{b\PYZus{}eq}\PY{o}{=}\PY{n}{b\PYZus{}eq}\PY{p}{,}
    \PY{n}{bounds}\PY{o}{=}\PY{n}{bounds}\PY{p}{,}
    \PY{n}{method}\PY{o}{=}\PY{l+s+s2}{\PYZdq{}}\PY{l+s+s2}{highs}\PY{l+s+s2}{\PYZdq{}}
\PY{p}{)}
\end{Verbatim}
\end{tcolorbox}

    \begin{tcolorbox}[breakable, size=fbox, boxrule=1pt, pad at break*=1mm,colback=cellbackground, colframe=cellborder]
\prompt{In}{incolor}{179}{\boxspacing}
\begin{Verbatim}[commandchars=\\\{\}]
\PY{n}{solucion\PYZus{}io} \PY{o}{=} \PY{n}{pd}\PY{o}{.}\PY{n}{DataFrame}\PY{p}{(}\PY{p}{\PYZob{}}
    \PY{l+s+s2}{\PYZdq{}}\PY{l+s+s2}{Mes}\PY{l+s+s2}{\PYZdq{}}\PY{p}{:} \PY{n}{pronostico\PYZus{}io}\PY{p}{[}\PY{l+s+s2}{\PYZdq{}}\PY{l+s+s2}{Mes}\PY{l+s+s2}{\PYZdq{}}\PY{p}{]}\PY{p}{,}
    \PY{l+s+s2}{\PYZdq{}}\PY{l+s+s2}{Precio}\PY{l+s+s2}{\PYZdq{}}\PY{p}{:} \PY{n}{costos}\PY{p}{,}
    \PY{l+s+s2}{\PYZdq{}}\PY{l+s+s2}{Compra\PYZus{}Optima\PYZus{}kg}\PY{l+s+s2}{\PYZdq{}}\PY{p}{:} \PY{n}{resultado\PYZus{}io}\PY{o}{.}\PY{n}{x}
\PY{p}{\PYZcb{}}\PY{p}{)}

\PY{n}{solucion\PYZus{}io}
\end{Verbatim}
\end{tcolorbox}

            \begin{tcolorbox}[breakable, size=fbox, boxrule=.5pt, pad at break*=1mm, opacityfill=0]
\prompt{Out}{outcolor}{179}{\boxspacing}
\begin{Verbatim}[commandchars=\\\{\}]
   Mes    Precio  Compra\_Optima\_kg
0    1  9.665025               0.0
1    2  9.836343               0.0
2    3  9.733932               0.0
3    4  9.673063               0.0
4    5  9.566773               0.0
5    6  9.546923           10000.0
\end{Verbatim}
\end{tcolorbox}
        
    \begin{tcolorbox}[breakable, size=fbox, boxrule=1pt, pad at break*=1mm,colback=cellbackground, colframe=cellborder]
\prompt{In}{incolor}{180}{\boxspacing}
\begin{Verbatim}[commandchars=\\\{\}]
\PY{n}{costo\PYZus{}total} \PY{o}{=} \PY{p}{(}
    \PY{n}{solucion\PYZus{}io}\PY{p}{[}\PY{l+s+s2}{\PYZdq{}}\PY{l+s+s2}{Precio}\PY{l+s+s2}{\PYZdq{}}\PY{p}{]}
    \PY{o}{*} \PY{n}{solucion\PYZus{}io}\PY{p}{[}\PY{l+s+s2}{\PYZdq{}}\PY{l+s+s2}{Compra\PYZus{}Optima\PYZus{}kg}\PY{l+s+s2}{\PYZdq{}}\PY{p}{]}
\PY{p}{)}\PY{o}{.}\PY{n}{sum}\PY{p}{(}\PY{p}{)}

\PY{n+nb}{print}\PY{p}{(}
    \PY{l+s+sa}{f}\PY{l+s+s2}{\PYZdq{}}\PY{l+s+s2}{Costo mínimo esperado: \PYZdl{}}\PY{l+s+si}{\PYZob{}}\PY{n}{costo\PYZus{}total}\PY{l+s+si}{:}\PY{l+s+s2}{,.2f}\PY{l+s+si}{\PYZcb{}}\PY{l+s+s2}{\PYZdq{}}
\PY{p}{)}
\end{Verbatim}
\end{tcolorbox}

    \begin{Verbatim}[commandchars=\\\{\}]
Costo mínimo esperado: \$95,469.23
    \end{Verbatim}

    \hypertarget{resultados-de-la-optimizaciuxf3n}{%
\subsection{Resultados de la
optimización}\label{resultados-de-la-optimizaciuxf3n}}

Utilizando los precios proyectados por el modelo SARIMAX para los
próximos seis meses, se formuló un problema de programación lineal cuyo
objetivo consistió en minimizar el costo total de abastecimiento para
una demanda fija de 10,000 kilogramos de maíz.

Los precios pronosticados fueron:

\begin{longtable}[]{@{}lr@{}}
\toprule\noalign{}
Mes & Precio esperado (MXN/kg) \\
\midrule\noalign{}
\endhead
\bottomrule\noalign{}
\endlastfoot
1 & 9.665 \\
2 & 9.836 \\
3 & 9.734 \\
4 & 9.673 \\
5 & 9.567 \\
6 & 9.547 \\
\end{longtable}

La solución óptima obtenida indicó que la totalidad de la demanda
debería adquirirse durante el sexto mes del horizonte de planeación.

\begin{longtable}[]{@{}lr@{}}
\toprule\noalign{}
Mes & Compra óptima (kg) \\
\midrule\noalign{}
\endhead
\bottomrule\noalign{}
\endlastfoot
1 & 0 \\
2 & 0 \\
3 & 0 \\
4 & 0 \\
5 & 0 \\
6 & 10,000 \\
\end{longtable}

El costo mínimo esperado asociado a esta estrategia fue:

\[
Costo_{mínimo}
=
\$95,469.23
\]

\begin{center}\rule{0.5\linewidth}{0.5pt}\end{center}

El modelo de optimización identificó que el menor precio esperado dentro
del horizonte analizado corresponde al sexto mes.

Por esta razón, la estrategia que minimiza el costo total consiste en
posponer la compra hasta dicho periodo.

Desde la perspectiva de Investigación de Operaciones, esta decisión
surge de manera automática al combinar:

\begin{itemize}
\tightlist
\item
  El modelo predictivo SARIMAX desarrollado previamente.
\item
  Un modelo de programación lineal orientado a la minimización de
  costos.
\end{itemize}

En consecuencia, el proyecto no sólo permite anticipar el comportamiento
futuro de los precios, sino también generar recomendaciones concretas
para la toma de decisiones de compra e inventario.

    \hypertarget{conclusiuxf3n-de-la-aplicaciuxf3n-de-investigaciuxf3n-de-operaciones}{%
\subsection{Conclusión de la aplicación de Investigación de
Operaciones}\label{conclusiuxf3n-de-la-aplicaciuxf3n-de-investigaciuxf3n-de-operaciones}}

La incorporación de técnicas de Investigación de Operaciones permitió
transformar los resultados predictivos obtenidos mediante SARIMAX en una
herramienta práctica para la toma de decisiones.

Mientras que el modelo de series de tiempo responde a la pregunta:

\begin{quote}
¿Cuál será el precio esperado del maíz?
\end{quote}

el modelo de optimización responde:

\begin{quote}
¿Cuál es la estrategia de compra que minimiza el costo esperado?
\end{quote}

Los resultados mostraron que, bajo el escenario base considerado, la
alternativa óptima consiste en diferir la adquisición del maíz hasta el
periodo con el menor precio proyectado.

De esta manera, el proyecto integra tres componentes complementarios:

\begin{enumerate}
\def\labelenumi{\arabic{enumi}.}
\tightlist
\item
  Minería de Datos para identificar patrones y relaciones relevantes.
\item
  Modelos Estocásticos para pronosticar el comportamiento futuro del
  mercado.
\item
  Investigación de Operaciones para optimizar decisiones de
  abastecimiento.
\end{enumerate}

Esta integración permite convertir información histórica en conocimiento
accionable que puede apoyar procesos de planeación dentro de la cadena
agroalimentaria mexicana.

    \hypertarget{formulaciuxf3n-del-problema-de-optimizaciuxf3n}{%
\subsection{40.3 Formulación del Problema de
Optimización}\label{formulaciuxf3n-del-problema-de-optimizaciuxf3n}}

Una vez obtenidos los pronósticos de precio mediante el modelo SARIMAX,
se formuló un problema de programación lineal cuyo propósito consistió
en determinar la estrategia de compra que minimizara el costo esperado
de abastecimiento.

La situación planteada corresponde a una organización que requiere
adquirir una cantidad fija de maíz blanco para satisfacer su demanda
futura.

\begin{center}\rule{0.5\linewidth}{0.5pt}\end{center}

\hypertarget{variables-de-decisiuxf3n}{%
\subsubsection{Variables de decisión}\label{variables-de-decisiuxf3n}}

Se definieron las siguientes variables:

\begin{itemize}
\tightlist
\item
  (x\_1): cantidad de maíz a comprar en el mes 1.
\item
  (x\_2): cantidad de maíz a comprar en el mes 2.
\item
  (x\_3): cantidad de maíz a comprar en el mes 3.
\item
  (x\_4): cantidad de maíz a comprar en el mes 4.
\item
  (x\_5): cantidad de maíz a comprar en el mes 5.
\item
  (x\_6): cantidad de maíz a comprar en el mes 6.
\end{itemize}

Cada variable representa la cantidad de kilogramos adquiridos en el
periodo correspondiente.

\begin{center}\rule{0.5\linewidth}{0.5pt}\end{center}

\hypertarget{funciuxf3n-objetivo}{%
\subsubsection{Función objetivo}\label{funciuxf3n-objetivo}}

El objetivo consiste en minimizar el costo total esperado de
adquisición.

A partir de los precios pronosticados por el modelo SARIMAX, la función
objetivo quedó definida como:

\[
\min Z =
9.665x_1 +
9.836x_2 +
9.734x_3 +
9.673x_4 +
9.567x_5 +
9.547x_6
\]

donde:

\begin{itemize}
\tightlist
\item
  (Z) representa el costo total esperado de compra.
\item
  Los coeficientes corresponden a los precios pronosticados para cada
  mes del horizonte de planeación.
\end{itemize}

\begin{center}\rule{0.5\linewidth}{0.5pt}\end{center}

\hypertarget{restricciuxf3n-de-demanda}{%
\subsubsection{Restricción de demanda}\label{restricciuxf3n-de-demanda}}

La organización requiere adquirir exactamente 10,000 kilogramos de maíz.

Por lo tanto, la suma de las compras realizadas durante el horizonte de
planeación debe satisfacer la demanda total:

\[
x_1+x_2+x_3+x_4+x_5+x_6=10,000
\]

\begin{center}\rule{0.5\linewidth}{0.5pt}\end{center}

\hypertarget{restricciuxf3n-de-no-negatividad}{%
\subsubsection{Restricción de no
negatividad}\label{restricciuxf3n-de-no-negatividad}}

Las cantidades adquiridas no pueden tomar valores negativos:

\[
x_i \geq 0
\qquad
\forall i
\]

\begin{center}\rule{0.5\linewidth}{0.5pt}\end{center}

\hypertarget{modelo-matemuxe1tico-completo}{%
\subsubsection{Modelo matemático
completo}\label{modelo-matemuxe1tico-completo}}

La formulación final del problema puede expresarse como:

\[
\min Z =
9.665x_1 +
9.836x_2 +
9.734x_3 +
9.673x_4 +
9.567x_5 +
9.547x_6
\]

sujeto a:

\[
x_1+x_2+x_3+x_4+x_5+x_6=10,000
\]

\[
x_i \geq 0
\qquad
\forall i
\]

\begin{center}\rule{0.5\linewidth}{0.5pt}\end{center}

Desde una perspectiva operativa, el modelo evalúa todas las
combinaciones posibles de compra que permitan satisfacer la demanda
requerida y selecciona aquella que produce el menor costo total
esperado.

De esta manera, los pronósticos generados mediante Ciencia de Datos y
Modelos Estocásticos se convierten en información accionable para apoyar
decisiones de abastecimiento dentro de la cadena agroalimentaria.

En términos prácticos, el problema responde a la siguiente pregunta:

\begin{quote}
\textbf{¿En qué momento conviene realizar las compras de maíz para
minimizar el gasto esperado de adquisición?}
\end{quote}

    \hypertarget{vinculaciuxf3n-con-la-uca-de-mercadotecnia-digital}{%
\section{41. Vinculación con la UCA de Mercadotecnia
Digital}\label{vinculaciuxf3n-con-la-uca-de-mercadotecnia-digital}}

\hypertarget{de-la-predicciuxf3n-al-posicionamiento-comercial}{%
\subsection{De la predicción al posicionamiento
comercial}\label{de-la-predicciuxf3n-al-posicionamiento-comercial}}

La UCA de Mercadotecnia Digital se integra al proyecto al transformar
los resultados técnicos del modelo predictivo en información accionable
para productores, cooperativas, comercializadores e industria
alimentaria.

El valor del modelo no se limita a estimar el precio futuro del maíz
blanco. Su utilidad aparece cuando el pronóstico se convierte en una
herramienta de inteligencia comercial capaz de responder preguntas como:

\begin{itemize}
\tightlist
\item
  ¿Cuándo conviene vender?
\item
  ¿Cuándo conviene comprar?
\item
  ¿Qué tan fuerte es la presión del maíz sobre la tortilla?
\item
  ¿Qué mensaje comercial puede construir una cooperativa frente a la
  volatilidad?
\item
  ¿Cómo puede una organización reducir su dependencia de intermediarios
  usando información predictiva?
\end{itemize}

Desde esta perspectiva, la mercadotecnia digital no se entiende
únicamente como publicidad o presencia en redes sociales, sino como un
sistema de toma de decisiones basado en datos, analítica predictiva,
segmentación de usuarios y comunicación estratégica.

    \begin{tcolorbox}[breakable, size=fbox, boxrule=1pt, pad at break*=1mm,colback=cellbackground, colframe=cellborder]
\prompt{In}{incolor}{299}{\boxspacing}
\begin{Verbatim}[commandchars=\\\{\}]
\PY{c+c1}{\PYZsh{}Definiendo público objetivo}

\PY{n}{segmentos\PYZus{}mercado} \PY{o}{=} \PY{n}{pd}\PY{o}{.}\PY{n}{DataFrame}\PY{p}{(}\PY{p}{\PYZob{}}
    \PY{l+s+s2}{\PYZdq{}}\PY{l+s+s2}{Actor}\PY{l+s+s2}{\PYZdq{}}\PY{p}{:} \PY{p}{[}
        \PY{l+s+s2}{\PYZdq{}}\PY{l+s+s2}{Productores}\PY{l+s+s2}{\PYZdq{}}\PY{p}{,}
        \PY{l+s+s2}{\PYZdq{}}\PY{l+s+s2}{Cooperativas agrícolas}\PY{l+s+s2}{\PYZdq{}}\PY{p}{,}
        \PY{l+s+s2}{\PYZdq{}}\PY{l+s+s2}{Comercializadores}\PY{l+s+s2}{\PYZdq{}}\PY{p}{,}
        \PY{l+s+s2}{\PYZdq{}}\PY{l+s+s2}{Industria tortillera}\PY{l+s+s2}{\PYZdq{}}\PY{p}{,}
        \PY{l+s+s2}{\PYZdq{}}\PY{l+s+s2}{Consumidores}\PY{l+s+s2}{\PYZdq{}}\PY{p}{,}
        \PY{l+s+s2}{\PYZdq{}}\PY{l+s+s2}{Sector público}\PY{l+s+s2}{\PYZdq{}}
    \PY{p}{]}\PY{p}{,}
    \PY{l+s+s2}{\PYZdq{}}\PY{l+s+s2}{Necesidad principal}\PY{l+s+s2}{\PYZdq{}}\PY{p}{:} \PY{p}{[}
        \PY{l+s+s2}{\PYZdq{}}\PY{l+s+s2}{Saber cuándo vender para evitar precios bajos}\PY{l+s+s2}{\PYZdq{}}\PY{p}{,}
        \PY{l+s+s2}{\PYZdq{}}\PY{l+s+s2}{Planear almacenamiento, compra colectiva y negociación}\PY{l+s+s2}{\PYZdq{}}\PY{p}{,}
        \PY{l+s+s2}{\PYZdq{}}\PY{l+s+s2}{Optimizar compras e inventarios}\PY{l+s+s2}{\PYZdq{}}\PY{p}{,}
        \PY{l+s+s2}{\PYZdq{}}\PY{l+s+s2}{Anticipar costos del insumo principal}\PY{l+s+s2}{\PYZdq{}}\PY{p}{,}
        \PY{l+s+s2}{\PYZdq{}}\PY{l+s+s2}{Comprender por qué cambia el precio de la tortilla}\PY{l+s+s2}{\PYZdq{}}\PY{p}{,}
        \PY{l+s+s2}{\PYZdq{}}\PY{l+s+s2}{Monitorear riesgos de seguridad alimentaria}\PY{l+s+s2}{\PYZdq{}}
    \PY{p}{]}\PY{p}{,}
    \PY{l+s+s2}{\PYZdq{}}\PY{l+s+s2}{Uso del modelo}\PY{l+s+s2}{\PYZdq{}}\PY{p}{:} \PY{p}{[}
        \PY{l+s+s2}{\PYZdq{}}\PY{l+s+s2}{Pronóstico del precio del maíz}\PY{l+s+s2}{\PYZdq{}}\PY{p}{,}
        \PY{l+s+s2}{\PYZdq{}}\PY{l+s+s2}{Escenarios de comercialización justa}\PY{l+s+s2}{\PYZdq{}}\PY{p}{,}
        \PY{l+s+s2}{\PYZdq{}}\PY{l+s+s2}{Optimización de abastecimiento}\PY{l+s+s2}{\PYZdq{}}\PY{p}{,}
        \PY{l+s+s2}{\PYZdq{}}\PY{l+s+s2}{Planeación de costos}\PY{l+s+s2}{\PYZdq{}}\PY{p}{,}
        \PY{l+s+s2}{\PYZdq{}}\PY{l+s+s2}{Comunicación clara del impacto del maíz}\PY{l+s+s2}{\PYZdq{}}\PY{p}{,}
        \PY{l+s+s2}{\PYZdq{}}\PY{l+s+s2}{Alertas tempranas de presión alimentaria}\PY{l+s+s2}{\PYZdq{}}
    \PY{p}{]}\PY{p}{,}
    \PY{l+s+s2}{\PYZdq{}}\PY{l+s+s2}{Mensaje digital sugerido}\PY{l+s+s2}{\PYZdq{}}\PY{p}{:} \PY{p}{[}
        \PY{l+s+s2}{\PYZdq{}}\PY{l+s+s2}{Vende con información, no con incertidumbre}\PY{l+s+s2}{\PYZdq{}}\PY{p}{,}
        \PY{l+s+s2}{\PYZdq{}}\PY{l+s+s2}{Negocia mejor usando datos del mercado}\PY{l+s+s2}{\PYZdq{}}\PY{p}{,}
        \PY{l+s+s2}{\PYZdq{}}\PY{l+s+s2}{Compra en el momento de menor costo esperado}\PY{l+s+s2}{\PYZdq{}}\PY{p}{,}
        \PY{l+s+s2}{\PYZdq{}}\PY{l+s+s2}{Anticipa tus costos antes de que suban}\PY{l+s+s2}{\PYZdq{}}\PY{p}{,}
        \PY{l+s+s2}{\PYZdq{}}\PY{l+s+s2}{El precio de la tortilla no depende sólo del maíz}\PY{l+s+s2}{\PYZdq{}}\PY{p}{,}
        \PY{l+s+s2}{\PYZdq{}}\PY{l+s+s2}{Datos para proteger la seguridad alimentaria}\PY{l+s+s2}{\PYZdq{}}
    \PY{p}{]}
\PY{p}{\PYZcb{}}\PY{p}{)}

\PY{n}{segmentos\PYZus{}mercado}
\end{Verbatim}
\end{tcolorbox}

            \begin{tcolorbox}[breakable, size=fbox, boxrule=.5pt, pad at break*=1mm, opacityfill=0]
\prompt{Out}{outcolor}{299}{\boxspacing}
\begin{Verbatim}[commandchars=\\\{\}]
                    Actor                                Necesidad principal  \textbackslash{}
0             Productores      Saber cuándo vender para evitar precios bajos
1  Cooperativas agrícolas  Planear almacenamiento, compra colectiva y neg{\ldots}
2       Comercializadores                    Optimizar compras e inventarios
3    Industria tortillera              Anticipar costos del insumo principal
4            Consumidores  Comprender por qué cambia el precio de la tort{\ldots}
5          Sector público        Monitorear riesgos de seguridad alimentaria

                             Uso del modelo  \textbackslash{}
0            Pronóstico del precio del maíz
1      Escenarios de comercialización justa
2            Optimización de abastecimiento
3                      Planeación de costos
4   Comunicación clara del impacto del maíz
5  Alertas tempranas de presión alimentaria

                            Mensaje digital sugerido
0        Vende con información, no con incertidumbre
1             Negocia mejor usando datos del mercado
2       Compra en el momento de menor costo esperado
3             Anticipa tus costos antes de que suban
4  El precio de la tortilla no depende sólo del maíz
5       Datos para proteger la seguridad alimentaria
\end{Verbatim}
\end{tcolorbox}
        
    \hypertarget{segmentaciuxf3n-estratuxe9gica-de-usuarios}{%
\subsection{Segmentación estratégica de
usuarios}\label{segmentaciuxf3n-estratuxe9gica-de-usuarios}}

La segmentación anterior permite traducir el modelo predictivo en una
estrategia de comunicación diferenciada. No todos los actores requieren
el mismo tipo de información.

Para un productor, el valor está en anticipar si conviene vender o
esperar. Para una cooperativa, el valor está en negociar mejor, planear
almacenamiento o coordinar compras. Para la industria tortillera, el
pronóstico permite anticipar presiones de costo. Para consumidores y
sector público, el análisis ayuda a comprender que el precio de la
tortilla no depende exclusivamente del maíz, sino de una cadena de
costos más amplia.

Esta segmentación convierte el proyecto en una herramienta de
mercadotecnia digital basada en datos: cada público recibe información
útil, comprensible y orientada a decisiones.

    \begin{tcolorbox}[breakable, size=fbox, boxrule=1pt, pad at break*=1mm,colback=cellbackground, colframe=cellborder]
\prompt{In}{incolor}{302}{\boxspacing}
\begin{Verbatim}[commandchars=\\\{\}]
\PY{n}{maiz\PYZus{}tortilla}\PY{p}{[}\PY{l+s+s2}{\PYZdq{}}\PY{l+s+s2}{Participacion\PYZus{}maiz}\PY{l+s+s2}{\PYZdq{}}\PY{p}{]} \PY{o}{=} \PY{p}{(}
    \PY{n}{maiz\PYZus{}tortilla}\PY{p}{[}\PY{l+s+s2}{\PYZdq{}}\PY{l+s+s2}{Precio\PYZus{}maiz}\PY{l+s+s2}{\PYZdq{}}\PY{p}{]}
    \PY{o}{/} \PY{n}{maiz\PYZus{}tortilla}\PY{p}{[}\PY{l+s+s2}{\PYZdq{}}\PY{l+s+s2}{Precio\PYZus{}tortilla}\PY{l+s+s2}{\PYZdq{}}\PY{p}{]}
\PY{p}{)} \PY{o}{*} \PY{l+m+mi}{100}
\end{Verbatim}
\end{tcolorbox}

    \begin{tcolorbox}[breakable, size=fbox, boxrule=1pt, pad at break*=1mm,colback=cellbackground, colframe=cellborder]
\prompt{In}{incolor}{303}{\boxspacing}
\begin{Verbatim}[commandchars=\\\{\}]
\PY{n}{maiz\PYZus{}tortilla}\PY{p}{[}\PY{l+s+s2}{\PYZdq{}}\PY{l+s+s2}{Diferencia\PYZus{}pesos}\PY{l+s+s2}{\PYZdq{}}\PY{p}{]} \PY{o}{=} \PY{p}{(}
    \PY{n}{maiz\PYZus{}tortilla}\PY{p}{[}\PY{l+s+s2}{\PYZdq{}}\PY{l+s+s2}{Precio\PYZus{}tortilla}\PY{l+s+s2}{\PYZdq{}}\PY{p}{]}
    \PY{o}{\PYZhy{}} \PY{n}{maiz\PYZus{}tortilla}\PY{p}{[}\PY{l+s+s2}{\PYZdq{}}\PY{l+s+s2}{Precio\PYZus{}maiz}\PY{l+s+s2}{\PYZdq{}}\PY{p}{]}
\PY{p}{)}
\end{Verbatim}
\end{tcolorbox}

    \begin{tcolorbox}[breakable, size=fbox, boxrule=1pt, pad at break*=1mm,colback=cellbackground, colframe=cellborder]
\prompt{In}{incolor}{304}{\boxspacing}
\begin{Verbatim}[commandchars=\\\{\}]
\PY{n}{participacion\PYZus{}prom} \PY{o}{=} \PY{n}{maiz\PYZus{}tortilla}\PY{p}{[}\PY{l+s+s2}{\PYZdq{}}\PY{l+s+s2}{Participacion\PYZus{}maiz}\PY{l+s+s2}{\PYZdq{}}\PY{p}{]}\PY{o}{.}\PY{n}{mean}\PY{p}{(}\PY{p}{)}
\PY{n}{diferencia\PYZus{}prom} \PY{o}{=} \PY{n}{maiz\PYZus{}tortilla}\PY{p}{[}\PY{l+s+s2}{\PYZdq{}}\PY{l+s+s2}{Diferencia\PYZus{}pesos}\PY{l+s+s2}{\PYZdq{}}\PY{p}{]}\PY{o}{.}\PY{n}{mean}\PY{p}{(}\PY{p}{)}
\end{Verbatim}
\end{tcolorbox}

    \begin{tcolorbox}[breakable, size=fbox, boxrule=1pt, pad at break*=1mm,colback=cellbackground, colframe=cellborder]
\prompt{In}{incolor}{305}{\boxspacing}
\begin{Verbatim}[commandchars=\\\{\}]
\PY{n}{maiz\PYZus{}tortilla}\PY{p}{[}
    \PY{p}{[}\PY{l+s+s2}{\PYZdq{}}\PY{l+s+s2}{Precio\PYZus{}maiz}\PY{l+s+s2}{\PYZdq{}}\PY{p}{,}
     \PY{l+s+s2}{\PYZdq{}}\PY{l+s+s2}{Precio\PYZus{}tortilla}\PY{l+s+s2}{\PYZdq{}}\PY{p}{,}
     \PY{l+s+s2}{\PYZdq{}}\PY{l+s+s2}{Participacion\PYZus{}maiz}\PY{l+s+s2}{\PYZdq{}}\PY{p}{,}
     \PY{l+s+s2}{\PYZdq{}}\PY{l+s+s2}{Diferencia\PYZus{}pesos}\PY{l+s+s2}{\PYZdq{}}\PY{p}{]}
\PY{p}{]}\PY{o}{.}\PY{n}{head}\PY{p}{(}\PY{p}{)}
\end{Verbatim}
\end{tcolorbox}

            \begin{tcolorbox}[breakable, size=fbox, boxrule=.5pt, pad at break*=1mm, opacityfill=0]
\prompt{Out}{outcolor}{305}{\boxspacing}
\begin{Verbatim}[commandchars=\\\{\}]
   Precio\_maiz  Precio\_tortilla  Participacion\_maiz  Diferencia\_pesos
0         2.05             4.00           51.250000              1.95
1         2.06             4.04           50.990099              1.98
2         2.09             4.07           51.351351              1.98
3         2.12             4.10           51.707317              1.98
4         2.14             4.14           51.690821              2.00
\end{Verbatim}
\end{tcolorbox}
        
    \begin{tcolorbox}[breakable, size=fbox, boxrule=1pt, pad at break*=1mm,colback=cellbackground, colframe=cellborder]
\prompt{In}{incolor}{306}{\boxspacing}
\begin{Verbatim}[commandchars=\\\{\}]
\PY{n}{maiz\PYZus{}tortilla}\PY{p}{[}\PY{l+s+s2}{\PYZdq{}}\PY{l+s+s2}{Participacion\PYZus{}otros\PYZus{}costos}\PY{l+s+s2}{\PYZdq{}}\PY{p}{]} \PY{o}{=} \PY{p}{(}
    \PY{l+m+mi}{100} \PY{o}{\PYZhy{}} \PY{n}{maiz\PYZus{}tortilla}\PY{p}{[}\PY{l+s+s2}{\PYZdq{}}\PY{l+s+s2}{Participacion\PYZus{}maiz}\PY{l+s+s2}{\PYZdq{}}\PY{p}{]}
\PY{p}{)}
\end{Verbatim}
\end{tcolorbox}

    \begin{tcolorbox}[breakable, size=fbox, boxrule=1pt, pad at break*=1mm,colback=cellbackground, colframe=cellborder]
\prompt{In}{incolor}{307}{\boxspacing}
\begin{Verbatim}[commandchars=\\\{\}]
\PY{n}{otros\PYZus{}costos\PYZus{}prom} \PY{o}{=} \PY{p}{(}
    \PY{n}{maiz\PYZus{}tortilla}\PY{p}{[}\PY{l+s+s2}{\PYZdq{}}\PY{l+s+s2}{Participacion\PYZus{}otros\PYZus{}costos}\PY{l+s+s2}{\PYZdq{}}\PY{p}{]}
    \PY{o}{.}\PY{n}{mean}\PY{p}{(}\PY{p}{)}
\PY{p}{)}
\end{Verbatim}
\end{tcolorbox}

    \begin{tcolorbox}[breakable, size=fbox, boxrule=1pt, pad at break*=1mm,colback=cellbackground, colframe=cellborder]
\prompt{In}{incolor}{309}{\boxspacing}
\begin{Verbatim}[commandchars=\\\{\}]
\PY{n+nb}{print}\PY{p}{(}
    \PY{l+s+sa}{f}\PY{l+s+s2}{\PYZdq{}}\PY{l+s+s2}{Participación promedio del maíz: }\PY{l+s+si}{\PYZob{}}\PY{n}{participacion\PYZus{}prom}\PY{l+s+si}{:}\PY{l+s+s2}{.2f}\PY{l+s+si}{\PYZcb{}}\PY{l+s+s2}{\PYZpc{}}\PY{l+s+s2}{\PYZdq{}}
\PY{p}{)}

\PY{n+nb}{print}\PY{p}{(}
    \PY{l+s+sa}{f}\PY{l+s+s2}{\PYZdq{}}\PY{l+s+s2}{Diferencia promedio tortilla\PYZhy{}maíz: \PYZdl{}}\PY{l+s+si}{\PYZob{}}\PY{n}{diferencia\PYZus{}prom}\PY{l+s+si}{:}\PY{l+s+s2}{.2f}\PY{l+s+si}{\PYZcb{}}\PY{l+s+s2}{\PYZdq{}}
\PY{p}{)}
\end{Verbatim}
\end{tcolorbox}

    \begin{Verbatim}[commandchars=\\\{\}]
Participación promedio del maíz: 39.05\%
Diferencia promedio tortilla-maíz: \$7.52
    \end{Verbatim}

    \begin{tcolorbox}[breakable, size=fbox, boxrule=1pt, pad at break*=1mm,colback=cellbackground, colframe=cellborder]
\prompt{In}{incolor}{310}{\boxspacing}
\begin{Verbatim}[commandchars=\\\{\}]
\PY{n}{participacion\PYZus{}global} \PY{o}{=} \PY{p}{(}
    \PY{n}{maiz\PYZus{}tortilla}\PY{p}{[}\PY{l+s+s2}{\PYZdq{}}\PY{l+s+s2}{Precio\PYZus{}maiz}\PY{l+s+s2}{\PYZdq{}}\PY{p}{]}\PY{o}{.}\PY{n}{sum}\PY{p}{(}\PY{p}{)}
    \PY{o}{/} \PY{n}{maiz\PYZus{}tortilla}\PY{p}{[}\PY{l+s+s2}{\PYZdq{}}\PY{l+s+s2}{Precio\PYZus{}tortilla}\PY{l+s+s2}{\PYZdq{}}\PY{p}{]}\PY{o}{.}\PY{n}{sum}\PY{p}{(}\PY{p}{)}
\PY{p}{)} \PY{o}{*} \PY{l+m+mi}{100}
\end{Verbatim}
\end{tcolorbox}

    \begin{tcolorbox}[breakable, size=fbox, boxrule=1pt, pad at break*=1mm,colback=cellbackground, colframe=cellborder]
\prompt{In}{incolor}{311}{\boxspacing}
\begin{Verbatim}[commandchars=\\\{\}]
\PY{n}{participacion\PYZus{}promedio} \PY{o}{=} \PY{p}{(}
    \PY{n}{maiz\PYZus{}tortilla}\PY{p}{[}\PY{l+s+s2}{\PYZdq{}}\PY{l+s+s2}{Participacion\PYZus{}maiz}\PY{l+s+s2}{\PYZdq{}}\PY{p}{]}
    \PY{o}{.}\PY{n}{mean}\PY{p}{(}\PY{p}{)}
\PY{p}{)}

\PY{n}{participacion\PYZus{}global} \PY{o}{=} \PY{p}{(}
    \PY{n}{maiz\PYZus{}tortilla}\PY{p}{[}\PY{l+s+s2}{\PYZdq{}}\PY{l+s+s2}{Precio\PYZus{}maiz}\PY{l+s+s2}{\PYZdq{}}\PY{p}{]}\PY{o}{.}\PY{n}{sum}\PY{p}{(}\PY{p}{)}
    \PY{o}{/} \PY{n}{maiz\PYZus{}tortilla}\PY{p}{[}\PY{l+s+s2}{\PYZdq{}}\PY{l+s+s2}{Precio\PYZus{}tortilla}\PY{l+s+s2}{\PYZdq{}}\PY{p}{]}\PY{o}{.}\PY{n}{sum}\PY{p}{(}\PY{p}{)}
\PY{p}{)} \PY{o}{*} \PY{l+m+mi}{100}

\PY{n+nb}{print}\PY{p}{(}\PY{l+s+s2}{\PYZdq{}}\PY{l+s+s2}{Promedio mensual:}\PY{l+s+s2}{\PYZdq{}}\PY{p}{,} \PY{n}{participacion\PYZus{}promedio}\PY{p}{)}
\PY{n+nb}{print}\PY{p}{(}\PY{l+s+s2}{\PYZdq{}}\PY{l+s+s2}{Participación global:}\PY{l+s+s2}{\PYZdq{}}\PY{p}{,} \PY{n}{participacion\PYZus{}global}\PY{p}{)}
\end{Verbatim}
\end{tcolorbox}

    \begin{Verbatim}[commandchars=\\\{\}]
Promedio mensual: 39.045728264478505
Participación global: 38.08971716823781
    \end{Verbatim}

    \begin{tcolorbox}[breakable, size=fbox, boxrule=1pt, pad at break*=1mm,colback=cellbackground, colframe=cellborder]
\prompt{In}{incolor}{314}{\boxspacing}
\begin{Verbatim}[commandchars=\\\{\}]
\PY{n}{indicadores\PYZus{}marketing} \PY{o}{=} \PY{n}{pd}\PY{o}{.}\PY{n}{DataFrame}\PY{p}{(}\PY{p}{\PYZob{}}
    \PY{l+s+s2}{\PYZdq{}}\PY{l+s+s2}{Indicador}\PY{l+s+s2}{\PYZdq{}}\PY{p}{:} \PY{p}{[}
        \PY{l+s+s2}{\PYZdq{}}\PY{l+s+s2}{Precio promedio del maíz}\PY{l+s+s2}{\PYZdq{}}\PY{p}{,}
        \PY{l+s+s2}{\PYZdq{}}\PY{l+s+s2}{Precio promedio de la tortilla}\PY{l+s+s2}{\PYZdq{}}\PY{p}{,}
        \PY{l+s+s2}{\PYZdq{}}\PY{l+s+s2}{Participación promedio del maíz}\PY{l+s+s2}{\PYZdq{}}\PY{p}{,}
        \PY{l+s+s2}{\PYZdq{}}\PY{l+s+s2}{Participación global del maíz}\PY{l+s+s2}{\PYZdq{}}\PY{p}{,}
        \PY{l+s+s2}{\PYZdq{}}\PY{l+s+s2}{Diferencia promedio tortilla\PYZhy{}maíz}\PY{l+s+s2}{\PYZdq{}}\PY{p}{,}
        \PY{l+s+s2}{\PYZdq{}}\PY{l+s+s2}{Correlación maíz\PYZhy{}tortilla (niveles)}\PY{l+s+s2}{\PYZdq{}}\PY{p}{,}
        \PY{l+s+s2}{\PYZdq{}}\PY{l+s+s2}{Correlación maíz\PYZhy{}tortilla (variaciones)}\PY{l+s+s2}{\PYZdq{}}
    \PY{p}{]}\PY{p}{,}

    \PY{l+s+s2}{\PYZdq{}}\PY{l+s+s2}{Valor}\PY{l+s+s2}{\PYZdq{}}\PY{p}{:} \PY{p}{[}
        \PY{n+nb}{float}\PY{p}{(}\PY{n}{precio\PYZus{}maiz\PYZus{}prom}\PY{p}{)}\PY{p}{,}
        \PY{n+nb}{float}\PY{p}{(}\PY{n}{precio\PYZus{}tortilla\PYZus{}prom}\PY{p}{)}\PY{p}{,}
        \PY{n+nb}{float}\PY{p}{(}\PY{n}{participacion\PYZus{}prom}\PY{p}{)}\PY{p}{,}
        \PY{n+nb}{float}\PY{p}{(}\PY{n}{participacion\PYZus{}global}\PY{p}{)}\PY{p}{,}
        \PY{n+nb}{float}\PY{p}{(}\PY{n}{diferencia\PYZus{}prom}\PY{p}{)}\PY{p}{,}
        \PY{l+m+mf}{0.972}\PY{p}{,}
        \PY{l+m+mf}{0.266}
    \PY{p}{]}\PY{p}{,}

    \PY{l+s+s2}{\PYZdq{}}\PY{l+s+s2}{Interpretación comercial}\PY{l+s+s2}{\PYZdq{}}\PY{p}{:} \PY{p}{[}
        \PY{l+s+s2}{\PYZdq{}}\PY{l+s+s2}{Costo histórico promedio del insumo agrícola}\PY{l+s+s2}{\PYZdq{}}\PY{p}{,}
        \PY{l+s+s2}{\PYZdq{}}\PY{l+s+s2}{Precio promedio pagado por consumidores}\PY{l+s+s2}{\PYZdq{}}\PY{p}{,}
        \PY{l+s+s2}{\PYZdq{}}\PY{l+s+s2}{Porcentaje del precio de la tortilla explicado por el maíz}\PY{l+s+s2}{\PYZdq{}}\PY{p}{,}
        \PY{l+s+s2}{\PYZdq{}}\PY{l+s+s2}{Participación histórica agregada del maíz en la cadena}\PY{l+s+s2}{\PYZdq{}}\PY{p}{,}
        \PY{l+s+s2}{\PYZdq{}}\PY{l+s+s2}{Valor económico asociado a otros costos y márgenes}\PY{l+s+s2}{\PYZdq{}}\PY{p}{,}
        \PY{l+s+s2}{\PYZdq{}}\PY{l+s+s2}{Relación estructural de largo plazo entre ambos productos}\PY{l+s+s2}{\PYZdq{}}\PY{p}{,}
        \PY{l+s+s2}{\PYZdq{}}\PY{l+s+s2}{Transmisión mensual parcial y no proporcional}\PY{l+s+s2}{\PYZdq{}}
    \PY{p}{]}
\PY{p}{\PYZcb{}}\PY{p}{)}

\PY{n}{indicadores\PYZus{}marketing}\PY{p}{[}\PY{l+s+s2}{\PYZdq{}}\PY{l+s+s2}{Valor}\PY{l+s+s2}{\PYZdq{}}\PY{p}{]} \PY{o}{=} \PY{n}{indicadores\PYZus{}marketing}\PY{p}{[}\PY{l+s+s2}{\PYZdq{}}\PY{l+s+s2}{Valor}\PY{l+s+s2}{\PYZdq{}}\PY{p}{]}\PY{o}{.}\PY{n}{round}\PY{p}{(}\PY{l+m+mi}{2}\PY{p}{)}

\PY{n}{indicadores\PYZus{}marketing}
\end{Verbatim}
\end{tcolorbox}

            \begin{tcolorbox}[breakable, size=fbox, boxrule=.5pt, pad at break*=1mm, opacityfill=0]
\prompt{Out}{outcolor}{314}{\boxspacing}
\begin{Verbatim}[commandchars=\\\{\}]
                                 Indicador  Valor  \textbackslash{}
0                 Precio promedio del maíz   4.63
1           Precio promedio de la tortilla  12.14
2          Participación promedio del maíz  39.05
3            Participación global del maíz  38.09
4        Diferencia promedio tortilla-maíz   7.52
5      Correlación maíz-tortilla (niveles)   0.97
6  Correlación maíz-tortilla (variaciones)   0.27

                            Interpretación comercial
0       Costo histórico promedio del insumo agrícola
1            Precio promedio pagado por consumidores
2  Porcentaje del precio de la tortilla explicado{\ldots}
3  Participación histórica agregada del maíz en l{\ldots}
4  Valor económico asociado a otros costos y márg{\ldots}
5  Relación estructural de largo plazo entre ambo{\ldots}
6      Transmisión mensual parcial y no proporcional
\end{Verbatim}
\end{tcolorbox}
        
    \begin{tcolorbox}[breakable, size=fbox, boxrule=1pt, pad at break*=1mm,colback=cellbackground, colframe=cellborder]
\prompt{In}{incolor}{316}{\boxspacing}
\begin{Verbatim}[commandchars=\\\{\}]
\PY{k+kn}{import} \PY{n+nn}{matplotlib}\PY{n+nn}{.}\PY{n+nn}{pyplot} \PY{k}{as} \PY{n+nn}{plt}

\PY{n}{participacion\PYZus{}otros} \PY{o}{=} \PY{l+m+mi}{100} \PY{o}{\PYZhy{}} \PY{n}{participacion\PYZus{}prom}

\PY{n}{fig}\PY{p}{,} \PY{n}{ax} \PY{o}{=} \PY{n}{plt}\PY{o}{.}\PY{n}{subplots}\PY{p}{(}\PY{n}{figsize}\PY{o}{=}\PY{p}{(}\PY{l+m+mi}{8}\PY{p}{,} \PY{l+m+mi}{8}\PY{p}{)}\PY{p}{)}

\PY{n}{ax}\PY{o}{.}\PY{n}{pie}\PY{p}{(}
    \PY{p}{[}\PY{n}{participacion\PYZus{}prom}\PY{p}{,} \PY{n}{participacion\PYZus{}otros}\PY{p}{]}\PY{p}{,}
    \PY{n}{labels}\PY{o}{=}\PY{p}{[}
        \PY{l+s+sa}{f}\PY{l+s+s2}{\PYZdq{}}\PY{l+s+s2}{Maíz}\PY{l+s+se}{\PYZbs{}n}\PY{l+s+s2}{(}\PY{l+s+si}{\PYZob{}}\PY{n}{participacion\PYZus{}prom}\PY{l+s+si}{:}\PY{l+s+s2}{.2f}\PY{l+s+si}{\PYZcb{}}\PY{l+s+s2}{\PYZpc{})}\PY{l+s+s2}{\PYZdq{}}\PY{p}{,}
        \PY{l+s+sa}{f}\PY{l+s+s2}{\PYZdq{}}\PY{l+s+s2}{Otros costos}\PY{l+s+se}{\PYZbs{}n}\PY{l+s+s2}{(}\PY{l+s+si}{\PYZob{}}\PY{n}{participacion\PYZus{}otros}\PY{l+s+si}{:}\PY{l+s+s2}{.2f}\PY{l+s+si}{\PYZcb{}}\PY{l+s+s2}{\PYZpc{})}\PY{l+s+s2}{\PYZdq{}}
    \PY{p}{]}\PY{p}{,}
    \PY{n}{autopct}\PY{o}{=}\PY{l+s+s2}{\PYZdq{}}\PY{l+s+si}{\PYZpc{}1.1f}\PY{l+s+si}{\PYZpc{}\PYZpc{}}\PY{l+s+s2}{\PYZdq{}}\PY{p}{,}
    \PY{n}{startangle}\PY{o}{=}\PY{l+m+mi}{90}
\PY{p}{)}

\PY{n}{ax}\PY{o}{.}\PY{n}{set\PYZus{}title}\PY{p}{(}
    \PY{l+s+s2}{\PYZdq{}}\PY{l+s+s2}{Composición promedio del precio de la tortilla}\PY{l+s+se}{\PYZbs{}n}\PY{l+s+s2}{México 2000\PYZhy{}2026}\PY{l+s+s2}{\PYZdq{}}\PY{p}{,}
    \PY{n}{fontsize}\PY{o}{=}\PY{l+m+mi}{14}
\PY{p}{)}

\PY{n}{plt}\PY{o}{.}\PY{n}{show}\PY{p}{(}\PY{p}{)}
\end{Verbatim}
\end{tcolorbox}

    \begin{center}
    \adjustimage{max size={0.9\linewidth}{0.9\paperheight}}{output_233_0.png}
    \end{center}
    { \hspace*{\fill} \\}
    
    \begin{tcolorbox}[breakable, size=fbox, boxrule=1pt, pad at break*=1mm,colback=cellbackground, colframe=cellborder]
\prompt{In}{incolor}{317}{\boxspacing}
\begin{Verbatim}[commandchars=\\\{\}]
\PY{n}{pronostico\PYZus{}io}\PY{o}{.}\PY{n}{head}\PY{p}{(}\PY{p}{)}
\end{Verbatim}
\end{tcolorbox}

            \begin{tcolorbox}[breakable, size=fbox, boxrule=.5pt, pad at break*=1mm, opacityfill=0]
\prompt{Out}{outcolor}{317}{\boxspacing}
\begin{Verbatim}[commandchars=\\\{\}]
   Mes  Precio\_Pronosticado
0    1             9.665025
1    2             9.836343
2    3             9.733932
3    4             9.673063
4    5             9.566773
\end{Verbatim}
\end{tcolorbox}
        
    \begin{tcolorbox}[breakable, size=fbox, boxrule=1pt, pad at break*=1mm,colback=cellbackground, colframe=cellborder]
\prompt{In}{incolor}{318}{\boxspacing}
\begin{Verbatim}[commandchars=\\\{\}]
\PY{n}{pronostico\PYZus{}marketing} \PY{o}{=} \PY{n}{pronostico\PYZus{}io}\PY{o}{.}\PY{n}{copy}\PY{p}{(}\PY{p}{)}

\PY{n}{precio\PYZus{}prom\PYZus{}pronosticado} \PY{o}{=} \PY{p}{(}
    \PY{n}{pronostico\PYZus{}marketing}\PY{p}{[}\PY{l+s+s2}{\PYZdq{}}\PY{l+s+s2}{Precio\PYZus{}Pronosticado}\PY{l+s+s2}{\PYZdq{}}\PY{p}{]}
    \PY{o}{.}\PY{n}{mean}\PY{p}{(}\PY{p}{)}
\PY{p}{)}

\PY{n}{pronostico\PYZus{}marketing}\PY{p}{[}\PY{l+s+s2}{\PYZdq{}}\PY{l+s+s2}{Diferencia\PYZus{}vs\PYZus{}promedio}\PY{l+s+s2}{\PYZdq{}}\PY{p}{]} \PY{o}{=} \PY{p}{(}
    \PY{n}{pronostico\PYZus{}marketing}\PY{p}{[}\PY{l+s+s2}{\PYZdq{}}\PY{l+s+s2}{Precio\PYZus{}Pronosticado}\PY{l+s+s2}{\PYZdq{}}\PY{p}{]}
    \PY{o}{\PYZhy{}} \PY{n}{precio\PYZus{}prom\PYZus{}pronosticado}
\PY{p}{)}

\PY{n}{pronostico\PYZus{}marketing}\PY{p}{[}\PY{l+s+s2}{\PYZdq{}}\PY{l+s+s2}{Senal\PYZus{}comercial}\PY{l+s+s2}{\PYZdq{}}\PY{p}{]} \PY{o}{=} \PY{p}{(}
    \PY{n}{pronostico\PYZus{}marketing}\PY{p}{[}\PY{l+s+s2}{\PYZdq{}}\PY{l+s+s2}{Diferencia\PYZus{}vs\PYZus{}promedio}\PY{l+s+s2}{\PYZdq{}}\PY{p}{]}
    \PY{o}{.}\PY{n}{apply}\PY{p}{(}
        \PY{k}{lambda} \PY{n}{x}\PY{p}{:}
        \PY{l+s+s2}{\PYZdq{}}\PY{l+s+s2}{Oportunidad de compra}\PY{l+s+s2}{\PYZdq{}}
        \PY{k}{if} \PY{n}{x} \PY{o}{\PYZlt{}} \PY{l+m+mi}{0}
        \PY{k}{else} \PY{l+s+s2}{\PYZdq{}}\PY{l+s+s2}{Precio superior al promedio}\PY{l+s+s2}{\PYZdq{}}
    \PY{p}{)}
\PY{p}{)}

\PY{n}{pronostico\PYZus{}marketing}
\end{Verbatim}
\end{tcolorbox}

            \begin{tcolorbox}[breakable, size=fbox, boxrule=.5pt, pad at break*=1mm, opacityfill=0]
\prompt{Out}{outcolor}{318}{\boxspacing}
\begin{Verbatim}[commandchars=\\\{\}]
   Mes  Precio\_Pronosticado  Diferencia\_vs\_promedio  \textbackslash{}
0    1             9.665025               -0.005318
1    2             9.836343                0.166000
2    3             9.733932                0.063589
3    4             9.673063                0.002720
4    5             9.566773               -0.103570
5    6             9.546923               -0.123420

               Senal\_comercial
0        Oportunidad de compra
1  Precio superior al promedio
2  Precio superior al promedio
3  Precio superior al promedio
4        Oportunidad de compra
5        Oportunidad de compra
\end{Verbatim}
\end{tcolorbox}
        
    \begin{tcolorbox}[breakable, size=fbox, boxrule=1pt, pad at break*=1mm,colback=cellbackground, colframe=cellborder]
\prompt{In}{incolor}{319}{\boxspacing}
\begin{Verbatim}[commandchars=\\\{\}]
\PY{k+kn}{import} \PY{n+nn}{matplotlib}\PY{n+nn}{.}\PY{n+nn}{pyplot} \PY{k}{as} \PY{n+nn}{plt}

\PY{n}{plt}\PY{o}{.}\PY{n}{figure}\PY{p}{(}\PY{n}{figsize}\PY{o}{=}\PY{p}{(}\PY{l+m+mi}{10}\PY{p}{,}\PY{l+m+mi}{5}\PY{p}{)}\PY{p}{)}

\PY{n}{plt}\PY{o}{.}\PY{n}{plot}\PY{p}{(}
    \PY{n}{pronostico\PYZus{}marketing}\PY{p}{[}\PY{l+s+s2}{\PYZdq{}}\PY{l+s+s2}{Mes}\PY{l+s+s2}{\PYZdq{}}\PY{p}{]}\PY{p}{,}
    \PY{n}{pronostico\PYZus{}marketing}\PY{p}{[}\PY{l+s+s2}{\PYZdq{}}\PY{l+s+s2}{Precio\PYZus{}Pronosticado}\PY{l+s+s2}{\PYZdq{}}\PY{p}{]}\PY{p}{,}
    \PY{n}{marker}\PY{o}{=}\PY{l+s+s2}{\PYZdq{}}\PY{l+s+s2}{o}\PY{l+s+s2}{\PYZdq{}}\PY{p}{,}
    \PY{n}{linewidth}\PY{o}{=}\PY{l+m+mi}{2}\PY{p}{,}
    \PY{n}{label}\PY{o}{=}\PY{l+s+s2}{\PYZdq{}}\PY{l+s+s2}{Pronóstico SARIMAX}\PY{l+s+s2}{\PYZdq{}}
\PY{p}{)}

\PY{n}{plt}\PY{o}{.}\PY{n}{axhline}\PY{p}{(}
    \PY{n}{y}\PY{o}{=}\PY{n}{precio\PYZus{}prom\PYZus{}pronosticado}\PY{p}{,}
    \PY{n}{linestyle}\PY{o}{=}\PY{l+s+s2}{\PYZdq{}}\PY{l+s+s2}{\PYZhy{}\PYZhy{}}\PY{l+s+s2}{\PYZdq{}}\PY{p}{,}
    \PY{n}{linewidth}\PY{o}{=}\PY{l+m+mi}{2}\PY{p}{,}
    \PY{n}{label}\PY{o}{=}\PY{l+s+s2}{\PYZdq{}}\PY{l+s+s2}{Promedio pronosticado}\PY{l+s+s2}{\PYZdq{}}
\PY{p}{)}

\PY{n}{plt}\PY{o}{.}\PY{n}{title}\PY{p}{(}
    \PY{l+s+s2}{\PYZdq{}}\PY{l+s+s2}{Identificación de oportunidades de compra}\PY{l+s+se}{\PYZbs{}n}\PY{l+s+s2}{basadas en el pronóstico SARIMAX}\PY{l+s+s2}{\PYZdq{}}
\PY{p}{)}

\PY{n}{plt}\PY{o}{.}\PY{n}{xlabel}\PY{p}{(}\PY{l+s+s2}{\PYZdq{}}\PY{l+s+s2}{Mes}\PY{l+s+s2}{\PYZdq{}}\PY{p}{)}
\PY{n}{plt}\PY{o}{.}\PY{n}{ylabel}\PY{p}{(}\PY{l+s+s2}{\PYZdq{}}\PY{l+s+s2}{Precio proyectado (MXN/Kg)}\PY{l+s+s2}{\PYZdq{}}\PY{p}{)}

\PY{n}{plt}\PY{o}{.}\PY{n}{grid}\PY{p}{(}\PY{k+kc}{True}\PY{p}{)}
\PY{n}{plt}\PY{o}{.}\PY{n}{legend}\PY{p}{(}\PY{p}{)}

\PY{n}{plt}\PY{o}{.}\PY{n}{show}\PY{p}{(}\PY{p}{)}
\end{Verbatim}
\end{tcolorbox}

    \begin{center}
    \adjustimage{max size={0.9\linewidth}{0.9\paperheight}}{output_236_0.png}
    \end{center}
    { \hspace*{\fill} \\}
    
    \hypertarget{seuxf1ales-comerciales-derivadas-del-pronuxf3stico}{%
\subsection{Señales comerciales derivadas del
pronóstico}\label{seuxf1ales-comerciales-derivadas-del-pronuxf3stico}}

La capacidad predictiva del modelo SARIMAX permite transformar un
pronóstico estadístico en una herramienta de apoyo comercial.

Al comparar cada precio proyectado contra el promedio esperado del
horizonte de planeación, es posible identificar periodos relativamente
favorables para la compra o el almacenamiento del maíz.

Esta lógica constituye un ejemplo de inteligencia comercial basada en
datos, ya que permite convertir resultados analíticos en acciones
concretas para productores, cooperativas y comercializadores.

Desde la perspectiva de Mercadotecnia Digital, el valor no se encuentra
únicamente en mostrar una predicción, sino en traducirla a mensajes
simples, comprensibles y accionables para distintos tipos de usuarios.

    \begin{tcolorbox}[breakable, size=fbox, boxrule=1pt, pad at break*=1mm,colback=cellbackground, colframe=cellborder]
\prompt{In}{incolor}{320}{\boxspacing}
\begin{Verbatim}[commandchars=\\\{\}]
\PY{n}{segmentacion\PYZus{}usuarios} \PY{o}{=} \PY{n}{pd}\PY{o}{.}\PY{n}{DataFrame}\PY{p}{(}\PY{p}{\PYZob{}}

    \PY{l+s+s2}{\PYZdq{}}\PY{l+s+s2}{Usuario}\PY{l+s+s2}{\PYZdq{}}\PY{p}{:} \PY{p}{[}
        \PY{l+s+s2}{\PYZdq{}}\PY{l+s+s2}{Productor}\PY{l+s+s2}{\PYZdq{}}\PY{p}{,}
        \PY{l+s+s2}{\PYZdq{}}\PY{l+s+s2}{Cooperativa agrícola}\PY{l+s+s2}{\PYZdq{}}\PY{p}{,}
        \PY{l+s+s2}{\PYZdq{}}\PY{l+s+s2}{Comercializador}\PY{l+s+s2}{\PYZdq{}}\PY{p}{,}
        \PY{l+s+s2}{\PYZdq{}}\PY{l+s+s2}{Industria tortillera}\PY{l+s+s2}{\PYZdq{}}\PY{p}{,}
        \PY{l+s+s2}{\PYZdq{}}\PY{l+s+s2}{Gobierno}\PY{l+s+s2}{\PYZdq{}}
    \PY{p}{]}\PY{p}{,}

    \PY{l+s+s2}{\PYZdq{}}\PY{l+s+s2}{Información recibida}\PY{l+s+s2}{\PYZdq{}}\PY{p}{:} \PY{p}{[}
        \PY{l+s+s2}{\PYZdq{}}\PY{l+s+s2}{Pronóstico SARIMAX del maíz}\PY{l+s+s2}{\PYZdq{}}\PY{p}{,}
        \PY{l+s+s2}{\PYZdq{}}\PY{l+s+s2}{Señales comerciales y pronóstico}\PY{l+s+s2}{\PYZdq{}}\PY{p}{,}
        \PY{l+s+s2}{\PYZdq{}}\PY{l+s+s2}{Valle de oportunidad de compra}\PY{l+s+s2}{\PYZdq{}}\PY{p}{,}
        \PY{l+s+s2}{\PYZdq{}}\PY{l+s+s2}{Relación maíz\PYZhy{}tortilla}\PY{l+s+s2}{\PYZdq{}}\PY{p}{,}
        \PY{l+s+s2}{\PYZdq{}}\PY{l+s+s2}{Indicadores de riesgo alimentario}\PY{l+s+s2}{\PYZdq{}}
    \PY{p}{]}\PY{p}{,}

    \PY{l+s+s2}{\PYZdq{}}\PY{l+s+s2}{Decisión apoyada}\PY{l+s+s2}{\PYZdq{}}\PY{p}{:} \PY{p}{[}
        \PY{l+s+s2}{\PYZdq{}}\PY{l+s+s2}{Cuándo vender}\PY{l+s+s2}{\PYZdq{}}\PY{p}{,}
        \PY{l+s+s2}{\PYZdq{}}\PY{l+s+s2}{Cuándo almacenar}\PY{l+s+s2}{\PYZdq{}}\PY{p}{,}
        \PY{l+s+s2}{\PYZdq{}}\PY{l+s+s2}{Cuándo comprar}\PY{l+s+s2}{\PYZdq{}}\PY{p}{,}
        \PY{l+s+s2}{\PYZdq{}}\PY{l+s+s2}{Planeación de costos}\PY{l+s+s2}{\PYZdq{}}\PY{p}{,}
        \PY{l+s+s2}{\PYZdq{}}\PY{l+s+s2}{Monitoreo y prevención}\PY{l+s+s2}{\PYZdq{}}
    \PY{p}{]}\PY{p}{,}

    \PY{l+s+s2}{\PYZdq{}}\PY{l+s+s2}{Beneficio esperado}\PY{l+s+s2}{\PYZdq{}}\PY{p}{:} \PY{p}{[}
        \PY{l+s+s2}{\PYZdq{}}\PY{l+s+s2}{Mejor precio de venta}\PY{l+s+s2}{\PYZdq{}}\PY{p}{,}
        \PY{l+s+s2}{\PYZdq{}}\PY{l+s+s2}{Mayor capacidad de negociación}\PY{l+s+s2}{\PYZdq{}}\PY{p}{,}
        \PY{l+s+s2}{\PYZdq{}}\PY{l+s+s2}{Menor costo de abastecimiento}\PY{l+s+s2}{\PYZdq{}}\PY{p}{,}
        \PY{l+s+s2}{\PYZdq{}}\PY{l+s+s2}{Control de márgenes}\PY{l+s+s2}{\PYZdq{}}\PY{p}{,}
        \PY{l+s+s2}{\PYZdq{}}\PY{l+s+s2}{Reducción de riesgos}\PY{l+s+s2}{\PYZdq{}}
    \PY{p}{]}
\PY{p}{\PYZcb{}}\PY{p}{)}

\PY{n}{segmentacion\PYZus{}usuarios}
\end{Verbatim}
\end{tcolorbox}

            \begin{tcolorbox}[breakable, size=fbox, boxrule=.5pt, pad at break*=1mm, opacityfill=0]
\prompt{Out}{outcolor}{320}{\boxspacing}
\begin{Verbatim}[commandchars=\\\{\}]
                Usuario               Información recibida  \textbackslash{}
0             Productor        Pronóstico SARIMAX del maíz
1  Cooperativa agrícola   Señales comerciales y pronóstico
2       Comercializador     Valle de oportunidad de compra
3  Industria tortillera             Relación maíz-tortilla
4              Gobierno  Indicadores de riesgo alimentario

         Decisión apoyada              Beneficio esperado
0           Cuándo vender           Mejor precio de venta
1        Cuándo almacenar  Mayor capacidad de negociación
2          Cuándo comprar   Menor costo de abastecimiento
3    Planeación de costos             Control de márgenes
4  Monitoreo y prevención            Reducción de riesgos
\end{Verbatim}
\end{tcolorbox}
        
    \hypertarget{segmentaciuxf3n-de-usuarios-y-propuesta-de-valor}{%
\subsection{Segmentación de usuarios y propuesta de
valor}\label{segmentaciuxf3n-de-usuarios-y-propuesta-de-valor}}

Una característica fundamental de los sistemas de mercadotecnia digital
es que distintos usuarios requieren distintos tipos de información.

El proyecto Mercado Justo transforma los resultados del modelo
predictivo en información específica para cada actor de la cadena
agroalimentaria.

De esta forma, el sistema deja de ser únicamente un ejercicio de
análisis de datos y se convierte en una herramienta digital de
inteligencia comercial capaz de apoyar decisiones concretas para
productores, cooperativas, comercializadores, industria alimentaria y
organismos públicos.

    \begin{tcolorbox}[breakable, size=fbox, boxrule=1pt, pad at break*=1mm,colback=cellbackground, colframe=cellborder]
\prompt{In}{incolor}{321}{\boxspacing}
\begin{Verbatim}[commandchars=\\\{\}]
\PY{n}{propuesta\PYZus{}valor} \PY{o}{=} \PY{n}{pd}\PY{o}{.}\PY{n}{DataFrame}\PY{p}{(}\PY{p}{\PYZob{}}

    \PY{l+s+s2}{\PYZdq{}}\PY{l+s+s2}{Componente}\PY{l+s+s2}{\PYZdq{}}\PY{p}{:} \PY{p}{[}
        \PY{l+s+s2}{\PYZdq{}}\PY{l+s+s2}{Predicción de precios}\PY{l+s+s2}{\PYZdq{}}\PY{p}{,}
        \PY{l+s+s2}{\PYZdq{}}\PY{l+s+s2}{Explicación de la tortilla}\PY{l+s+s2}{\PYZdq{}}\PY{p}{,}
        \PY{l+s+s2}{\PYZdq{}}\PY{l+s+s2}{Señales comerciales}\PY{l+s+s2}{\PYZdq{}}\PY{p}{,}
        \PY{l+s+s2}{\PYZdq{}}\PY{l+s+s2}{Optimización de compras}\PY{l+s+s2}{\PYZdq{}}\PY{p}{,}
        \PY{l+s+s2}{\PYZdq{}}\PY{l+s+s2}{Dashboard estratégico}\PY{l+s+s2}{\PYZdq{}}
    \PY{p}{]}\PY{p}{,}

    \PY{l+s+s2}{\PYZdq{}}\PY{l+s+s2}{Valor generado}\PY{l+s+s2}{\PYZdq{}}\PY{p}{:} \PY{p}{[}
        \PY{l+s+s2}{\PYZdq{}}\PY{l+s+s2}{Reduce incertidumbre}\PY{l+s+s2}{\PYZdq{}}\PY{p}{,}
        \PY{l+s+s2}{\PYZdq{}}\PY{l+s+s2}{Explica formación de precios}\PY{l+s+s2}{\PYZdq{}}\PY{p}{,}
        \PY{l+s+s2}{\PYZdq{}}\PY{l+s+s2}{Identifica oportunidades}\PY{l+s+s2}{\PYZdq{}}\PY{p}{,}
        \PY{l+s+s2}{\PYZdq{}}\PY{l+s+s2}{Minimiza costos esperados}\PY{l+s+s2}{\PYZdq{}}\PY{p}{,}
        \PY{l+s+s2}{\PYZdq{}}\PY{l+s+s2}{Facilita toma de decisiones}\PY{l+s+s2}{\PYZdq{}}
    \PY{p}{]}\PY{p}{,}

    \PY{l+s+s2}{\PYZdq{}}\PY{l+s+s2}{Beneficiario principal}\PY{l+s+s2}{\PYZdq{}}\PY{p}{:} \PY{p}{[}
        \PY{l+s+s2}{\PYZdq{}}\PY{l+s+s2}{Productores}\PY{l+s+s2}{\PYZdq{}}\PY{p}{,}
        \PY{l+s+s2}{\PYZdq{}}\PY{l+s+s2}{Consumidores}\PY{l+s+s2}{\PYZdq{}}\PY{p}{,}
        \PY{l+s+s2}{\PYZdq{}}\PY{l+s+s2}{Comercializadores}\PY{l+s+s2}{\PYZdq{}}\PY{p}{,}
        \PY{l+s+s2}{\PYZdq{}}\PY{l+s+s2}{Cooperativas}\PY{l+s+s2}{\PYZdq{}}\PY{p}{,}
        \PY{l+s+s2}{\PYZdq{}}\PY{l+s+s2}{Toda la cadena}\PY{l+s+s2}{\PYZdq{}}
    \PY{p}{]}
\PY{p}{\PYZcb{}}\PY{p}{)}

\PY{n}{propuesta\PYZus{}valor}
\end{Verbatim}
\end{tcolorbox}

            \begin{tcolorbox}[breakable, size=fbox, boxrule=.5pt, pad at break*=1mm, opacityfill=0]
\prompt{Out}{outcolor}{321}{\boxspacing}
\begin{Verbatim}[commandchars=\\\{\}]
                   Componente                Valor generado  \textbackslash{}
0       Predicción de precios          Reduce incertidumbre
1  Explicación de la tortilla  Explica formación de precios
2         Señales comerciales      Identifica oportunidades
3     Optimización de compras     Minimiza costos esperados
4       Dashboard estratégico   Facilita toma de decisiones

  Beneficiario principal
0            Productores
1           Consumidores
2      Comercializadores
3           Cooperativas
4         Toda la cadena
\end{Verbatim}
\end{tcolorbox}
        
    \begin{tcolorbox}[breakable, size=fbox, boxrule=1pt, pad at break*=1mm,colback=cellbackground, colframe=cellborder]
\prompt{In}{incolor}{322}{\boxspacing}
\begin{Verbatim}[commandchars=\\\{\}]
\PY{n}{kpis\PYZus{}mercado\PYZus{}justo} \PY{o}{=} \PY{n}{pd}\PY{o}{.}\PY{n}{DataFrame}\PY{p}{(}\PY{p}{\PYZob{}}

    \PY{l+s+s2}{\PYZdq{}}\PY{l+s+s2}{Indicador}\PY{l+s+s2}{\PYZdq{}}\PY{p}{:} \PY{p}{[}
        \PY{l+s+s2}{\PYZdq{}}\PY{l+s+s2}{MAPE modelo SARIMAX}\PY{l+s+s2}{\PYZdq{}}\PY{p}{,}
        \PY{l+s+s2}{\PYZdq{}}\PY{l+s+s2}{Participación del maíz en tortilla}\PY{l+s+s2}{\PYZdq{}}\PY{p}{,}
        \PY{l+s+s2}{\PYZdq{}}\PY{l+s+s2}{Correlación maíz\PYZhy{}tortilla}\PY{l+s+s2}{\PYZdq{}}\PY{p}{,}
        \PY{l+s+s2}{\PYZdq{}}\PY{l+s+s2}{Horizonte de pronóstico}\PY{l+s+s2}{\PYZdq{}}\PY{p}{,}
        \PY{l+s+s2}{\PYZdq{}}\PY{l+s+s2}{Fuentes integradas}\PY{l+s+s2}{\PYZdq{}}
    \PY{p}{]}\PY{p}{,}

    \PY{l+s+s2}{\PYZdq{}}\PY{l+s+s2}{Valor}\PY{l+s+s2}{\PYZdq{}}\PY{p}{:} \PY{p}{[}
        \PY{l+s+s2}{\PYZdq{}}\PY{l+s+s2}{15.83}\PY{l+s+s2}{\PYZpc{}}\PY{l+s+s2}{\PYZdq{}}\PY{p}{,}
        \PY{l+s+s2}{\PYZdq{}}\PY{l+s+s2}{39.05}\PY{l+s+s2}{\PYZpc{}}\PY{l+s+s2}{\PYZdq{}}\PY{p}{,}
        \PY{l+s+s2}{\PYZdq{}}\PY{l+s+s2}{0.97}\PY{l+s+s2}{\PYZdq{}}\PY{p}{,}
        \PY{l+s+s2}{\PYZdq{}}\PY{l+s+s2}{6 meses}\PY{l+s+s2}{\PYZdq{}}\PY{p}{,}
        \PY{l+s+s2}{\PYZdq{}}\PY{l+s+s2}{9 fuentes oficiales}\PY{l+s+s2}{\PYZdq{}}
    \PY{p}{]}\PY{p}{,}

    \PY{l+s+s2}{\PYZdq{}}\PY{l+s+s2}{Contribución}\PY{l+s+s2}{\PYZdq{}}\PY{p}{:} \PY{p}{[}
        \PY{l+s+s2}{\PYZdq{}}\PY{l+s+s2}{Precisión predictiva}\PY{l+s+s2}{\PYZdq{}}\PY{p}{,}
        \PY{l+s+s2}{\PYZdq{}}\PY{l+s+s2}{Comprensión de precios}\PY{l+s+s2}{\PYZdq{}}\PY{p}{,}
        \PY{l+s+s2}{\PYZdq{}}\PY{l+s+s2}{Relación estructural}\PY{l+s+s2}{\PYZdq{}}\PY{p}{,}
        \PY{l+s+s2}{\PYZdq{}}\PY{l+s+s2}{Planeación comercial}\PY{l+s+s2}{\PYZdq{}}\PY{p}{,}
        \PY{l+s+s2}{\PYZdq{}}\PY{l+s+s2}{Inteligencia de mercado}\PY{l+s+s2}{\PYZdq{}}
    \PY{p}{]}
\PY{p}{\PYZcb{}}\PY{p}{)}

\PY{n}{kpis\PYZus{}mercado\PYZus{}justo}
\end{Verbatim}
\end{tcolorbox}

            \begin{tcolorbox}[breakable, size=fbox, boxrule=.5pt, pad at break*=1mm, opacityfill=0]
\prompt{Out}{outcolor}{322}{\boxspacing}
\begin{Verbatim}[commandchars=\\\{\}]
                            Indicador                Valor  \textbackslash{}
0                 MAPE modelo SARIMAX               15.83\%
1  Participación del maíz en tortilla               39.05\%
2           Correlación maíz-tortilla                 0.97
3             Horizonte de pronóstico              6 meses
4                  Fuentes integradas  9 fuentes oficiales

              Contribución
0     Precisión predictiva
1   Comprensión de precios
2     Relación estructural
3     Planeación comercial
4  Inteligencia de mercado
\end{Verbatim}
\end{tcolorbox}
        
    \hypertarget{mercadotecnia-digital-basada-en-ciencia-de-datos}{%
\subsection{Mercadotecnia Digital basada en Ciencia de
Datos}\label{mercadotecnia-digital-basada-en-ciencia-de-datos}}

Tradicionalmente, la mercadotecnia digital se asocia con campañas
publicitarias, redes sociales o posicionamiento de productos.

Sin embargo, en mercados complejos como el agroalimentario, la
generación de valor también puede surgir mediante la democratización de
información estratégica.

El sistema desarrollado permite transformar grandes volúmenes de datos
en conocimiento accionable, reduciendo las asimetrías de información
entre actores con distintos niveles de acceso a información de mercado.

Desde esta perspectiva, el Modelo de Mercado Justo, análisis y
predicicón de precios del maíz y su impacto en la tortilla puede
entenderse como un producto digital de inteligencia comercial sustentado
en analítica predictiva.

    La aportación de Mercadotecnia Digital dentro del proyecto no se limita
a la difusión de información, sino a la creación de una plataforma de
inteligencia comercial basada en datos.

Los modelos predictivos desarrollados permiten anticipar escenarios de
mercado, mientras que las herramientas de visualización y segmentación
convierten dichos resultados en información útil para distintos
usuarios.

De esta manera, Mercado Justo integra analítica predictiva, comunicación
estratégica y apoyo a la toma de decisiones dentro de una misma solución
digital.

La principal contribución consiste en transformar datos complejos en
información accesible, comprensible y accionable, favoreciendo mercados
más transparentes y reduciendo las asimetrías de información entre
productores, comercializadores y consumidores.

    \hypertarget{del-productor-al-consumidor}{%
\subsubsection{Del productor al
consumidor}\label{del-productor-al-consumidor}}

Una de las aportaciones más relevantes del proyecto consiste en extender
el análisis más allá de la producción agrícola.

Mientras que el modelo SARIMAX del maíz permite apoyar decisiones de
productores, cooperativas y comercializadores, el modelo SARIMAX de la
tortilla traslada el análisis hacia el consumidor final.

Esta extensión permite comprender cómo los cambios observados en el
mercado agrícola terminan reflejándose en uno de los alimentos más
representativos de la dieta mexicana.

Desde la perspectiva de Mercadotecnia Digital, esta ampliación del
alcance incrementa el valor, ya que permite ofrecer información
relevante tanto a actores productivos como a consumidores, reduciendo
asimetrías de información a lo largo de toda la cadena agroalimentaria.

    \begin{tcolorbox}[breakable, size=fbox, boxrule=1pt, pad at break*=1mm,colback=cellbackground, colframe=cellborder]
\prompt{In}{incolor}{324}{\boxspacing}
\begin{Verbatim}[commandchars=\\\{\}]
\PY{n}{arquitectura\PYZus{}producto} \PY{o}{=} \PY{n}{pd}\PY{o}{.}\PY{n}{DataFrame}\PY{p}{(}\PY{p}{\PYZob{}}

    \PY{l+s+s2}{\PYZdq{}}\PY{l+s+s2}{Módulo}\PY{l+s+s2}{\PYZdq{}}\PY{p}{:} \PY{p}{[}
        \PY{l+s+s2}{\PYZdq{}}\PY{l+s+s2}{Mercado del Maíz}\PY{l+s+s2}{\PYZdq{}}\PY{p}{,}
        \PY{l+s+s2}{\PYZdq{}}\PY{l+s+s2}{Mercado de la Tortilla}\PY{l+s+s2}{\PYZdq{}}\PY{p}{,}
        \PY{l+s+s2}{\PYZdq{}}\PY{l+s+s2}{Alertas Comerciales}\PY{l+s+s2}{\PYZdq{}}\PY{p}{,}
        \PY{l+s+s2}{\PYZdq{}}\PY{l+s+s2}{Planeación de Compras}\PY{l+s+s2}{\PYZdq{}}\PY{p}{,}
        \PY{l+s+s2}{\PYZdq{}}\PY{l+s+s2}{Panel de Indicadores}\PY{l+s+s2}{\PYZdq{}}
    \PY{p}{]}\PY{p}{,}

    \PY{l+s+s2}{\PYZdq{}}\PY{l+s+s2}{Funcionalidad}\PY{l+s+s2}{\PYZdq{}}\PY{p}{:} \PY{p}{[}
        \PY{l+s+s2}{\PYZdq{}}\PY{l+s+s2}{Pronóstico SARIMAX del maíz}\PY{l+s+s2}{\PYZdq{}}\PY{p}{,}
        \PY{l+s+s2}{\PYZdq{}}\PY{l+s+s2}{Pronóstico SARIMAX de tortilla}\PY{l+s+s2}{\PYZdq{}}\PY{p}{,}
        \PY{l+s+s2}{\PYZdq{}}\PY{l+s+s2}{Señales de oportunidad}\PY{l+s+s2}{\PYZdq{}}\PY{p}{,}
        \PY{l+s+s2}{\PYZdq{}}\PY{l+s+s2}{Optimización de abastecimiento}\PY{l+s+s2}{\PYZdq{}}\PY{p}{,}
        \PY{l+s+s2}{\PYZdq{}}\PY{l+s+s2}{KPIs y métricas estratégicas}\PY{l+s+s2}{\PYZdq{}}
    \PY{p}{]}\PY{p}{,}

    \PY{l+s+s2}{\PYZdq{}}\PY{l+s+s2}{Usuario principal}\PY{l+s+s2}{\PYZdq{}}\PY{p}{:} \PY{p}{[}
        \PY{l+s+s2}{\PYZdq{}}\PY{l+s+s2}{Productores}\PY{l+s+s2}{\PYZdq{}}\PY{p}{,}
        \PY{l+s+s2}{\PYZdq{}}\PY{l+s+s2}{Consumidores}\PY{l+s+s2}{\PYZdq{}}\PY{p}{,}
        \PY{l+s+s2}{\PYZdq{}}\PY{l+s+s2}{Comercializadores}\PY{l+s+s2}{\PYZdq{}}\PY{p}{,}
        \PY{l+s+s2}{\PYZdq{}}\PY{l+s+s2}{Cooperativas}\PY{l+s+s2}{\PYZdq{}}\PY{p}{,}
        \PY{l+s+s2}{\PYZdq{}}\PY{l+s+s2}{Todos los usuarios}\PY{l+s+s2}{\PYZdq{}}
    \PY{p}{]}
\PY{p}{\PYZcb{}}\PY{p}{)}

\PY{n}{arquitectura\PYZus{}producto}
\end{Verbatim}
\end{tcolorbox}

            \begin{tcolorbox}[breakable, size=fbox, boxrule=.5pt, pad at break*=1mm, opacityfill=0]
\prompt{Out}{outcolor}{324}{\boxspacing}
\begin{Verbatim}[commandchars=\\\{\}]
                   Módulo                   Funcionalidad   Usuario principal
0        Mercado del Maíz     Pronóstico SARIMAX del maíz         Productores
1  Mercado de la Tortilla  Pronóstico SARIMAX de tortilla        Consumidores
2     Alertas Comerciales          Señales de oportunidad   Comercializadores
3   Planeación de Compras  Optimización de abastecimiento        Cooperativas
4    Panel de Indicadores    KPIs y métricas estratégicas  Todos los usuarios
\end{Verbatim}
\end{tcolorbox}
        
    \begin{tcolorbox}[breakable, size=fbox, boxrule=1pt, pad at break*=1mm,colback=cellbackground, colframe=cellborder]
\prompt{In}{incolor}{325}{\boxspacing}
\begin{Verbatim}[commandchars=\\\{\}]
\PY{k+kn}{import} \PY{n+nn}{networkx} \PY{k}{as} \PY{n+nn}{nx}
\PY{k+kn}{import} \PY{n+nn}{matplotlib}\PY{n+nn}{.}\PY{n+nn}{pyplot} \PY{k}{as} \PY{n+nn}{plt}

\PY{n}{G} \PY{o}{=} \PY{n}{nx}\PY{o}{.}\PY{n}{DiGraph}\PY{p}{(}\PY{p}{)}

\PY{n}{G}\PY{o}{.}\PY{n}{add\PYZus{}edges\PYZus{}from}\PY{p}{(}\PY{p}{[}

    \PY{p}{(}\PY{l+s+s2}{\PYZdq{}}\PY{l+s+s2}{Datos oficiales}\PY{l+s+s2}{\PYZdq{}}\PY{p}{,} \PY{l+s+s2}{\PYZdq{}}\PY{l+s+s2}{SARIMAX Maíz}\PY{l+s+s2}{\PYZdq{}}\PY{p}{)}\PY{p}{,}

    \PY{p}{(}\PY{l+s+s2}{\PYZdq{}}\PY{l+s+s2}{SARIMAX Maíz}\PY{l+s+s2}{\PYZdq{}}\PY{p}{,} \PY{l+s+s2}{\PYZdq{}}\PY{l+s+s2}{Pronóstico Maíz}\PY{l+s+s2}{\PYZdq{}}\PY{p}{)}\PY{p}{,}

    \PY{p}{(}\PY{l+s+s2}{\PYZdq{}}\PY{l+s+s2}{Pronóstico Maíz}\PY{l+s+s2}{\PYZdq{}}\PY{p}{,} \PY{l+s+s2}{\PYZdq{}}\PY{l+s+s2}{Optimización}\PY{l+s+s2}{\PYZdq{}}\PY{p}{)}\PY{p}{,}

    \PY{p}{(}\PY{l+s+s2}{\PYZdq{}}\PY{l+s+s2}{Pronóstico Maíz}\PY{l+s+s2}{\PYZdq{}}\PY{p}{,} \PY{l+s+s2}{\PYZdq{}}\PY{l+s+s2}{SARIMAX Tortilla}\PY{l+s+s2}{\PYZdq{}}\PY{p}{)}\PY{p}{,}

    \PY{p}{(}\PY{l+s+s2}{\PYZdq{}}\PY{l+s+s2}{SARIMAX Tortilla}\PY{l+s+s2}{\PYZdq{}}\PY{p}{,} \PY{l+s+s2}{\PYZdq{}}\PY{l+s+s2}{Pronóstico Tortilla}\PY{l+s+s2}{\PYZdq{}}\PY{p}{)}\PY{p}{,}

    \PY{p}{(}\PY{l+s+s2}{\PYZdq{}}\PY{l+s+s2}{Optimización}\PY{l+s+s2}{\PYZdq{}}\PY{p}{,} \PY{l+s+s2}{\PYZdq{}}\PY{l+s+s2}{Productores}\PY{l+s+s2}{\PYZdq{}}\PY{p}{)}\PY{p}{,}
    \PY{p}{(}\PY{l+s+s2}{\PYZdq{}}\PY{l+s+s2}{Optimización}\PY{l+s+s2}{\PYZdq{}}\PY{p}{,} \PY{l+s+s2}{\PYZdq{}}\PY{l+s+s2}{Cooperativas}\PY{l+s+s2}{\PYZdq{}}\PY{p}{)}\PY{p}{,}
    \PY{p}{(}\PY{l+s+s2}{\PYZdq{}}\PY{l+s+s2}{Optimización}\PY{l+s+s2}{\PYZdq{}}\PY{p}{,} \PY{l+s+s2}{\PYZdq{}}\PY{l+s+s2}{Comercializadores}\PY{l+s+s2}{\PYZdq{}}\PY{p}{)}\PY{p}{,}

    \PY{p}{(}\PY{l+s+s2}{\PYZdq{}}\PY{l+s+s2}{Pronóstico Tortilla}\PY{l+s+s2}{\PYZdq{}}\PY{p}{,} \PY{l+s+s2}{\PYZdq{}}\PY{l+s+s2}{Consumidores}\PY{l+s+s2}{\PYZdq{}}\PY{p}{)}\PY{p}{,}
    \PY{p}{(}\PY{l+s+s2}{\PYZdq{}}\PY{l+s+s2}{Pronóstico Tortilla}\PY{l+s+s2}{\PYZdq{}}\PY{p}{,} \PY{l+s+s2}{\PYZdq{}}\PY{l+s+s2}{Gobierno}\PY{l+s+s2}{\PYZdq{}}\PY{p}{)}\PY{p}{,}
    \PY{p}{(}\PY{l+s+s2}{\PYZdq{}}\PY{l+s+s2}{Pronóstico Tortilla}\PY{l+s+s2}{\PYZdq{}}\PY{p}{,} \PY{l+s+s2}{\PYZdq{}}\PY{l+s+s2}{Industria Tortillera}\PY{l+s+s2}{\PYZdq{}}\PY{p}{)}

\PY{p}{]}\PY{p}{)}

\PY{n}{plt}\PY{o}{.}\PY{n}{figure}\PY{p}{(}\PY{n}{figsize}\PY{o}{=}\PY{p}{(}\PY{l+m+mi}{14}\PY{p}{,}\PY{l+m+mi}{10}\PY{p}{)}\PY{p}{)}

\PY{n}{pos} \PY{o}{=} \PY{n}{nx}\PY{o}{.}\PY{n}{spring\PYZus{}layout}\PY{p}{(}\PY{n}{G}\PY{p}{,} \PY{n}{seed}\PY{o}{=}\PY{l+m+mi}{42}\PY{p}{)}

\PY{n}{nx}\PY{o}{.}\PY{n}{draw}\PY{p}{(}
    \PY{n}{G}\PY{p}{,}
    \PY{n}{pos}\PY{p}{,}
    \PY{n}{with\PYZus{}labels}\PY{o}{=}\PY{k+kc}{True}\PY{p}{,}
    \PY{n}{node\PYZus{}size}\PY{o}{=}\PY{l+m+mi}{5000}\PY{p}{,}
    \PY{n}{font\PYZus{}size}\PY{o}{=}\PY{l+m+mi}{8}
\PY{p}{)}

\PY{n}{plt}\PY{o}{.}\PY{n}{title}\PY{p}{(}
    \PY{l+s+s2}{\PYZdq{}}\PY{l+s+s2}{Generación de valor de la plataforma Mercado Justo}\PY{l+s+s2}{\PYZdq{}}
\PY{p}{)}

\PY{n}{plt}\PY{o}{.}\PY{n}{show}\PY{p}{(}\PY{p}{)}
\end{Verbatim}
\end{tcolorbox}

    \begin{center}
    \adjustimage{max size={0.9\linewidth}{0.9\paperheight}}{output_246_0.png}
    \end{center}
    { \hspace*{\fill} \\}
    
    \begin{tcolorbox}[breakable, size=fbox, boxrule=1pt, pad at break*=1mm,colback=cellbackground, colframe=cellborder]
\prompt{In}{incolor}{330}{\boxspacing}
\begin{Verbatim}[commandchars=\\\{\}]
\PY{n}{precio\PYZus{}actual} \PY{o}{=} \PY{n}{maiz\PYZus{}tortilla}\PY{p}{[}\PY{l+s+s2}{\PYZdq{}}\PY{l+s+s2}{Precio\PYZus{}maiz}\PY{l+s+s2}{\PYZdq{}}\PY{p}{]}\PY{o}{.}\PY{n}{iloc}\PY{p}{[}\PY{o}{\PYZhy{}}\PY{l+m+mi}{1}\PY{p}{]}

\PY{n}{precio\PYZus{}futuro} \PY{o}{=} \PY{n}{pronostico\PYZus{}io}\PY{p}{[}\PY{l+s+s2}{\PYZdq{}}\PY{l+s+s2}{Precio\PYZus{}Pronosticado}\PY{l+s+s2}{\PYZdq{}}\PY{p}{]}\PY{o}{.}\PY{n}{iloc}\PY{p}{[}\PY{o}{\PYZhy{}}\PY{l+m+mi}{1}\PY{p}{]}

\PY{n+nb}{print}\PY{p}{(}\PY{l+s+s2}{\PYZdq{}}\PY{l+s+s2}{Precio actual:}\PY{l+s+s2}{\PYZdq{}}\PY{p}{,} \PY{n}{precio\PYZus{}actual}\PY{p}{)}
\PY{n+nb}{print}\PY{p}{(}\PY{l+s+s2}{\PYZdq{}}\PY{l+s+s2}{Precio futuro:}\PY{l+s+s2}{\PYZdq{}}\PY{p}{,} \PY{n}{precio\PYZus{}futuro}\PY{p}{)}

\PY{n}{variacion} \PY{o}{=} \PY{p}{(}
    \PY{p}{(}\PY{n}{precio\PYZus{}futuro} \PY{o}{\PYZhy{}} \PY{n}{precio\PYZus{}actual}\PY{p}{)}
    \PY{o}{/} \PY{n}{precio\PYZus{}actual}
\PY{p}{)} \PY{o}{*} \PY{l+m+mi}{100}

\PY{n+nb}{print}\PY{p}{(}\PY{l+s+s2}{\PYZdq{}}\PY{l+s+s2}{Variación esperada:}\PY{l+s+s2}{\PYZdq{}}\PY{p}{,} \PY{n}{variacion}\PY{p}{)}
\end{Verbatim}
\end{tcolorbox}

    \begin{Verbatim}[commandchars=\\\{\}]
Precio actual: 9.82
Precio futuro: 9.546922715569542
Variación esperada: -2.780827743691024
    \end{Verbatim}

    \begin{tcolorbox}[breakable, size=fbox, boxrule=1pt, pad at break*=1mm,colback=cellbackground, colframe=cellborder]
\prompt{In}{incolor}{341}{\boxspacing}
\begin{Verbatim}[commandchars=\\\{\}]
\PY{k+kn}{import} \PY{n+nn}{matplotlib}\PY{n+nn}{.}\PY{n+nn}{pyplot} \PY{k}{as} \PY{n+nn}{plt}

\PY{c+c1}{\PYZsh{} =====================================================}
\PY{c+c1}{\PYZsh{} VARIABLES MAÍZ}
\PY{c+c1}{\PYZsh{} =====================================================}

\PY{n}{precio\PYZus{}actual\PYZus{}maiz} \PY{o}{=} \PY{n}{maiz\PYZus{}tortilla}\PY{p}{[}\PY{l+s+s2}{\PYZdq{}}\PY{l+s+s2}{Precio\PYZus{}maiz}\PY{l+s+s2}{\PYZdq{}}\PY{p}{]}\PY{o}{.}\PY{n}{iloc}\PY{p}{[}\PY{o}{\PYZhy{}}\PY{l+m+mi}{1}\PY{p}{]}

\PY{n}{precio\PYZus{}futuro\PYZus{}maiz} \PY{o}{=} \PY{p}{(}
    \PY{n}{pronostico\PYZus{}io}\PY{p}{[}\PY{l+s+s2}{\PYZdq{}}\PY{l+s+s2}{Precio\PYZus{}Pronosticado}\PY{l+s+s2}{\PYZdq{}}\PY{p}{]}
    \PY{o}{.}\PY{n}{iloc}\PY{p}{[}\PY{o}{\PYZhy{}}\PY{l+m+mi}{1}\PY{p}{]}
\PY{p}{)}

\PY{n}{variacion\PYZus{}maiz} \PY{o}{=} \PY{p}{(}
    \PY{p}{(}\PY{n}{precio\PYZus{}futuro\PYZus{}maiz} \PY{o}{\PYZhy{}} \PY{n}{precio\PYZus{}actual\PYZus{}maiz}\PY{p}{)}
    \PY{o}{/} \PY{n}{precio\PYZus{}actual\PYZus{}maiz}
\PY{p}{)} \PY{o}{*} \PY{l+m+mi}{100}

\PY{n}{mejor\PYZus{}mes} \PY{o}{=} \PY{p}{(}
    \PY{n}{pronostico\PYZus{}io}\PY{o}{.}\PY{n}{loc}\PY{p}{[}
        \PY{n}{pronostico\PYZus{}io}\PY{p}{[}\PY{l+s+s2}{\PYZdq{}}\PY{l+s+s2}{Precio\PYZus{}Pronosticado}\PY{l+s+s2}{\PYZdq{}}\PY{p}{]}\PY{o}{.}\PY{n}{idxmin}\PY{p}{(}\PY{p}{)}\PY{p}{,}
        \PY{l+s+s2}{\PYZdq{}}\PY{l+s+s2}{Mes}\PY{l+s+s2}{\PYZdq{}}
    \PY{p}{]}
\PY{p}{)}

\PY{n}{mejor\PYZus{}precio} \PY{o}{=} \PY{p}{(}
    \PY{n}{pronostico\PYZus{}io}\PY{p}{[}\PY{l+s+s2}{\PYZdq{}}\PY{l+s+s2}{Precio\PYZus{}Pronosticado}\PY{l+s+s2}{\PYZdq{}}\PY{p}{]}
    \PY{o}{.}\PY{n}{min}\PY{p}{(}\PY{p}{)}
\PY{p}{)}

\PY{c+c1}{\PYZsh{} =====================================================}
\PY{c+c1}{\PYZsh{} VARIABLES TORTILLA}
\PY{c+c1}{\PYZsh{} =====================================================}

\PY{n}{precio\PYZus{}actual\PYZus{}tortilla} \PY{o}{=} \PY{l+m+mf}{25.00}

\PY{n}{precio\PYZus{}futuro\PYZus{}tortilla} \PY{o}{=} \PY{p}{(}
    \PY{n}{tabla\PYZus{}pronostico\PYZus{}tortilla}\PY{p}{[}
        \PY{l+s+s2}{\PYZdq{}}\PY{l+s+s2}{Precio Pronosticado Tortilla (MXN/Kg)}\PY{l+s+s2}{\PYZdq{}}
    \PY{p}{]}
    \PY{o}{.}\PY{n}{iloc}\PY{p}{[}\PY{o}{\PYZhy{}}\PY{l+m+mi}{1}\PY{p}{]}
\PY{p}{)}

\PY{n}{variacion\PYZus{}tortilla} \PY{o}{=} \PY{p}{(}
    \PY{p}{(}\PY{n}{precio\PYZus{}futuro\PYZus{}tortilla} \PY{o}{\PYZhy{}} \PY{n}{precio\PYZus{}actual\PYZus{}tortilla}\PY{p}{)}
    \PY{o}{/} \PY{n}{precio\PYZus{}actual\PYZus{}tortilla}
\PY{p}{)} \PY{o}{*} \PY{l+m+mi}{100}

\PY{c+c1}{\PYZsh{} =====================================================}
\PY{c+c1}{\PYZsh{} FIGURA PRINCIPAL}
\PY{c+c1}{\PYZsh{} =====================================================}

\PY{n}{fig} \PY{o}{=} \PY{n}{plt}\PY{o}{.}\PY{n}{figure}\PY{p}{(}\PY{n}{figsize}\PY{o}{=}\PY{p}{(}\PY{l+m+mi}{18}\PY{p}{,}\PY{l+m+mi}{12}\PY{p}{)}\PY{p}{)}

\PY{n}{fig}\PY{o}{.}\PY{n}{suptitle}\PY{p}{(}
    \PY{l+s+s2}{\PYZdq{}}\PY{l+s+s2}{MERCADO JUSTO}\PY{l+s+se}{\PYZbs{}n}\PY{l+s+s2}{Promovemos la toma de decisiones basadas en datos}\PY{l+s+s2}{\PYZdq{}}\PY{p}{,}
    \PY{n}{fontsize}\PY{o}{=}\PY{l+m+mi}{22}\PY{p}{,}
    \PY{n}{fontweight}\PY{o}{=}\PY{l+s+s2}{\PYZdq{}}\PY{l+s+s2}{bold}\PY{l+s+s2}{\PYZdq{}}
\PY{p}{)}

\PY{c+c1}{\PYZsh{} =====================================================}
\PY{c+c1}{\PYZsh{} TARJETA 1}
\PY{c+c1}{\PYZsh{} =====================================================}

\PY{n}{ax1} \PY{o}{=} \PY{n}{plt}\PY{o}{.}\PY{n}{subplot2grid}\PY{p}{(}\PY{p}{(}\PY{l+m+mi}{4}\PY{p}{,}\PY{l+m+mi}{4}\PY{p}{)}\PY{p}{,}\PY{p}{(}\PY{l+m+mi}{0}\PY{p}{,}\PY{l+m+mi}{0}\PY{p}{)}\PY{p}{)}
\PY{n}{ax1}\PY{o}{.}\PY{n}{axis}\PY{p}{(}\PY{l+s+s2}{\PYZdq{}}\PY{l+s+s2}{off}\PY{l+s+s2}{\PYZdq{}}\PY{p}{)}

\PY{n}{ax1}\PY{o}{.}\PY{n}{text}\PY{p}{(}
    \PY{l+m+mf}{0.5}\PY{p}{,}
    \PY{l+m+mf}{0.60}\PY{p}{,}
    \PY{l+s+sa}{f}\PY{l+s+s2}{\PYZdq{}}\PY{l+s+s2}{\PYZdl{}}\PY{l+s+si}{\PYZob{}}\PY{n}{precio\PYZus{}actual\PYZus{}maiz}\PY{l+s+si}{:}\PY{l+s+s2}{.2f}\PY{l+s+si}{\PYZcb{}}\PY{l+s+s2}{\PYZdq{}}\PY{p}{,}
    \PY{n}{ha}\PY{o}{=}\PY{l+s+s2}{\PYZdq{}}\PY{l+s+s2}{center}\PY{l+s+s2}{\PYZdq{}}\PY{p}{,}
    \PY{n}{fontsize}\PY{o}{=}\PY{l+m+mi}{24}\PY{p}{,}
    \PY{n}{fontweight}\PY{o}{=}\PY{l+s+s2}{\PYZdq{}}\PY{l+s+s2}{bold}\PY{l+s+s2}{\PYZdq{}}
\PY{p}{)}

\PY{n}{ax1}\PY{o}{.}\PY{n}{text}\PY{p}{(}
    \PY{l+m+mf}{0.5}\PY{p}{,}
    \PY{l+m+mf}{0.25}\PY{p}{,}
    \PY{l+s+s2}{\PYZdq{}}\PY{l+s+s2}{Precio actual del maíz}\PY{l+s+s2}{\PYZdq{}}\PY{p}{,}
    \PY{n}{ha}\PY{o}{=}\PY{l+s+s2}{\PYZdq{}}\PY{l+s+s2}{center}\PY{l+s+s2}{\PYZdq{}}\PY{p}{,}
    \PY{n}{fontsize}\PY{o}{=}\PY{l+m+mi}{11}
\PY{p}{)}

\PY{c+c1}{\PYZsh{} =====================================================}
\PY{c+c1}{\PYZsh{} TARJETA 2}
\PY{c+c1}{\PYZsh{} =====================================================}

\PY{n}{ax2} \PY{o}{=} \PY{n}{plt}\PY{o}{.}\PY{n}{subplot2grid}\PY{p}{(}\PY{p}{(}\PY{l+m+mi}{4}\PY{p}{,}\PY{l+m+mi}{4}\PY{p}{)}\PY{p}{,}\PY{p}{(}\PY{l+m+mi}{0}\PY{p}{,}\PY{l+m+mi}{1}\PY{p}{)}\PY{p}{)}
\PY{n}{ax2}\PY{o}{.}\PY{n}{axis}\PY{p}{(}\PY{l+s+s2}{\PYZdq{}}\PY{l+s+s2}{off}\PY{l+s+s2}{\PYZdq{}}\PY{p}{)}

\PY{n}{ax2}\PY{o}{.}\PY{n}{text}\PY{p}{(}
    \PY{l+m+mf}{0.5}\PY{p}{,}
    \PY{l+m+mf}{0.60}\PY{p}{,}
    \PY{l+s+sa}{f}\PY{l+s+s2}{\PYZdq{}}\PY{l+s+s2}{\PYZdl{}}\PY{l+s+si}{\PYZob{}}\PY{n}{precio\PYZus{}actual\PYZus{}tortilla}\PY{l+s+si}{:}\PY{l+s+s2}{.2f}\PY{l+s+si}{\PYZcb{}}\PY{l+s+s2}{\PYZdq{}}\PY{p}{,}
    \PY{n}{ha}\PY{o}{=}\PY{l+s+s2}{\PYZdq{}}\PY{l+s+s2}{center}\PY{l+s+s2}{\PYZdq{}}\PY{p}{,}
    \PY{n}{fontsize}\PY{o}{=}\PY{l+m+mi}{24}\PY{p}{,}
    \PY{n}{fontweight}\PY{o}{=}\PY{l+s+s2}{\PYZdq{}}\PY{l+s+s2}{bold}\PY{l+s+s2}{\PYZdq{}}
\PY{p}{)}

\PY{n}{ax2}\PY{o}{.}\PY{n}{text}\PY{p}{(}
    \PY{l+m+mf}{0.5}\PY{p}{,}
    \PY{l+m+mf}{0.25}\PY{p}{,}
    \PY{l+s+s2}{\PYZdq{}}\PY{l+s+s2}{Precio actual de la tortilla}\PY{l+s+s2}{\PYZdq{}}\PY{p}{,}
    \PY{n}{ha}\PY{o}{=}\PY{l+s+s2}{\PYZdq{}}\PY{l+s+s2}{center}\PY{l+s+s2}{\PYZdq{}}\PY{p}{,}
    \PY{n}{fontsize}\PY{o}{=}\PY{l+m+mi}{11}
\PY{p}{)}

\PY{c+c1}{\PYZsh{} =====================================================}
\PY{c+c1}{\PYZsh{} TARJETA 3}
\PY{c+c1}{\PYZsh{} =====================================================}

\PY{n}{ax3} \PY{o}{=} \PY{n}{plt}\PY{o}{.}\PY{n}{subplot2grid}\PY{p}{(}\PY{p}{(}\PY{l+m+mi}{4}\PY{p}{,}\PY{l+m+mi}{4}\PY{p}{)}\PY{p}{,}\PY{p}{(}\PY{l+m+mi}{0}\PY{p}{,}\PY{l+m+mi}{2}\PY{p}{)}\PY{p}{)}
\PY{n}{ax3}\PY{o}{.}\PY{n}{axis}\PY{p}{(}\PY{l+s+s2}{\PYZdq{}}\PY{l+s+s2}{off}\PY{l+s+s2}{\PYZdq{}}\PY{p}{)}

\PY{n}{ax3}\PY{o}{.}\PY{n}{text}\PY{p}{(}
    \PY{l+m+mf}{0.5}\PY{p}{,}
    \PY{l+m+mf}{0.60}\PY{p}{,}
    \PY{l+s+sa}{f}\PY{l+s+s2}{\PYZdq{}}\PY{l+s+s2}{Mes }\PY{l+s+si}{\PYZob{}}\PY{n+nb}{int}\PY{p}{(}\PY{n}{mejor\PYZus{}mes}\PY{p}{)}\PY{l+s+si}{\PYZcb{}}\PY{l+s+s2}{\PYZdq{}}\PY{p}{,}
    \PY{n}{ha}\PY{o}{=}\PY{l+s+s2}{\PYZdq{}}\PY{l+s+s2}{center}\PY{l+s+s2}{\PYZdq{}}\PY{p}{,}
    \PY{n}{fontsize}\PY{o}{=}\PY{l+m+mi}{22}\PY{p}{,}
    \PY{n}{fontweight}\PY{o}{=}\PY{l+s+s2}{\PYZdq{}}\PY{l+s+s2}{bold}\PY{l+s+s2}{\PYZdq{}}
\PY{p}{)}

\PY{n}{ax3}\PY{o}{.}\PY{n}{text}\PY{p}{(}
    \PY{l+m+mf}{0.5}\PY{p}{,}
    \PY{l+m+mf}{0.25}\PY{p}{,}
    \PY{l+s+s2}{\PYZdq{}}\PY{l+s+s2}{Mejor momento para comprar maíz}\PY{l+s+s2}{\PYZdq{}}\PY{p}{,}
    \PY{n}{ha}\PY{o}{=}\PY{l+s+s2}{\PYZdq{}}\PY{l+s+s2}{center}\PY{l+s+s2}{\PYZdq{}}\PY{p}{,}
    \PY{n}{fontsize}\PY{o}{=}\PY{l+m+mi}{11}
\PY{p}{)}

\PY{c+c1}{\PYZsh{} =====================================================}
\PY{c+c1}{\PYZsh{} TARJETA 4}
\PY{c+c1}{\PYZsh{} =====================================================}

\PY{n}{ax4} \PY{o}{=} \PY{n}{plt}\PY{o}{.}\PY{n}{subplot2grid}\PY{p}{(}\PY{p}{(}\PY{l+m+mi}{4}\PY{p}{,}\PY{l+m+mi}{4}\PY{p}{)}\PY{p}{,}\PY{p}{(}\PY{l+m+mi}{0}\PY{p}{,}\PY{l+m+mi}{3}\PY{p}{)}\PY{p}{)}
\PY{n}{ax4}\PY{o}{.}\PY{n}{axis}\PY{p}{(}\PY{l+s+s2}{\PYZdq{}}\PY{l+s+s2}{off}\PY{l+s+s2}{\PYZdq{}}\PY{p}{)}

\PY{n}{ax4}\PY{o}{.}\PY{n}{text}\PY{p}{(}
    \PY{l+m+mf}{0.5}\PY{p}{,}
    \PY{l+m+mf}{0.60}\PY{p}{,}
    \PY{l+s+s2}{\PYZdq{}}\PY{l+s+s2}{Estable}\PY{l+s+s2}{\PYZdq{}}\PY{p}{,}
    \PY{n}{ha}\PY{o}{=}\PY{l+s+s2}{\PYZdq{}}\PY{l+s+s2}{center}\PY{l+s+s2}{\PYZdq{}}\PY{p}{,}
    \PY{n}{fontsize}\PY{o}{=}\PY{l+m+mi}{22}\PY{p}{,}
    \PY{n}{fontweight}\PY{o}{=}\PY{l+s+s2}{\PYZdq{}}\PY{l+s+s2}{bold}\PY{l+s+s2}{\PYZdq{}}
\PY{p}{)}

\PY{n}{ax4}\PY{o}{.}\PY{n}{text}\PY{p}{(}
    \PY{l+m+mf}{0.5}\PY{p}{,}
    \PY{l+m+mf}{0.25}\PY{p}{,}
    \PY{l+s+s2}{\PYZdq{}}\PY{l+s+s2}{Perspectiva de la tortilla}\PY{l+s+s2}{\PYZdq{}}\PY{p}{,}
    \PY{n}{ha}\PY{o}{=}\PY{l+s+s2}{\PYZdq{}}\PY{l+s+s2}{center}\PY{l+s+s2}{\PYZdq{}}\PY{p}{,}
    \PY{n}{fontsize}\PY{o}{=}\PY{l+m+mi}{11}
\PY{p}{)}

\PY{c+c1}{\PYZsh{} =====================================================}
\PY{c+c1}{\PYZsh{} PRONÓSTICO MAÍZ}
\PY{c+c1}{\PYZsh{} =====================================================}

\PY{n}{ax5} \PY{o}{=} \PY{n}{plt}\PY{o}{.}\PY{n}{subplot2grid}\PY{p}{(}\PY{p}{(}\PY{l+m+mi}{4}\PY{p}{,}\PY{l+m+mi}{4}\PY{p}{)}\PY{p}{,}\PY{p}{(}\PY{l+m+mi}{1}\PY{p}{,}\PY{l+m+mi}{0}\PY{p}{)}\PY{p}{,} \PY{n}{colspan}\PY{o}{=}\PY{l+m+mi}{2}\PY{p}{)}

\PY{n}{ax5}\PY{o}{.}\PY{n}{plot}\PY{p}{(}
    \PY{n}{pronostico\PYZus{}io}\PY{p}{[}\PY{l+s+s2}{\PYZdq{}}\PY{l+s+s2}{Mes}\PY{l+s+s2}{\PYZdq{}}\PY{p}{]}\PY{p}{,}
    \PY{n}{pronostico\PYZus{}io}\PY{p}{[}\PY{l+s+s2}{\PYZdq{}}\PY{l+s+s2}{Precio\PYZus{}Pronosticado}\PY{l+s+s2}{\PYZdq{}}\PY{p}{]}\PY{p}{,}
    \PY{n}{marker}\PY{o}{=}\PY{l+s+s2}{\PYZdq{}}\PY{l+s+s2}{o}\PY{l+s+s2}{\PYZdq{}}\PY{p}{,}
    \PY{n}{linewidth}\PY{o}{=}\PY{l+m+mi}{2}
\PY{p}{)}

\PY{n}{ax5}\PY{o}{.}\PY{n}{set\PYZus{}title}\PY{p}{(}
    \PY{l+s+s2}{\PYZdq{}}\PY{l+s+s2}{¿Qué podría pasar con el precio del maíz?}\PY{l+s+s2}{\PYZdq{}}
\PY{p}{)}

\PY{n}{ax5}\PY{o}{.}\PY{n}{set\PYZus{}xlabel}\PY{p}{(}\PY{l+s+s2}{\PYZdq{}}\PY{l+s+s2}{Mes}\PY{l+s+s2}{\PYZdq{}}\PY{p}{)}

\PY{n}{ax5}\PY{o}{.}\PY{n}{set\PYZus{}ylabel}\PY{p}{(}\PY{l+s+s2}{\PYZdq{}}\PY{l+s+s2}{MXN/kg}\PY{l+s+s2}{\PYZdq{}}\PY{p}{)}

\PY{n}{ax5}\PY{o}{.}\PY{n}{grid}\PY{p}{(}\PY{k+kc}{True}\PY{p}{)}

\PY{c+c1}{\PYZsh{} =====================================================}
\PY{c+c1}{\PYZsh{} PRONÓSTICO TORTILLA}
\PY{c+c1}{\PYZsh{} =====================================================}

\PY{n}{ax6} \PY{o}{=} \PY{n}{plt}\PY{o}{.}\PY{n}{subplot2grid}\PY{p}{(}\PY{p}{(}\PY{l+m+mi}{4}\PY{p}{,}\PY{l+m+mi}{4}\PY{p}{)}\PY{p}{,}\PY{p}{(}\PY{l+m+mi}{1}\PY{p}{,}\PY{l+m+mi}{2}\PY{p}{)}\PY{p}{,} \PY{n}{colspan}\PY{o}{=}\PY{l+m+mi}{2}\PY{p}{)}

\PY{n}{ax6}\PY{o}{.}\PY{n}{plot}\PY{p}{(}
    \PY{n}{tabla\PYZus{}pronostico\PYZus{}tortilla}\PY{o}{.}\PY{n}{index}\PY{p}{,}
    \PY{n}{tabla\PYZus{}pronostico\PYZus{}tortilla}\PY{p}{[}
        \PY{l+s+s2}{\PYZdq{}}\PY{l+s+s2}{Precio Pronosticado Tortilla (MXN/Kg)}\PY{l+s+s2}{\PYZdq{}}
    \PY{p}{]}\PY{p}{,}
    \PY{n}{marker}\PY{o}{=}\PY{l+s+s2}{\PYZdq{}}\PY{l+s+s2}{o}\PY{l+s+s2}{\PYZdq{}}\PY{p}{,}
    \PY{n}{linewidth}\PY{o}{=}\PY{l+m+mi}{2}
\PY{p}{)}

\PY{n}{ax6}\PY{o}{.}\PY{n}{set\PYZus{}title}\PY{p}{(}
    \PY{l+s+s2}{\PYZdq{}}\PY{l+s+s2}{¿Qué podría pasar con el precio de la tortilla?}\PY{l+s+s2}{\PYZdq{}}
\PY{p}{)}

\PY{n}{ax6}\PY{o}{.}\PY{n}{set\PYZus{}ylabel}\PY{p}{(}\PY{l+s+s2}{\PYZdq{}}\PY{l+s+s2}{MXN/kg}\PY{l+s+s2}{\PYZdq{}}\PY{p}{)}

\PY{n}{ax6}\PY{o}{.}\PY{n}{grid}\PY{p}{(}\PY{k+kc}{True}\PY{p}{)}

\PY{c+c1}{\PYZsh{} =====================================================}
\PY{c+c1}{\PYZsh{} RECOMENDACIÓN MAÍZ}
\PY{c+c1}{\PYZsh{} =====================================================}

\PY{n}{ax7} \PY{o}{=} \PY{n}{plt}\PY{o}{.}\PY{n}{subplot2grid}\PY{p}{(}\PY{p}{(}\PY{l+m+mi}{4}\PY{p}{,}\PY{l+m+mi}{4}\PY{p}{)}\PY{p}{,}\PY{p}{(}\PY{l+m+mi}{2}\PY{p}{,}\PY{l+m+mi}{0}\PY{p}{)}\PY{p}{,} \PY{n}{colspan}\PY{o}{=}\PY{l+m+mi}{2}\PY{p}{)}

\PY{n}{ax7}\PY{o}{.}\PY{n}{axis}\PY{p}{(}\PY{l+s+s2}{\PYZdq{}}\PY{l+s+s2}{off}\PY{l+s+s2}{\PYZdq{}}\PY{p}{)}

\PY{n}{texto\PYZus{}maiz} \PY{o}{=} \PY{l+s+sa}{f}\PY{l+s+s2}{\PYZdq{}\PYZdq{}\PYZdq{}}
\PY{l+s+s2}{¿QUÉ NOS DICEN LOS DATOS DEL MAÍZ?}

\PY{l+s+s2}{No se esperan cambios importantes}
\PY{l+s+s2}{en el precio del maíz durante los próximos meses.}

\PY{l+s+s2}{Si puede elegir cuándo comprar,}
\PY{l+s+s2}{los últimos meses muestran los precios}
\PY{l+s+s2}{más favorables del periodo analizado.}

\PY{l+s+s2}{El mes }\PY{l+s+si}{\PYZob{}}\PY{n+nb}{int}\PY{p}{(}\PY{n}{mejor\PYZus{}mes}\PY{p}{)}\PY{l+s+si}{\PYZcb{}}
\PY{l+s+s2}{presenta la estimación más baja.}
\PY{l+s+s2}{\PYZdq{}\PYZdq{}\PYZdq{}}

\PY{n}{ax7}\PY{o}{.}\PY{n}{text}\PY{p}{(}
    \PY{l+m+mf}{0.5}\PY{p}{,}
    \PY{l+m+mf}{0.5}\PY{p}{,}
    \PY{n}{texto\PYZus{}maiz}\PY{p}{,}
    \PY{n}{ha}\PY{o}{=}\PY{l+s+s2}{\PYZdq{}}\PY{l+s+s2}{center}\PY{l+s+s2}{\PYZdq{}}\PY{p}{,}
    \PY{n}{va}\PY{o}{=}\PY{l+s+s2}{\PYZdq{}}\PY{l+s+s2}{center}\PY{l+s+s2}{\PYZdq{}}\PY{p}{,}
    \PY{n}{fontsize}\PY{o}{=}\PY{l+m+mi}{12}\PY{p}{,}
    \PY{n}{bbox}\PY{o}{=}\PY{n+nb}{dict}\PY{p}{(}
        \PY{n}{boxstyle}\PY{o}{=}\PY{l+s+s2}{\PYZdq{}}\PY{l+s+s2}{round}\PY{l+s+s2}{\PYZdq{}}\PY{p}{,}
        \PY{n}{alpha}\PY{o}{=}\PY{l+m+mf}{0.15}
    \PY{p}{)}
\PY{p}{)}
\PY{c+c1}{\PYZsh{} =====================================================}
\PY{c+c1}{\PYZsh{} RECOMENDACIÓN TORTILLA}
\PY{c+c1}{\PYZsh{} =====================================================}

\PY{n}{ax8} \PY{o}{=} \PY{n}{plt}\PY{o}{.}\PY{n}{subplot2grid}\PY{p}{(}\PY{p}{(}\PY{l+m+mi}{4}\PY{p}{,}\PY{l+m+mi}{4}\PY{p}{)}\PY{p}{,}\PY{p}{(}\PY{l+m+mi}{2}\PY{p}{,}\PY{l+m+mi}{2}\PY{p}{)}\PY{p}{,} \PY{n}{colspan}\PY{o}{=}\PY{l+m+mi}{2}\PY{p}{)}

\PY{n}{ax8}\PY{o}{.}\PY{n}{axis}\PY{p}{(}\PY{l+s+s2}{\PYZdq{}}\PY{l+s+s2}{off}\PY{l+s+s2}{\PYZdq{}}\PY{p}{)}

\PY{n}{texto\PYZus{}tortilla} \PY{o}{=} \PY{l+s+s2}{\PYZdq{}\PYZdq{}\PYZdq{}}
\PY{l+s+s2}{¿QUÉ ESPERAR DEL PRECIO DE LA TORTILLA?}

\PY{l+s+s2}{Los datos indican que el precio}
\PY{l+s+s2}{de la tortilla podría mantenerse estable.}

\PY{l+s+s2}{No se observan aumentos fuertes}
\PY{l+s+s2}{en el escenario más probable.}

\PY{l+s+s2}{Por ahora no existen señales}
\PY{l+s+s2}{de cambios importantes para quienes}
\PY{l+s+s2}{compran tortilla regularmente.}
\PY{l+s+s2}{\PYZdq{}\PYZdq{}\PYZdq{}}

\PY{n}{ax8}\PY{o}{.}\PY{n}{text}\PY{p}{(}
    \PY{l+m+mf}{0.5}\PY{p}{,}
    \PY{l+m+mf}{0.5}\PY{p}{,}
    \PY{n}{texto\PYZus{}tortilla}\PY{p}{,}
    \PY{n}{ha}\PY{o}{=}\PY{l+s+s2}{\PYZdq{}}\PY{l+s+s2}{center}\PY{l+s+s2}{\PYZdq{}}\PY{p}{,}
    \PY{n}{va}\PY{o}{=}\PY{l+s+s2}{\PYZdq{}}\PY{l+s+s2}{center}\PY{l+s+s2}{\PYZdq{}}\PY{p}{,}
    \PY{n}{fontsize}\PY{o}{=}\PY{l+m+mi}{12}\PY{p}{,}
    \PY{n}{bbox}\PY{o}{=}\PY{n+nb}{dict}\PY{p}{(}
        \PY{n}{boxstyle}\PY{o}{=}\PY{l+s+s2}{\PYZdq{}}\PY{l+s+s2}{round}\PY{l+s+s2}{\PYZdq{}}\PY{p}{,}
        \PY{n}{alpha}\PY{o}{=}\PY{l+m+mf}{0.15}
    \PY{p}{)}
\PY{p}{)}

\PY{n}{plt}\PY{o}{.}\PY{n}{tight\PYZus{}layout}\PY{p}{(}\PY{p}{)}

\PY{n}{plt}\PY{o}{.}\PY{n}{show}\PY{p}{(}\PY{p}{)}
\end{Verbatim}
\end{tcolorbox}

    \begin{center}
    \adjustimage{max size={0.9\linewidth}{0.9\paperheight}}{output_248_0.png}
    \end{center}
    { \hspace*{\fill} \\}
    
    \begin{tcolorbox}[breakable, size=fbox, boxrule=1pt, pad at break*=1mm,colback=cellbackground, colframe=cellborder]
\prompt{In}{incolor}{ }{\boxspacing}
\begin{Verbatim}[commandchars=\\\{\}]

\end{Verbatim}
\end{tcolorbox}

    \hypertarget{modelo-xgboost-para-integraciuxf3n-de-aprendizaje-de-muxe1quina}{%
\section{42. Modelo XGBoost para Integración de Aprendizaje de
Máquina}\label{modelo-xgboost-para-integraciuxf3n-de-aprendizaje-de-muxe1quina}}

\hypertarget{objetivo}{%
\subsection{Objetivo}\label{objetivo}}

Con el propósito de integrar la UCA de Aprendizaje de Máquina, se
construye un modelo XGBoost para predecir el precio nacional del maíz
blanco.

A diferencia de SARIMAX, que modela explícitamente la estructura
temporal, XGBoost permite capturar relaciones no lineales entre
variables económicas, climáticas, logísticas y rezagos históricos del
precio.

Este modelo se utiliza como comparación complementaria frente al modelo
SARIMAX reducido.

    \begin{tcolorbox}[breakable, size=fbox, boxrule=1pt, pad at break*=1mm,colback=cellbackground, colframe=cellborder]
\prompt{In}{incolor}{182}{\boxspacing}
\begin{Verbatim}[commandchars=\\\{\}]
\PY{n}{df\PYZus{}ml} \PY{o}{=} \PY{n}{pdf\PYZus{}completo}\PY{o}{.}\PY{n}{copy}\PY{p}{(}\PY{p}{)}

\PY{n}{df\PYZus{}ml}\PY{p}{[}\PY{l+s+s2}{\PYZdq{}}\PY{l+s+s2}{fecha}\PY{l+s+s2}{\PYZdq{}}\PY{p}{]} \PY{o}{=} \PY{n}{pd}\PY{o}{.}\PY{n}{to\PYZus{}datetime}\PY{p}{(}\PY{n}{df\PYZus{}ml}\PY{p}{[}\PY{l+s+s2}{\PYZdq{}}\PY{l+s+s2}{fecha}\PY{l+s+s2}{\PYZdq{}}\PY{p}{]}\PY{p}{)}

\PY{n}{df\PYZus{}ml} \PY{o}{=} \PY{n}{df\PYZus{}ml}\PY{o}{.}\PY{n}{sort\PYZus{}values}\PY{p}{(}\PY{l+s+s2}{\PYZdq{}}\PY{l+s+s2}{fecha}\PY{l+s+s2}{\PYZdq{}}\PY{p}{)}
\end{Verbatim}
\end{tcolorbox}

    \begin{tcolorbox}[breakable, size=fbox, boxrule=1pt, pad at break*=1mm,colback=cellbackground, colframe=cellborder]
\prompt{In}{incolor}{183}{\boxspacing}
\begin{Verbatim}[commandchars=\\\{\}]
\PY{n}{df\PYZus{}ml}\PY{p}{[}\PY{l+s+s2}{\PYZdq{}}\PY{l+s+s2}{lag\PYZus{}1}\PY{l+s+s2}{\PYZdq{}}\PY{p}{]} \PY{o}{=} \PY{n}{df\PYZus{}ml}\PY{p}{[}\PY{l+s+s2}{\PYZdq{}}\PY{l+s+s2}{Precio\PYZus{}maizblanco\PYZus{}mxpkilo\PYZus{}prom\PYZus{}nacional}\PY{l+s+s2}{\PYZdq{}}\PY{p}{]}\PY{o}{.}\PY{n}{shift}\PY{p}{(}\PY{l+m+mi}{1}\PY{p}{)}
\PY{n}{df\PYZus{}ml}\PY{p}{[}\PY{l+s+s2}{\PYZdq{}}\PY{l+s+s2}{lag\PYZus{}2}\PY{l+s+s2}{\PYZdq{}}\PY{p}{]} \PY{o}{=} \PY{n}{df\PYZus{}ml}\PY{p}{[}\PY{l+s+s2}{\PYZdq{}}\PY{l+s+s2}{Precio\PYZus{}maizblanco\PYZus{}mxpkilo\PYZus{}prom\PYZus{}nacional}\PY{l+s+s2}{\PYZdq{}}\PY{p}{]}\PY{o}{.}\PY{n}{shift}\PY{p}{(}\PY{l+m+mi}{2}\PY{p}{)}
\PY{n}{df\PYZus{}ml}\PY{p}{[}\PY{l+s+s2}{\PYZdq{}}\PY{l+s+s2}{lag\PYZus{}3}\PY{l+s+s2}{\PYZdq{}}\PY{p}{]} \PY{o}{=} \PY{n}{df\PYZus{}ml}\PY{p}{[}\PY{l+s+s2}{\PYZdq{}}\PY{l+s+s2}{Precio\PYZus{}maizblanco\PYZus{}mxpkilo\PYZus{}prom\PYZus{}nacional}\PY{l+s+s2}{\PYZdq{}}\PY{p}{]}\PY{o}{.}\PY{n}{shift}\PY{p}{(}\PY{l+m+mi}{3}\PY{p}{)}
\PY{n}{df\PYZus{}ml}\PY{p}{[}\PY{l+s+s2}{\PYZdq{}}\PY{l+s+s2}{lag\PYZus{}6}\PY{l+s+s2}{\PYZdq{}}\PY{p}{]} \PY{o}{=} \PY{n}{df\PYZus{}ml}\PY{p}{[}\PY{l+s+s2}{\PYZdq{}}\PY{l+s+s2}{Precio\PYZus{}maizblanco\PYZus{}mxpkilo\PYZus{}prom\PYZus{}nacional}\PY{l+s+s2}{\PYZdq{}}\PY{p}{]}\PY{o}{.}\PY{n}{shift}\PY{p}{(}\PY{l+m+mi}{6}\PY{p}{)}
\PY{n}{df\PYZus{}ml}\PY{p}{[}\PY{l+s+s2}{\PYZdq{}}\PY{l+s+s2}{lag\PYZus{}12}\PY{l+s+s2}{\PYZdq{}}\PY{p}{]} \PY{o}{=} \PY{n}{df\PYZus{}ml}\PY{p}{[}\PY{l+s+s2}{\PYZdq{}}\PY{l+s+s2}{Precio\PYZus{}maizblanco\PYZus{}mxpkilo\PYZus{}prom\PYZus{}nacional}\PY{l+s+s2}{\PYZdq{}}\PY{p}{]}\PY{o}{.}\PY{n}{shift}\PY{p}{(}\PY{l+m+mi}{12}\PY{p}{)}

\PY{n}{df\PYZus{}ml}\PY{p}{[}\PY{l+s+s2}{\PYZdq{}}\PY{l+s+s2}{media\PYZus{}movil\PYZus{}3}\PY{l+s+s2}{\PYZdq{}}\PY{p}{]} \PY{o}{=} \PY{p}{(}
    \PY{n}{df\PYZus{}ml}\PY{p}{[}\PY{l+s+s2}{\PYZdq{}}\PY{l+s+s2}{Precio\PYZus{}maizblanco\PYZus{}mxpkilo\PYZus{}prom\PYZus{}nacional}\PY{l+s+s2}{\PYZdq{}}\PY{p}{]}
    \PY{o}{.}\PY{n}{rolling}\PY{p}{(}\PY{l+m+mi}{3}\PY{p}{)}
    \PY{o}{.}\PY{n}{mean}\PY{p}{(}\PY{p}{)}
\PY{p}{)}

\PY{n}{df\PYZus{}ml}\PY{p}{[}\PY{l+s+s2}{\PYZdq{}}\PY{l+s+s2}{media\PYZus{}movil\PYZus{}6}\PY{l+s+s2}{\PYZdq{}}\PY{p}{]} \PY{o}{=} \PY{p}{(}
    \PY{n}{df\PYZus{}ml}\PY{p}{[}\PY{l+s+s2}{\PYZdq{}}\PY{l+s+s2}{Precio\PYZus{}maizblanco\PYZus{}mxpkilo\PYZus{}prom\PYZus{}nacional}\PY{l+s+s2}{\PYZdq{}}\PY{p}{]}
    \PY{o}{.}\PY{n}{rolling}\PY{p}{(}\PY{l+m+mi}{6}\PY{p}{)}
    \PY{o}{.}\PY{n}{mean}\PY{p}{(}\PY{p}{)}
\PY{p}{)}

\PY{n}{df\PYZus{}ml}\PY{p}{[}\PY{l+s+s2}{\PYZdq{}}\PY{l+s+s2}{media\PYZus{}movil\PYZus{}12}\PY{l+s+s2}{\PYZdq{}}\PY{p}{]} \PY{o}{=} \PY{p}{(}
    \PY{n}{df\PYZus{}ml}\PY{p}{[}\PY{l+s+s2}{\PYZdq{}}\PY{l+s+s2}{Precio\PYZus{}maizblanco\PYZus{}mxpkilo\PYZus{}prom\PYZus{}nacional}\PY{l+s+s2}{\PYZdq{}}\PY{p}{]}
    \PY{o}{.}\PY{n}{rolling}\PY{p}{(}\PY{l+m+mi}{12}\PY{p}{)}
    \PY{o}{.}\PY{n}{mean}\PY{p}{(}\PY{p}{)}
\PY{p}{)}

\PY{n}{df\PYZus{}ml} \PY{o}{=} \PY{n}{df\PYZus{}ml}\PY{o}{.}\PY{n}{dropna}\PY{p}{(}\PY{p}{)}
\end{Verbatim}
\end{tcolorbox}

    \begin{tcolorbox}[breakable, size=fbox, boxrule=1pt, pad at break*=1mm,colback=cellbackground, colframe=cellborder]
\prompt{In}{incolor}{184}{\boxspacing}
\begin{Verbatim}[commandchars=\\\{\}]
\PY{n}{features\PYZus{}xgb} \PY{o}{=} \PY{p}{[}
    \PY{l+s+s2}{\PYZdq{}}\PY{l+s+s2}{Precio\PYZus{}prom\PYZus{}CBOT\PYZus{}mxpkilo\PYZus{}maiz\PYZus{}internacional}\PY{l+s+s2}{\PYZdq{}}\PY{p}{,}
    \PY{l+s+s2}{\PYZdq{}}\PY{l+s+s2}{Costo\PYZus{}transporte\PYZus{}maiz\PYZus{}mxpkilo\PYZus{}prom\PYZus{}nac}\PY{l+s+s2}{\PYZdq{}}\PY{p}{,}
    \PY{l+s+s2}{\PYZdq{}}\PY{l+s+s2}{Temperatura\PYZus{}celsius\PYZus{}prom\PYZus{}nacional}\PY{l+s+s2}{\PYZdq{}}\PY{p}{,}
    \PY{l+s+s2}{\PYZdq{}}\PY{l+s+s2}{Indice\PYZus{}severidad\PYZus{}ponderada\PYZus{}sequia\PYZus{}prom\PYZus{}nacional}\PY{l+s+s2}{\PYZdq{}}\PY{p}{,}
    \PY{l+s+s2}{\PYZdq{}}\PY{l+s+s2}{lag\PYZus{}1}\PY{l+s+s2}{\PYZdq{}}\PY{p}{,}
    \PY{l+s+s2}{\PYZdq{}}\PY{l+s+s2}{lag\PYZus{}2}\PY{l+s+s2}{\PYZdq{}}\PY{p}{,}
    \PY{l+s+s2}{\PYZdq{}}\PY{l+s+s2}{lag\PYZus{}3}\PY{l+s+s2}{\PYZdq{}}\PY{p}{,}
    \PY{l+s+s2}{\PYZdq{}}\PY{l+s+s2}{lag\PYZus{}6}\PY{l+s+s2}{\PYZdq{}}\PY{p}{,}
    \PY{l+s+s2}{\PYZdq{}}\PY{l+s+s2}{lag\PYZus{}12}\PY{l+s+s2}{\PYZdq{}}\PY{p}{,}
    \PY{l+s+s2}{\PYZdq{}}\PY{l+s+s2}{media\PYZus{}movil\PYZus{}3}\PY{l+s+s2}{\PYZdq{}}\PY{p}{,}
    \PY{l+s+s2}{\PYZdq{}}\PY{l+s+s2}{media\PYZus{}movil\PYZus{}6}\PY{l+s+s2}{\PYZdq{}}\PY{p}{,}
    \PY{l+s+s2}{\PYZdq{}}\PY{l+s+s2}{media\PYZus{}movil\PYZus{}12}\PY{l+s+s2}{\PYZdq{}}
\PY{p}{]}

\PY{n}{target} \PY{o}{=} \PY{l+s+s2}{\PYZdq{}}\PY{l+s+s2}{Precio\PYZus{}maizblanco\PYZus{}mxpkilo\PYZus{}prom\PYZus{}nacional}\PY{l+s+s2}{\PYZdq{}}
\end{Verbatim}
\end{tcolorbox}

    \begin{tcolorbox}[breakable, size=fbox, boxrule=1pt, pad at break*=1mm,colback=cellbackground, colframe=cellborder]
\prompt{In}{incolor}{185}{\boxspacing}
\begin{Verbatim}[commandchars=\\\{\}]
\PY{n}{train\PYZus{}size\PYZus{}ml} \PY{o}{=} \PY{n+nb}{int}\PY{p}{(}\PY{n+nb}{len}\PY{p}{(}\PY{n}{df\PYZus{}ml}\PY{p}{)} \PY{o}{*} \PY{l+m+mf}{0.80}\PY{p}{)}

\PY{n}{train\PYZus{}ml} \PY{o}{=} \PY{n}{df\PYZus{}ml}\PY{o}{.}\PY{n}{iloc}\PY{p}{[}\PY{p}{:}\PY{n}{train\PYZus{}size\PYZus{}ml}\PY{p}{]}
\PY{n}{test\PYZus{}ml} \PY{o}{=} \PY{n}{df\PYZus{}ml}\PY{o}{.}\PY{n}{iloc}\PY{p}{[}\PY{n}{train\PYZus{}size\PYZus{}ml}\PY{p}{:}\PY{p}{]}

\PY{n}{X\PYZus{}train\PYZus{}xgb} \PY{o}{=} \PY{n}{train\PYZus{}ml}\PY{p}{[}\PY{n}{features\PYZus{}xgb}\PY{p}{]}
\PY{n}{y\PYZus{}train\PYZus{}xgb} \PY{o}{=} \PY{n}{train\PYZus{}ml}\PY{p}{[}\PY{n}{target}\PY{p}{]}

\PY{n}{X\PYZus{}test\PYZus{}xgb} \PY{o}{=} \PY{n}{test\PYZus{}ml}\PY{p}{[}\PY{n}{features\PYZus{}xgb}\PY{p}{]}
\PY{n}{y\PYZus{}test\PYZus{}xgb} \PY{o}{=} \PY{n}{test\PYZus{}ml}\PY{p}{[}\PY{n}{target}\PY{p}{]}

\PY{n+nb}{print}\PY{p}{(}\PY{n}{X\PYZus{}train\PYZus{}xgb}\PY{o}{.}\PY{n}{shape}\PY{p}{)}
\PY{n+nb}{print}\PY{p}{(}\PY{n}{X\PYZus{}test\PYZus{}xgb}\PY{o}{.}\PY{n}{shape}\PY{p}{)}
\end{Verbatim}
\end{tcolorbox}

    \begin{Verbatim}[commandchars=\\\{\}]
(244, 12)
(61, 12)
    \end{Verbatim}

    \begin{tcolorbox}[breakable, size=fbox, boxrule=1pt, pad at break*=1mm,colback=cellbackground, colframe=cellborder]
\prompt{In}{incolor}{188}{\boxspacing}
\begin{Verbatim}[commandchars=\\\{\}]
\PY{n}{modelo\PYZus{}rf} \PY{o}{=} \PY{n}{RandomForestRegressor}\PY{p}{(}
    \PY{n}{n\PYZus{}estimators}\PY{o}{=}\PY{l+m+mi}{500}\PY{p}{,}
    \PY{n}{max\PYZus{}depth}\PY{o}{=}\PY{l+m+mi}{8}\PY{p}{,}
    \PY{n}{random\PYZus{}state}\PY{o}{=}\PY{l+m+mi}{42}\PY{p}{,}
    \PY{n}{n\PYZus{}jobs}\PY{o}{=}\PY{o}{\PYZhy{}}\PY{l+m+mi}{1}
\PY{p}{)}

\PY{n}{modelo\PYZus{}rf}\PY{o}{.}\PY{n}{fit}\PY{p}{(}
    \PY{n}{X\PYZus{}train\PYZus{}xgb}\PY{p}{,}
    \PY{n}{y\PYZus{}train\PYZus{}xgb}
\PY{p}{)}
\end{Verbatim}
\end{tcolorbox}

            \begin{tcolorbox}[breakable, size=fbox, boxrule=.5pt, pad at break*=1mm, opacityfill=0]
\prompt{Out}{outcolor}{188}{\boxspacing}
\begin{Verbatim}[commandchars=\\\{\}]
RandomForestRegressor(max\_depth=8, n\_estimators=500, n\_jobs=-1, random\_state=42)
\end{Verbatim}
\end{tcolorbox}
        
    \begin{tcolorbox}[breakable, size=fbox, boxrule=1pt, pad at break*=1mm,colback=cellbackground, colframe=cellborder]
\prompt{In}{incolor}{189}{\boxspacing}
\begin{Verbatim}[commandchars=\\\{\}]
\PY{n}{pred\PYZus{}rf} \PY{o}{=} \PY{n}{modelo\PYZus{}rf}\PY{o}{.}\PY{n}{predict}\PY{p}{(}\PY{n}{X\PYZus{}test\PYZus{}xgb}\PY{p}{)}
\end{Verbatim}
\end{tcolorbox}

    \begin{tcolorbox}[breakable, size=fbox, boxrule=1pt, pad at break*=1mm,colback=cellbackground, colframe=cellborder]
\prompt{In}{incolor}{190}{\boxspacing}
\begin{Verbatim}[commandchars=\\\{\}]
\PY{k+kn}{from} \PY{n+nn}{sklearn}\PY{n+nn}{.}\PY{n+nn}{metrics} \PY{k+kn}{import} \PY{p}{(}
    \PY{n}{mean\PYZus{}absolute\PYZus{}error}\PY{p}{,}
    \PY{n}{mean\PYZus{}squared\PYZus{}error}
\PY{p}{)}

\PY{n}{mae\PYZus{}rf} \PY{o}{=} \PY{n}{mean\PYZus{}absolute\PYZus{}error}\PY{p}{(}
    \PY{n}{y\PYZus{}test\PYZus{}xgb}\PY{p}{,}
    \PY{n}{pred\PYZus{}rf}
\PY{p}{)}

\PY{n}{rmse\PYZus{}rf} \PY{o}{=} \PY{n}{np}\PY{o}{.}\PY{n}{sqrt}\PY{p}{(}
    \PY{n}{mean\PYZus{}squared\PYZus{}error}\PY{p}{(}
        \PY{n}{y\PYZus{}test\PYZus{}xgb}\PY{p}{,}
        \PY{n}{pred\PYZus{}rf}
    \PY{p}{)}
\PY{p}{)}

\PY{n}{mape\PYZus{}rf} \PY{o}{=} \PY{p}{(}
    \PY{n}{np}\PY{o}{.}\PY{n}{mean}\PY{p}{(}
        \PY{n}{np}\PY{o}{.}\PY{n}{abs}\PY{p}{(}
            \PY{p}{(}\PY{n}{y\PYZus{}test\PYZus{}xgb} \PY{o}{\PYZhy{}} \PY{n}{pred\PYZus{}rf}\PY{p}{)}
            \PY{o}{/} \PY{n}{y\PYZus{}test\PYZus{}xgb}
        \PY{p}{)}
    \PY{p}{)}
    \PY{o}{*} \PY{l+m+mi}{100}
\PY{p}{)}

\PY{n+nb}{print}\PY{p}{(}\PY{l+s+s2}{\PYZdq{}}\PY{l+s+s2}{MAE :}\PY{l+s+s2}{\PYZdq{}}\PY{p}{,} \PY{n}{mae\PYZus{}rf}\PY{p}{)}
\PY{n+nb}{print}\PY{p}{(}\PY{l+s+s2}{\PYZdq{}}\PY{l+s+s2}{RMSE:}\PY{l+s+s2}{\PYZdq{}}\PY{p}{,} \PY{n}{rmse\PYZus{}rf}\PY{p}{)}
\PY{n+nb}{print}\PY{p}{(}\PY{l+s+s2}{\PYZdq{}}\PY{l+s+s2}{MAPE:}\PY{l+s+s2}{\PYZdq{}}\PY{p}{,} \PY{n}{mape\PYZus{}rf}\PY{p}{)}
\end{Verbatim}
\end{tcolorbox}

    \begin{Verbatim}[commandchars=\\\{\}]
MAE : 2.790027409836063
RMSE: 3.010637640646865
MAPE: 29.792979974003636
    \end{Verbatim}

    \begin{tcolorbox}[breakable, size=fbox, boxrule=1pt, pad at break*=1mm,colback=cellbackground, colframe=cellborder]
\prompt{In}{incolor}{191}{\boxspacing}
\begin{Verbatim}[commandchars=\\\{\}]
\PY{n}{importancias\PYZus{}rf} \PY{o}{=} \PY{n}{pd}\PY{o}{.}\PY{n}{DataFrame}\PY{p}{(}\PY{p}{\PYZob{}}
    \PY{l+s+s2}{\PYZdq{}}\PY{l+s+s2}{Variable}\PY{l+s+s2}{\PYZdq{}}\PY{p}{:} \PY{n}{features\PYZus{}xgb}\PY{p}{,}
    \PY{l+s+s2}{\PYZdq{}}\PY{l+s+s2}{Importancia}\PY{l+s+s2}{\PYZdq{}}\PY{p}{:} \PY{n}{modelo\PYZus{}rf}\PY{o}{.}\PY{n}{feature\PYZus{}importances\PYZus{}}
\PY{p}{\PYZcb{}}\PY{p}{)}\PY{o}{.}\PY{n}{sort\PYZus{}values}\PY{p}{(}
    \PY{n}{by}\PY{o}{=}\PY{l+s+s2}{\PYZdq{}}\PY{l+s+s2}{Importancia}\PY{l+s+s2}{\PYZdq{}}\PY{p}{,}
    \PY{n}{ascending}\PY{o}{=}\PY{k+kc}{False}
\PY{p}{)}

\PY{n}{importancias\PYZus{}rf}
\end{Verbatim}
\end{tcolorbox}

            \begin{tcolorbox}[breakable, size=fbox, boxrule=.5pt, pad at break*=1mm, opacityfill=0]
\prompt{Out}{outcolor}{191}{\boxspacing}
\begin{Verbatim}[commandchars=\\\{\}]
                                           Variable  Importancia
9                                     media\_movil\_3     0.328038
8                                            lag\_12     0.326474
11                                   media\_movil\_12     0.146844
4                                             lag\_1     0.067694
1            Costo\_transporte\_maiz\_mxpkilo\_prom\_nac     0.056751
10                                    media\_movil\_6     0.044030
5                                             lag\_2     0.015341
6                                             lag\_3     0.009295
0       Precio\_prom\_CBOT\_mxpkilo\_maiz\_internacional     0.004307
7                                             lag\_6     0.000559
3   Indice\_severidad\_ponderada\_sequia\_prom\_nacional     0.000336
2                 Temperatura\_celsius\_prom\_nacional     0.000330
\end{Verbatim}
\end{tcolorbox}
        
    \begin{tcolorbox}[breakable, size=fbox, boxrule=1pt, pad at break*=1mm,colback=cellbackground, colframe=cellborder]
\prompt{In}{incolor}{192}{\boxspacing}
\begin{Verbatim}[commandchars=\\\{\}]
\PY{n}{plt}\PY{o}{.}\PY{n}{figure}\PY{p}{(}\PY{n}{figsize}\PY{o}{=}\PY{p}{(}\PY{l+m+mi}{14}\PY{p}{,}\PY{l+m+mi}{6}\PY{p}{)}\PY{p}{)}

\PY{n}{plt}\PY{o}{.}\PY{n}{plot}\PY{p}{(}
    \PY{n}{test\PYZus{}ml}\PY{p}{[}\PY{l+s+s2}{\PYZdq{}}\PY{l+s+s2}{fecha}\PY{l+s+s2}{\PYZdq{}}\PY{p}{]}\PY{p}{,}
    \PY{n}{y\PYZus{}test\PYZus{}xgb}\PY{p}{,}
    \PY{n}{label}\PY{o}{=}\PY{l+s+s2}{\PYZdq{}}\PY{l+s+s2}{Real}\PY{l+s+s2}{\PYZdq{}}\PY{p}{,}
    \PY{n}{linewidth}\PY{o}{=}\PY{l+m+mi}{2}
\PY{p}{)}

\PY{n}{plt}\PY{o}{.}\PY{n}{plot}\PY{p}{(}
    \PY{n}{test\PYZus{}ml}\PY{p}{[}\PY{l+s+s2}{\PYZdq{}}\PY{l+s+s2}{fecha}\PY{l+s+s2}{\PYZdq{}}\PY{p}{]}\PY{p}{,}
    \PY{n}{pred\PYZus{}rf}\PY{p}{,}
    \PY{n}{label}\PY{o}{=}\PY{l+s+s2}{\PYZdq{}}\PY{l+s+s2}{Random Forest}\PY{l+s+s2}{\PYZdq{}}\PY{p}{,}
    \PY{n}{linestyle}\PY{o}{=}\PY{l+s+s2}{\PYZdq{}}\PY{l+s+s2}{\PYZhy{}\PYZhy{}}\PY{l+s+s2}{\PYZdq{}}
\PY{p}{)}

\PY{n}{plt}\PY{o}{.}\PY{n}{title}\PY{p}{(}
    \PY{l+s+s2}{\PYZdq{}}\PY{l+s+s2}{Validación del Modelo Random Forest}\PY{l+s+s2}{\PYZdq{}}
\PY{p}{)}

\PY{n}{plt}\PY{o}{.}\PY{n}{xlabel}\PY{p}{(}\PY{l+s+s2}{\PYZdq{}}\PY{l+s+s2}{Fecha}\PY{l+s+s2}{\PYZdq{}}\PY{p}{)}
\PY{n}{plt}\PY{o}{.}\PY{n}{ylabel}\PY{p}{(}\PY{l+s+s2}{\PYZdq{}}\PY{l+s+s2}{Precio del Maíz (MXN/Kg)}\PY{l+s+s2}{\PYZdq{}}\PY{p}{)}

\PY{n}{plt}\PY{o}{.}\PY{n}{legend}\PY{p}{(}\PY{p}{)}
\PY{n}{plt}\PY{o}{.}\PY{n}{grid}\PY{p}{(}\PY{k+kc}{True}\PY{p}{)}

\PY{n}{plt}\PY{o}{.}\PY{n}{show}\PY{p}{(}\PY{p}{)}
\end{Verbatim}
\end{tcolorbox}

    \begin{center}
    \adjustimage{max size={0.9\linewidth}{0.9\paperheight}}{output_259_0.png}
    \end{center}
    { \hspace*{\fill} \\}
    
    \hypertarget{comparaciuxf3n-entre-sarimax-y-random-forest}{%
\subsection{Comparación entre SARIMAX y Random
Forest}\label{comparaciuxf3n-entre-sarimax-y-random-forest}}

Con el objetivo de integrar la UCA de Aprendizaje de Máquina, se
construyó un modelo Random Forest utilizando variables económicas,
climáticas y rezagos históricos del precio del maíz.

Los resultados obtenidos muestran que el desempeño del Random Forest fue
inferior al alcanzado por el modelo SARIMAX reducido.

Mientras que el modelo SARIMAX reducido alcanzó un MAPE de 15.83\%, el
Random Forest obtuvo un MAPE de 29.79\%.

La inspección visual de las predicciones indica que el Random Forest
tendió a generar estimaciones excesivamente conservadoras, sin
reproducir adecuadamente los incrementos observados durante el periodo
2021--2024.

Este comportamiento puede explicarse porque los algoritmos basados en
árboles presentan limitaciones para extrapolar valores fuera de los
rangos observados durante el entrenamiento. En contraste, los modelos
SARIMAX están específicamente diseñados para capturar tendencias,
estacionalidad y dependencia temporal, características fundamentales de
las series de tiempo económicas.

Por lo tanto, para el caso del precio nacional del maíz blanco, los
resultados sugieren que el enfoque basado en series de tiempo
proporciona una representación más adecuada de la dinámica observada en
el mercado.

    \hypertarget{anexo-linea-de-investigaciuxf3n-futura}{%
\section{Anexo: Linea de investigación
futura}\label{anexo-linea-de-investigaciuxf3n-futura}}

\hypertarget{modelo-de-predicciuxf3n-de-precio-nacional-de-tortilla}{%
\subsection{Modelo de predicción de precio nacional de
tortilla}\label{modelo-de-predicciuxf3n-de-precio-nacional-de-tortilla}}

    \begin{tcolorbox}[breakable, size=fbox, boxrule=1pt, pad at break*=1mm,colback=cellbackground, colframe=cellborder]
\prompt{In}{incolor}{263}{\boxspacing}
\begin{Verbatim}[commandchars=\\\{\}]
\PY{n}{modelo\PYZus{}tortilla} \PY{o}{=} \PY{p}{(}
    \PY{n}{pdf\PYZus{}completo}\PY{p}{[}
        \PY{p}{[}
            \PY{l+s+s2}{\PYZdq{}}\PY{l+s+s2}{fecha}\PY{l+s+s2}{\PYZdq{}}\PY{p}{,}
            \PY{l+s+s2}{\PYZdq{}}\PY{l+s+s2}{Precio\PYZus{}maizblanco\PYZus{}mxpkilo\PYZus{}prom\PYZus{}nacional}\PY{l+s+s2}{\PYZdq{}}\PY{p}{,}
            \PY{l+s+s2}{\PYZdq{}}\PY{l+s+s2}{Precio\PYZus{}prom\PYZus{}CBOT\PYZus{}mxpkilo\PYZus{}maiz\PYZus{}internacional}\PY{l+s+s2}{\PYZdq{}}\PY{p}{,}
            \PY{l+s+s2}{\PYZdq{}}\PY{l+s+s2}{Costo\PYZus{}transporte\PYZus{}maiz\PYZus{}mxpkilo\PYZus{}prom\PYZus{}nac}\PY{l+s+s2}{\PYZdq{}}\PY{p}{,}
            \PY{l+s+s2}{\PYZdq{}}\PY{l+s+s2}{Temperatura\PYZus{}celsius\PYZus{}prom\PYZus{}nacional}\PY{l+s+s2}{\PYZdq{}}\PY{p}{,}
            \PY{l+s+s2}{\PYZdq{}}\PY{l+s+s2}{Indice\PYZus{}severidad\PYZus{}ponderada\PYZus{}sequia\PYZus{}prom\PYZus{}nacional}\PY{l+s+s2}{\PYZdq{}}
        \PY{p}{]}
    \PY{p}{]}
    \PY{o}{.}\PY{n}{merge}\PY{p}{(}
        \PY{n}{pdf\PYZus{}tortilla}\PY{p}{[}\PY{p}{[}\PY{l+s+s2}{\PYZdq{}}\PY{l+s+s2}{fecha}\PY{l+s+s2}{\PYZdq{}}\PY{p}{,} \PY{l+s+s2}{\PYZdq{}}\PY{l+s+s2}{Nacional}\PY{l+s+s2}{\PYZdq{}}\PY{p}{]}\PY{p}{]}\PY{p}{,}
        \PY{n}{on}\PY{o}{=}\PY{l+s+s2}{\PYZdq{}}\PY{l+s+s2}{fecha}\PY{l+s+s2}{\PYZdq{}}\PY{p}{,}
        \PY{n}{how}\PY{o}{=}\PY{l+s+s2}{\PYZdq{}}\PY{l+s+s2}{inner}\PY{l+s+s2}{\PYZdq{}}
    \PY{p}{)}
\PY{p}{)}

\PY{n}{modelo\PYZus{}tortilla} \PY{o}{=} \PY{n}{modelo\PYZus{}tortilla}\PY{o}{.}\PY{n}{rename}\PY{p}{(}
    \PY{n}{columns}\PY{o}{=}\PY{p}{\PYZob{}}
        \PY{l+s+s2}{\PYZdq{}}\PY{l+s+s2}{Nacional}\PY{l+s+s2}{\PYZdq{}}\PY{p}{:} \PY{l+s+s2}{\PYZdq{}}\PY{l+s+s2}{Precio\PYZus{}tortilla}\PY{l+s+s2}{\PYZdq{}}
    \PY{p}{\PYZcb{}}
\PY{p}{)}

\PY{n}{modelo\PYZus{}tortilla}\PY{p}{[}\PY{l+s+s2}{\PYZdq{}}\PY{l+s+s2}{fecha}\PY{l+s+s2}{\PYZdq{}}\PY{p}{]} \PY{o}{=} \PY{n}{pd}\PY{o}{.}\PY{n}{to\PYZus{}datetime}\PY{p}{(}\PY{n}{modelo\PYZus{}tortilla}\PY{p}{[}\PY{l+s+s2}{\PYZdq{}}\PY{l+s+s2}{fecha}\PY{l+s+s2}{\PYZdq{}}\PY{p}{]}\PY{p}{)}
\PY{n}{modelo\PYZus{}tortilla} \PY{o}{=} \PY{n}{modelo\PYZus{}tortilla}\PY{o}{.}\PY{n}{sort\PYZus{}values}\PY{p}{(}\PY{l+s+s2}{\PYZdq{}}\PY{l+s+s2}{fecha}\PY{l+s+s2}{\PYZdq{}}\PY{p}{)}

\PY{n}{modelo\PYZus{}tortilla}\PY{o}{.}\PY{n}{head}\PY{p}{(}\PY{p}{)}
\end{Verbatim}
\end{tcolorbox}

            \begin{tcolorbox}[breakable, size=fbox, boxrule=.5pt, pad at break*=1mm, opacityfill=0]
\prompt{Out}{outcolor}{263}{\boxspacing}
\begin{Verbatim}[commandchars=\\\{\}]
       fecha  Precio\_maizblanco\_mxpkilo\_prom\_nacional  \textbackslash{}
0 2000-01-01                                     2.05
1 2000-02-01                                     2.06
2 2000-03-01                                     2.09
3 2000-04-01                                     2.12
4 2000-05-01                                     2.14

   Precio\_prom\_CBOT\_mxpkilo\_maiz\_internacional  \textbackslash{}
0                                         0.92
1                                         0.91
2                                         0.89
3                                         0.92
4                                         1.00

   Costo\_transporte\_maiz\_mxpkilo\_prom\_nac  Temperatura\_celsius\_prom\_nacional  \textbackslash{}
0                                    0.22                               16.6
1                                    0.22                               18.3
2                                    0.23                               21.8
3                                    0.23                               24.4
4                                    0.23                               26.7

   Indice\_severidad\_ponderada\_sequia\_prom\_nacional  Precio\_tortilla
0                                            43.12             4.00
1                                            44.28             4.04
2                                            46.24             4.07
3                                            56.12             4.10
4                                            53.83             4.14
\end{Verbatim}
\end{tcolorbox}
        
    \begin{tcolorbox}[breakable, size=fbox, boxrule=1pt, pad at break*=1mm,colback=cellbackground, colframe=cellborder]
\prompt{In}{incolor}{264}{\boxspacing}
\begin{Verbatim}[commandchars=\\\{\}]
\PY{n}{corr\PYZus{}tortilla} \PY{o}{=} \PY{p}{(}
    \PY{n}{modelo\PYZus{}tortilla}
    \PY{o}{.}\PY{n}{drop}\PY{p}{(}\PY{n}{columns}\PY{o}{=}\PY{p}{[}\PY{l+s+s2}{\PYZdq{}}\PY{l+s+s2}{fecha}\PY{l+s+s2}{\PYZdq{}}\PY{p}{]}\PY{p}{)}
    \PY{o}{.}\PY{n}{corr}\PY{p}{(}\PY{p}{)}\PY{p}{[}\PY{l+s+s2}{\PYZdq{}}\PY{l+s+s2}{Precio\PYZus{}tortilla}\PY{l+s+s2}{\PYZdq{}}\PY{p}{]}
    \PY{o}{.}\PY{n}{sort\PYZus{}values}\PY{p}{(}\PY{n}{ascending}\PY{o}{=}\PY{k+kc}{False}\PY{p}{)}
\PY{p}{)}

\PY{n}{corr\PYZus{}tortilla}
\end{Verbatim}
\end{tcolorbox}

            \begin{tcolorbox}[breakable, size=fbox, boxrule=.5pt, pad at break*=1mm, opacityfill=0]
\prompt{Out}{outcolor}{264}{\boxspacing}
\begin{Verbatim}[commandchars=\\\{\}]
Precio\_tortilla                                    1.000000
Costo\_transporte\_maiz\_mxpkilo\_prom\_nac             0.981050
Precio\_maizblanco\_mxpkilo\_prom\_nacional            0.972386
Precio\_prom\_CBOT\_mxpkilo\_maiz\_internacional        0.810711
Indice\_severidad\_ponderada\_sequia\_prom\_nacional    0.442966
Temperatura\_celsius\_prom\_nacional                  0.069833
Name: Precio\_tortilla, dtype: float64
\end{Verbatim}
\end{tcolorbox}
        
    \begin{tcolorbox}[breakable, size=fbox, boxrule=1pt, pad at break*=1mm,colback=cellbackground, colframe=cellborder]
\prompt{In}{incolor}{265}{\boxspacing}
\begin{Verbatim}[commandchars=\\\{\}]
\PY{n}{serie\PYZus{}tortilla} \PY{o}{=} \PY{p}{(}
    \PY{n}{modelo\PYZus{}tortilla}
    \PY{o}{.}\PY{n}{set\PYZus{}index}\PY{p}{(}\PY{l+s+s2}{\PYZdq{}}\PY{l+s+s2}{fecha}\PY{l+s+s2}{\PYZdq{}}\PY{p}{)}
    \PY{p}{[}\PY{l+s+s2}{\PYZdq{}}\PY{l+s+s2}{Precio\PYZus{}tortilla}\PY{l+s+s2}{\PYZdq{}}\PY{p}{]}
\PY{p}{)}

\PY{n}{variables\PYZus{}tortilla} \PY{o}{=} \PY{p}{[}
    \PY{l+s+s2}{\PYZdq{}}\PY{l+s+s2}{Precio\PYZus{}maizblanco\PYZus{}mxpkilo\PYZus{}prom\PYZus{}nacional}\PY{l+s+s2}{\PYZdq{}}\PY{p}{,}
    \PY{l+s+s2}{\PYZdq{}}\PY{l+s+s2}{Costo\PYZus{}transporte\PYZus{}maiz\PYZus{}mxpkilo\PYZus{}prom\PYZus{}nac}\PY{l+s+s2}{\PYZdq{}}\PY{p}{,}
    \PY{l+s+s2}{\PYZdq{}}\PY{l+s+s2}{Temperatura\PYZus{}celsius\PYZus{}prom\PYZus{}nacional}\PY{l+s+s2}{\PYZdq{}}\PY{p}{,}
    \PY{l+s+s2}{\PYZdq{}}\PY{l+s+s2}{Indice\PYZus{}severidad\PYZus{}ponderada\PYZus{}sequia\PYZus{}prom\PYZus{}nacional}\PY{l+s+s2}{\PYZdq{}}
\PY{p}{]}

\PY{n}{X\PYZus{}tortilla} \PY{o}{=} \PY{p}{(}
    \PY{n}{modelo\PYZus{}tortilla}
    \PY{o}{.}\PY{n}{set\PYZus{}index}\PY{p}{(}\PY{l+s+s2}{\PYZdq{}}\PY{l+s+s2}{fecha}\PY{l+s+s2}{\PYZdq{}}\PY{p}{)}
    \PY{p}{[}\PY{n}{variables\PYZus{}tortilla}\PY{p}{]}
\PY{p}{)}
\end{Verbatim}
\end{tcolorbox}

    \begin{tcolorbox}[breakable, size=fbox, boxrule=1pt, pad at break*=1mm,colback=cellbackground, colframe=cellborder]
\prompt{In}{incolor}{266}{\boxspacing}
\begin{Verbatim}[commandchars=\\\{\}]
\PY{k+kn}{from} \PY{n+nn}{statsmodels}\PY{n+nn}{.}\PY{n+nn}{tsa}\PY{n+nn}{.}\PY{n+nn}{stattools} \PY{k+kn}{import} \PY{n}{adfuller}

\PY{n+nb}{print}\PY{p}{(}\PY{l+s+s2}{\PYZdq{}}\PY{l+s+s2}{Serie original}\PY{l+s+s2}{\PYZdq{}}\PY{p}{)}
\PY{n}{adf\PYZus{}original} \PY{o}{=} \PY{n}{adfuller}\PY{p}{(}\PY{n}{serie\PYZus{}tortilla}\PY{p}{)}

\PY{n+nb}{print}\PY{p}{(}\PY{l+s+s2}{\PYZdq{}}\PY{l+s+s2}{ADF:}\PY{l+s+s2}{\PYZdq{}}\PY{p}{,} \PY{n}{adf\PYZus{}original}\PY{p}{[}\PY{l+m+mi}{0}\PY{p}{]}\PY{p}{)}
\PY{n+nb}{print}\PY{p}{(}\PY{l+s+s2}{\PYZdq{}}\PY{l+s+s2}{p\PYZhy{}value:}\PY{l+s+s2}{\PYZdq{}}\PY{p}{,} \PY{n}{adf\PYZus{}original}\PY{p}{[}\PY{l+m+mi}{1}\PY{p}{]}\PY{p}{)}

\PY{n+nb}{print}\PY{p}{(}\PY{l+s+s2}{\PYZdq{}}\PY{l+s+se}{\PYZbs{}n}\PY{l+s+s2}{Serie diferenciada}\PY{l+s+s2}{\PYZdq{}}\PY{p}{)}
\PY{n}{serie\PYZus{}tortilla\PYZus{}diff} \PY{o}{=} \PY{n}{serie\PYZus{}tortilla}\PY{o}{.}\PY{n}{diff}\PY{p}{(}\PY{p}{)}\PY{o}{.}\PY{n}{dropna}\PY{p}{(}\PY{p}{)}

\PY{n}{adf\PYZus{}diff} \PY{o}{=} \PY{n}{adfuller}\PY{p}{(}\PY{n}{serie\PYZus{}tortilla\PYZus{}diff}\PY{p}{)}

\PY{n+nb}{print}\PY{p}{(}\PY{l+s+s2}{\PYZdq{}}\PY{l+s+s2}{ADF:}\PY{l+s+s2}{\PYZdq{}}\PY{p}{,} \PY{n}{adf\PYZus{}diff}\PY{p}{[}\PY{l+m+mi}{0}\PY{p}{]}\PY{p}{)}
\PY{n+nb}{print}\PY{p}{(}\PY{l+s+s2}{\PYZdq{}}\PY{l+s+s2}{p\PYZhy{}value:}\PY{l+s+s2}{\PYZdq{}}\PY{p}{,} \PY{n}{adf\PYZus{}diff}\PY{p}{[}\PY{l+m+mi}{1}\PY{p}{]}\PY{p}{)}
\end{Verbatim}
\end{tcolorbox}

    \begin{Verbatim}[commandchars=\\\{\}]
Serie original
ADF: 1.778498766226762
p-value: 0.9983054057793194

Serie diferenciada
ADF: -16.611544415196967
p-value: 1.72912995198354e-29
    \end{Verbatim}

    \begin{tcolorbox}[breakable, size=fbox, boxrule=1pt, pad at break*=1mm,colback=cellbackground, colframe=cellborder]
\prompt{In}{incolor}{267}{\boxspacing}
\begin{Verbatim}[commandchars=\\\{\}]
\PY{k+kn}{from} \PY{n+nn}{statsmodels}\PY{n+nn}{.}\PY{n+nn}{tsa}\PY{n+nn}{.}\PY{n+nn}{seasonal} \PY{k+kn}{import} \PY{n}{seasonal\PYZus{}decompose}

\PY{n}{descomposicion\PYZus{}tortilla} \PY{o}{=} \PY{n}{seasonal\PYZus{}decompose}\PY{p}{(}
    \PY{n}{serie\PYZus{}tortilla}\PY{p}{,}
    \PY{n}{model}\PY{o}{=}\PY{l+s+s2}{\PYZdq{}}\PY{l+s+s2}{additive}\PY{l+s+s2}{\PYZdq{}}\PY{p}{,}
    \PY{n}{period}\PY{o}{=}\PY{l+m+mi}{12}
\PY{p}{)}

\PY{n}{descomposicion\PYZus{}tortilla}\PY{o}{.}\PY{n}{plot}\PY{p}{(}\PY{p}{)}
\PY{n}{plt}\PY{o}{.}\PY{n}{show}\PY{p}{(}\PY{p}{)}
\end{Verbatim}
\end{tcolorbox}

    \begin{center}
    \adjustimage{max size={0.9\linewidth}{0.9\paperheight}}{output_266_0.png}
    \end{center}
    { \hspace*{\fill} \\}
    
    \begin{tcolorbox}[breakable, size=fbox, boxrule=1pt, pad at break*=1mm,colback=cellbackground, colframe=cellborder]
\prompt{In}{incolor}{268}{\boxspacing}
\begin{Verbatim}[commandchars=\\\{\}]
\PY{k+kn}{from} \PY{n+nn}{statsmodels}\PY{n+nn}{.}\PY{n+nn}{graphics}\PY{n+nn}{.}\PY{n+nn}{tsaplots} \PY{k+kn}{import} \PY{n}{plot\PYZus{}acf}\PY{p}{,} \PY{n}{plot\PYZus{}pacf}

\PY{n}{fig}\PY{p}{,} \PY{n}{ax} \PY{o}{=} \PY{n}{plt}\PY{o}{.}\PY{n}{subplots}\PY{p}{(}\PY{l+m+mi}{2}\PY{p}{,} \PY{l+m+mi}{1}\PY{p}{,} \PY{n}{figsize}\PY{o}{=}\PY{p}{(}\PY{l+m+mi}{12}\PY{p}{,} \PY{l+m+mi}{8}\PY{p}{)}\PY{p}{)}

\PY{n}{plot\PYZus{}acf}\PY{p}{(}
    \PY{n}{serie\PYZus{}tortilla\PYZus{}diff}\PY{p}{,}
    \PY{n}{lags}\PY{o}{=}\PY{l+m+mi}{36}\PY{p}{,}
    \PY{n}{ax}\PY{o}{=}\PY{n}{ax}\PY{p}{[}\PY{l+m+mi}{0}\PY{p}{]}
\PY{p}{)}

\PY{n}{plot\PYZus{}pacf}\PY{p}{(}
    \PY{n}{serie\PYZus{}tortilla\PYZus{}diff}\PY{p}{,}
    \PY{n}{lags}\PY{o}{=}\PY{l+m+mi}{36}\PY{p}{,}
    \PY{n}{ax}\PY{o}{=}\PY{n}{ax}\PY{p}{[}\PY{l+m+mi}{1}\PY{p}{]}
\PY{p}{)}

\PY{n}{plt}\PY{o}{.}\PY{n}{tight\PYZus{}layout}\PY{p}{(}\PY{p}{)}
\PY{n}{plt}\PY{o}{.}\PY{n}{show}\PY{p}{(}\PY{p}{)}
\end{Verbatim}
\end{tcolorbox}

    \begin{center}
    \adjustimage{max size={0.9\linewidth}{0.9\paperheight}}{output_267_0.png}
    \end{center}
    { \hspace*{\fill} \\}
    
    \begin{tcolorbox}[breakable, size=fbox, boxrule=1pt, pad at break*=1mm,colback=cellbackground, colframe=cellborder]
\prompt{In}{incolor}{269}{\boxspacing}
\begin{Verbatim}[commandchars=\\\{\}]
\PY{n}{train\PYZus{}size\PYZus{}tortilla} \PY{o}{=} \PY{n+nb}{int}\PY{p}{(}\PY{n+nb}{len}\PY{p}{(}\PY{n}{serie\PYZus{}tortilla}\PY{p}{)} \PY{o}{*} \PY{l+m+mf}{0.80}\PY{p}{)}

\PY{n}{y\PYZus{}train\PYZus{}tortilla} \PY{o}{=} \PY{n}{serie\PYZus{}tortilla}\PY{o}{.}\PY{n}{iloc}\PY{p}{[}\PY{p}{:}\PY{n}{train\PYZus{}size\PYZus{}tortilla}\PY{p}{]}
\PY{n}{y\PYZus{}test\PYZus{}tortilla} \PY{o}{=} \PY{n}{serie\PYZus{}tortilla}\PY{o}{.}\PY{n}{iloc}\PY{p}{[}\PY{n}{train\PYZus{}size\PYZus{}tortilla}\PY{p}{:}\PY{p}{]}

\PY{n}{X\PYZus{}train\PYZus{}tortilla} \PY{o}{=} \PY{n}{X\PYZus{}tortilla}\PY{o}{.}\PY{n}{iloc}\PY{p}{[}\PY{p}{:}\PY{n}{train\PYZus{}size\PYZus{}tortilla}\PY{p}{]}
\PY{n}{X\PYZus{}test\PYZus{}tortilla} \PY{o}{=} \PY{n}{X\PYZus{}tortilla}\PY{o}{.}\PY{n}{iloc}\PY{p}{[}\PY{n}{train\PYZus{}size\PYZus{}tortilla}\PY{p}{:}\PY{p}{]}

\PY{n+nb}{print}\PY{p}{(}\PY{l+s+s2}{\PYZdq{}}\PY{l+s+s2}{Entrenamiento:}\PY{l+s+s2}{\PYZdq{}}\PY{p}{,} \PY{n}{y\PYZus{}train\PYZus{}tortilla}\PY{o}{.}\PY{n}{index}\PY{o}{.}\PY{n}{min}\PY{p}{(}\PY{p}{)}\PY{p}{,} \PY{n}{y\PYZus{}train\PYZus{}tortilla}\PY{o}{.}\PY{n}{index}\PY{o}{.}\PY{n}{max}\PY{p}{(}\PY{p}{)}\PY{p}{)}
\PY{n+nb}{print}\PY{p}{(}\PY{l+s+s2}{\PYZdq{}}\PY{l+s+s2}{Prueba:}\PY{l+s+s2}{\PYZdq{}}\PY{p}{,} \PY{n}{y\PYZus{}test\PYZus{}tortilla}\PY{o}{.}\PY{n}{index}\PY{o}{.}\PY{n}{min}\PY{p}{(}\PY{p}{)}\PY{p}{,} \PY{n}{y\PYZus{}test\PYZus{}tortilla}\PY{o}{.}\PY{n}{index}\PY{o}{.}\PY{n}{max}\PY{p}{(}\PY{p}{)}\PY{p}{)}
\PY{n+nb}{print}\PY{p}{(}\PY{n+nb}{len}\PY{p}{(}\PY{n}{y\PYZus{}train\PYZus{}tortilla}\PY{p}{)}\PY{p}{,} \PY{n+nb}{len}\PY{p}{(}\PY{n}{y\PYZus{}test\PYZus{}tortilla}\PY{p}{)}\PY{p}{)}
\end{Verbatim}
\end{tcolorbox}

    \begin{Verbatim}[commandchars=\\\{\}]
Entrenamiento: 2000-01-01 00:00:00 2021-01-01 00:00:00
Prueba: 2021-02-01 00:00:00 2026-05-01 00:00:00
253 64
    \end{Verbatim}

    \begin{tcolorbox}[breakable, size=fbox, boxrule=1pt, pad at break*=1mm,colback=cellbackground, colframe=cellborder]
\prompt{In}{incolor}{270}{\boxspacing}
\begin{Verbatim}[commandchars=\\\{\}]
\PY{k+kn}{from} \PY{n+nn}{statsmodels}\PY{n+nn}{.}\PY{n+nn}{tsa}\PY{n+nn}{.}\PY{n+nn}{statespace}\PY{n+nn}{.}\PY{n+nn}{sarimax} \PY{k+kn}{import} \PY{n}{SARIMAX}

\PY{n}{modelo\PYZus{}sarimax\PYZus{}tortilla} \PY{o}{=} \PY{n}{SARIMAX}\PY{p}{(}
    \PY{n}{y\PYZus{}train\PYZus{}tortilla}\PY{p}{,}
    \PY{n}{exog}\PY{o}{=}\PY{n}{X\PYZus{}train\PYZus{}tortilla}\PY{p}{,}
    \PY{n}{order}\PY{o}{=}\PY{p}{(}\PY{l+m+mi}{1}\PY{p}{,} \PY{l+m+mi}{1}\PY{p}{,} \PY{l+m+mi}{0}\PY{p}{)}\PY{p}{,}
    \PY{n}{seasonal\PYZus{}order}\PY{o}{=}\PY{p}{(}\PY{l+m+mi}{1}\PY{p}{,} \PY{l+m+mi}{0}\PY{p}{,} \PY{l+m+mi}{0}\PY{p}{,} \PY{l+m+mi}{12}\PY{p}{)}\PY{p}{,}
    \PY{n}{enforce\PYZus{}stationarity}\PY{o}{=}\PY{k+kc}{False}\PY{p}{,}
    \PY{n}{enforce\PYZus{}invertibility}\PY{o}{=}\PY{k+kc}{False}
\PY{p}{)}

\PY{n}{resultado\PYZus{}tortilla} \PY{o}{=} \PY{n}{modelo\PYZus{}sarimax\PYZus{}tortilla}\PY{o}{.}\PY{n}{fit}\PY{p}{(}\PY{p}{)}

\PY{n+nb}{print}\PY{p}{(}\PY{n}{resultado\PYZus{}tortilla}\PY{o}{.}\PY{n}{summary}\PY{p}{(}\PY{p}{)}\PY{p}{)}
\end{Verbatim}
\end{tcolorbox}

    \begin{Verbatim}[commandchars=\\\{\}]
/opt/conda/lib/python3.11/site-packages/statsmodels/tsa/base/tsa\_model.py:473:
ValueWarning: No frequency information was provided, so inferred frequency MS
will be used.
  self.\_init\_dates(dates, freq)
/opt/conda/lib/python3.11/site-packages/statsmodels/tsa/base/tsa\_model.py:473:
ValueWarning: No frequency information was provided, so inferred frequency MS
will be used.
  self.\_init\_dates(dates, freq)
    \end{Verbatim}

    \begin{Verbatim}[commandchars=\\\{\}]
                                     SARIMAX Results
================================================================================
==========
Dep. Variable:                    Precio\_tortilla   No. Observations:
253
Model:             SARIMAX(1, 1, 0)x(1, 0, 0, 12)   Log Likelihood
32.495
Date:                            Thu, 04 Jun 2026   AIC
-50.990
Time:                                    09:37:10   BIC
-26.655
Sample:                                01-01-2000   HQIC
-41.184
                                     - 01-01-2021
Covariance Type:                              opg
================================================================================
===================================
                                                      coef    std err          z
P>|z|      [0.025      0.975]
--------------------------------------------------------------------------------
-----------------------------------
Precio\_maizblanco\_mxpkilo\_prom\_nacional             1.2916      0.087     14.825
0.000       1.121       1.462
Costo\_transporte\_maiz\_mxpkilo\_prom\_nac              1.9596      2.939      0.667
0.505      -3.800       7.719
Temperatura\_celsius\_prom\_nacional                  -0.0212      0.020     -1.074
0.283      -0.060       0.018
Indice\_severidad\_ponderada\_sequia\_prom\_nacional    -0.0011      0.001     -1.804
0.071      -0.002    9.81e-05
ar.L1                                              -0.0321      0.062     -0.522
0.602      -0.153       0.088
ar.S.L12                                            0.2438      0.038      6.390
0.000       0.169       0.319
sigma2                                              0.0446      0.002     22.096
0.000       0.041       0.049
================================================================================
===
Ljung-Box (L1) (Q):                   0.02   Jarque-Bera (JB):
16398.26
Prob(Q):                              0.90   Prob(JB):
0.00
Heteroskedasticity (H):               1.12   Skew:
2.35
Prob(H) (two-sided):                  0.63   Kurtosis:
43.31
================================================================================
===

Warnings:
[1] Covariance matrix calculated using the outer product of gradients (complex-
step).
    \end{Verbatim}

    \begin{tcolorbox}[breakable, size=fbox, boxrule=1pt, pad at break*=1mm,colback=cellbackground, colframe=cellborder]
\prompt{In}{incolor}{271}{\boxspacing}
\begin{Verbatim}[commandchars=\\\{\}]
\PY{n}{pred\PYZus{}tortilla} \PY{o}{=} \PY{n}{resultado\PYZus{}tortilla}\PY{o}{.}\PY{n}{get\PYZus{}prediction}\PY{p}{(}
    \PY{n}{start}\PY{o}{=}\PY{n}{y\PYZus{}test\PYZus{}tortilla}\PY{o}{.}\PY{n}{index}\PY{p}{[}\PY{l+m+mi}{0}\PY{p}{]}\PY{p}{,}
    \PY{n}{end}\PY{o}{=}\PY{n}{y\PYZus{}test\PYZus{}tortilla}\PY{o}{.}\PY{n}{index}\PY{p}{[}\PY{o}{\PYZhy{}}\PY{l+m+mi}{1}\PY{p}{]}\PY{p}{,}
    \PY{n}{exog}\PY{o}{=}\PY{n}{X\PYZus{}test\PYZus{}tortilla}
\PY{p}{)}

\PY{n}{pred\PYZus{}tortilla} \PY{o}{=} \PY{n}{pred\PYZus{}tortilla}\PY{o}{.}\PY{n}{predicted\PYZus{}mean}
\end{Verbatim}
\end{tcolorbox}

    \begin{tcolorbox}[breakable, size=fbox, boxrule=1pt, pad at break*=1mm,colback=cellbackground, colframe=cellborder]
\prompt{In}{incolor}{272}{\boxspacing}
\begin{Verbatim}[commandchars=\\\{\}]
\PY{k+kn}{from} \PY{n+nn}{sklearn}\PY{n+nn}{.}\PY{n+nn}{metrics} \PY{k+kn}{import} \PY{n}{mean\PYZus{}absolute\PYZus{}error}\PY{p}{,} \PY{n}{mean\PYZus{}squared\PYZus{}error}

\PY{n}{mae\PYZus{}tortilla} \PY{o}{=} \PY{n}{mean\PYZus{}absolute\PYZus{}error}\PY{p}{(}
    \PY{n}{y\PYZus{}test\PYZus{}tortilla}\PY{p}{,}
    \PY{n}{pred\PYZus{}tortilla}
\PY{p}{)}

\PY{n}{rmse\PYZus{}tortilla} \PY{o}{=} \PY{n}{np}\PY{o}{.}\PY{n}{sqrt}\PY{p}{(}
    \PY{n}{mean\PYZus{}squared\PYZus{}error}\PY{p}{(}
        \PY{n}{y\PYZus{}test\PYZus{}tortilla}\PY{p}{,}
        \PY{n}{pred\PYZus{}tortilla}
    \PY{p}{)}
\PY{p}{)}

\PY{n}{mape\PYZus{}tortilla} \PY{o}{=} \PY{p}{(}
    \PY{n}{np}\PY{o}{.}\PY{n}{mean}\PY{p}{(}
        \PY{n}{np}\PY{o}{.}\PY{n}{abs}\PY{p}{(}
            \PY{p}{(}\PY{n}{y\PYZus{}test\PYZus{}tortilla} \PY{o}{\PYZhy{}} \PY{n}{pred\PYZus{}tortilla}\PY{p}{)}
            \PY{o}{/} \PY{n}{y\PYZus{}test\PYZus{}tortilla}
        \PY{p}{)}
    \PY{p}{)}
    \PY{o}{*} \PY{l+m+mi}{100}
\PY{p}{)}

\PY{n+nb}{print}\PY{p}{(}\PY{l+s+s2}{\PYZdq{}}\PY{l+s+s2}{MAE :}\PY{l+s+s2}{\PYZdq{}}\PY{p}{,} \PY{n}{mae\PYZus{}tortilla}\PY{p}{)}
\PY{n+nb}{print}\PY{p}{(}\PY{l+s+s2}{\PYZdq{}}\PY{l+s+s2}{RMSE:}\PY{l+s+s2}{\PYZdq{}}\PY{p}{,} \PY{n}{rmse\PYZus{}tortilla}\PY{p}{)}
\PY{n+nb}{print}\PY{p}{(}\PY{l+s+s2}{\PYZdq{}}\PY{l+s+s2}{MAPE:}\PY{l+s+s2}{\PYZdq{}}\PY{p}{,} \PY{n}{mape\PYZus{}tortilla}\PY{p}{)}
\end{Verbatim}
\end{tcolorbox}

    \begin{Verbatim}[commandchars=\\\{\}]
MAE : 1.6500135174060606
RMSE: 1.8170114271777669
MAPE: 7.368909481674122
    \end{Verbatim}

    \begin{tcolorbox}[breakable, size=fbox, boxrule=1pt, pad at break*=1mm,colback=cellbackground, colframe=cellborder]
\prompt{In}{incolor}{273}{\boxspacing}
\begin{Verbatim}[commandchars=\\\{\}]
\PY{n}{plt}\PY{o}{.}\PY{n}{figure}\PY{p}{(}\PY{n}{figsize}\PY{o}{=}\PY{p}{(}\PY{l+m+mi}{14}\PY{p}{,} \PY{l+m+mi}{6}\PY{p}{)}\PY{p}{)}

\PY{n}{plt}\PY{o}{.}\PY{n}{plot}\PY{p}{(}
    \PY{n}{y\PYZus{}test\PYZus{}tortilla}\PY{o}{.}\PY{n}{index}\PY{p}{,}
    \PY{n}{y\PYZus{}test\PYZus{}tortilla}\PY{p}{,}
    \PY{n}{label}\PY{o}{=}\PY{l+s+s2}{\PYZdq{}}\PY{l+s+s2}{Precio real}\PY{l+s+s2}{\PYZdq{}}\PY{p}{,}
    \PY{n}{linewidth}\PY{o}{=}\PY{l+m+mi}{2}
\PY{p}{)}

\PY{n}{plt}\PY{o}{.}\PY{n}{plot}\PY{p}{(}
    \PY{n}{pred\PYZus{}tortilla}\PY{o}{.}\PY{n}{index}\PY{p}{,}
    \PY{n}{pred\PYZus{}tortilla}\PY{p}{,}
    \PY{n}{label}\PY{o}{=}\PY{l+s+s2}{\PYZdq{}}\PY{l+s+s2}{Predicción SARIMAX}\PY{l+s+s2}{\PYZdq{}}\PY{p}{,}
    \PY{n}{linestyle}\PY{o}{=}\PY{l+s+s2}{\PYZdq{}}\PY{l+s+s2}{\PYZhy{}\PYZhy{}}\PY{l+s+s2}{\PYZdq{}}\PY{p}{,}
    \PY{n}{linewidth}\PY{o}{=}\PY{l+m+mi}{2}
\PY{p}{)}

\PY{n}{plt}\PY{o}{.}\PY{n}{title}\PY{p}{(}\PY{l+s+s2}{\PYZdq{}}\PY{l+s+s2}{Validación del Modelo SARIMAX para Precio Nacional de la Tortilla}\PY{l+s+s2}{\PYZdq{}}\PY{p}{)}
\PY{n}{plt}\PY{o}{.}\PY{n}{xlabel}\PY{p}{(}\PY{l+s+s2}{\PYZdq{}}\PY{l+s+s2}{Fecha}\PY{l+s+s2}{\PYZdq{}}\PY{p}{)}
\PY{n}{plt}\PY{o}{.}\PY{n}{ylabel}\PY{p}{(}\PY{l+s+s2}{\PYZdq{}}\PY{l+s+s2}{MXN/Kg}\PY{l+s+s2}{\PYZdq{}}\PY{p}{)}

\PY{n}{plt}\PY{o}{.}\PY{n}{legend}\PY{p}{(}\PY{p}{)}
\PY{n}{plt}\PY{o}{.}\PY{n}{grid}\PY{p}{(}\PY{k+kc}{True}\PY{p}{)}

\PY{n}{plt}\PY{o}{.}\PY{n}{show}\PY{p}{(}\PY{p}{)}
\end{Verbatim}
\end{tcolorbox}

    \begin{center}
    \adjustimage{max size={0.9\linewidth}{0.9\paperheight}}{output_272_0.png}
    \end{center}
    { \hspace*{\fill} \\}
    
    \begin{tcolorbox}[breakable, size=fbox, boxrule=1pt, pad at break*=1mm,colback=cellbackground, colframe=cellborder]
\prompt{In}{incolor}{274}{\boxspacing}
\begin{Verbatim}[commandchars=\\\{\}]
\PY{n}{residuos\PYZus{}tortilla} \PY{o}{=} \PY{n}{resultado\PYZus{}tortilla}\PY{o}{.}\PY{n}{resid}

\PY{n}{plt}\PY{o}{.}\PY{n}{figure}\PY{p}{(}\PY{n}{figsize}\PY{o}{=}\PY{p}{(}\PY{l+m+mi}{12}\PY{p}{,} \PY{l+m+mi}{5}\PY{p}{)}\PY{p}{)}

\PY{n}{plt}\PY{o}{.}\PY{n}{plot}\PY{p}{(}\PY{n}{residuos\PYZus{}tortilla}\PY{p}{)}

\PY{n}{plt}\PY{o}{.}\PY{n}{title}\PY{p}{(}\PY{l+s+s2}{\PYZdq{}}\PY{l+s+s2}{Residuos del Modelo SARIMAX de Tortilla}\PY{l+s+s2}{\PYZdq{}}\PY{p}{)}
\PY{n}{plt}\PY{o}{.}\PY{n}{grid}\PY{p}{(}\PY{k+kc}{True}\PY{p}{)}

\PY{n}{plt}\PY{o}{.}\PY{n}{show}\PY{p}{(}\PY{p}{)}
\end{Verbatim}
\end{tcolorbox}

    \begin{center}
    \adjustimage{max size={0.9\linewidth}{0.9\paperheight}}{output_273_0.png}
    \end{center}
    { \hspace*{\fill} \\}
    
    \begin{tcolorbox}[breakable, size=fbox, boxrule=1pt, pad at break*=1mm,colback=cellbackground, colframe=cellborder]
\prompt{In}{incolor}{275}{\boxspacing}
\begin{Verbatim}[commandchars=\\\{\}]
\PY{n}{plot\PYZus{}acf}\PY{p}{(}
    \PY{n}{residuos\PYZus{}tortilla}\PY{p}{,}
    \PY{n}{lags}\PY{o}{=}\PY{l+m+mi}{36}
\PY{p}{)}

\PY{n}{plt}\PY{o}{.}\PY{n}{title}\PY{p}{(}\PY{l+s+s2}{\PYZdq{}}\PY{l+s+s2}{Autocorrelación de los Residuos}\PY{l+s+s2}{\PYZdq{}}\PY{p}{)}
\PY{n}{plt}\PY{o}{.}\PY{n}{show}\PY{p}{(}\PY{p}{)}
\end{Verbatim}
\end{tcolorbox}

    \begin{center}
    \adjustimage{max size={0.9\linewidth}{0.9\paperheight}}{output_274_0.png}
    \end{center}
    { \hspace*{\fill} \\}
    
    \begin{tcolorbox}[breakable, size=fbox, boxrule=1pt, pad at break*=1mm,colback=cellbackground, colframe=cellborder]
\prompt{In}{incolor}{276}{\boxspacing}
\begin{Verbatim}[commandchars=\\\{\}]
\PY{n}{modelo\PYZus{}tortilla\PYZus{}final} \PY{o}{=} \PY{n}{SARIMAX}\PY{p}{(}
    \PY{n}{serie\PYZus{}tortilla}\PY{p}{,}
    \PY{n}{exog}\PY{o}{=}\PY{n}{X\PYZus{}tortilla}\PY{p}{,}
    \PY{n}{order}\PY{o}{=}\PY{p}{(}\PY{l+m+mi}{1}\PY{p}{,} \PY{l+m+mi}{1}\PY{p}{,} \PY{l+m+mi}{0}\PY{p}{)}\PY{p}{,}
    \PY{n}{seasonal\PYZus{}order}\PY{o}{=}\PY{p}{(}\PY{l+m+mi}{1}\PY{p}{,} \PY{l+m+mi}{0}\PY{p}{,} \PY{l+m+mi}{0}\PY{p}{,} \PY{l+m+mi}{12}\PY{p}{)}\PY{p}{,}
    \PY{n}{enforce\PYZus{}stationarity}\PY{o}{=}\PY{k+kc}{False}\PY{p}{,}
    \PY{n}{enforce\PYZus{}invertibility}\PY{o}{=}\PY{k+kc}{False}
\PY{p}{)}

\PY{n}{resultado\PYZus{}tortilla\PYZus{}final} \PY{o}{=} \PY{n}{modelo\PYZus{}tortilla\PYZus{}final}\PY{o}{.}\PY{n}{fit}\PY{p}{(}\PY{p}{)}

\PY{n+nb}{print}\PY{p}{(}\PY{n}{resultado\PYZus{}tortilla\PYZus{}final}\PY{o}{.}\PY{n}{summary}\PY{p}{(}\PY{p}{)}\PY{p}{)}
\end{Verbatim}
\end{tcolorbox}

    \begin{Verbatim}[commandchars=\\\{\}]
/opt/conda/lib/python3.11/site-packages/statsmodels/tsa/base/tsa\_model.py:473:
ValueWarning: No frequency information was provided, so inferred frequency MS
will be used.
  self.\_init\_dates(dates, freq)
/opt/conda/lib/python3.11/site-packages/statsmodels/tsa/base/tsa\_model.py:473:
ValueWarning: No frequency information was provided, so inferred frequency MS
will be used.
  self.\_init\_dates(dates, freq)
    \end{Verbatim}

    \begin{Verbatim}[commandchars=\\\{\}]
                                     SARIMAX Results
================================================================================
==========
Dep. Variable:                    Precio\_tortilla   No. Observations:
317
Model:             SARIMAX(1, 1, 0)x(1, 0, 0, 12)   Log Likelihood
27.584
Date:                            Thu, 04 Jun 2026   AIC
-41.167
Time:                                    09:37:51   BIC
-15.171
Sample:                                01-01-2000   HQIC
-30.767
                                     - 05-01-2026
Covariance Type:                              opg
================================================================================
===================================
                                                      coef    std err          z
P>|z|      [0.025      0.975]
--------------------------------------------------------------------------------
-----------------------------------
Precio\_maizblanco\_mxpkilo\_prom\_nacional             0.3497      0.060      5.855
0.000       0.233       0.467
Costo\_transporte\_maiz\_mxpkilo\_prom\_nac              6.5995      2.191      3.013
0.003       2.306      10.893
Temperatura\_celsius\_prom\_nacional                   0.0018      0.014      0.130
0.897      -0.025       0.028
Indice\_severidad\_ponderada\_sequia\_prom\_nacional    -0.0004      0.001     -0.693
0.488      -0.002       0.001
ar.L1                                              -0.0028      0.037     -0.076
0.940      -0.075       0.070
ar.S.L12                                            0.0643      0.041      1.576
0.115      -0.016       0.144
sigma2                                              0.0488      0.002     30.648
0.000       0.046       0.052
================================================================================
===
Ljung-Box (L1) (Q):                   0.05   Jarque-Bera (JB):
55557.55
Prob(Q):                              0.83   Prob(JB):
0.00
Heteroskedasticity (H):               0.94   Skew:
4.38
Prob(H) (two-sided):                  0.74   Kurtosis:
68.76
================================================================================
===

Warnings:
[1] Covariance matrix calculated using the outer product of gradients (complex-
step).
    \end{Verbatim}

    \begin{tcolorbox}[breakable, size=fbox, boxrule=1pt, pad at break*=1mm,colback=cellbackground, colframe=cellborder]
\prompt{In}{incolor}{277}{\boxspacing}
\begin{Verbatim}[commandchars=\\\{\}]
\PY{n}{fechas\PYZus{}tortilla\PYZus{}futuras} \PY{o}{=} \PY{n}{pd}\PY{o}{.}\PY{n}{date\PYZus{}range}\PY{p}{(}
    \PY{n}{start}\PY{o}{=}\PY{n}{serie\PYZus{}tortilla}\PY{o}{.}\PY{n}{index}\PY{o}{.}\PY{n}{max}\PY{p}{(}\PY{p}{)} \PY{o}{+} \PY{n}{pd}\PY{o}{.}\PY{n}{DateOffset}\PY{p}{(}\PY{n}{months}\PY{o}{=}\PY{l+m+mi}{1}\PY{p}{)}\PY{p}{,}
    \PY{n}{periods}\PY{o}{=}\PY{l+m+mi}{12}\PY{p}{,}
    \PY{n}{freq}\PY{o}{=}\PY{l+s+s2}{\PYZdq{}}\PY{l+s+s2}{MS}\PY{l+s+s2}{\PYZdq{}}
\PY{p}{)}

\PY{n}{X\PYZus{}future\PYZus{}tortilla} \PY{o}{=} \PY{n}{pd}\PY{o}{.}\PY{n}{DataFrame}\PY{p}{(}
    \PY{n}{index}\PY{o}{=}\PY{n}{fechas\PYZus{}tortilla\PYZus{}futuras}
\PY{p}{)}

\PY{n}{X\PYZus{}future\PYZus{}tortilla}\PY{p}{[}
    \PY{l+s+s2}{\PYZdq{}}\PY{l+s+s2}{Precio\PYZus{}maizblanco\PYZus{}mxpkilo\PYZus{}prom\PYZus{}nacional}\PY{l+s+s2}{\PYZdq{}}
\PY{p}{]} \PY{o}{=} \PY{n}{precio\PYZus{}pronosticado}\PY{o}{.}\PY{n}{values}

\PY{n}{X\PYZus{}future\PYZus{}tortilla}\PY{p}{[}
    \PY{l+s+s2}{\PYZdq{}}\PY{l+s+s2}{Costo\PYZus{}transporte\PYZus{}maiz\PYZus{}mxpkilo\PYZus{}prom\PYZus{}nac}\PY{l+s+s2}{\PYZdq{}}
\PY{p}{]} \PY{o}{=} \PY{n}{X\PYZus{}tortilla}\PY{p}{[}
    \PY{l+s+s2}{\PYZdq{}}\PY{l+s+s2}{Costo\PYZus{}transporte\PYZus{}maiz\PYZus{}mxpkilo\PYZus{}prom\PYZus{}nac}\PY{l+s+s2}{\PYZdq{}}
\PY{p}{]}\PY{o}{.}\PY{n}{iloc}\PY{p}{[}\PY{o}{\PYZhy{}}\PY{l+m+mi}{1}\PY{p}{]}

\PY{n}{X\PYZus{}future\PYZus{}tortilla}\PY{p}{[}
    \PY{l+s+s2}{\PYZdq{}}\PY{l+s+s2}{Temperatura\PYZus{}celsius\PYZus{}prom\PYZus{}nacional}\PY{l+s+s2}{\PYZdq{}}
\PY{p}{]} \PY{o}{=} \PY{n}{X\PYZus{}tortilla}\PY{p}{[}
    \PY{l+s+s2}{\PYZdq{}}\PY{l+s+s2}{Temperatura\PYZus{}celsius\PYZus{}prom\PYZus{}nacional}\PY{l+s+s2}{\PYZdq{}}
\PY{p}{]}\PY{o}{.}\PY{n}{iloc}\PY{p}{[}\PY{o}{\PYZhy{}}\PY{l+m+mi}{1}\PY{p}{]}

\PY{n}{X\PYZus{}future\PYZus{}tortilla}\PY{p}{[}
    \PY{l+s+s2}{\PYZdq{}}\PY{l+s+s2}{Indice\PYZus{}severidad\PYZus{}ponderada\PYZus{}sequia\PYZus{}prom\PYZus{}nacional}\PY{l+s+s2}{\PYZdq{}}
\PY{p}{]} \PY{o}{=} \PY{n}{X\PYZus{}tortilla}\PY{p}{[}
    \PY{l+s+s2}{\PYZdq{}}\PY{l+s+s2}{Indice\PYZus{}severidad\PYZus{}ponderada\PYZus{}sequia\PYZus{}prom\PYZus{}nacional}\PY{l+s+s2}{\PYZdq{}}
\PY{p}{]}\PY{o}{.}\PY{n}{iloc}\PY{p}{[}\PY{o}{\PYZhy{}}\PY{l+m+mi}{1}\PY{p}{]}

\PY{n}{X\PYZus{}future\PYZus{}tortilla}
\end{Verbatim}
\end{tcolorbox}

            \begin{tcolorbox}[breakable, size=fbox, boxrule=.5pt, pad at break*=1mm, opacityfill=0]
\prompt{Out}{outcolor}{277}{\boxspacing}
\begin{Verbatim}[commandchars=\\\{\}]
            Precio\_maizblanco\_mxpkilo\_prom\_nacional  \textbackslash{}
2026-06-01                                 9.665025
2026-07-01                                 9.836343
2026-08-01                                 9.733932
2026-09-01                                 9.673063
2026-10-01                                 9.566773
2026-11-01                                 9.546923
2026-12-01                                 9.517849
2027-01-01                                 9.575653
2027-02-01                                 9.559805
2027-03-01                                 9.477932
2027-04-01                                 9.529332
2027-05-01                                 9.495355

            Costo\_transporte\_maiz\_mxpkilo\_prom\_nac  \textbackslash{}
2026-06-01                                    1.87
2026-07-01                                    1.87
2026-08-01                                    1.87
2026-09-01                                    1.87
2026-10-01                                    1.87
2026-11-01                                    1.87
2026-12-01                                    1.87
2027-01-01                                    1.87
2027-02-01                                    1.87
2027-03-01                                    1.87
2027-04-01                                    1.87
2027-05-01                                    1.87

            Temperatura\_celsius\_prom\_nacional  \textbackslash{}
2026-06-01                               27.3
2026-07-01                               27.3
2026-08-01                               27.3
2026-09-01                               27.3
2026-10-01                               27.3
2026-11-01                               27.3
2026-12-01                               27.3
2027-01-01                               27.3
2027-02-01                               27.3
2027-03-01                               27.3
2027-04-01                               27.3
2027-05-01                               27.3

            Indice\_severidad\_ponderada\_sequia\_prom\_nacional
2026-06-01                                           236.82
2026-07-01                                           236.82
2026-08-01                                           236.82
2026-09-01                                           236.82
2026-10-01                                           236.82
2026-11-01                                           236.82
2026-12-01                                           236.82
2027-01-01                                           236.82
2027-02-01                                           236.82
2027-03-01                                           236.82
2027-04-01                                           236.82
2027-05-01                                           236.82
\end{Verbatim}
\end{tcolorbox}
        
    \begin{tcolorbox}[breakable, size=fbox, boxrule=1pt, pad at break*=1mm,colback=cellbackground, colframe=cellborder]
\prompt{In}{incolor}{278}{\boxspacing}
\begin{Verbatim}[commandchars=\\\{\}]
\PY{n}{forecast\PYZus{}tortilla\PYZus{}final} \PY{o}{=} \PY{n}{resultado\PYZus{}tortilla\PYZus{}final}\PY{o}{.}\PY{n}{get\PYZus{}forecast}\PY{p}{(}
    \PY{n}{steps}\PY{o}{=}\PY{l+m+mi}{12}\PY{p}{,}
    \PY{n}{exog}\PY{o}{=}\PY{n}{X\PYZus{}future\PYZus{}tortilla}
\PY{p}{)}

\PY{n}{precio\PYZus{}tortilla\PYZus{}futuro} \PY{o}{=} \PY{n}{forecast\PYZus{}tortilla\PYZus{}final}\PY{o}{.}\PY{n}{predicted\PYZus{}mean}
\PY{n}{precio\PYZus{}tortilla\PYZus{}futuro}\PY{o}{.}\PY{n}{index} \PY{o}{=} \PY{n}{fechas\PYZus{}tortilla\PYZus{}futuras}

\PY{n}{intervalos\PYZus{}tortilla} \PY{o}{=} \PY{n}{forecast\PYZus{}tortilla\PYZus{}final}\PY{o}{.}\PY{n}{conf\PYZus{}int}\PY{p}{(}\PY{p}{)}
\PY{n}{intervalos\PYZus{}tortilla}\PY{o}{.}\PY{n}{index} \PY{o}{=} \PY{n}{fechas\PYZus{}tortilla\PYZus{}futuras}
\end{Verbatim}
\end{tcolorbox}

    \begin{tcolorbox}[breakable, size=fbox, boxrule=1pt, pad at break*=1mm,colback=cellbackground, colframe=cellborder]
\prompt{In}{incolor}{279}{\boxspacing}
\begin{Verbatim}[commandchars=\\\{\}]
\PY{n}{tabla\PYZus{}pronostico\PYZus{}tortilla} \PY{o}{=} \PY{n}{pd}\PY{o}{.}\PY{n}{DataFrame}\PY{p}{(}\PY{p}{\PYZob{}}
    \PY{l+s+s2}{\PYZdq{}}\PY{l+s+s2}{Precio Pronosticado Tortilla (MXN/Kg)}\PY{l+s+s2}{\PYZdq{}}\PY{p}{:} \PY{n}{precio\PYZus{}tortilla\PYZus{}futuro}\PY{o}{.}\PY{n}{round}\PY{p}{(}\PY{l+m+mi}{2}\PY{p}{)}\PY{p}{,}
    \PY{l+s+s2}{\PYZdq{}}\PY{l+s+s2}{Límite Inferior (MXN)}\PY{l+s+s2}{\PYZdq{}}\PY{p}{:} \PY{n}{intervalos\PYZus{}tortilla}\PY{o}{.}\PY{n}{iloc}\PY{p}{[}\PY{p}{:}\PY{p}{,} \PY{l+m+mi}{0}\PY{p}{]}\PY{o}{.}\PY{n}{round}\PY{p}{(}\PY{l+m+mi}{2}\PY{p}{)}\PY{p}{,}
    \PY{l+s+s2}{\PYZdq{}}\PY{l+s+s2}{Límite Superior (MXN)}\PY{l+s+s2}{\PYZdq{}}\PY{p}{:} \PY{n}{intervalos\PYZus{}tortilla}\PY{o}{.}\PY{n}{iloc}\PY{p}{[}\PY{p}{:}\PY{p}{,} \PY{l+m+mi}{1}\PY{p}{]}\PY{o}{.}\PY{n}{round}\PY{p}{(}\PY{l+m+mi}{2}\PY{p}{)}
\PY{p}{\PYZcb{}}\PY{p}{)}

\PY{n}{tabla\PYZus{}pronostico\PYZus{}tortilla}
\end{Verbatim}
\end{tcolorbox}

            \begin{tcolorbox}[breakable, size=fbox, boxrule=.5pt, pad at break*=1mm, opacityfill=0]
\prompt{Out}{outcolor}{279}{\boxspacing}
\begin{Verbatim}[commandchars=\\\{\}]
            Precio Pronosticado Tortilla (MXN/Kg)  Límite Inferior (MXN)  \textbackslash{}
2026-06-01                                  24.96                  24.53
2026-07-01                                  25.02                  24.41
2026-08-01                                  24.98                  24.23
2026-09-01                                  24.95                  24.09
2026-10-01                                  24.92                  23.95
2026-11-01                                  24.91                  23.85
2026-12-01                                  24.89                  23.75
2027-01-01                                  24.91                  23.69
2027-02-01                                  24.91                  23.62
2027-03-01                                  24.88                  23.51
2027-04-01                                  24.91                  23.47
2027-05-01                                  24.93                  23.44

            Límite Superior (MXN)
2026-06-01                  25.40
2026-07-01                  25.63
2026-08-01                  25.73
2026-09-01                  25.82
2026-10-01                  25.88
2026-11-01                  25.97
2026-12-01                  26.04
2027-01-01                  26.14
2027-02-01                  26.21
2027-03-01                  26.24
2027-04-01                  26.34
2027-05-01                  26.43
\end{Verbatim}
\end{tcolorbox}
        
    \begin{tcolorbox}[breakable, size=fbox, boxrule=1pt, pad at break*=1mm,colback=cellbackground, colframe=cellborder]
\prompt{In}{incolor}{280}{\boxspacing}
\begin{Verbatim}[commandchars=\\\{\}]
\PY{n}{plt}\PY{o}{.}\PY{n}{figure}\PY{p}{(}\PY{n}{figsize}\PY{o}{=}\PY{p}{(}\PY{l+m+mi}{14}\PY{p}{,} \PY{l+m+mi}{7}\PY{p}{)}\PY{p}{)}

\PY{n}{historico\PYZus{}tortilla\PYZus{}reciente} \PY{o}{=} \PY{n}{serie\PYZus{}tortilla}\PY{o}{.}\PY{n}{loc}\PY{p}{[}\PY{l+s+s2}{\PYZdq{}}\PY{l+s+s2}{2023\PYZhy{}01\PYZhy{}01}\PY{l+s+s2}{\PYZdq{}}\PY{p}{:}\PY{p}{]}

\PY{n}{plt}\PY{o}{.}\PY{n}{plot}\PY{p}{(}
    \PY{n}{historico\PYZus{}tortilla\PYZus{}reciente}\PY{o}{.}\PY{n}{index}\PY{p}{,}
    \PY{n}{historico\PYZus{}tortilla\PYZus{}reciente}\PY{p}{,}
    \PY{n}{label}\PY{o}{=}\PY{l+s+s2}{\PYZdq{}}\PY{l+s+s2}{Precio Histórico de la Tortilla}\PY{l+s+s2}{\PYZdq{}}\PY{p}{,}
    \PY{n}{linewidth}\PY{o}{=}\PY{l+m+mi}{2}
\PY{p}{)}

\PY{n}{plt}\PY{o}{.}\PY{n}{plot}\PY{p}{(}
    \PY{n}{precio\PYZus{}tortilla\PYZus{}futuro}\PY{o}{.}\PY{n}{index}\PY{p}{,}
    \PY{n}{precio\PYZus{}tortilla\PYZus{}futuro}\PY{p}{,}
    \PY{n}{label}\PY{o}{=}\PY{l+s+s2}{\PYZdq{}}\PY{l+s+s2}{Pronóstico SARIMAX Tortilla}\PY{l+s+s2}{\PYZdq{}}\PY{p}{,}
    \PY{n}{linestyle}\PY{o}{=}\PY{l+s+s2}{\PYZdq{}}\PY{l+s+s2}{\PYZhy{}\PYZhy{}}\PY{l+s+s2}{\PYZdq{}}\PY{p}{,}
    \PY{n}{linewidth}\PY{o}{=}\PY{l+m+mf}{2.5}
\PY{p}{)}

\PY{n}{plt}\PY{o}{.}\PY{n}{fill\PYZus{}between}\PY{p}{(}
    \PY{n}{intervalos\PYZus{}tortilla}\PY{o}{.}\PY{n}{index}\PY{p}{,}
    \PY{n}{intervalos\PYZus{}tortilla}\PY{o}{.}\PY{n}{iloc}\PY{p}{[}\PY{p}{:}\PY{p}{,} \PY{l+m+mi}{0}\PY{p}{]}\PY{p}{,}
    \PY{n}{intervalos\PYZus{}tortilla}\PY{o}{.}\PY{n}{iloc}\PY{p}{[}\PY{p}{:}\PY{p}{,} \PY{l+m+mi}{1}\PY{p}{]}\PY{p}{,}
    \PY{n}{alpha}\PY{o}{=}\PY{l+m+mf}{0.2}\PY{p}{,}
    \PY{n}{label}\PY{o}{=}\PY{l+s+s2}{\PYZdq{}}\PY{l+s+s2}{Intervalo de Confianza 95}\PY{l+s+s2}{\PYZpc{}}\PY{l+s+s2}{\PYZdq{}}
\PY{p}{)}

\PY{n}{plt}\PY{o}{.}\PY{n}{axvline}\PY{p}{(}
    \PY{n}{serie\PYZus{}tortilla}\PY{o}{.}\PY{n}{index}\PY{o}{.}\PY{n}{max}\PY{p}{(}\PY{p}{)}\PY{p}{,}
    \PY{n}{linestyle}\PY{o}{=}\PY{l+s+s2}{\PYZdq{}}\PY{l+s+s2}{:}\PY{l+s+s2}{\PYZdq{}}\PY{p}{,}
    \PY{n}{linewidth}\PY{o}{=}\PY{l+m+mi}{2}\PY{p}{,}
    \PY{n}{label}\PY{o}{=}\PY{l+s+sa}{f}\PY{l+s+s2}{\PYZdq{}}\PY{l+s+s2}{Frontera de datos actuales (}\PY{l+s+si}{\PYZob{}}\PY{n}{serie\PYZus{}tortilla}\PY{o}{.}\PY{n}{index}\PY{o}{.}\PY{n}{max}\PY{p}{(}\PY{p}{)}\PY{o}{.}\PY{n}{strftime}\PY{p}{(}\PY{l+s+s1}{\PYZsq{}}\PY{l+s+s1}{\PYZpc{}}\PY{l+s+s1}{B }\PY{l+s+s1}{\PYZpc{}}\PY{l+s+s1}{Y}\PY{l+s+s1}{\PYZsq{}}\PY{p}{)}\PY{l+s+si}{\PYZcb{}}\PY{l+s+s2}{)}\PY{l+s+s2}{\PYZdq{}}
\PY{p}{)}

\PY{n}{plt}\PY{o}{.}\PY{n}{title}\PY{p}{(}
    \PY{l+s+s2}{\PYZdq{}}\PY{l+s+s2}{Proyección del Precio Nacional de la Tortilla}\PY{l+s+s2}{\PYZdq{}}\PY{p}{,}
    \PY{n}{fontsize}\PY{o}{=}\PY{l+m+mi}{14}\PY{p}{,}
    \PY{n}{fontweight}\PY{o}{=}\PY{l+s+s2}{\PYZdq{}}\PY{l+s+s2}{bold}\PY{l+s+s2}{\PYZdq{}}
\PY{p}{)}

\PY{n}{plt}\PY{o}{.}\PY{n}{xlabel}\PY{p}{(}\PY{l+s+s2}{\PYZdq{}}\PY{l+s+s2}{Fecha}\PY{l+s+s2}{\PYZdq{}}\PY{p}{)}
\PY{n}{plt}\PY{o}{.}\PY{n}{ylabel}\PY{p}{(}\PY{l+s+s2}{\PYZdq{}}\PY{l+s+s2}{MXN/Kg}\PY{l+s+s2}{\PYZdq{}}\PY{p}{)}

\PY{n}{plt}\PY{o}{.}\PY{n}{legend}\PY{p}{(}\PY{p}{)}
\PY{n}{plt}\PY{o}{.}\PY{n}{grid}\PY{p}{(}\PY{k+kc}{True}\PY{p}{)}

\PY{n}{plt}\PY{o}{.}\PY{n}{show}\PY{p}{(}\PY{p}{)}
\end{Verbatim}
\end{tcolorbox}

    \begin{center}
    \adjustimage{max size={0.9\linewidth}{0.9\paperheight}}{output_279_0.png}
    \end{center}
    { \hspace*{\fill} \\}
    
    \begin{tcolorbox}[breakable, size=fbox, boxrule=1pt, pad at break*=1mm,colback=cellbackground, colframe=cellborder]
\prompt{In}{incolor}{281}{\boxspacing}
\begin{Verbatim}[commandchars=\\\{\}]
\PY{n}{residuos\PYZus{}tortilla}\PY{o}{.}\PY{n}{describe}\PY{p}{(}\PY{p}{)}
\end{Verbatim}
\end{tcolorbox}

            \begin{tcolorbox}[breakable, size=fbox, boxrule=.5pt, pad at break*=1mm, opacityfill=0]
\prompt{Out}{outcolor}{281}{\boxspacing}
\begin{Verbatim}[commandchars=\\\{\}]
count    253.000000
mean       0.020199
std        0.221706
min       -1.385263
25\%       -0.039654
50\%        0.019637
75\%        0.057507
max        1.997657
dtype: float64
\end{Verbatim}
\end{tcolorbox}
        
    \begin{tcolorbox}[breakable, size=fbox, boxrule=1pt, pad at break*=1mm,colback=cellbackground, colframe=cellborder]
\prompt{In}{incolor}{285}{\boxspacing}
\begin{Verbatim}[commandchars=\\\{\}]
\PY{k+kn}{from} \PY{n+nn}{statsmodels}\PY{n+nn}{.}\PY{n+nn}{stats}\PY{n+nn}{.}\PY{n+nn}{diagnostic} \PY{k+kn}{import} \PY{n}{acorr\PYZus{}ljungbox}

\PY{n}{ljung} \PY{o}{=} \PY{n}{acorr\PYZus{}ljungbox}\PY{p}{(}
    \PY{n}{residuos\PYZus{}tortilla}\PY{p}{,}
    \PY{n}{lags}\PY{o}{=}\PY{p}{[}\PY{l+m+mi}{12}\PY{p}{]}\PY{p}{,}
    \PY{n}{return\PYZus{}df}\PY{o}{=}\PY{k+kc}{True}
\PY{p}{)}

\PY{n}{ljung}
\end{Verbatim}
\end{tcolorbox}

            \begin{tcolorbox}[breakable, size=fbox, boxrule=.5pt, pad at break*=1mm, opacityfill=0]
\prompt{Out}{outcolor}{285}{\boxspacing}
\begin{Verbatim}[commandchars=\\\{\}]
     lb\_stat  lb\_pvalue
12  4.145355   0.980668
\end{Verbatim}
\end{tcolorbox}
        
    \begin{tcolorbox}[breakable, size=fbox, boxrule=1pt, pad at break*=1mm,colback=cellbackground, colframe=cellborder]
\prompt{In}{incolor}{287}{\boxspacing}
\begin{Verbatim}[commandchars=\\\{\}]
\PY{n}{residuos\PYZus{}tortilla\PYZus{}final} \PY{o}{=} \PY{n}{resultado\PYZus{}tortilla\PYZus{}final}\PY{o}{.}\PY{n}{resid}

\PY{n}{media\PYZus{}res} \PY{o}{=} \PY{n}{residuos\PYZus{}tortilla\PYZus{}final}\PY{o}{.}\PY{n}{mean}\PY{p}{(}\PY{p}{)}
\PY{n}{std\PYZus{}res} \PY{o}{=} \PY{n}{residuos\PYZus{}tortilla\PYZus{}final}\PY{o}{.}\PY{n}{std}\PY{p}{(}\PY{p}{)}
\end{Verbatim}
\end{tcolorbox}

    \begin{tcolorbox}[breakable, size=fbox, boxrule=1pt, pad at break*=1mm,colback=cellbackground, colframe=cellborder]
\prompt{In}{incolor}{288}{\boxspacing}
\begin{Verbatim}[commandchars=\\\{\}]
\PY{k+kn}{from} \PY{n+nn}{statsmodels}\PY{n+nn}{.}\PY{n+nn}{graphics}\PY{n+nn}{.}\PY{n+nn}{tsaplots} \PY{k+kn}{import} \PY{n}{plot\PYZus{}acf}
\PY{k+kn}{from} \PY{n+nn}{scipy}\PY{n+nn}{.}\PY{n+nn}{stats} \PY{k+kn}{import} \PY{n}{norm}\PY{p}{,} \PY{n}{probplot}

\PY{n}{fig}\PY{p}{,} \PY{n}{ax} \PY{o}{=} \PY{n}{plt}\PY{o}{.}\PY{n}{subplots}\PY{p}{(}
    \PY{l+m+mi}{2}\PY{p}{,} \PY{l+m+mi}{2}\PY{p}{,}
    \PY{n}{figsize}\PY{o}{=}\PY{p}{(}\PY{l+m+mi}{15}\PY{p}{,}\PY{l+m+mi}{10}\PY{p}{)}
\PY{p}{)}

\PY{n}{fig}\PY{o}{.}\PY{n}{suptitle}\PY{p}{(}
    \PY{l+s+s2}{\PYZdq{}}\PY{l+s+s2}{Diagnóstico Estadístico de los Residuos (Modelo SARIMAX Tortilla)}\PY{l+s+s2}{\PYZdq{}}\PY{p}{,}
    \PY{n}{fontsize}\PY{o}{=}\PY{l+m+mi}{20}\PY{p}{,}
    \PY{n}{fontweight}\PY{o}{=}\PY{l+s+s2}{\PYZdq{}}\PY{l+s+s2}{bold}\PY{l+s+s2}{\PYZdq{}}
\PY{p}{)}

\PY{c+c1}{\PYZsh{} \PYZhy{}\PYZhy{}\PYZhy{}\PYZhy{}\PYZhy{}\PYZhy{}\PYZhy{}\PYZhy{}\PYZhy{}\PYZhy{}\PYZhy{}\PYZhy{}\PYZhy{}\PYZhy{}\PYZhy{}\PYZhy{}\PYZhy{}\PYZhy{}\PYZhy{}\PYZhy{}\PYZhy{}\PYZhy{}\PYZhy{}\PYZhy{}\PYZhy{}\PYZhy{}\PYZhy{}\PYZhy{}\PYZhy{}\PYZhy{}\PYZhy{}\PYZhy{}\PYZhy{}\PYZhy{}\PYZhy{}\PYZhy{}\PYZhy{}\PYZhy{}\PYZhy{}\PYZhy{}\PYZhy{}\PYZhy{}\PYZhy{}\PYZhy{}\PYZhy{}\PYZhy{}\PYZhy{}\PYZhy{}\PYZhy{}\PYZhy{}}
\PY{c+c1}{\PYZsh{} Evolución temporal}
\PY{c+c1}{\PYZsh{} \PYZhy{}\PYZhy{}\PYZhy{}\PYZhy{}\PYZhy{}\PYZhy{}\PYZhy{}\PYZhy{}\PYZhy{}\PYZhy{}\PYZhy{}\PYZhy{}\PYZhy{}\PYZhy{}\PYZhy{}\PYZhy{}\PYZhy{}\PYZhy{}\PYZhy{}\PYZhy{}\PYZhy{}\PYZhy{}\PYZhy{}\PYZhy{}\PYZhy{}\PYZhy{}\PYZhy{}\PYZhy{}\PYZhy{}\PYZhy{}\PYZhy{}\PYZhy{}\PYZhy{}\PYZhy{}\PYZhy{}\PYZhy{}\PYZhy{}\PYZhy{}\PYZhy{}\PYZhy{}\PYZhy{}\PYZhy{}\PYZhy{}\PYZhy{}\PYZhy{}\PYZhy{}\PYZhy{}\PYZhy{}\PYZhy{}\PYZhy{}}

\PY{n}{ax}\PY{p}{[}\PY{l+m+mi}{0}\PY{p}{,}\PY{l+m+mi}{0}\PY{p}{]}\PY{o}{.}\PY{n}{plot}\PY{p}{(}
    \PY{n}{residuos\PYZus{}tortilla\PYZus{}final}\PY{p}{,}
    \PY{n}{color}\PY{o}{=}\PY{l+s+s2}{\PYZdq{}}\PY{l+s+s2}{navy}\PY{l+s+s2}{\PYZdq{}}
\PY{p}{)}

\PY{n}{ax}\PY{p}{[}\PY{l+m+mi}{0}\PY{p}{,}\PY{l+m+mi}{0}\PY{p}{]}\PY{o}{.}\PY{n}{axhline}\PY{p}{(}
    \PY{l+m+mi}{0}\PY{p}{,}
    \PY{n}{color}\PY{o}{=}\PY{l+s+s2}{\PYZdq{}}\PY{l+s+s2}{red}\PY{l+s+s2}{\PYZdq{}}\PY{p}{,}
    \PY{n}{linestyle}\PY{o}{=}\PY{l+s+s2}{\PYZdq{}}\PY{l+s+s2}{\PYZhy{}\PYZhy{}}\PY{l+s+s2}{\PYZdq{}}
\PY{p}{)}

\PY{n}{ax}\PY{p}{[}\PY{l+m+mi}{0}\PY{p}{,}\PY{l+m+mi}{0}\PY{p}{]}\PY{o}{.}\PY{n}{set\PYZus{}title}\PY{p}{(}
    \PY{l+s+s2}{\PYZdq{}}\PY{l+s+s2}{Evolución Temporal de los Residuos}\PY{l+s+s2}{\PYZdq{}}
\PY{p}{)}

\PY{n}{ax}\PY{p}{[}\PY{l+m+mi}{0}\PY{p}{,}\PY{l+m+mi}{0}\PY{p}{]}\PY{o}{.}\PY{n}{set\PYZus{}ylabel}\PY{p}{(}
    \PY{l+s+s2}{\PYZdq{}}\PY{l+s+s2}{Error (MXN)}\PY{l+s+s2}{\PYZdq{}}
\PY{p}{)}

\PY{c+c1}{\PYZsh{} \PYZhy{}\PYZhy{}\PYZhy{}\PYZhy{}\PYZhy{}\PYZhy{}\PYZhy{}\PYZhy{}\PYZhy{}\PYZhy{}\PYZhy{}\PYZhy{}\PYZhy{}\PYZhy{}\PYZhy{}\PYZhy{}\PYZhy{}\PYZhy{}\PYZhy{}\PYZhy{}\PYZhy{}\PYZhy{}\PYZhy{}\PYZhy{}\PYZhy{}\PYZhy{}\PYZhy{}\PYZhy{}\PYZhy{}\PYZhy{}\PYZhy{}\PYZhy{}\PYZhy{}\PYZhy{}\PYZhy{}\PYZhy{}\PYZhy{}\PYZhy{}\PYZhy{}\PYZhy{}\PYZhy{}\PYZhy{}\PYZhy{}\PYZhy{}\PYZhy{}\PYZhy{}\PYZhy{}\PYZhy{}\PYZhy{}\PYZhy{}}
\PY{c+c1}{\PYZsh{} Histograma + normal teórica}
\PY{c+c1}{\PYZsh{} \PYZhy{}\PYZhy{}\PYZhy{}\PYZhy{}\PYZhy{}\PYZhy{}\PYZhy{}\PYZhy{}\PYZhy{}\PYZhy{}\PYZhy{}\PYZhy{}\PYZhy{}\PYZhy{}\PYZhy{}\PYZhy{}\PYZhy{}\PYZhy{}\PYZhy{}\PYZhy{}\PYZhy{}\PYZhy{}\PYZhy{}\PYZhy{}\PYZhy{}\PYZhy{}\PYZhy{}\PYZhy{}\PYZhy{}\PYZhy{}\PYZhy{}\PYZhy{}\PYZhy{}\PYZhy{}\PYZhy{}\PYZhy{}\PYZhy{}\PYZhy{}\PYZhy{}\PYZhy{}\PYZhy{}\PYZhy{}\PYZhy{}\PYZhy{}\PYZhy{}\PYZhy{}\PYZhy{}\PYZhy{}\PYZhy{}\PYZhy{}}

\PY{n}{ax}\PY{p}{[}\PY{l+m+mi}{0}\PY{p}{,}\PY{l+m+mi}{1}\PY{p}{]}\PY{o}{.}\PY{n}{hist}\PY{p}{(}
    \PY{n}{residuos\PYZus{}tortilla\PYZus{}final}\PY{p}{,}
    \PY{n}{bins}\PY{o}{=}\PY{l+m+mi}{40}\PY{p}{,}
    \PY{n}{density}\PY{o}{=}\PY{k+kc}{True}\PY{p}{,}
    \PY{n}{alpha}\PY{o}{=}\PY{l+m+mf}{0.7}\PY{p}{,}
    \PY{n}{color}\PY{o}{=}\PY{l+s+s2}{\PYZdq{}}\PY{l+s+s2}{forestgreen}\PY{l+s+s2}{\PYZdq{}}\PY{p}{,}
    \PY{n}{edgecolor}\PY{o}{=}\PY{l+s+s2}{\PYZdq{}}\PY{l+s+s2}{black}\PY{l+s+s2}{\PYZdq{}}
\PY{p}{)}

\PY{n}{x} \PY{o}{=} \PY{n}{np}\PY{o}{.}\PY{n}{linspace}\PY{p}{(}
    \PY{n}{residuos\PYZus{}tortilla\PYZus{}final}\PY{o}{.}\PY{n}{min}\PY{p}{(}\PY{p}{)}\PY{p}{,}
    \PY{n}{residuos\PYZus{}tortilla\PYZus{}final}\PY{o}{.}\PY{n}{max}\PY{p}{(}\PY{p}{)}\PY{p}{,}
    \PY{l+m+mi}{100}
\PY{p}{)}

\PY{n}{ax}\PY{p}{[}\PY{l+m+mi}{0}\PY{p}{,}\PY{l+m+mi}{1}\PY{p}{]}\PY{o}{.}\PY{n}{plot}\PY{p}{(}
    \PY{n}{x}\PY{p}{,}
    \PY{n}{norm}\PY{o}{.}\PY{n}{pdf}\PY{p}{(}\PY{n}{x}\PY{p}{,} \PY{n}{media\PYZus{}res}\PY{p}{,} \PY{n}{std\PYZus{}res}\PY{p}{)}\PY{p}{,}
    \PY{n}{color}\PY{o}{=}\PY{l+s+s2}{\PYZdq{}}\PY{l+s+s2}{red}\PY{l+s+s2}{\PYZdq{}}\PY{p}{,}
    \PY{n}{linewidth}\PY{o}{=}\PY{l+m+mi}{2}\PY{p}{,}
    \PY{n}{label}\PY{o}{=}\PY{l+s+s2}{\PYZdq{}}\PY{l+s+s2}{Normal Teórica}\PY{l+s+s2}{\PYZdq{}}
\PY{p}{)}

\PY{n}{ax}\PY{p}{[}\PY{l+m+mi}{0}\PY{p}{,}\PY{l+m+mi}{1}\PY{p}{]}\PY{o}{.}\PY{n}{set\PYZus{}title}\PY{p}{(}
    \PY{l+s+s2}{\PYZdq{}}\PY{l+s+s2}{Distribución de los Residuos vs. Curva Normal}\PY{l+s+s2}{\PYZdq{}}
\PY{p}{)}

\PY{n}{ax}\PY{p}{[}\PY{l+m+mi}{0}\PY{p}{,}\PY{l+m+mi}{1}\PY{p}{]}\PY{o}{.}\PY{n}{legend}\PY{p}{(}\PY{p}{)}

\PY{c+c1}{\PYZsh{} \PYZhy{}\PYZhy{}\PYZhy{}\PYZhy{}\PYZhy{}\PYZhy{}\PYZhy{}\PYZhy{}\PYZhy{}\PYZhy{}\PYZhy{}\PYZhy{}\PYZhy{}\PYZhy{}\PYZhy{}\PYZhy{}\PYZhy{}\PYZhy{}\PYZhy{}\PYZhy{}\PYZhy{}\PYZhy{}\PYZhy{}\PYZhy{}\PYZhy{}\PYZhy{}\PYZhy{}\PYZhy{}\PYZhy{}\PYZhy{}\PYZhy{}\PYZhy{}\PYZhy{}\PYZhy{}\PYZhy{}\PYZhy{}\PYZhy{}\PYZhy{}\PYZhy{}\PYZhy{}\PYZhy{}\PYZhy{}\PYZhy{}\PYZhy{}\PYZhy{}\PYZhy{}\PYZhy{}\PYZhy{}\PYZhy{}\PYZhy{}}
\PY{c+c1}{\PYZsh{} ACF}
\PY{c+c1}{\PYZsh{} \PYZhy{}\PYZhy{}\PYZhy{}\PYZhy{}\PYZhy{}\PYZhy{}\PYZhy{}\PYZhy{}\PYZhy{}\PYZhy{}\PYZhy{}\PYZhy{}\PYZhy{}\PYZhy{}\PYZhy{}\PYZhy{}\PYZhy{}\PYZhy{}\PYZhy{}\PYZhy{}\PYZhy{}\PYZhy{}\PYZhy{}\PYZhy{}\PYZhy{}\PYZhy{}\PYZhy{}\PYZhy{}\PYZhy{}\PYZhy{}\PYZhy{}\PYZhy{}\PYZhy{}\PYZhy{}\PYZhy{}\PYZhy{}\PYZhy{}\PYZhy{}\PYZhy{}\PYZhy{}\PYZhy{}\PYZhy{}\PYZhy{}\PYZhy{}\PYZhy{}\PYZhy{}\PYZhy{}\PYZhy{}\PYZhy{}\PYZhy{}}

\PY{n}{plot\PYZus{}acf}\PY{p}{(}
    \PY{n}{residuos\PYZus{}tortilla\PYZus{}final}\PY{p}{,}
    \PY{n}{lags}\PY{o}{=}\PY{l+m+mi}{24}\PY{p}{,}
    \PY{n}{ax}\PY{o}{=}\PY{n}{ax}\PY{p}{[}\PY{l+m+mi}{1}\PY{p}{,}\PY{l+m+mi}{0}\PY{p}{]}
\PY{p}{)}

\PY{n}{ax}\PY{p}{[}\PY{l+m+mi}{1}\PY{p}{,}\PY{l+m+mi}{0}\PY{p}{]}\PY{o}{.}\PY{n}{set\PYZus{}title}\PY{p}{(}
    \PY{l+s+s2}{\PYZdq{}}\PY{l+s+s2}{Función de Autocorrelación (ACF) de los Residuos}\PY{l+s+s2}{\PYZdq{}}
\PY{p}{)}

\PY{c+c1}{\PYZsh{} \PYZhy{}\PYZhy{}\PYZhy{}\PYZhy{}\PYZhy{}\PYZhy{}\PYZhy{}\PYZhy{}\PYZhy{}\PYZhy{}\PYZhy{}\PYZhy{}\PYZhy{}\PYZhy{}\PYZhy{}\PYZhy{}\PYZhy{}\PYZhy{}\PYZhy{}\PYZhy{}\PYZhy{}\PYZhy{}\PYZhy{}\PYZhy{}\PYZhy{}\PYZhy{}\PYZhy{}\PYZhy{}\PYZhy{}\PYZhy{}\PYZhy{}\PYZhy{}\PYZhy{}\PYZhy{}\PYZhy{}\PYZhy{}\PYZhy{}\PYZhy{}\PYZhy{}\PYZhy{}\PYZhy{}\PYZhy{}\PYZhy{}\PYZhy{}\PYZhy{}\PYZhy{}\PYZhy{}\PYZhy{}\PYZhy{}\PYZhy{}}
\PY{c+c1}{\PYZsh{} QQ Plot}
\PY{c+c1}{\PYZsh{} \PYZhy{}\PYZhy{}\PYZhy{}\PYZhy{}\PYZhy{}\PYZhy{}\PYZhy{}\PYZhy{}\PYZhy{}\PYZhy{}\PYZhy{}\PYZhy{}\PYZhy{}\PYZhy{}\PYZhy{}\PYZhy{}\PYZhy{}\PYZhy{}\PYZhy{}\PYZhy{}\PYZhy{}\PYZhy{}\PYZhy{}\PYZhy{}\PYZhy{}\PYZhy{}\PYZhy{}\PYZhy{}\PYZhy{}\PYZhy{}\PYZhy{}\PYZhy{}\PYZhy{}\PYZhy{}\PYZhy{}\PYZhy{}\PYZhy{}\PYZhy{}\PYZhy{}\PYZhy{}\PYZhy{}\PYZhy{}\PYZhy{}\PYZhy{}\PYZhy{}\PYZhy{}\PYZhy{}\PYZhy{}\PYZhy{}\PYZhy{}}

\PY{n}{probplot}\PY{p}{(}
    \PY{n}{residuos\PYZus{}tortilla\PYZus{}final}\PY{p}{,}
    \PY{n}{dist}\PY{o}{=}\PY{l+s+s2}{\PYZdq{}}\PY{l+s+s2}{norm}\PY{l+s+s2}{\PYZdq{}}\PY{p}{,}
    \PY{n}{plot}\PY{o}{=}\PY{n}{ax}\PY{p}{[}\PY{l+m+mi}{1}\PY{p}{,}\PY{l+m+mi}{1}\PY{p}{]}
\PY{p}{)}

\PY{n}{ax}\PY{p}{[}\PY{l+m+mi}{1}\PY{p}{,}\PY{l+m+mi}{1}\PY{p}{]}\PY{o}{.}\PY{n}{set\PYZus{}title}\PY{p}{(}
    \PY{l+s+s2}{\PYZdq{}}\PY{l+s+s2}{Gráfico Q\PYZhy{}Q de Normalidad}\PY{l+s+s2}{\PYZdq{}}
\PY{p}{)}

\PY{n}{plt}\PY{o}{.}\PY{n}{tight\PYZus{}layout}\PY{p}{(}\PY{p}{)}

\PY{n}{plt}\PY{o}{.}\PY{n}{show}\PY{p}{(}\PY{p}{)}
\end{Verbatim}
\end{tcolorbox}

    \begin{center}
    \adjustimage{max size={0.9\linewidth}{0.9\paperheight}}{output_283_0.png}
    \end{center}
    { \hspace*{\fill} \\}
    
    El objetivo principal de este proyecto consistió en analizar y
pronosticar el comportamiento del precio nacional del maíz blanco en
México. Sin embargo, durante el desarrollo de la investigación surgió
una línea adicional de análisis derivada de la fuerte relación observada
entre el precio del maíz y el precio nacional de la tortilla.

La tortilla constituye el principal producto alimenticio derivado del
maíz dentro de la dieta mexicana, por lo que resulta relevante explorar
si los factores que explican la evolución del precio del maíz también
permiten modelar el comportamiento de los precios al consumidor final.

Esta línea de investigación no forma parte del objetivo central del
proyecto, pero representa una extensión natural del modelo desarrollado.

\begin{center}\rule{0.5\linewidth}{0.5pt}\end{center}

\hypertarget{relaciuxf3n-entre-el-precio-del-mauxedz-y-la-tortilla}{%
\subsection{Relación entre el Precio del Maíz y la
Tortilla}\label{relaciuxf3n-entre-el-precio-del-mauxedz-y-la-tortilla}}

El análisis de correlación mostró una asociación positiva muy elevada
entre ambas variables.

\begin{longtable}[]{@{}lr@{}}
\toprule\noalign{}
Variable & Correlación \\
\midrule\noalign{}
\endhead
\bottomrule\noalign{}
\endlastfoot
Precio del maíz blanco nacional & 0.972 \\
Costo de transporte & 0.981 \\
Precio internacional del maíz (CBOT) & 0.811 \\
Índice de severidad de sequía & 0.443 \\
Temperatura promedio nacional & 0.070 \\
\end{longtable}

Los resultados sugieren que los movimientos observados en el precio de
la tortilla guardan una estrecha relación con la evolución de los costos
de producción y distribución del maíz.

\begin{center}\rule{0.5\linewidth}{0.5pt}\end{center}

\hypertarget{estacionariedad-y-estructura-temporal}{%
\subsection{Estacionariedad y Estructura
Temporal}\label{estacionariedad-y-estructura-temporal}}

La prueba Dickey-Fuller aumentada (ADF) aplicada a la serie original del
precio de la tortilla arrojó:

\begin{itemize}
\tightlist
\item
  Estadístico ADF = 1.778
\item
  p-value = 0.998
\end{itemize}

Lo anterior indica que la serie original no es estacionaria.

Posteriormente se aplicó una diferenciación de primer orden, obteniendo:

\begin{itemize}
\tightlist
\item
  Estadístico ADF = -16.612
\item
  p-value \textless{} 0.001
\end{itemize}

Estos resultados confirman que la serie diferenciada es estacionaria y
cumple los supuestos requeridos para la construcción de modelos SARIMAX.

\begin{center}\rule{0.5\linewidth}{0.5pt}\end{center}

\hypertarget{modelo-sarimax-para-el-precio-nacional-de-la-tortilla}{%
\subsection{Modelo SARIMAX para el Precio Nacional de la
Tortilla}\label{modelo-sarimax-para-el-precio-nacional-de-la-tortilla}}

Se estimó un modelo SARIMAX con variables exógenas relacionadas con la
cadena de suministro agrícola:

\begin{itemize}
\tightlist
\item
  Precio nacional del maíz blanco.
\item
  Costo de transporte del maíz.
\item
  Temperatura promedio nacional.
\item
  Índice de severidad de sequía.
\end{itemize}

Los coeficientes obtenidos en el modelo final muestran que las variables
con mayor capacidad explicativa son:

\begin{longtable}[]{@{}lrr@{}}
\toprule\noalign{}
Variable & Coeficiente & Significancia \\
\midrule\noalign{}
\endhead
\bottomrule\noalign{}
\endlastfoot
Precio nacional del maíz blanco & 0.350 & p \textless{} 0.001 \\
Costo de transporte & 6.600 & p = 0.003 \\
\end{longtable}

Las variables climáticas no mostraron significancia estadística directa
dentro del modelo, sugiriendo que su efecto podría encontrarse mediado
por el comportamiento del mercado agrícola y los costos logísticos.

\begin{center}\rule{0.5\linewidth}{0.5pt}\end{center}

\hypertarget{evaluaciuxf3n-del-modelo}{%
\subsection{Evaluación del Modelo}\label{evaluaciuxf3n-del-modelo}}

Las métricas obtenidas durante la etapa de validación fueron:

\begin{longtable}[]{@{}lr@{}}
\toprule\noalign{}
Métrica & Valor \\
\midrule\noalign{}
\endhead
\bottomrule\noalign{}
\endlastfoot
MAE & 1.650 \\
RMSE & 1.817 \\
MAPE & 7.37 \% \\
\end{longtable}

Un error porcentual absoluto medio cercano al 7 \% indica una capacidad
predictiva aceptable para una serie económica de largo plazo.

\begin{center}\rule{0.5\linewidth}{0.5pt}\end{center}

\hypertarget{diagnuxf3stico-de-residuos}{%
\subsection{Diagnóstico de Residuos}\label{diagnuxf3stico-de-residuos}}

Los residuos generados por el modelo fueron sometidos a pruebas
estadísticas con el fin de verificar que no permaneciera información
temporal sin explicar.

Los principales hallazgos fueron:

\begin{itemize}
\tightlist
\item
  Media residual cercana a cero.
\item
  Ausencia de patrones sistemáticos en el tiempo.
\item
  Autocorrelaciones residuales dentro de las bandas de confianza.
\item
  Prueba de Ljung-Box con p-value = 0.981.
\end{itemize}

Estos resultados sugieren que los residuos se comportan como ruido
blanco y que el modelo captura adecuadamente la estructura temporal de
la serie.

Aunque las pruebas de normalidad muestran desviaciones en las colas de
la distribución, este comportamiento es común en series económicas
sujetas a choques externos extraordinarios.

\begin{center}\rule{0.5\linewidth}{0.5pt}\end{center}

\hypertarget{proyecciuxf3n-exploratoria-del-precio-nacional-de-la-tortilla-20262027}{%
\subsection{Proyección Exploratoria del Precio Nacional de la Tortilla
(2026--2027)}\label{proyecciuxf3n-exploratoria-del-precio-nacional-de-la-tortilla-20262027}}

Utilizando los pronósticos obtenidos previamente para el precio nacional
del maíz blanco, se generó una proyección exploratoria del precio
nacional de la tortilla para el periodo comprendido entre junio de 2026
y mayo de 2027.

Los resultados sugieren un escenario de estabilidad relativa, con
precios esperados cercanos a:

\begin{itemize}
\tightlist
\item
  Valor mínimo proyectado: 24.88 MXN/kg.
\item
  Valor máximo proyectado: 25.02 MXN/kg.
\end{itemize}

El intervalo de confianza al 95 \% indica que podrían presentarse
fluctuaciones moderadas alrededor de dichos valores, principalmente
derivadas de incertidumbre económica, logística o climática.

\begin{center}\rule{0.5\linewidth}{0.5pt}\end{center}

\hypertarget{implicaciones-para-la-cadena-agroalimentaria}{%
\subsection{Implicaciones para la Cadena
Agroalimentaria}\label{implicaciones-para-la-cadena-agroalimentaria}}

Los resultados obtenidos permiten plantear la hipótesis de que el
comportamiento del precio nacional de la tortilla puede ser explicado en
gran medida por los mismos factores que afectan al mercado del maíz
blanco.

Desde una perspectiva de política pública y planeación económica, esto
implica que acciones orientadas a estabilizar:

\begin{itemize}
\tightlist
\item
  los precios del maíz,
\item
  los costos de transporte,
\item
  la disponibilidad de producción agrícola,
\end{itemize}

podrían contribuir indirectamente a reducir la volatilidad de uno de los
productos básicos más importantes de la canasta alimentaria mexicana.

\begin{center}\rule{0.5\linewidth}{0.5pt}\end{center}

\hypertarget{trabajo-futuro}{%
\subsection{Trabajo Futuro}\label{trabajo-futuro}}

Como continuación de esta investigación se propone:

\begin{enumerate}
\def\labelenumi{\arabic{enumi}.}
\tightlist
\item
  Incorporar precios regionales de tortilla y maíz.
\item
  Integrar variables energéticas adicionales (diesel, gasolina y
  electricidad).
\item
  Analizar rezagos temporales entre el aumento del precio del maíz y su
  transmisión al precio de la tortilla.
\item
  Comparar modelos SARIMAX con enfoques de Machine Learning como Random
  Forest, XGBoost y LightGBM.
\item
  Construir modelos de optimización que permitan estimar impactos
  potenciales de cambios en costos agrícolas sobre el consumidor final.
\end{enumerate}

Esta línea de trabajo permitiría evolucionar desde el análisis del
mercado primario del maíz hacia un modelo integral de la cadena
agroalimentaria mexicana.

    \begin{tcolorbox}[breakable, size=fbox, boxrule=1pt, pad at break*=1mm,colback=cellbackground, colframe=cellborder]
\prompt{In}{incolor}{ }{\boxspacing}
\begin{Verbatim}[commandchars=\\\{\}]

\end{Verbatim}
\end{tcolorbox}


    % Add a bibliography block to the postdoc
    
    
    
\end{document}
